% Reconstructed text serializer; not generated by the unavailable supplied builder.

\chapter{Seminar records 597 and 599--639}

\label{ch:GY1983}

\section{Overview and editorial principles}

This chapter covers catalogue records 597 and 599--639, spanning the 1983--84
and 1984--85 academic-year headings and scan sequence P001113--P001198.
The catalogue contains 42 records, including documentary headings.
The occurrence inventory has 413 entries. An item is an individual recorded
definition, claim, construction, exercise, question, or programme; repetition
is preserved. Distinct canonical problems are indexed separately.

The catalogue page ranges overlap. Physical headings, the Russian
transcription, and the scans actually inspected determine the attribution.
The September Deligne report starts on P001115; the preceding April material
on P001113--P001114 is linked to earlier-year entries.
The missing-image Beilinson record 603 is reconstructed from the explicit
heading on P001124. The written date 3/X-84 is retained as evidence while the
session is placed on 3 October 1983 from the surrounding sequence.
On 14 November the physical order is Gelfand, Chulaevskii, Zelevinsky
(records 613, 612, 611); on 23 April it is Serre then Margulis
(records 633, 632). Unattributed interpolations remain unattributed.

Every occurrence separates the notebook's content from a qualified statement
and from the editor's additions. Literal correctness, publication-match
strength, evidence depth, and development status are independent fields.
A faithful mathematical error remains visible in the correction ledger.
A title-only identification or inaccessible original is not described as a
read theorem. Classification E denotes work already established or circulating
before the talk, even when its journal publication was later. Classification A
requires an identifiable near-term speaker or Gelfand continuation; B records
later work by others. C includes a precise residual after partial development.
D would require a documented search supporting absence of development.
F records an underspecified or unidentified item and supplies no inference of
openness. No D assessment is made merely because a search failed.

Several chronological distinctions materially change the account.
Deligne--Mostow circulated in July 1982, Vinberg's compact-reflection bound
was announced in 1981, and the Zumino--Wu--Zee preprint was available in
July 1983 \citep{delignemostow1986monodromy,vinberg1981absence,zuminowuzee1984anomalies}.
The O(4) Polyakov--Wiegmann publication and Baxter's tetrahedron verification
predate their seminar occurrences
\citep{polyakovwiegmann1983goldstone,baxter1983tetrahedron}.
The corresponding general SU(N), elliptic-fusion, hidden-symmetry, Toda,
superflag, affine singular-vector, absolute-Hodge, and dressing-action
continuations are evaluated individually rather than grouped by publication
year \citep{wiegmann1984sun,cherednik1985finite,gelfandzelevinsky1984models,
berezanskii1985integration,penkov1988borel,malikovfeiginfuchs1986singular,
beilinson1986absolute,semenov1985dressing}.

The research status is 12 November 2026. The checked Aggarwal--Elboim preprint
proves existence, with a quantitative remainder, of the maximum-dimension
exponential constant; its explicit evaluation is a distinct remaining task
\citep{aggarwalelboim2026maximal}. The direct Dyson-rank bijection question is
bounded by the checked 2019 formulation, not certified exhaustively open
through 2026 \citep{stanton2019rankcrank}. In the critical Ising discussion,
nearest-neighbour magnetization-field convergence is separated from a
precisely reconstructed arbitrary-strength finite-range universality question
\citep{camia2015magnetization,cavagiulianigreenblatt2025boundary}.
An inaccessible September 2026 bulk preprint is retained only as an
unverified lead.

The mathematical development proves a conditional covariance-scaling theorem
for the Ising residual, without claiming field convergence or priority.
Separate annotations give a finite-measure Radon identity, a counterexample
to a pole-count-only algebraic-integrability converse, a counterexample to
necessity of coisotropy for a fixed action, and two symbol repairs.
Independent reviewers checked the analytic arguments and constants; finite
JavaScript calculations are retained with source and output.
The execution environment failed during the run and the original supplied
builder became unavailable. This text therefore uses an independently
reconstructed, included serializer. Compilation and visual-inspection status
are recorded separately in the report and build-status file; the recovery
serializer is not represented as the supplied builder.


\section{Talks and mathematical occurrences}

\subsection{Record 599: 1983--1984 year heading and overlapping catalogue provenance}

\label{talk:599}

\noindent\textbf{Date:} 1983. \textbf{Speaker:} Editorial year heading. \textbf{Pages:} 1113--1117 (catalogue span).\par

This is an annual heading, not a mathematical lecture, and has zero items. P001113--P001114 contain adjacent April 1983 material belonging to the prior volume: moment-map/Duistermaat--Heckman material should be linked to1982-596.1--3, and zeta-function continuation to1982-595.1--7. The broad catalogue span overlaps later autumn lectures; actual source headings control their boundaries. This record is kept to preserve the supplied inventory.

\subsection{Record 597: Monodromy, families of abelian varieties, and hypergeometric lattices}

\label{talk:597}

\noindent\textbf{Date:} 1983-09-05. \textbf{Speaker:} P. Deligne. \textbf{Pages:} 1115--1116.\par

The explicit5 September Deligne heading is on P001115; the talk continues into the upper portion of P001116 and ends before the 12 September Arnold heading. The catalogue's broader P001113--P001115 range contains prior-April material and must not determine mathematical attribution. The notes mix local monodromy, semistable abelian degenerations, a degree bound, Deligne--Mostow examples and rational monodromy finiteness. The fixed-base hypotheses and equivalence relation are essential.

\subsubsection{1983-597.1 --- Semistable degenerations and semiabelian identity components}

\label{item:1983-597.1}

\noindent\textbf{Kind:} construction. \textbf{Fate:} E. \textbf{Pages:} 1115.\par

\paragraph{Notebook content.} The talk discusses extending a family of abelian varieties across a missing point by allowing a torus in the special fibre.

\paragraph{Precise statement.} For an abelian variety over the fraction field of a characteristic-zero discrete valuation ring, semistable reduction means that the identity component of the special fibre of its Neron model is an extension \(1\to T\to G\to B\to1\) of an abelian variety \(B\) by a torus \(T\). After a suitable finite extension of the base field, semistable reduction exists.

\paragraph{Editorial reconstruction.} This is a statement about the identity component of a group-scheme model. It does not say that the total proper family remains an abelian scheme at a genuinely degenerating fibre. A finite base change may be necessary.

\paragraph{Literature and fate.} The corrected formulation belongs to established mathematics available before the talk. Locator and verification scope: The opened publications verify context and metadata; the group-scheme formulation is an editorial clarification. \cite{faltings1983arakelov,schmid1973variation}.

\paragraph{Present status.} As of 12 November 2026: the stated qualified formulation is established within the cited and explicitly checked scope.

\paragraph{Audit.} Literal correctness: correct with missing hypotheses. Publication match: partial. Verification: publisher metadata checked. Development: annotated. Confidence: high.

\noindent References: \citep{faltings1983arakelov,schmid1973variation}.\par

\subsubsection{1983-597.2 --- Level structures and the torsion-free congruence condition}

\label{item:1983-597.2}

\noindent\textbf{Kind:} method. \textbf{Fate:} E. \textbf{Pages:} 1115.\par

\paragraph{Notebook content.} The notebook uses torsion points and a level condition to control local monodromy.

\paragraph{Precise statement.} If an integral matrix \(T\) is quasi-unipotent and \(T\equiv I\pmod m\) for an integer \(m\ge3\), then \(T\) is unipotent. Thus a full level-\(m\) structure makes a quasi-unipotent local monodromy action unipotent in the corresponding integral local system.

\paragraph{Editorial reconstruction.} The quasi-unipotence hypothesis remains indispensable. The principal congruence condition removes nontrivial finite-order eigenvalue factors, but it is not itself a proof of geometric quasi-unipotence. Characteristic-zero integral monodromy and level at least 3 are the intended setting.

\paragraph{Literature and fate.} The corrected formulation belongs to established mathematics available before the talk. Locator and verification scope: Classical monodromy context; congruence implication separated from the false converse. \cite{schmid1973variation}.

\paragraph{Present status.} As of 12 November 2026: the stated qualified formulation is established within the cited and explicitly checked scope.

\paragraph{Audit.} Literal correctness: correct with missing hypotheses. Publication match: partial. Verification: publisher metadata checked. Development: annotated. Confidence: high.

\noindent References: \citep{schmid1973variation}.\par

\subsubsection{1983-597.3 --- The literal congruence/quasi-unipotence equivalence fails}

\label{item:1983-597.3}

\noindent\textbf{Kind:} claim. \textbf{Fate:} E. \textbf{Pages:} 1115.\par

\paragraph{Notebook content.} The scan associates quasi-unipotence and a level-three congruence without retaining the needed one-way qualification.

\paragraph{Precise statement.} The condition \(T\equiv I\pmod3\) does not imply that an arbitrary integral invertible matrix \(T\) is quasi-unipotent. For example \(T=\begin{pmatrix}4&3\\9&7\end{pmatrix}\) has determinant1 and trace11; its eigenvalues are real reciprocal numbers different from roots of unity.

\paragraph{Editorial reconstruction.} Conversely, a unipotent matrix need not itself be congruent to the identity modulo3, as \(\begin{pmatrix}1&1\\0&1\end{pmatrix}\) shows. The valid geometric implication is the preceding item with both hypotheses. The counterexamples are exact elementary checks, not claims about geometric realizability of these matrices.

\paragraph{Literature and fate.} The corrected formulation belongs to established mathematics available before the talk. Locator and verification scope: Direct determinant, trace and congruence calculation; source context checked at metadata level. \cite{schmid1973variation}.

\paragraph{Present status.} As of 12 November 2026: the stated qualified formulation is established within the cited and explicitly checked scope.

\paragraph{Audit.} Literal correctness: false as written. Publication match: partial. Verification: publisher metadata checked. Development: annotated. Confidence: high.

\noindent References: \citep{schmid1973variation}.\par

\subsubsection{1983-597.4 --- A Hodge-degree bound and its missing data}

\label{item:1983-597.4}

\noindent\textbf{Kind:} claim. \textbf{Fate:} A. \textbf{Pages:} 1115.\par

\paragraph{Notebook content.} A bound for the degree of a family of abelian varieties is invoked in the finiteness discussion.

\paragraph{Precise statement.} For a polarized weight-one variation arising from a nonisotrivial abelian scheme on \(U=C\setminus D\), with \(C\) a fixed smooth projective complex curve and unipotent local monodromy, the degree of the canonical extension of its Hodge bundle is bounded in terms of its rank and \(2g(C)-2+|D|\).

\paragraph{Editorial reconstruction.} A precise inequality requires the Hodge-bundle convention, rank, base curve, punctures and semistable extension. The notebook does not supply enough data to fix the displayed constant. The modern formulation records only the supported boundedness content; it does not invent an exact numerical inequality.

\paragraph{Literature and fate.} Faltings's identified paper was published in October 1983 after the September talk. It is a close boundedness match, but the notebook's exact normalization and constant are not independently verified. Locator and verification scope: Faltings's October 1983 article, publisher metadata and title; full inequality not read. \cite{faltings1983arakelov}.

\paragraph{Present status.} As of 12 November 2026: the Hodge-degree boundedness principle is established, while this occurrence's exact constant is unverified.

\paragraph{Audit.} Literal correctness: correct with missing hypotheses. Publication match: partial. Verification: publisher metadata checked. Development: annotated. Confidence: medium.

\noindent References: \citep{faltings1983arakelov}.\par

\subsubsection{1983-597.5 --- Cyclic covers underlying hypergeometric monodromy}

\label{item:1983-597.5}

\noindent\textbf{Kind:} construction. \textbf{Fate:} E. \textbf{Pages:} 1115--1116.\par

\paragraph{Notebook content.} The notes give a branched-cover construction in the Deligne--Mostow discussion.

\paragraph{Precise statement.} For rational weights \(0<\mu_i<1\) with \(\sum_{i=1}^N\mu_i=2\), choose \(d\) with \(d\mu_i\in\mathbb Z\). A cyclic cover modeled by \(y^d=\prod_i(x-z_i)^{d\mu_i}\), normalized and compactified, realizes the relevant rank-\(N-2\) hypergeometric local system as a suitable eigenspace of its first cohomology as the distinct branch points \(z_i\) vary.

\paragraph{Editorial reconstruction.} The cover must be normalized; repeated branch points belong to boundary degenerations and are excluded from the open configuration space. Character choice can replace \(\mu_i\) by complementary weights, so the eigenvalue convention must be fixed.

\paragraph{Literature and fate.} The primary1986 publication contains the construction, and the IHES preprint is dated July 1982. It was therefore circulated before the September 1983 talk; classE records that chronology. Locator and verification scope: Section 12.9, cyclic covers and eigenspace realization. \cite{delignemostow1986monodromy}.

\paragraph{Present status.} As of 12 November 2026: the stated qualified formulation is established within the cited and explicitly checked scope.

\paragraph{Audit.} Literal correctness: correct with missing hypotheses. Publication match: exact. Verification: read the result. Development: annotated. Confidence: high.

\noindent References: \citep{delignemostow1986monodromy}.\par

\subsubsection{1983-597.6 --- The Deligne--Mostow lattice criterion}

\label{item:1983-597.6}

\noindent\textbf{Kind:} theorem. \textbf{Fate:} E. \textbf{Pages:} 1116, upper part.\par

\paragraph{Notebook content.} The discussion produces complex hyperbolic monodromy groups from weighted points.

\paragraph{Precise statement.} For rational weights \(0<\mu_i<1\), \(\sum_i\mu_i=2\), the Hermitian form on the relevant hypergeometric space has signature \((1,N-3)\). If \((1-\mu_i-\mu_j)^{-1}\in\mathbb Z\) whenever \(\mu_i+\mu_j<1\), the INT criterion yields a lattice in \(PU(1,N-3)\). Five points give a complex2-ball; six give a complex3-ball.

\paragraph{Editorial reconstruction.} An unrestricted choice of weights does not yield a lattice. The weaker symmetry-enhanced half-integrality criterion is a separate refinement and is not silently substituted for INT. Neither every such lattice nor every hypergeometric monodromy group is nonarithmetic.

\paragraph{Literature and fate.} The exact primary theorem is in the 1986 article, but the corresponding IHES preprint circulated in July 1982. ClassE therefore reflects prior circulation, not the later journal date. Locator and verification scope: Theorem 10.19 and section 11.4, especially p. 66. \cite{delignemostow1986monodromy}.

\paragraph{Present status.} As of 12 November 2026: the stated qualified formulation is established within the cited and explicitly checked scope.

\paragraph{Audit.} Literal correctness: correct with missing hypotheses. Publication match: exact. Verification: read the result. Development: annotated. Confidence: high.

\noindent References: \citep{delignemostow1986monodromy}.\par

\subsubsection{1983-597.7 --- Finiteness of rational geometric monodromy on a fixed base}

\label{item:1983-597.7}

\noindent\textbf{Kind:} theorem. \textbf{Fate:} A. \textbf{Pages:} 1116, upper part.\par

\paragraph{Notebook content.} Deligne announces a finiteness statement for monodromy representations.

\paragraph{Precise statement.} Fix a smooth connected complex algebraic variety \(S\) and an integer \(r\). There are only finitely many isomorphism classes of rank-\(r\) rational local systems of geometric origin, where geometric origin means a direct summand of \(R^jf_*\mathbb Q\) for a smooth projective morphism \(f:X\to S\), on \(S\) in Deligne's finiteness theorem. The equivalence is rational isomorphism, hence conjugacy over \(\mathbb Q\), not an unqualified integral-conjugacy assertion.

\paragraph{Editorial reconstruction.} The geometric-origin and fixed-base hypotheses cannot be omitted: arbitrary rank-one complex characters already form continuous families. The precise category and semisimplicity conditions are those of the cited theorem; no claim is made for every local system carrying only a formal Hodge filtration. The smooth-projective/direct-summand definition excludes arbitrary mixed or unipotent extensions; this hypothesis is part of the theorem.

\paragraph{Literature and fate.} Deligne published the finiteness theorem in 1987. The primary publication was opened; the exact rational formulation was read in Achinger's author notes, rather than inferred from thematic similarity. Locator and verification scope: Achinger Theorem 1 and Theorem 2; Deligne's1987 primary bibliographic record. \cite{deligne1987finitude,achinger2018monodromy}.

\paragraph{Present status.} As of 12 November 2026: the fixed-base rational finiteness theorem is established. Direct verification of the 1987 proof is limited because the opened scan had no usable text.

\paragraph{Audit.} Literal correctness: correct with missing hypotheses. Publication match: exact. Verification: read the result. Development: annotated. Confidence: high.

\noindent References: \citep{deligne1987finitude,achinger2018monodromy}.\par

\subsubsection{1983-597.8 --- A finiteness formulation with an unspecified varying base}

\label{item:1983-597.8}

\noindent\textbf{Kind:} claim. \textbf{Fate:} F. \textbf{Pages:} 1116, upper part.\par

\paragraph{Notebook content.} A further shorthand finiteness sentence does not clearly fix the base or the equivalence relation.

\paragraph{Precise statement.} No unique theorem can be reconstructed from a claim of finitely many monodromies while both the base and the notion of equivalence remain unspecified. Fixing dimension, rank or genus alone is not a substitute for specifying the algebraic family and the allowed identifications.

\paragraph{Editorial reconstruction.} The preceding fixed-base theorem is a candidate clarification, not an exact identification of every broader sentence. The occurrence is retained separately so that a stronger claim is not silently imported into Deligne's published theorem.

\paragraph{Literature and fate.} The surviving wording does not identify a unique result or publication. Locator and verification scope: Comparison with Achinger Theorem 1; no exact match for the broader fragment. \cite{deligne1987finitude,achinger2018monodromy}.

\paragraph{Present status.} As of 12 November 2026: the precise content remains unidentified within the documented search scope; this does not establish that a well-posed problem is open.

\paragraph{Audit.} Literal correctness: not determinable. Publication match: partial. Verification: indirect. Development: annotated. Confidence: low.

\noindent References: \citep{deligne1987finitude,achinger2018monodromy}.\par

\subsubsection{1983-597.9 --- Rational geometric local systems as the general formulation}

\label{item:1983-597.9}

\noindent\textbf{Kind:} remark. \textbf{Fate:} A. \textbf{Pages:} 1116, end of Deligne section.\par

\paragraph{Notebook content.} The abelian-family discussion is generalized to geometric local systems.

\paragraph{Precise statement.} The monodromy-finiteness framework applies to rational local systems of geometric origin, where geometric origin means a direct summand of \(R^jf_*\mathbb Q\) for a smooth projective morphism \(f:X\to S\), of fixed rank on a fixed smooth complex algebraic base, beyond the weight-one local systems of abelian schemes.

\paragraph{Editorial reconstruction.} This is a distinct generalization occurrence, although it overlaps the preceding theorem mathematically. It is not a statement about all integral lattices inside a rational representation. Earlier quasi-unipotent monodromy occurs at 1969-182.1. The smooth-projective/direct-summand definition excludes arbitrary mixed or unipotent extensions; this hypothesis is part of the theorem.

\paragraph{Literature and fate.} The general rational statement is identified with Deligne's1987 finiteness theorem; the exact theorem formulation was checked in the opened author notes. Locator and verification scope: Achinger Theorem 1, with variation-of-Hodge-structure discussion in Theorem 2. \cite{deligne1987finitude,achinger2018monodromy}.

\paragraph{Present status.} As of 12 November 2026: the stated qualified formulation is established within the cited and explicitly checked scope.

\paragraph{Audit.} Literal correctness: correct with missing hypotheses. Publication match: exact. Verification: read the result. Development: annotated. Confidence: high.

\noindent References: \citep{deligne1987finitude,achinger2018monodromy}.\par

\subsection{Record 600: Sperner theory, polytopes, reflection groups, and quadratic dynamics}

\label{talk:600}

\noindent\textbf{Date:} 1983-09-12. \textbf{Speaker:} V. I. Arnold, with seminar comments. \textbf{Pages:} 1116--1119.\par

The12 September heading in the lower part of P001116 follows Deligne. Arnold discusses rectangular partition posets and hard Lefschetz, the polytope g-theorem, Vinberg's dimension bound and the quadratic map \(z^2+i\). Bare notices of Sato, MacPherson and Yau do not identify new mathematical problems. The parameter-density remark lacks an ambient set and must not be silently converted into the density-of-hyperbolicity conjecture.

\subsubsection{1983-600.1 --- The notice of Sato's work}

\label{item:1983-600.1}

\noindent\textbf{Kind:} remark. \textbf{Fate:} F. \textbf{Pages:} 1116, after12September heading.\par

\paragraph{Notebook content.} Sato is named in a brief notice without a complete theorem or publication title.

\paragraph{Precise statement.} The notice does not determine a unique mathematical statement.

\paragraph{Editorial reconstruction.} The surrounding subjects include representation theory and integrable systems, but neither context identifies a particular Sato result. No question is invented from the name alone.

\paragraph{Literature and fate.} The surviving wording does not identify a unique result or publication. Locator and verification scope: No uniquely identifying result locator survives.

\paragraph{Present status.} As of 12 November 2026: the precise content remains unidentified within the documented search scope; this does not establish that a well-posed problem is open.

\paragraph{Audit.} Literal correctness: not determinable. Publication match: none found. Verification: unresolved lead. Development: annotated. Confidence: low.

\subsubsection{1983-600.2 --- The notice of MacPherson's work}

\label{item:1983-600.2}

\noindent\textbf{Kind:} remark. \textbf{Fate:} F. \textbf{Pages:} 1116, opening notices.\par

\paragraph{Notebook content.} MacPherson is named in a separate brief report without a precise assertion.

\paragraph{Precise statement.} The notice does not specify a unique theorem, invariant, or problem.

\paragraph{Editorial reconstruction.} The hard-Lefschetz and intersection-homology context supplies possible leads only. The attribution is retained separately from the Sato notice, and is not treated as an exact publication identification.

\paragraph{Literature and fate.} The surviving wording does not identify a unique result or publication. Locator and verification scope: No uniquely identifying result locator survives.

\paragraph{Present status.} As of 12 November 2026: the precise content remains unidentified within the documented search scope; this does not establish that a well-posed problem is open.

\paragraph{Audit.} Literal correctness: not determinable. Publication match: none found. Verification: unresolved lead. Development: annotated. Confidence: low.

\subsubsection{1983-600.3 --- The rectangular partition poset}

\label{item:1983-600.3}

\noindent\textbf{Kind:} definition. \textbf{Fate:} E. \textbf{Pages:} 1116--1117.\par

\paragraph{Notebook content.} Arnold introduces Young diagrams fitting in a rectangle.

\paragraph{Precise statement.} For positive integers \(m,n\), let \(L(m,n)\) be the set of partitions \(\lambda=(\lambda_1\ge\cdots\ge\lambda_m\ge0)\) with \(\lambda_1\le n\), ordered by componentwise inclusion. Its rank is \(|\lambda|=\sum_i\lambda_i\), ranging from 0 to \(mn\).

\paragraph{Editorial reconstruction.} Interchanging the two rectangle dimensions transposes diagrams and gives an isomorphic poset. Rank counts boxes, not rows or the number of nonzero parts.

\paragraph{Literature and fate.} The corrected formulation belongs to established mathematics available before the talk. Locator and verification scope: Section 4, rectangular partition lattices. \cite{stanley1980weyl}.

\paragraph{Present status.} As of 12 November 2026: the stated qualified formulation is established within the cited and explicitly checked scope.

\paragraph{Audit.} Literal correctness: correct with missing hypotheses. Publication match: exact. Verification: read the result. Development: annotated. Confidence: high.

\noindent References: \citep{stanley1980weyl}.\par

\subsubsection{1983-600.4 --- Width and antichains}

\label{item:1983-600.4}

\noindent\textbf{Kind:} definition. \textbf{Fate:} E. \textbf{Pages:} 1117.\par

\paragraph{Notebook content.} The maximum number of incomparable diagrams is considered.

\paragraph{Precise statement.} An antichain in a partially ordered set is a subset whose distinct elements are pairwise incomparable. Its width is the maximum cardinality of an antichain; a rank level in \(L(m,n)\) is always an antichain.

\paragraph{Editorial reconstruction.} The nontrivial claim is that a largest rank level has maximum possible cardinality, not merely that rank levels are antichains.

\paragraph{Literature and fate.} The corrected formulation belongs to established mathematics available before the talk. Locator and verification scope: Introduction and Theorem 3.1. \cite{stanley1980weyl}.

\paragraph{Present status.} As of 12 November 2026: the stated qualified formulation is established within the cited and explicitly checked scope.

\paragraph{Audit.} Literal correctness: correct. Publication match: exact. Verification: read the result. Development: annotated. Confidence: high.

\noindent References: \citep{stanley1980weyl}.\par

\subsubsection{1983-600.5 --- Sperner and strong Sperner properties of L(m,n)}

\label{item:1983-600.5}

\noindent\textbf{Kind:} theorem. \textbf{Fate:} E. \textbf{Pages:} 1117.\par

\paragraph{Notebook content.} The largest rank level is asserted to solve the width problem.

\paragraph{Precise statement.} The poset \(L(m,n)\) is rank symmetric, rank unimodal and strongly Sperner. In particular its width is the coefficient of \(q^{\lfloor mn/2\rfloor}\) in the Gaussian binomial \({m+n\brack m}_q=\sum_{\lambda\in L(m,n)}q^{|\lambda|}\).

\paragraph{Editorial reconstruction.} Strong Sperner means that a union of \(k\) antichains has cardinality at most the sum of the \(k\) largest rank sizes. This stronger formulation is supplied by the identified theorem; the width assertion is its \(k=1\) case.

\paragraph{Literature and fate.} The corrected formulation belongs to established mathematics available before the talk. Locator and verification scope: Stanley 1980 Theorem 3.1 and section 4;1984 section 4, Theorem 1. \cite{stanley1980weyl,stanley1984lefschetz}.

\paragraph{Present status.} As of 12 November 2026: the stated qualified formulation is established within the cited and explicitly checked scope.

\paragraph{Audit.} Literal correctness: correct with missing hypotheses. Publication match: exact. Verification: read the result. Development: annotated. Confidence: high.

\noindent Canonical problem: GP-1983-07.\par

\noindent References: \citep{stanley1980weyl,stanley1984lefschetz}.\par

\subsubsection{1983-600.6 --- Grassmannians, Schubert classes, and hard Lefschetz}

\label{item:1983-600.6}

\noindent\textbf{Kind:} method. \textbf{Fate:} E. \textbf{Pages:} 1117.\par

\paragraph{Notebook content.} The combinatorial theorem is related to a projective embedding and Schubert cells.

\paragraph{Precise statement.} Schubert varieties in \(\operatorname{Gr}(m,m+n)\) are indexed by the rectangle partitions. Their cohomology classes form a basis compatible with the incidence order, and cup product with the ample Plucker hyperplane class gives the raising operator used with hard Lefschetz to prove rank unimodality and Sperner statements.

\paragraph{Editorial reconstruction.} Cells are disjoint, so inclusion refers to their closures, the Schubert varieties. The standard Grassmannian embedding is Plucker; calling it Veronese confuses two distinct constructions. Hard Lefschetz is applied to the smooth projective Grassmannian, not to arbitrary singular Schubert varieties without an intersection-cohomology qualification.

\paragraph{Literature and fate.} The corrected formulation belongs to established mathematics available before the talk. Locator and verification scope: Theorem 2.4, Theorem 3.1 and section 4. \cite{stanley1980weyl}.

\paragraph{Present status.} As of 12 November 2026: the stated qualified formulation is established within the cited and explicitly checked scope.

\paragraph{Audit.} Literal correctness: correct with missing hypotheses. Publication match: exact. Verification: read the result. Development: annotated. Confidence: high.

\noindent References: \citep{stanley1980weyl}.\par

\subsubsection{1983-600.7 --- The Horrocks attribution}

\label{item:1983-600.7}

\noindent\textbf{Kind:} remark. \textbf{Fate:} E. \textbf{Pages:} 1117.\par

\paragraph{Notebook content.} Horrocks is cited in connection with Schubert varieties and symmetric-function calculations.

\paragraph{Precise statement.} Horrocks's1957 paper \emph{On the relation of S-functions to Schubert varieties} is a plausible bibliographic match for this connection.

\paragraph{Editorial reconstruction.} The publisher issue contents verifies the work's existence and metadata, but the original theorem was not accessible. The notebook's specific sketch therefore remains only partially matched; it cannot be used to assign a precise theorem to Horrocks from title similarity alone.

\paragraph{Literature and fate.} The corrected formulation belongs to established mathematics available before the talk. Locator and verification scope: Publisher issue contents, pp. 265--280; direct article page inaccessible. \cite{horrocks1957schubert}.

\paragraph{Present status.} As of 12 November 2026: the cited1957 work is identified bibliographically, but the exact result intended by this occurrence is unverified.

\paragraph{Audit.} Literal correctness: correct with missing hypotheses. Publication match: partial. Verification: publisher metadata checked. Development: annotated. Confidence: medium.

\noindent References: \citep{horrocks1957schubert}.\par

\subsubsection{1983-600.8 --- Euler's relation and the simple/simplicial distinction}

\label{item:1983-600.8}

\noindent\textbf{Kind:} definition and identity. \textbf{Fate:} E. \textbf{Pages:} 1118.\par

\paragraph{Notebook content.} The face-number discussion invokes Euler's formula and a special class of polytopes.

\paragraph{Precise statement.} For a convex \(d\)-polytope \(P\), let \(f_i\) count its \(i\)-faces and put \(f_d=1\). Then \(\sum_{i=0}^d(-1)^if_i=1\). A simple \(d\)-polytope has exactly \(d\) facets through each vertex; a simplicial polytope has every proper face a simplex. Polarity interchanges these classes.

\paragraph{Editorial reconstruction.} The face-polynomial convention in the following items is that of a simple polytope, equivalently the dual simplicial polytope. Interchanging the adjectives without dualizing changes the interpretation of the coefficients. The scan of P001118 was inspected for the terminology.

\paragraph{Literature and fate.} The corrected formulation belongs to established mathematics available before the talk. Locator and verification scope: Stanley p. 236 and Danilov's toric cohomology discussion. \cite{stanley1980faces,danilov1978toric}.

\paragraph{Present status.} As of 12 November 2026: the stated qualified formulation is established within the cited and explicitly checked scope.

\paragraph{Audit.} Literal correctness: correct with missing hypotheses. Publication match: exact. Verification: read the result. Development: annotated. Confidence: high.

\noindent References: \citep{stanley1980faces,danilov1978toric}.\par

\subsubsection{1983-600.9 --- The cube face-polynomial identity}

\label{item:1983-600.9}

\noindent\textbf{Kind:} example. \textbf{Fate:} E. \textbf{Pages:} 1118.\par

\paragraph{Notebook content.} The \(d\)-cube is used to illustrate the shifted face polynomial.

\paragraph{Precise statement.} For the \(d\)-cube, \(f_i=2^{d-i}\binom di\), so \(F(t)=\sum_{i=0}^df_it^i=(2+t)^d\). Hence the shifted polynomial \(H(t)=F(t-1)\) equals \((1+t)^d\), with coefficients \(\binom di\).

\paragraph{Editorial reconstruction.} This exact calculation fixes whether the notebook's variable shift is \(t-1\) or \(t+1\) and prevents confusing the face vector with the h-vector. It is a direct binomial computation.

\paragraph{Literature and fate.} The corrected formulation belongs to established mathematics available before the talk. Locator and verification scope: Direct verification under the simple-polytope convention. \cite{stanley1980faces}.

\paragraph{Present status.} As of 12 November 2026: the stated qualified formulation is established within the cited and explicitly checked scope.

\paragraph{Audit.} Literal correctness: correct. Publication match: exact. Verification: read the result. Development: annotated. Confidence: high.

\noindent References: \citep{stanley1980faces}.\par

\subsubsection{1983-600.10 --- Dehn--Sommerville symmetry and h-vector unimodality}

\label{item:1983-600.10}

\noindent\textbf{Kind:} theorem. \textbf{Fate:} E. \textbf{Pages:} 1118.\par

\paragraph{Notebook content.} The shifted coefficients are asserted to be symmetric and unimodal.

\paragraph{Precise statement.} For a simple convex \(d\)-polytope, define \(H(t)=\sum_{i=0}^df_i(t-1)^i=\sum_{i=0}^dh_it^i\). Then \(h_i=h_{d-i}\) and \(1=h_0\le h_1\le\cdots\le h_{\lfloor d/2\rfloor}\).

\paragraph{Editorial reconstruction.} The symmetry is Dehn--Sommerville; unimodality uses the deeper g-theorem/hard-Lefschetz input. This is not a claim that the unshifted face numbers have the same symmetry. For a simplicial polytope use the dual h-vector convention.

\paragraph{Literature and fate.} The corrected formulation belongs to established mathematics available before the talk. Locator and verification scope: Stanley pp. 236--237. \cite{stanley1980faces,danilov1978toric}.

\paragraph{Present status.} As of 12 November 2026: the stated qualified formulation is established within the cited and explicitly checked scope.

\paragraph{Audit.} Literal correctness: correct with missing hypotheses. Publication match: exact. Verification: read the result. Development: annotated. Confidence: high.

\noindent References: \citep{stanley1980faces,danilov1978toric}.\par

\subsubsection{1983-600.11 --- Characterizing polytope h-vectors: the g-theorem}

\label{item:1983-600.11}

\noindent\textbf{Kind:} question and theorem. \textbf{Fate:} E. \textbf{Pages:} 1118.\par

\paragraph{Notebook content.} The notebook asks which coefficient lists arise and mentions Billera--Lee and Stanley.

\paragraph{Precise statement.} An integer sequence \(h=(h_0,\ldots,h_d)\) is the h-vector of a simplicial convex \(d\)-polytope, equivalently of a dual simple polytope, precisely when \(h_0=1\), \(h_i=h_{d-i}\), and \(g=(1,h_1-h_0,\ldots,h_{\lfloor d/2\rfloor}-h_{\lfloor d/2\rfloor-1})\) is an M-sequence, that is the Hilbert function of a standard graded Artinian algebra through these degrees.

\paragraph{Editorial reconstruction.} Nonnegativity and unimodality alone are insufficient; the Macaulay growth conditions encoded by an M-sequence are essential. Billera--Lee proved sufficiency and Stanley necessity, rather than each supplying an independent proof of the full equivalence.

\paragraph{Literature and fate.} Both directions were published before the 1983 talk: Billera--Lee sufficiency and Stanley's1980 necessity. This is classE, a retrospective discussion of the solved characterization. Locator and verification scope: Stanley 1980 p. 236 states the criterion and direction attribution;1984 references the Billera--Lee1980 announcement and 1981 proof. \cite{stanley1980faces,stanley1984lefschetz}.

\paragraph{Present status.} As of 12 November 2026: the stated qualified formulation is established within the cited and explicitly checked scope.

\paragraph{Audit.} Literal correctness: correct with missing hypotheses. Publication match: exact. Verification: read the result. Development: annotated. Confidence: high.

\noindent Canonical problem: GP-1983-01.\par

\noindent References: \citep{stanley1980faces,stanley1984lefschetz}.\par

\subsubsection{1983-600.12 --- The toric hard-Lefschetz proof mechanism}

\label{item:1983-600.12}

\noindent\textbf{Kind:} method. \textbf{Fate:} E. \textbf{Pages:} 1118.\par

\paragraph{Notebook content.} The notes pass from a polytope to a toric space and invoke Lefschetz.

\paragraph{Precise statement.} A rational simple polytope defines a projective toric variety with at worst finite quotient singularities. Its rational cohomology realizes the h-numbers, and the ample class supplies hard Lefschetz. A general real simple polytope can be perturbed to a rational one with the same face lattice.

\paragraph{Editorial reconstruction.} Simple does not imply smooth toric: smoothness additionally requires primitive facet normals at each vertex to form a lattice basis. Rational cohomology and the quotient-singularity version of the argument suffice for the necessity proof.

\paragraph{Literature and fate.} The corrected formulation belongs to established mathematics available before the talk. Locator and verification scope: Stanley p. 237; Danilov sections6 and 10. \cite{stanley1980faces,danilov1978toric}.

\paragraph{Present status.} As of 12 November 2026: the stated qualified formulation is established within the cited and explicitly checked scope.

\paragraph{Audit.} Literal correctness: correct with missing hypotheses. Publication match: exact. Verification: read the result. Development: annotated. Confidence: high.

\noindent References: \citep{stanley1980faces,danilov1978toric}.\par

\subsubsection{1983-600.13 --- Yau's unnamed problem list}

\label{item:1983-600.13}

\noindent\textbf{Kind:} remark. \textbf{Fate:} F. \textbf{Pages:} 1118.\par

\paragraph{Notebook content.} Yau's problems are mentioned without a specific statement.

\paragraph{Precise statement.} The reference supplies no unique problem, invariant or geometric hypothesis.

\paragraph{Editorial reconstruction.} No item from a later published Yau problem list is assigned merely because the author and subject are plausible. This is an unresolved bibliographic notice, not an open-problem classification.

\paragraph{Literature and fate.} The surviving wording does not identify a unique result or publication. Locator and verification scope: No exact problem locator identified.

\paragraph{Present status.} As of 12 November 2026: the precise content remains unidentified within the documented search scope; this does not establish that a well-posed problem is open.

\paragraph{Audit.} Literal correctness: not determinable. Publication match: none found. Verification: unresolved lead. Development: annotated. Confidence: low.

\subsubsection{1983-600.14 --- Vinberg's dimension bound for compact reflection chambers}

\label{item:1983-600.14}

\noindent\textbf{Kind:} theorem. \textbf{Fate:} E. \textbf{Pages:} 1118--1119.\par

\paragraph{Notebook content.} The talk reports the absence of compact hyperbolic reflection chambers in sufficiently high dimension.

\paragraph{Precise statement.} There is no discrete group generated by hyperplane reflections in real hyperbolic space \(\mathbb H^d\) with a bounded fundamental polyhedron when \(d\ge30\). Equivalently there is no compact hyperbolic Coxeter polytope in these dimensions.

\paragraph{Editorial reconstruction.} Compactness is essential to this exact theorem. An unrestricted statement for noncompact finite-volume nonarithmetic chambers is different and is not justified by the bound30.

\paragraph{Literature and fate.} The exact bound is already Theorem 1 of the 1981 primary announcement, received 4 December1980. The1984 full paper and 1985 survey are later accounts. The survey's own elementary proof gives62 and explicitly cites the sharper bound30 elsewhere. Locator and verification scope: Vinberg 1981 Theorem 1;1985 section 6.4 and references17--19. \cite{vinberg1981absence,vinberg1985hyperbolic}.

\paragraph{Present status.} As of 12 November 2026: the stated qualified formulation is established within the cited and explicitly checked scope.

\paragraph{Audit.} Literal correctness: correct with missing hypotheses. Publication match: exact. Verification: read the result. Development: annotated. Confidence: high.

\noindent References: \citep{vinberg1981absence,vinberg1985hyperbolic}.\par

\subsubsection{1983-600.15 --- The critical orbit of z squared plus i}

\label{item:1983-600.15}

\noindent\textbf{Kind:} example. \textbf{Fate:} E. \textbf{Pages:} 1119.\par

\paragraph{Notebook content.} The map \(f(z)=z^2+i\) is used as a concrete quadratic example.

\paragraph{Precise statement.} Its critical orbit is \(0\mapsto i\mapsto-1+i\mapsto-i\mapsto-1+i\). Thus the unique critical point is strictly preperiodic to a two-cycle, and the polynomial is postcritically finite.

\paragraph{Editorial reconstruction.} Squaring verifies every arrow exactly. The critical point does not lie on the eventual cycle, so this is not a superattracting critical cycle.

\paragraph{Literature and fate.} The corrected formulation belongs to established mathematics available before the talk. Locator and verification scope: Direct complex-arithmetic check; Orsay notes for the postcritically finite setting. \cite{douadyhubbard1985polynomes}.

\paragraph{Present status.} As of 12 November 2026: the stated qualified formulation is established within the cited and explicitly checked scope.

\paragraph{Audit.} Literal correctness: correct. Publication match: exact. Verification: publisher metadata checked. Development: annotated. Confidence: high.

\noindent References: \citep{douadyhubbard1985polynomes}.\par

\subsubsection{1983-600.16 --- The parameter-density question lacks an ambient set}

\label{item:1983-600.16}

\noindent\textbf{Kind:} question. \textbf{Fate:} F. \textbf{Pages:} 1119.\par

\paragraph{Notebook content.} The notes ask whether parameters for which the critical orbit is finite are dense, without an unambiguous ambient parameter set.

\paragraph{Precise statement.} For \(f_c(z)=z^2+c\), the set of parameters with finite critical orbit is not dense in all of \(\mathbb C\): if \(|c|>2\), the orbit of 0 escapes, so the open set \(|c|>2\) contains no such parameter. Density relative to a bifurcation boundary is a different precise question.

\paragraph{Editorial reconstruction.} Neither the Mandelbrot boundary nor a particular moduli topology is specified clearly enough to replace the notebook's wording by that better-known problem. It must also not be confused with density of hyperbolicity, which concerns open hyperbolic components rather than the countable set of postcritically finite parameters.

\paragraph{Literature and fate.} The surviving wording does not identify a unique result or publication. Locator and verification scope: Direct escape estimate; Orsay notes as background only. \cite{douadyhubbard1985polynomes}.

\paragraph{Present status.} As of 12 November 2026: the literal whole-plane interpretation is false; the intended restricted-density question remains unidentified within the source and verified search scope.

\paragraph{Audit.} Literal correctness: not determinable. Publication match: thematic only. Verification: indirect. Development: annotated. Confidence: low.

\noindent Canonical problem: GP-1983-02.\par

\noindent References: \citep{douadyhubbard1985polynomes}.\par

\subsubsection{1983-600.17 --- Hyperbolic uniformization and contracting inverse branches}

\label{item:1983-600.17}

\noindent\textbf{Kind:} method. \textbf{Fate:} E. \textbf{Pages:} 1119.\par

\paragraph{Notebook content.} A uniformization argument is used to explain expansion near the Julia set.

\paragraph{Precise statement.} Remove the finite postcritical set \(P\) of a postcritically finite rational map from the sphere, and use the hyperbolic metric when the resulting surface is hyperbolic. On the full preimage of that punctured surface the map is an unbranched covering, hence locally an isometry for the corresponding covering metric; comparison with the smaller punctured domain yields contraction of inverse branches where the inclusion is strict.

\paragraph{Editorial reconstruction.} The full preimage, rather than the same punctured surface without qualification, is the covering domain. Strict contraction requires the proper-inclusion or hyperbolic-orbifold hypotheses; the Euclidean/parabolic exceptions must not be suppressed.

\paragraph{Literature and fate.} The corrected formulation belongs to established mathematics available before the talk. Locator and verification scope: Orsay notes, postcritical/uniformization framework; exact notebook proof not fully matched. \cite{douadyhubbard1985polynomes}.

\paragraph{Present status.} As of 12 November 2026: the stated qualified formulation is established within the cited and explicitly checked scope.

\paragraph{Audit.} Literal correctness: correct with missing hypotheses. Publication match: partial. Verification: indirect. Development: annotated. Confidence: medium.

\noindent References: \citep{douadyhubbard1985polynomes}.\par

\subsubsection{1983-600.18 --- Julia sets, dendrites, and the nonwandering-set shorthand}

\label{item:1983-600.18}

\noindent\textbf{Kind:} claim and remark. \textbf{Fate:} E. \textbf{Pages:} 1119.\par

\paragraph{Notebook content.} The final discussion refers to the invariant fractal set and to \ru{рац. числа} in the transcription.

\paragraph{Precise statement.} For \(z^2+i\), the connected locally connected Julia set is a dendrite: it contains no simple closed curve and no bounded Fatou component. A Julia set is not generally equal to the full nonwandering set of the rational map, which also includes nonwandering Fatou dynamics such as attracting cycles, including infinity for a polynomial on the sphere.

\paragraph{Editorial reconstruction.} The literal phrase \ru{рац. числа} means rational numbers. In context rational maps may have been intended, but without a decisive scan reading the transcription is retained and any emendation is marked conjectural. The Julia/nonwandering distinction is a mathematical correction independent of that lexical uncertainty.

\paragraph{Literature and fate.} The corrected formulation belongs to established mathematics available before the talk. Locator and verification scope: Postcritically finite polynomial setting; direct distinction using the fixed point at infinity. \cite{douadyhubbard1985polynomes}.

\paragraph{Present status.} As of 12 November 2026: the stated qualified formulation is established within the cited and explicitly checked scope.

\paragraph{Audit.} Literal correctness: false as written. Publication match: partial. Verification: indirect. Development: annotated. Confidence: medium.

\noindent References: \citep{douadyhubbard1985polynomes}.\par

\subsection{Record 601: Chiral models, topology of solitons, and Witten's topological term}

\label{talk:601}

\noindent\textbf{Date:} 1983-09-26. \textbf{Speaker:} A. M. Perelomov. \textbf{Pages:} 1120--1123.\par

The lecture begins on P001120 and ends before Daletskii's3 October heading on P001123. It moves from symmetry and particle interactions to sigma models, coadjoint-orbit instantons, scaling and the Skyrme term, and the five-dimensional Wess--Zumino functional. Several shorthand assertions require missing compactness, simple-connectedness, first-order self-duality and period-normalization hypotheses. P001122 was inspected for the decisive homotopy and scaling formulas.

\subsubsection{1983-601.1 --- Charge selection rules from U(1) symmetry}

\label{item:1983-601.1}

\noindent\textbf{Kind:} claim. \textbf{Fate:} E. \textbf{Pages:} 1120.\par

\paragraph{Notebook content.} The opening particle discussion links charge conservation with a circle symmetry.

\paragraph{Precise statement.} If a field transforms under \(U(1)\) with weight \(q\), then an invariant polynomial interaction monomial has total weight zero, counting conjugate fields with opposite weight. Under the corresponding classical symmetry, Noether's theorem gives a conserved charge subject to the field equations and boundary conditions.

\paragraph{Editorial reconstruction.} At the quantum level a classical symmetry may be anomalous. The algebraic selection rule is therefore separated from an unconditional claim of quantum conservation.

\paragraph{Literature and fate.} The corrected formulation belongs to established mathematics available before the talk. Locator and verification scope: General symmetry context; the cited abstract concerns anomalies and topological terms. \cite{witten1983global}.

\paragraph{Present status.} As of 12 November 2026: the stated qualified formulation is established within the cited and explicitly checked scope.

\paragraph{Audit.} Literal correctness: correct with missing hypotheses. Publication match: thematic only. Verification: publisher metadata checked. Development: annotated. Confidence: high.

\noindent References: \citep{witten1983global}.\par

\subsubsection{1983-601.2 --- Particles as charge carriers: a heuristic}

\label{item:1983-601.2}

\noindent\textbf{Kind:} remark. \textbf{Fate:} E. \textbf{Pages:} 1120.\par

\paragraph{Notebook content.} The notes describe particles as carriers of interaction or charge in a motivating physical discussion.

\paragraph{Precise statement.} A particle species may transform in a representation of a symmetry group, and its representation determines the charges to which that symmetry couples. This does not imply that every symmetry generator or every representation corresponds to an independently existing stable particle.

\paragraph{Editorial reconstruction.} The statement is interpreted as motivation for choosing fields and symmetry representations, not as a universal mathematical classification theorem. The original particle wording does not determine a sharper claim.

\paragraph{Literature and fate.} The corrected formulation belongs to established mathematics available before the talk. Locator and verification scope: No specific theorem is asserted.

\paragraph{Present status.} As of 12 November 2026: the stated qualified formulation is established within the cited and explicitly checked scope.

\paragraph{Audit.} Literal correctness: not determinable. Publication match: none found. Verification: indirect. Development: annotated. Confidence: medium.

\subsubsection{1983-601.3 --- The QED interaction vertex and derivative notation}

\label{item:1983-601.3}

\noindent\textbf{Kind:} formula. \textbf{Fate:} E. \textbf{Pages:} 1120.\par

\paragraph{Notebook content.} The electron--photon interaction is used as a model vertex.

\paragraph{Precise statement.} With \(D_\mu=\partial_\mu+ieA_\mu\), the Dirac term \(\overline\psi(i\gamma^\mu D_\mu-m)\psi\) contains the interaction \(-e\overline\psi\gamma^\mu A_\mu\psi\). The derivative belongs to the kinetic/covariant derivative term, not to an additional derivative of the photon field in this minimal-coupling vertex.

\paragraph{Editorial reconstruction.} The overall sign changes with the convention for \(D_\mu\). A derivative displayed as though it acts on \(A_\mu\) in the minimal vertex is not repaired merely by changing the charge sign. This is an elementary expansion of the stated covariant derivative.

\paragraph{Literature and fate.} The corrected formulation belongs to established mathematics available before the talk. Locator and verification scope: Direct algebraic expansion.

\paragraph{Present status.} As of 12 November 2026: the stated qualified formulation is established within the cited and explicitly checked scope.

\paragraph{Audit.} Literal correctness: false as written. Publication match: none found. Verification: indirect. Development: annotated. Confidence: high.

\subsubsection{1983-601.4 --- The nonlinear sigma-model action}

\label{item:1983-601.4}

\noindent\textbf{Kind:} definition. \textbf{Fate:} E. \textbf{Pages:} 1120--1121.\par

\paragraph{Notebook content.} The target-space action is introduced for a field with values in a Lie group or homogeneous space.

\paragraph{Precise statement.} For a map \(\phi:(M,g)\to(X,h)\), define \(S(\phi)=\frac12\int_M g^{\mu\nu}h_{ab}(\phi)\partial_\mu\phi^a\partial_\nu\phi^b\,d\operatorname{vol}_g\). For a compact group target with bi-invariant metric this is the principal chiral action expressed by \(g^{-1}dg\).

\paragraph{Editorial reconstruction.} Euclidean signature gives a nonnegative Dirichlet energy for a positive target metric; Lorentzian signature gives a wave-map action and need not be bounded below. These settings must not be interchanged in a minimization argument.

\paragraph{Literature and fate.} The corrected formulation belongs to established mathematics available before the talk. Locator and verification scope: Section 1, chiral/sigma-model definition. \cite{perelomov1981instantons}.

\paragraph{Present status.} As of 12 November 2026: the stated qualified formulation is established within the cited and explicitly checked scope.

\paragraph{Audit.} Literal correctness: correct with missing hypotheses. Publication match: exact. Verification: read the result. Development: annotated. Confidence: high.

\noindent References: \citep{perelomov1981instantons}.\par

\subsubsection{1983-601.5 --- Euler--Lagrange equations: harmonic maps and wave maps}

\label{item:1983-601.5}

\noindent\textbf{Kind:} equation. \textbf{Fate:} E. \textbf{Pages:} 1121.\par

\paragraph{Notebook content.} The variational field equation is written for the sigma action.

\paragraph{Precise statement.} The equation is \(\Box_g\phi^a+\Gamma^a_{bc}(\phi)g^{\mu\nu}\partial_\mu\phi^b\partial_\nu\phi^c=0\), with the appropriate Laplace--Beltrami or wave operator. For a group with bi-invariant metric it can be written \(\partial_\mu(g^{-1}\partial^\mu g)=0\) in flat coordinates.

\paragraph{Editorial reconstruction.} Coordinate derivatives of the metric are contained in the Christoffel term. A free linear wave equation describes only a flat-target or linearized approximation.

\paragraph{Literature and fate.} The corrected formulation belongs to established mathematics available before the talk. Locator and verification scope: Section 1, field equations. \cite{perelomov1981instantons}.

\paragraph{Present status.} As of 12 November 2026: the stated qualified formulation is established within the cited and explicitly checked scope.

\paragraph{Audit.} Literal correctness: correct with missing hypotheses. Publication match: exact. Verification: read the result. Development: annotated. Confidence: high.

\noindent References: \citep{perelomov1981instantons}.\par

\subsubsection{1983-601.6 --- Spatial compactification and the topological sector}

\label{item:1983-601.6}

\noindent\textbf{Kind:} construction. \textbf{Fate:} E. \textbf{Pages:} 1121.\par

\paragraph{Notebook content.} A field tending to a constant at spatial infinity is interpreted as a map from a sphere.

\paragraph{Precise statement.} A continuous map \(g:\mathbb R^3\to G\) with a specified constant limit at infinity extends to a based map \(S^3\to G\). Homotopy classes of such fields are elements of \(\pi_3(G)\).

\paragraph{Editorial reconstruction.} Finite energy alone must not be used without a decay/limit condition to assert continuous extension at the compactification point. The three-dimensional spatial sphere is distinct from four-dimensional spacetime compactification.

\paragraph{Literature and fate.} The corrected formulation belongs to established mathematics available before the talk. Locator and verification scope: pp. 434--435. \cite{witten1983current}.

\paragraph{Present status.} As of 12 November 2026: the stated qualified formulation is established within the cited and explicitly checked scope.

\paragraph{Audit.} Literal correctness: correct with missing hypotheses. Publication match: exact. Verification: read the result. Development: annotated. Confidence: high.

\noindent References: \citep{witten1983current}.\par

\subsubsection{1983-601.7 --- An energy lower bound does not prove existence of a minimizer}

\label{item:1983-601.7}

\noindent\textbf{Kind:} claim. \textbf{Fate:} E. \textbf{Pages:} 1121.\par

\paragraph{Notebook content.} A topological lower bound is used to motivate stable solutions.

\paragraph{Precise statement.} A bound \(E(\phi)\ge c\) on a homotopy class does not imply that \(\inf E\) is attained. For example the function \(E(x)=\varepsilon+e^{-x}\) on \(\mathbb R\), with \(\varepsilon>0\), has positive lower bound \(\varepsilon\) and no minimizer.

\paragraph{Editorial reconstruction.} Existence requires compactness or a suitable direct-method argument, including control of bubbling, escape and singular limits. This logical gap remains even when the lower bound is topological.

\paragraph{Literature and fate.} The corrected formulation belongs to established mathematics available before the talk. Locator and verification scope: Section 2(b) first-order bounds; the counterexample is an independent elementary check. \cite{perelomov1981instantons}.

\paragraph{Present status.} As of 12 November 2026: the stated qualified formulation is established within the cited and explicitly checked scope.

\paragraph{Audit.} Literal correctness: false as written. Publication match: partial. Verification: read the result. Development: annotated. Confidence: high.

\noindent References: \citep{perelomov1981instantons}.\par

\subsubsection{1983-601.8 --- Principal group targets versus coadjoint-orbit targets}

\label{item:1983-601.8}

\noindent\textbf{Kind:} definition. \textbf{Fate:} E. \textbf{Pages:} 1121.\par

\paragraph{Notebook content.} The discussion changes from group-valued fields to orbits.

\paragraph{Precise statement.} A principal chiral model has target \(G\). A coadjoint-orbit sigma model has target \(\mathcal O_\xi=G/G_\xi\) with a chosen invariant metric and, for compact \(G\), a natural symplectic structure. These are different target geometries and can have different homotopy groups.

\paragraph{Editorial reconstruction.} For example \(\pi_2(G)=0\) for a compact Lie group, whereas a nontrivial compact coadjoint orbit typically has nonzero \(\pi_2\). Topological charges cannot be transferred between the two without changing the model.

\paragraph{Literature and fate.} The corrected formulation belongs to established mathematics available before the talk. Locator and verification scope: Section 2(c), homogeneous targets. \cite{perelomov1981instantons}.

\paragraph{Present status.} As of 12 November 2026: the stated qualified formulation is established within the cited and explicitly checked scope.

\paragraph{Audit.} Literal correctness: correct with missing hypotheses. Publication match: exact. Verification: read the result. Development: annotated. Confidence: high.

\noindent References: \citep{perelomov1981instantons}.\par

\subsubsection{1983-601.9 --- Compact coadjoint orbits are homogeneous Kahler manifolds}

\label{item:1983-601.9}

\noindent\textbf{Kind:} theorem. \textbf{Fate:} E. \textbf{Pages:} 1121.\par

\paragraph{Notebook content.} Coadjoint orbits are described as Kahler and rational.

\paragraph{Precise statement.} A coadjoint orbit of a compact connected Lie group admits a compatible invariant complex structure making its Kirillov symplectic form Kahler; as a complex homogeneous space it is a generalized flag variety and is projective and rational.

\paragraph{Editorial reconstruction.} Compactness and the compatible complex structure are essential. The same assertion is false for arbitrary real noncompact coadjoint orbits. The conclusion concerns the orbit, not every simply connected compact Kahler manifold.

\paragraph{Literature and fate.} The corrected formulation belongs to established mathematics available before the talk. Locator and verification scope: Perelomov section 2(c); Borel publisher metadata. \cite{borel1954kahler,perelomov1981instantons}.

\paragraph{Present status.} As of 12 November 2026: the stated qualified formulation is established within the cited and explicitly checked scope.

\paragraph{Audit.} Literal correctness: correct with missing hypotheses. Publication match: exact. Verification: read the result. Development: annotated. Confidence: high.

\noindent References: \citep{borel1954kahler,perelomov1981instantons}.\par

\subsubsection{1983-601.10 --- The homotopy exact sequence for G/H}

\label{item:1983-601.10}

\noindent\textbf{Kind:} formula. \textbf{Fate:} E. \textbf{Pages:} 1121.\par

\paragraph{Notebook content.} The second homotopy group of an orbit is identified with a fundamental group of its stabilizer.

\paragraph{Precise statement.} For connected compact \(G\) and connected closed \(H\), the fibration \(H\to G\to G/H\) gives \(\pi_2(G/H)\simeq\ker(\pi_1(H)\to\pi_1(G))\), since \(\pi_2(G)=0\). If \(G\) is simply connected, this reduces to \(\pi_2(G/H)\simeq\pi_1(H)\).

\paragraph{Editorial reconstruction.} The simpler equality requires the simple-connectedness hypothesis on \(G\). If \(H\) is disconnected, additional components enter the long exact sequence and the formulation must be modified.

\paragraph{Literature and fate.} The corrected formulation belongs to established mathematics available before the talk. Locator and verification scope: Fibration long exact sequence; Perelomov section 2(c). \cite{hatcher2002algebraic,perelomov1981instantons}.

\paragraph{Present status.} As of 12 November 2026: the stated qualified formulation is established within the cited and explicitly checked scope.

\paragraph{Audit.} Literal correctness: correct with missing hypotheses. Publication match: exact. Verification: read the result. Development: annotated. Confidence: high.

\noindent References: \citep{hatcher2002algebraic,perelomov1981instantons}.\par

\subsubsection{1983-601.11 --- Counting invariant complex structures}

\label{item:1983-601.11}

\noindent\textbf{Kind:} claim. \textbf{Fate:} E. \textbf{Pages:} 1121.\par

\paragraph{Notebook content.} The Weyl group is used to count invariant complex structures.

\paragraph{Precise statement.} For a compact connected semisimple group \(G\) and a maximal torus \(T\), the invariant integrable complex structures on the full flag manifold \(G/T\) correspond to choices of positive roots and hence to the Weyl chambers; their number is \(|W|\).

\paragraph{Editorial reconstruction.} This full-flag count cannot be applied unchanged to every coadjoint orbit \(G/H\). Nor is \(|W|/|W_H|\) a universal count of invariant complex structures: the quotient can index other geometric data. The notebook's general orbit wording therefore needs restriction.

\paragraph{Literature and fate.} The corrected formulation belongs to established mathematics available before the talk. Locator and verification scope: Perelomov section 2(c), full-flag discussion. \cite{perelomov1981instantons,borel1954kahler}.

\paragraph{Present status.} As of 12 November 2026: the stated qualified formulation is established within the cited and explicitly checked scope.

\paragraph{Audit.} Literal correctness: correct with missing hypotheses. Publication match: partial. Verification: read the result. Development: annotated. Confidence: high.

\noindent References: \citep{perelomov1981instantons,borel1954kahler}.\par

\subsubsection{1983-601.12 --- Triangular-unitary factorization}

\label{item:1983-601.12}

\noindent\textbf{Kind:} method. \textbf{Fate:} E. \textbf{Pages:} 1121.\par

\paragraph{Notebook content.} A triangular and unitary decomposition is used to describe complex coordinates.

\paragraph{Precise statement.} Every invertible complex matrix admits a QR factorization \(A=QR\), with \(Q\) unitary and \(R\) upper triangular with positive real diagonal; equivalently an RQ factorization is obtained with a consistent reversal of convention. The positivity condition gives uniqueness.

\paragraph{Editorial reconstruction.} This is a finite-dimensional Iwasawa/Gram--Schmidt construction. Upper versus lower triangular factors and their order must be fixed before using them as coordinates on a flag chart.

\paragraph{Literature and fate.} The corrected formulation belongs to established mathematics available before the talk. Locator and verification scope: Section 2(c), complex homogeneous-space coordinates; elementary Gram--Schmidt reconstruction. \cite{perelomov1981instantons}.

\paragraph{Present status.} As of 12 November 2026: the stated qualified formulation is established within the cited and explicitly checked scope.

\paragraph{Audit.} Literal correctness: correct with missing hypotheses. Publication match: partial. Verification: read the result. Development: annotated. Confidence: high.

\noindent References: \citep{perelomov1981instantons}.\par

\subsubsection{1983-601.13 --- Holomorphic instantons versus arbitrary harmonic maps}

\label{item:1983-601.13}

\noindent\textbf{Kind:} claim. \textbf{Fate:} E. \textbf{Pages:} 1121.\par

\paragraph{Notebook content.} Complex analyticity is associated with the instanton equations.

\paragraph{Precise statement.} For a Kahler target, a holomorphic or antiholomorphic map from an oriented Riemann surface solves the sigma-model harmonic-map equation and saturates the corresponding topological energy bound. The converse requires a first-order self-duality or energy-saturation assumption.

\paragraph{Editorial reconstruction.} Not every harmonic map into a Kahler target is holomorphic. Perelomov's review distinguishes the self-dual instanton sector from other stationary solutions; suppressing this distinction makes the notebook's strongest literal reading false.

\paragraph{Literature and fate.} The corrected formulation belongs to established mathematics available before the talk. Locator and verification scope: Section 2(b), equations2.39--2.48; section 3(c). \cite{perelomov1981instantons}.

\paragraph{Present status.} As of 12 November 2026: the stated qualified formulation is established within the cited and explicitly checked scope.

\paragraph{Audit.} Literal correctness: false as written. Publication match: exact. Verification: read the result. Development: annotated. Confidence: high.

\noindent References: \citep{perelomov1981instantons}.\par

\subsubsection{1983-601.14 --- Rational-map instantons for CP1}

\label{item:1983-601.14}

\noindent\textbf{Kind:} example. \textbf{Fate:} E. \textbf{Pages:} 1121--1122.\par

\paragraph{Notebook content.} The sphere target gives the Belavin--Polyakov rational solutions.

\paragraph{Precise statement.} A finite-energy self-dual map \(\mathbb R^2\to S^2\simeq\mathbb{CP}^1\) with a constant limit at infinity extends to a holomorphic map \(S^2\to\mathbb{CP}^1\), hence is rational. Its antiholomorphic counterpart gives the opposite topological charge.

\paragraph{Editorial reconstruction.} The rationality theorem uses compactification and the self-duality equations. It is not an unrestricted classification of all time-dependent or higher-dimensional chiral fields.

\paragraph{Literature and fate.} The corrected formulation belongs to established mathematics available before the talk. Locator and verification scope: Section 1, CP1 model. \cite{perelomov1981instantons}.

\paragraph{Present status.} As of 12 November 2026: the stated qualified formulation is established within the cited and explicitly checked scope.

\paragraph{Audit.} Literal correctness: correct with missing hypotheses. Publication match: exact. Verification: read the result. Development: annotated. Confidence: high.

\noindent References: \citep{perelomov1981instantons}.\par

\subsubsection{1983-601.15 --- Low homotopy groups of SU2 and SU3}

\label{item:1983-601.15}

\noindent\textbf{Kind:} claim. \textbf{Fate:} E. \textbf{Pages:} 1122.\par

\paragraph{Notebook content.} The scan contains a \(\mathbb Z/2\)-type reading in the discussion of \(\pi_3\) and related groups.

\paragraph{Precise statement.} \(\pi_3(SU(2))=\pi_3(SU(3))=\mathbb Z\), while \(\pi_4(SU(2))=\pi_5(SU(2))=\mathbb Z/2\) and \(\pi_4(SU(3))=0\). More generally a compact connected simply connected simple group has \(\pi_3\simeq\mathbb Z\); a semisimple product contributes one copy per simple factor.

\paragraph{Editorial reconstruction.} The integer baryon/winding charge belongs to \(\pi_3\), not \(\mathbb Z/2\). Simple connectedness is needed for claims about \(\pi_1\), and the different roles of \(\pi_3,\pi_4,\pi_5\) must not be conflated. P001122 was visually inspected.

\paragraph{Literature and fate.} The corrected formulation belongs to established mathematics available before the talk. Locator and verification scope: Witten pp. 434--436; Hatcher section 4.J. \cite{witten1983current,hatcher2002algebraic}.

\paragraph{Present status.} As of 12 November 2026: the stated qualified formulation is established within the cited and explicitly checked scope.

\paragraph{Audit.} Literal correctness: false as written. Publication match: exact. Verification: read the result. Development: annotated. Confidence: high.

\noindent References: \citep{witten1983current,hatcher2002algebraic}.\par

\subsubsection{1983-601.16 --- Derrick scaling depends on the dimension}

\label{item:1983-601.16}

\noindent\textbf{Kind:} method. \textbf{Fate:} E. \textbf{Pages:} 1122.\par

\paragraph{Notebook content.} Rescaling is used to show the instability of the pure quadratic model.

\paragraph{Precise statement.} For a Dirichlet energy in \(d\) Euclidean dimensions and \(g_\lambda(x)=g(x/\lambda)\), \(E_2(g_\lambda)=\lambda^{d-2}E_2(g)\). Thus a four-dimensional Euclidean action scales as \(\lambda^2\), while the static energy in three spatial dimensions scales as \(\lambda\).

\paragraph{Editorial reconstruction.} The two exponents correspond to different integrals. A nonconstant stationary finite-energy field of the pure quadratic three-dimensional model cannot minimize against all scalings when the rescaled fields remain admissible. This scaling argument alone is not a general theorem excluding every saddle or singular solution.

\paragraph{Literature and fate.} The corrected formulation belongs to established mathematics available before the talk. Locator and verification scope: Skyrme-model discussion; direct change of variables. \cite{witten1983current}.

\paragraph{Present status.} As of 12 November 2026: the stated qualified formulation is established within the cited and explicitly checked scope.

\paragraph{Audit.} Literal correctness: correct with missing hypotheses. Publication match: exact. Verification: read the result. Development: annotated. Confidence: high.

\noindent References: \citep{witten1983current}.\par

\subsubsection{1983-601.17 --- The quartic Skyrme term balances static scaling}

\label{item:1983-601.17}

\noindent\textbf{Kind:} construction. \textbf{Fate:} E. \textbf{Pages:} 1122.\par

\paragraph{Notebook content.} A fourth-order derivative term is introduced to stabilize the scale.

\paragraph{Precise statement.} For a compact matrix group with \(R_i=g^{-1}\partial_i g\), a positive static quartic energy is a positive multiple of \(-\int\operatorname{tr}([R_i,R_j]^2)\,d^3x\). It scales as \(\lambda^{-1}\), so \(E_2+E_4\) has opposing scale powers \(\lambda\) and \(\lambda^{-1}\).

\paragraph{Editorial reconstruction.} Removing Derrick's immediate shrinking obstruction does not establish existence or uniqueness of a smooth global minimizer. The matrix trace sign assumes anti-Hermitian \(R_i\).

\paragraph{Literature and fate.} The corrected formulation belongs to established mathematics available before the talk. Locator and verification scope: pp. 433--435. \cite{witten1983current}.

\paragraph{Present status.} As of 12 November 2026: the stated qualified formulation is established within the cited and explicitly checked scope.

\paragraph{Audit.} Literal correctness: correct with missing hypotheses. Publication match: exact. Verification: read the result. Development: annotated. Confidence: high.

\noindent References: \citep{witten1983current}.\par

\subsubsection{1983-601.18 --- Filling a periodic spatial history}

\label{item:1983-601.18}

\noindent\textbf{Kind:} construction. \textbf{Fate:} E. \textbf{Pages:} 1122--1123.\par

\paragraph{Notebook content.} A periodic field history is extended from \(S^3\times S^1\) to \(S^3\times D^2\).

\paragraph{Precise statement.} For target \(SU(3)\), both \(\pi_1(SU(3))=0\) and \(\pi_4(SU(3))=0\) remove the obstructions to filling the loop of based spatial maps in a fixed \(\pi_3\) component. A loop in the free mapping space also has its evaluation-loop obstruction.

\paragraph{Editorial reconstruction.} The extension is not a consequence of \(\pi_1=0\) alone. The relative cells of the pair \(S^3\times D^2,S^3\times S^1\) explain why the degree-four homotopy group enters.

\paragraph{Literature and fate.} The corrected formulation belongs to established mathematics available before the talk. Locator and verification scope: pp. 434--436. \cite{witten1983current}.

\paragraph{Present status.} As of 12 November 2026: the stated qualified formulation is established within the cited and explicitly checked scope.

\paragraph{Audit.} Literal correctness: correct with missing hypotheses. Publication match: exact. Verification: read the result. Development: annotated. Confidence: high.

\noindent References: \citep{witten1983current}.\par

\subsubsection{1983-601.19 --- The cohomology generators of SU3}

\label{item:1983-601.19}

\noindent\textbf{Kind:} theorem. \textbf{Fate:} E. \textbf{Pages:} 1123.\par

\paragraph{Notebook content.} Invariant forms in degrees3 and 5 are invoked.

\paragraph{Precise statement.} The ring \(H^*(SU(3);\mathbb Z)\) is the exterior algebra on generators in degrees3 and 5. Invariant closed forms represent their real images; after choosing a bi-invariant metric they can be chosen as harmonic representatives.

\paragraph{Editorial reconstruction.} Forms represent cohomology rather than homology. Their scalar normalization is independent of the existence of the cohomology generator and must be specified before discussing integral periods. An integral-cohomology representative is distinct from Witten's physical five-form, which can have half-integral periods after division by2pi on unrestricted cycles.

\paragraph{Literature and fate.} The corrected formulation belongs to established mathematics available before the talk. Locator and verification scope: Current-algebra paper's invariant5-form discussion; global paper abstract. \cite{witten1983current,witten1983global}.

\paragraph{Present status.} As of 12 November 2026: the stated qualified formulation is established within the cited and explicitly checked scope.

\paragraph{Audit.} Literal correctness: correct with missing hypotheses. Publication match: partial. Verification: read the result. Development: annotated. Confidence: high.

\noindent References: \citep{witten1983current,witten1983global}.\par

\subsubsection{1983-601.20 --- The Wess--Zumino phase and integer level}

\label{item:1983-601.20}

\noindent\textbf{Kind:} construction and claim. \textbf{Fate:} E. \textbf{Pages:} 1123.\par

\paragraph{Notebook content.} The extra action is an integral over a five-dimensional extension with an integer coefficient.

\paragraph{Precise statement.} In Witten's normalization, let \(\Gamma(U)=\int_Q\widetilde U^*\omega_5\), with spacetime \(S^3\times S^1\) and filling \(Q=S^3\times D^2\). The allowed differences of fillings have periods in \(2\pi\mathbb Z\), so \(\exp(iN\Gamma(U))\) is well defined for \(N\in\mathbb Z\). The action is defined modulo \(2\pi\), not integer-valued.

\paragraph{Editorial reconstruction.} The allowed cycle class is essential. Witten notes that some closed five-cycles in \(SU(3)\) have odd multiples of \(\pi\) as periods of his form. They do not arise from the stated fillings: classes represented by \(S^3\times S^2\) are spherical in this setting. Thus his physical normalization must not be identified with a primitive integral-cohomology normalization. A form with all periods in \(2\pi\mathbb Z\) gives a more general sufficient construction, but changes this normalization. The rotation phase \((-1)^N\) in 1983-601.24 uses Witten's \(N\).

\paragraph{Literature and fate.} This precise restriction is stated in Witten's1983 paper, p. 434, the technical paragraph immediately before baryon number, and predates the talk. The general multivalued-action comparison is Novikov 1982 section 5. \cite{witten1983current,novikov1982hamiltonian}

\paragraph{Present status.} As of 12 November 2026: the stated qualified formulation is established within the cited and explicitly checked scope.

\paragraph{Audit.} Literal correctness: correct with missing hypotheses. Publication match: exact. Verification: read the result. Development: annotated. Confidence: high.

\noindent References: \citep{witten1983current,novikov1982hamiltonian}.\par

\subsubsection{1983-601.21 --- Relative chains make the extension ambiguity precise}

\label{item:1983-601.21}

\noindent\textbf{Kind:} method. \textbf{Fate:} E. \textbf{Pages:} 1123.\par

\paragraph{Notebook content.} The extension integral is compared for two choices of filling.

\paragraph{Precise statement.} If \(B_1,B_2\) are oriented fillings of the same mapped boundary and the maps agree there, glue \(B_1\) to \(-B_2\). The difference of the two integrals of a closed form equals its integral over the resulting closed cycle, or equivalently its pairing with the pushed-forward homology class.

\paragraph{Editorial reconstruction.} The fundamental class is pushed forward by the map; a form is pulled back. Confusing these operations obscures the integrality argument. A smooth singular-chain formulation also works when a single embedded filling is unavailable. Only periods on realizable differences of the allowed fillings control the ambiguity; these are not all integral5-cycles in Witten's setting.

\paragraph{Literature and fate.} The corrected formulation belongs to established mathematics available before the talk. Locator and verification scope: Witten extension argument; Novikov period interpretation. \cite{witten1983current,novikov1982hamiltonian}.

\paragraph{Present status.} As of 12 November 2026: the stated qualified formulation is established within the cited and explicitly checked scope.

\paragraph{Audit.} Literal correctness: correct with missing hypotheses. Publication match: exact. Verification: read the result. Development: annotated. Confidence: high.

\noindent References: \citep{witten1983current,novikov1982hamiltonian}.\par

\subsubsection{1983-601.22 --- Zumino's higher-dimensional anomaly construction}

\label{item:1983-601.22}

\noindent\textbf{Kind:} remark. \textbf{Fate:} E. \textbf{Pages:} 1123.\par

\paragraph{Notebook content.} Zumino is credited with a higher-dimensional formulation.

\paragraph{Precise statement.} A closely matching construction derives consistent chiral gauge anomalies in even dimension \(2n\) from invariant differential forms and descent, with the Wess--Zumino consistency condition built in.

\paragraph{Editorial reconstruction.} The notebook names Zumino but gives no title. Zumino--Wu--Zee's LBL16443 preprint is therefore a strong partial match, not a uniquely proved bibliographic identification.

\paragraph{Literature and fate.} The opened primary preprint is dated July 1983 and bears a16 September1983 library stamp, both before the 26 September talk. The journal version appeared9 July1984. Prior circulation supports classE. Locator and verification scope: July 1983 preprint cover; abstract and sectionsII--III. \cite{zuminowuzee1984anomalies}.

\paragraph{Present status.} As of 12 November 2026: the stated qualified formulation is established within the cited and explicitly checked scope.

\paragraph{Audit.} Literal correctness: correct with missing hypotheses. Publication match: partial. Verification: read the result. Development: annotated. Confidence: medium.

\noindent References: \citep{zuminowuzee1984anomalies}.\par

\subsubsection{1983-601.23 --- The baryon winding number}

\label{item:1983-601.23}

\noindent\textbf{Kind:} formula. \textbf{Fate:} E. \textbf{Pages:} 1123.\par

\paragraph{Notebook content.} An integer assigned to the spatial field is interpreted as baryon number.

\paragraph{Precise statement.} For a sufficiently regular based map \(g:S^3\to SU(N)\), \(N\ge2\), its winding number is \(\frac1{24\pi^2}\int_{S^3}\operatorname{tr}((g^{-1}dg)^3)\), up to the global orientation/sign convention for the chosen fundamental trace. It represents the integer in \(\pi_3(SU(N))\).

\paragraph{Editorial reconstruction.} The normalization must be calibrated on a generator; an arbitrary invariant trace rescales it. This is a topological charge of a configuration, not by itself a theorem identifying all physical baryons.

\paragraph{Literature and fate.} The corrected formulation belongs to established mathematics available before the talk. Locator and verification scope: pp. 434--435. \cite{witten1983current}.

\paragraph{Present status.} As of 12 November 2026: the stated qualified formulation is established within the cited and explicitly checked scope.

\paragraph{Audit.} Literal correctness: correct with missing hypotheses. Publication match: exact. Verification: read the result. Development: annotated. Confidence: high.

\noindent References: \citep{witten1983current}.\par

\subsubsection{1983-601.24 --- Rotation of the soliton and the phase (-1) to the level}

\label{item:1983-601.24}

\noindent\textbf{Kind:} claim. \textbf{Fate:} E. \textbf{Pages:} 1123, end of lecture.\par

\paragraph{Notebook content.} The history of a rotating soliton is used to motivate fermionic behavior.

\paragraph{Precise statement.} For the normalized unit-baryon Skyrme configuration in Witten's \(SU(3)\) construction, a \(2\pi\) spatial rotation contributes the topological phase \((-1)^N\), where \(N\) is the integer Wess--Zumino level. Odd \(N\) therefore gives the fermionic rotation sign.

\paragraph{Editorial reconstruction.} The statement fixes the topological sector and normalization and concerns the phase under rotation. A general quantum spin-statistics theorem or construction of the interacting quantum theory is not supplied by this calculation. Here \(N\) and \(\Gamma\) use Witten's normalization, with rotating-history \(\Gamma=\pi\) modulo \(2\pi\), as stated on p. 435.

\paragraph{Literature and fate.} The exact calculation was published 22 August1983, before this September lecture; the notebook is reporting the contemporary result. Locator and verification scope: p. 435, rotation phase. \cite{witten1983current}.

\paragraph{Present status.} As of 12 November 2026: the stated qualified formulation is established within the cited and explicitly checked scope.

\paragraph{Audit.} Literal correctness: correct with missing hypotheses. Publication match: exact. Verification: read the result. Development: annotated. Confidence: high.

\noindent References: \citep{witten1983current}.\par

\subsection{Record 602: Lie--Poisson brackets, kinetic equations, and point-particle reductions}

\label{talk:602}

\noindent\textbf{Date:} 1983-10-03. \textbf{Speaker:} Yu. L. Daletskii. \textbf{Pages:} 1123--1124.\par

Daletskii's heading begins in the lower part of P001123 and continues on P001124 until the explicit Beilinson heading. The date line reads3/X-84 in the transcription and scan context, but the surrounding autumn1983 sequence supports3 October1983 as the editorial session date; the discrepancy is retained. The lecture realizes particle and Vlasov equations in a dual Poisson algebra and mentions a graded extension with Kadobianskii, whose exact publication remains unidentified.

\subsubsection{1983-602.1 --- The Lie--Poisson bracket on a dual algebra}

\label{item:1983-602.1}

\noindent\textbf{Kind:} construction. \textbf{Fate:} E. \textbf{Pages:} 1123.\par

\paragraph{Notebook content.} Daletskii places evolution on the dual of a Lie algebra.

\paragraph{Precise statement.} For a Lie algebra \(\mathfrak g\), sufficiently differentiable functions \(F,G\) on a chosen dual \(\mathfrak g^*\) have the Lie--Poisson bracket \(\{F,G\}(\xi)=\langle\xi,[dF(\xi),dG(\xi)]\rangle\), provided their differentials are represented by elements of \(\mathfrak g\) and the pairings are defined.

\paragraph{Editorial reconstruction.} In infinite dimensions the dual, topology and admissible functionals must be specified; the formula does not automatically define a smooth Poisson manifold on every algebraic dual. Jacobi follows algebraically for the admissible functionals.

\paragraph{Literature and fate.} The corrected formulation belongs to established mathematics available before the talk. Locator and verification scope: p. 385, equation 12 in the kinetic realization. \cite{morrison1980maxwell}.

\paragraph{Present status.} As of 12 November 2026: the stated qualified formulation is established within the cited and explicitly checked scope.

\paragraph{Audit.} Literal correctness: correct with missing hypotheses. Publication match: exact. Verification: read the result. Development: annotated. Confidence: high.

\noindent References: \citep{morrison1980maxwell}.\par

\subsubsection{1983-602.2 --- The coadjoint Hamiltonian equation and its sign}

\label{item:1983-602.2}

\noindent\textbf{Kind:} equation. \textbf{Fate:} E. \textbf{Pages:} 1123--1124.\par

\paragraph{Notebook content.} The Hamiltonian equation is expressed through a coadjoint action.

\paragraph{Precise statement.} With \(\langle\operatorname{ad}_u^*\xi,v\rangle=-\langle\xi,[u,v]\rangle\) and \(\dot F=\{F,H\}\), the evolution is \(\dot\xi=\operatorname{ad}_{dH(\xi)}^*\xi\).

\paragraph{Editorial reconstruction.} Testing against the linear functional \(F_v(\xi)=\langle\xi,v\rangle\) fixes the sign: \(\dot F_v=\langle\xi,[v,dH]\rangle\). A convention defining the coadjoint differential without the minus sign reverses the displayed formula.

\paragraph{Literature and fate.} The corrected formulation belongs to established mathematics available before the talk. Locator and verification scope: Lie--Poisson sign checked directly against the primary bracket. \cite{morrison1980maxwell}.

\paragraph{Present status.} As of 12 November 2026: the stated qualified formulation is established within the cited and explicitly checked scope.

\paragraph{Audit.} Literal correctness: correct with missing hypotheses. Publication match: exact. Verification: read the result. Development: annotated. Confidence: high.

\noindent References: \citep{morrison1980maxwell}.\par

\subsubsection{1983-602.3 --- A symmetric-kernel quadratic Hamiltonian}

\label{item:1983-602.3}

\noindent\textbf{Kind:} definition. \textbf{Fate:} E. \textbf{Pages:} 1124.\par

\paragraph{Notebook content.} A two-point kernel defines an energy on distributions or measures.

\paragraph{Precise statement.} On a compact symplectic manifold \(M\), let \(K\in C^\infty(M\times M)\) satisfy \(K(x,y)=K(y,x)\). For a finite measure \(f\), set \(H(f)=\frac12\iint K(x,y)\,df(x)\,df(y)\); then \(\delta H/\delta f(x)=\int K(x,y)\,df(y)\).

\paragraph{Editorial reconstruction.} Smoothness and compactness make the expression and its derivatives well-defined. On noncompact phase space impose appropriate decay and integrability. A singular physical interaction kernel requires a separate treatment, especially on the diagonal.

\paragraph{Literature and fate.} The corrected formulation belongs to established mathematics available before the talk. Locator and verification scope: Kinetic Hamiltonian framework; the smooth-kernel specialization is an editorial exact formulation. \cite{morrison1980maxwell}.

\paragraph{Present status.} As of 12 November 2026: the stated qualified formulation is established within the cited and explicitly checked scope.

\paragraph{Audit.} Literal correctness: correct with missing hypotheses. Publication match: partial. Verification: read the result. Development: annotated. Confidence: high.

\noindent References: \citep{morrison1980maxwell}.\par

\subsubsection{1983-602.4 --- Point-particle measures and the Klimontovich reduction}

\label{item:1983-602.4}

\noindent\textbf{Kind:} construction. \textbf{Fate:} E. \textbf{Pages:} 1124.\par

\paragraph{Notebook content.} A sum of point masses recovers the finite-particle system.

\paragraph{Precise statement.} For \(f=\sum_{i=1}^Na_i\delta_{z_i}\), with fixed nonzero weights, the smooth-kernel Hamiltonian becomes \(H_N=\frac12\sum_{i,j}a_ia_jK(z_i,z_j)\) on \(M^N\), with symplectic form \(\sum_i a_i\omega_i\). Its Hamiltonian trajectories push forward to solutions of the measure-valued kinetic equation.

\paragraph{Editorial reconstruction.} The diagonal terms are present for a general smooth \(K\). They can be omitted only with an explicit no-self-interaction convention or when they contribute no force. The normalization of weights and symplectic form must agree.

\paragraph{Literature and fate.} The corrected formulation belongs to established mathematics available before the talk. Locator and verification scope: Lie--Poisson kinetic bracket; finite-particle reduction checked by testing linear observables. \cite{morrison1980maxwell}.

\paragraph{Present status.} As of 12 November 2026: the stated qualified formulation is established within the cited and explicitly checked scope.

\paragraph{Audit.} Literal correctness: correct with missing hypotheses. Publication match: partial. Verification: read the result. Development: annotated. Confidence: high.

\noindent References: \citep{morrison1980maxwell}.\par

\subsubsection{1983-602.5 --- Invariance of the fixed-weight atomic-measure stratum}

\label{item:1983-602.5}

\noindent\textbf{Kind:} claim. \textbf{Fate:} E. \textbf{Pages:} 1124.\par

\paragraph{Notebook content.} The particle subspace is asserted to be preserved by the evolution.

\paragraph{Precise statement.} When the Hamiltonian vector field is smooth enough to generate a flow, a sum of finitely many Dirac masses with fixed weights remains such a sum under transport: the support points follow the characteristic equations.

\paragraph{Editorial reconstruction.} This assertion concerns the atomic-measure stratum. It does not say that every distribution is a finite sum of point masses, or that an arbitrary distributional solution has this form. Collisions and nonuniqueness for singular kernels lie outside the stated smooth setting.

\paragraph{Literature and fate.} The corrected formulation belongs to established mathematics available before the talk. Locator and verification scope: Characteristic interpretation of the kinetic equation. \cite{morrison1980maxwell}.

\paragraph{Present status.} As of 12 November 2026: the stated qualified formulation is established within the cited and explicitly checked scope.

\paragraph{Audit.} Literal correctness: correct with missing hypotheses. Publication match: partial. Verification: read the result. Development: annotated. Confidence: high.

\noindent References: \citep{morrison1980maxwell}.\par

\subsubsection{1983-602.6 --- The Vlasov equation in canonical coordinates}

\label{item:1983-602.6}

\noindent\textbf{Kind:} equation. \textbf{Fate:} E. \textbf{Pages:} 1124.\par

\paragraph{Notebook content.} The continuum equation is given in Poisson coordinates.

\paragraph{Precise statement.} For \(\{a,b\}=\sum_j(a_{q_j}b_{p_j}-a_{p_j}b_{q_j})\) and \(h_f=\delta H/\delta f\), the density equation is \(\partial_tf+\{f,h_f\}=0\), or \(\partial_tf+\sum_j((h_f)_{p_j}f_{q_j}-(h_f)_{q_j}f_{p_j})=0\). For measures the equation is interpreted weakly.

\paragraph{Editorial reconstruction.} Integration by parts with suitable decay or a closed phase space gives \(\frac d{dt}\langle f,a\rangle=\langle f,\{a,h_f\}\rangle\). The1981 Maxwell-bracket correction does not invalidate the pure Vlasov bracket used here.

\paragraph{Literature and fate.} The corrected formulation belongs to established mathematics available before the talk. Locator and verification scope: Morrison p. 385, equation 12; Weinstein--Morrison correction. \cite{morrison1980maxwell,weinsteinmorrison1981comments}.

\paragraph{Present status.} As of 12 November 2026: the stated qualified formulation is established within the cited and explicitly checked scope.

\paragraph{Audit.} Literal correctness: correct with missing hypotheses. Publication match: exact. Verification: read the result. Development: annotated. Confidence: high.

\noindent References: \citep{morrison1980maxwell,weinsteinmorrison1981comments}.\par

\subsubsection{1983-602.7 --- The graded Poisson extension and the unresolved joint-paper attribution}

\label{item:1983-602.7}

\noindent\textbf{Kind:} remark and method. \textbf{Fate:} E. \textbf{Pages:} 1124, before Beilinson heading.\par

\paragraph{Notebook content.} Daletskii mentions a super or graded extension with Kadobianskii.

\paragraph{Precise statement.} An even graded Poisson algebra has a graded Lie bracket, and the dual Lie--Poisson construction extends algebraically with the appropriate Koszul signs on homogeneous observables. Grassmann variables can encode classical spin degrees of freedom.

\paragraph{Editorial reconstruction.} Ordinary positive measures cannot simply be assigned odd Grassmann-valued masses and retain their probabilistic interpretation. Berezin--Marinov supplies published background, but searches in English and Russian did not identify the exact Daletskii--Kadobianskii work or its precise analytic claims.

\paragraph{Literature and fate.} The corrected formulation belongs to established mathematics available before the talk. Locator and verification scope: Marinov institutional bibliography; graded construction used as background only. \cite{berezinmarinov1977particle}.

\paragraph{Present status.} As of 12 November 2026: the algebraic graded construction is established; the named joint-paper attribution and any stronger measure-theoretic assertion remain unresolved.

\paragraph{Audit.} Literal correctness: correct with missing hypotheses. Publication match: thematic only. Verification: publisher metadata checked. Development: annotated. Confidence: medium.

\noindent References: \citep{berezinmarinov1977particle}.\par

\subsection{Record 603: Witten's SU(2) anomaly and the topological obstruction}

\label{talk:603}

\noindent\textbf{Date:} 1983-10-03 (reconstructed). \textbf{Speaker:} A. A. Beilinson. \textbf{Pages:} 1124--1125.\par

Although the catalogue supplies no page links, the lower part of P001124 explicitly labels Beilinson and leads into the anomaly discussion on P001125. These surviving notes are assigned to603 by source-heading evidence, with reconstructed attribution and the 3/X-84 versus1983 date discrepancy retained. The argument concerns an odd number of Weyl doublets, a Pfaffian sign, suspension topology and a real mod2 index; it does not define a rigorous interacting gauge-field path integral.

\subsubsection{1983-603.1 --- The odd-doublet SU2 anomaly}

\label{item:1983-603.1}

\noindent\textbf{Kind:} theorem. \textbf{Fate:} E. \textbf{Pages:} 1124, Beilinson heading.\par

\paragraph{Notebook content.} The Beilinson section reports Witten's anomaly for \(SU(2)\).

\paragraph{Precise statement.} On a four-dimensional spin manifold, an odd number of left-handed Weyl fermions in the fundamental doublet representation of \(SU(2)\) has the familiar global gauge anomaly: the fermion Pfaffian cannot be assigned a gauge-invariant sign over all gauge-equivalent fields. The generator of \(\pi_4(SU(2))=\mathbb Z/2\) produces the obstruction in the \(S^4\) test.

\paragraph{Editorial reconstruction.} The odd-doublet, representation and spin hypotheses are essential. This is the 1982 anomaly, not the distinct nonspin anomaly studied in 2019. The explicit Beilinson label supports assignment to603 despite the catalogue's missing page links.

\paragraph{Literature and fate.} Witten's paper was published 18 November1982, before the October 1983 lecture. The2019 primary article supplies the detailed modern formulation of the older anomaly. Locator and verification scope: Witten primary abstract; Wang--Wen--Witten sections2.1--2.3. \cite{witten1982anomaly,wangwenwitten2019new}.

\paragraph{Present status.} As of 12 November 2026: the stated qualified formulation is established within the cited and explicitly checked scope.

\paragraph{Audit.} Literal correctness: correct with missing hypotheses. Publication match: exact. Verification: read the result. Development: annotated. Confidence: high.

\noindent References: \citep{witten1982anomaly,wangwenwitten2019new}.\par

\subsubsection{1983-603.2 --- The formal Yang--Mills/fermion functional integral}

\label{item:1983-603.2}

\noindent\textbf{Kind:} construction. \textbf{Fate:} E. \textbf{Pages:} 1124--1125.\par

\paragraph{Notebook content.} The argument is introduced through a path integral over gauge and fermion fields.

\paragraph{Precise statement.} Formally integrating a quadratic fermionic action yields a determinant or Pfaffian depending on whether the fermion variables are complex paired or real/pseudoreal. The anomaly tests the transformation of this fermionic factor under gauge transformations of the background connection.

\paragraph{Editorial reconstruction.} This formal manipulation is not a construction of a countably additive measure for four-dimensional interacting Yang--Mills theory. The sign obstruction can be defined rigorously through families of elliptic operators without assigning a value to the complete path integral.

\paragraph{Literature and fate.} The corrected formulation belongs to established mathematics available before the talk. Locator and verification scope: Wang--Wen--Witten section 2.2. \cite{witten1982anomaly,wangwenwitten2019new}.

\paragraph{Present status.} As of 12 November 2026: the stated qualified formulation is established within the cited and explicitly checked scope.

\paragraph{Audit.} Literal correctness: correct with missing hypotheses. Publication match: exact. Verification: read the result. Development: annotated. Confidence: high.

\noindent References: \citep{witten1982anomaly,wangwenwitten2019new}.\par

\subsubsection{1983-603.3 --- The Polyakov--Wiegmann aside}

\label{item:1983-603.3}

\noindent\textbf{Kind:} remark. \textbf{Fate:} F. \textbf{Pages:} 1124.\par

\paragraph{Notebook content.} Polyakov and Wiegmann are named without an accompanying title or formula.

\paragraph{Precise statement.} The names and neighboring two-dimensional field-theory context do not determine a unique mathematical assertion.

\paragraph{Editorial reconstruction.} The1983 paper on nonabelian Goldstone bosons is a plausible lead, published 10 November1983. Its title and subject alone do not prove that it is the work meant in this October aside. No Polyakov--Wiegmann identity absent from the notes is inserted as though stated here.

\paragraph{Literature and fate.} The surviving wording does not identify a unique result or publication. Locator and verification scope: Princeton institutional publication metadata only. \cite{polyakovwiegmann1983goldstone}.

\paragraph{Present status.} As of 12 November 2026: the precise content remains unidentified within the documented search scope; this does not establish that a well-posed problem is open.

\paragraph{Audit.} Literal correctness: not determinable. Publication match: thematic only. Verification: publisher metadata checked. Development: annotated. Confidence: low.

\noindent References: \citep{polyakovwiegmann1983goldstone}.\par

\subsubsection{1983-603.4 --- Components of the gauge group on S4}

\label{item:1983-603.4}

\noindent\textbf{Kind:} topological claim. \textbf{Fate:} E. \textbf{Pages:} 1125.\par

\paragraph{Notebook content.} The gauge transformation group has a nontrivial component.

\paragraph{Precise statement.} For the trivial \(SU(2)\)-bundle over \(S^4\), based gauge transformations are based maps \(S^4\to SU(2)\), so their component group is \(\pi_4(SU(2))=\mathbb Z/2\). The same component classification holds for the free mapping group here because the connected target introduces no additional component.

\paragraph{Editorial reconstruction.} For a nontrivial principal bundle, the gauge group is a section space of the adjoint bundle and this identification cannot simply be copied. The notebook's elementary model uses the trivial bundle.

\paragraph{Literature and fate.} The corrected formulation belongs to established mathematics available before the talk. Locator and verification scope: Witten abstract; Hatcher section 4.J. \cite{witten1982anomaly,hatcher2002algebraic}.

\paragraph{Present status.} As of 12 November 2026: the stated qualified formulation is established within the cited and explicitly checked scope.

\paragraph{Audit.} Literal correctness: correct with missing hypotheses. Publication match: exact. Verification: read the result. Development: annotated. Confidence: high.

\noindent References: \citep{witten1982anomaly,hatcher2002algebraic}.\par

\subsubsection{1983-603.5 --- Suspension of the Hopf map and reduction modulo two}

\label{item:1983-603.5}

\noindent\textbf{Kind:} theorem. \textbf{Fate:} E. \textbf{Pages:} 1125.\par

\paragraph{Notebook content.} The nontrivial \(\pi_4(S^3)\) class is related to the Hopf map.

\paragraph{Precise statement.} The suspension homomorphism \(E:\pi_3(S^2)\simeq\mathbb Z\to\pi_4(S^3)\simeq\mathbb Z/2\) is surjective with kernel \(2\mathbb Z\). A Hopf-map generator suspends to the nonzero element.

\paragraph{Editorial reconstruction.} The kernel is generated by the Whitehead square \([\iota_2,\iota_2]=\pm2\eta\). Sign depends on orientation, while the index-two kernel does not.

\paragraph{Literature and fate.} The corrected formulation belongs to established mathematics available before the talk. Locator and verification scope: Section 4.J, pp. 473--474, Whitehead-square kernel and Corollary4J.4. \cite{hatcher2002algebraic}.

\paragraph{Present status.} As of 12 November 2026: the stated qualified formulation is established within the cited and explicitly checked scope.

\paragraph{Audit.} Literal correctness: correct with missing hypotheses. Publication match: exact. Verification: read the result. Development: annotated. Confidence: high.

\noindent References: \citep{hatcher2002algebraic}.\par

\subsubsection{1983-603.6 --- Exotic smooth structures on R4 and later uncountable families}

\label{item:1983-603.6}

\noindent\textbf{Kind:} claim. \textbf{Fate:} B. \textbf{Pages:} 1125.\par

\paragraph{Notebook content.} The topology discussion refers to exotic smooth4-space and an uncountable multiplicity claim.

\paragraph{Precise statement.} There exist uncountably many pairwise nondiffeomorphic smooth manifolds homeomorphic to \(\mathbb R^4\). This stronger cardinality assertion must be distinguished from the first existence consequences of four-manifold gauge theory.

\paragraph{Editorial reconstruction.} Donaldson's1983 definite-intersection-form theorem is not by itself a proof of uncountably many exotic structures. Taubes's1987 work and DeMichelis--Freedman's1992 construction supply later developments; the latter's primary abstract explicitly gives a continuum of exotic \(\mathbb R^4\)'s in standard4-space.

\paragraph{Literature and fate.} The uncountability assertion is supported by later1987/1992 work, not assigned to Donaldson1983. The opened 1992 primary abstract verifies the cardinality claim; its full proof and Taubes's original proof were not read. Locator and verification scope: Donaldson main theorem; DeMichelis--Freedman primary abstract. Taubes role verified indirectly only. \cite{donaldson1983selfdual,taubes1987periodic,demichelisfreedman1992uncountably}.

\paragraph{Present status.} As of 12 November 2026: uncountable families are established; historical attribution in this brief occurrence is only partially resolved.

\paragraph{Audit.} Literal correctness: correct with missing hypotheses. Publication match: partial. Verification: indirect. Development: annotated. Confidence: medium.

\noindent References: \citep{donaldson1983selfdual,taubes1987periodic,demichelisfreedman1992uncountably}.\par

\subsubsection{1983-603.7 --- The EHP edge and the Whitehead-square kernel}

\label{item:1983-603.7}

\noindent\textbf{Kind:} method. \textbf{Fate:} E. \textbf{Pages:} 1125.\par

\paragraph{Notebook content.} A more general homotopy-theoretic mechanism is sketched around the suspension calculation.

\paragraph{Precise statement.} For \(n\ge2\), the edge of the EHP sequence gives \(\pi_{2n+1}(S^{2n+1})\xrightarrow{P}\pi_{2n-1}(S^n)\xrightarrow{E}\pi_{2n}(S^{n+1})\to0\), where \(P\) sends a generator to the Whitehead square \([\iota_n,\iota_n]\).

\paragraph{Editorial reconstruction.} This is the relevant range statement, not an assertion that every suspension map has the same kernel or that the displayed pattern is exact in all dimensions without range restrictions. The \(n=2\) case produces the preceding modulo-two result.

\paragraph{Literature and fate.} The corrected formulation belongs to established mathematics available before the talk. Locator and verification scope: Hatcher section 4.J, pp. 473--475; Serre primary scan opened but not used as a read theorem. \cite{hatcher2002algebraic,serre1953classes}.

\paragraph{Present status.} As of 12 November 2026: the stated qualified formulation is established within the cited and explicitly checked scope.

\paragraph{Audit.} Literal correctness: correct with missing hypotheses. Publication match: exact. Verification: read the result. Development: annotated. Confidence: high.

\noindent References: \citep{hatcher2002algebraic,serre1953classes}.\par

\subsubsection{1983-603.8 --- Fermionic Gaussian integration gives a Pfaffian, not its reciprocal}

\label{item:1983-603.8}

\noindent\textbf{Kind:} formula. \textbf{Fate:} E. \textbf{Pages:} 1125.\par

\paragraph{Notebook content.} The displayed Grassmann Gaussian is written with a determinant half-power requiring a sign correction.

\paragraph{Precise statement.} For an antisymmetric \(2r\times2r\) matrix \(A\), \(\int d\theta_{2r}\cdots d\theta_1\,\exp(\frac12\theta^tA\theta)=\operatorname{Pf}(A)\) up to the fixed measure-order sign, and \(\operatorname{Pf}(A)^2=\det A\).

\paragraph{Editorial reconstruction.} A reciprocal determinant square root is the bosonic Gaussian expression. The finite-dimensional algebra explains the sign problem, while an infinite-dimensional Pfaffian requires an operator-family regularization and a Pfaffian line.

\paragraph{Literature and fate.} The corrected formulation belongs to established mathematics available before the talk. Locator and verification scope: Section 2.2, real-fermion Pfaffian interpretation. \cite{wangwenwitten2019new}.

\paragraph{Present status.} As of 12 November 2026: the stated qualified formulation is established within the cited and explicitly checked scope.

\paragraph{Audit.} Literal correctness: false as written. Publication match: exact. Verification: read the result. Development: annotated. Confidence: high.

\noindent References: \citep{wangwenwitten2019new}.\par

\subsubsection{1983-603.9 --- Mod-two spectral flow and the mapping-torus index}

\label{item:1983-603.9}

\noindent\textbf{Kind:} theorem and method. \textbf{Fate:} E. \textbf{Pages:} 1125.\par

\paragraph{Notebook content.} The anomaly sign is related to eigenvalues crossing zero along a family.

\paragraph{Precise statement.} Join a connection \(A\) to its gauge transform \(A^g\) by a path and identify the endpoints with \(g\). In the real structure relevant to the \(SU(2)\) Weyl doublet, the parity of the appropriate eigenvalue-pair crossings equals the mod2 index of the five-dimensional Dirac operator on the resulting mapping torus. The Pfaffian sign is \((-1)^\nu\).

\paragraph{Editorial reconstruction.} One counts crossings with the pseudoreal/real degeneracy accounted for, not the ordinary signed spectral flow of an arbitrary ungraded self-adjoint family. The bundle clutching and spin structure are part of the construction.

\paragraph{Literature and fate.} The corrected formulation belongs to established mathematics available before the talk. Locator and verification scope: Wang--Wen--Witten section 2.2, pp6--9; Atiyah--Singer section 2. \cite{witten1982anomaly,wangwenwitten2019new,atiyahsinger1971index5}.

\paragraph{Present status.} As of 12 November 2026: the stated qualified formulation is established within the cited and explicitly checked scope.

\paragraph{Audit.} Literal correctness: correct with missing hypotheses. Publication match: exact. Verification: read the result. Development: annotated. Confidence: high.

\noindent References: \citep{witten1982anomaly,wangwenwitten2019new,atiyahsinger1971index5}.\par

\subsubsection{1983-603.10 --- A general operator framework is requested but not specified}

\label{item:1983-603.10}

\noindent\textbf{Kind:} question. \textbf{Fate:} F. \textbf{Pages:} 1125, closing discussion.\par

\paragraph{Notebook content.} The discussion asks for a broad mathematical framework for the formal anomaly argument.

\paragraph{Precise statement.} The surviving question does not specify a category of field theories, a regularization, a class of operator families, or the exact object to be constructed.

\paragraph{Editorial reconstruction.} Determinant and Pfaffian line theories address a precise part of this motivation, developed further in 604. It would be misleading to classify the entire undefined program as solved or open. No new stable problem identifier is assigned to this bare formulation.

\paragraph{Literature and fate.} The surviving wording does not identify a unique result or publication. Locator and verification scope: Later determinant-line framework is a contextual lead only. \cite{daifreed1994eta}.

\paragraph{Present status.} As of 12 November 2026: the precise content remains unidentified within the documented search scope; this does not establish that a well-posed problem is open.

\paragraph{Audit.} Literal correctness: not determinable. Publication match: thematic only. Verification: indirect. Development: annotated. Confidence: low.

\noindent References: \citep{daifreed1994eta}.\par

\subsection{Record 604: Dirac families, mapping tori, and the anomaly sign}

\label{talk:604}

\noindent\textbf{Date:} 1983-10-17. \textbf{Speaker:} A. A. Beilinson, with I. M. Gelfand's question. \textbf{Pages:} 1126--1127.\par

The17 October Beilinson continuation uses a chiral Dirac family and a mapping torus to interpret spectral flow. The Pauli commutator normalization, real skew-adjoint structure and endpoint identification must be fixed. Gelfand asks for a mathematical framework; determinant/Pfaffian lines and later eta-invariant holonomy provide a precise partial realization, while the broad construction of an interacting quantum field theory is not claimed solved.

\subsubsection{1983-604.1 --- The chiral Dirac operator in four dimensions}

\label{item:1983-604.1}

\noindent\textbf{Kind:} construction. \textbf{Fate:} E. \textbf{Pages:} 1126.\par

\paragraph{Notebook content.} Beilinson writes a Dirac operator using two-component spinors.

\paragraph{Precise statement.} On an oriented Riemannian spin4-manifold, the complex chiral spinor bundles \(S^+\) and \(S^-\) have rank2. A connection on a Hermitian gauge bundle \(E\) defines \(D_A^+:\Gamma(S^+\otimes E)\to\Gamma(S^-\otimes E)\), with adjoint \(D_A^-\); the full Dirac operator has these as its off-diagonal blocks.

\paragraph{Editorial reconstruction.} The chiral operator alone is not an endomorphism of one spinor bundle and is not self-adjoint in that sense. Color/gauge indices are additional to the two spin indices. A spin structure or a specified replacement is required.

\paragraph{Literature and fate.} The corrected formulation belongs to established mathematics available before the talk. Locator and verification scope: Section 2.2, four-dimensional spinors and the relevant real structure. \cite{wangwenwitten2019new}.

\paragraph{Present status.} As of 12 November 2026: the stated qualified formulation is established within the cited and explicitly checked scope.

\paragraph{Audit.} Literal correctness: correct with missing hypotheses. Publication match: exact. Verification: read the result. Development: annotated. Confidence: high.

\noindent References: \citep{wangwenwitten2019new}.\par

\subsubsection{1983-604.2 --- Why a sign-changing formal integral is an obstruction}

\label{item:1983-604.2}

\noindent\textbf{Kind:} claim. \textbf{Fate:} E. \textbf{Pages:} 1126.\par

\paragraph{Notebook content.} A change of variables is used to argue that a formal partition function equals its own negative.

\paragraph{Precise statement.} If a formal fermionic measure changes sign under a gauge transformation while all bosonic terms and the domain are invariant, the formal integral would satisfy \(Z=-Z\). The mathematically useful conclusion is the obstruction to defining a gauge-invariant fermion measure or Pfaffian sign.

\paragraph{Editorial reconstruction.} One must not treat this as an ordinary numerical proof that a previously constructed infinite-dimensional integral is zero: such a measure has not been supplied. The line-valued formulation retains the obstruction without requiring that nonexistent numerical step.

\paragraph{Literature and fate.} The corrected formulation belongs to established mathematics available before the talk. Locator and verification scope: Witten abstract; Wang--Wen--Witten section 2.2. \cite{witten1982anomaly,wangwenwitten2019new}.

\paragraph{Present status.} As of 12 November 2026: the stated qualified formulation is established within the cited and explicitly checked scope.

\paragraph{Audit.} Literal correctness: correct with missing hypotheses. Publication match: exact. Verification: read the result. Development: annotated. Confidence: high.

\noindent References: \citep{witten1982anomaly,wangwenwitten2019new}.\par

\subsubsection{1983-604.3 --- The Pfaffian line over the space of connections}

\label{item:1983-604.3}

\noindent\textbf{Kind:} construction. \textbf{Fate:} B. \textbf{Pages:} 1126.\par

\paragraph{Notebook content.} The choice of sign for the fermion determinant square root is central.

\paragraph{Precise statement.} A smooth family of suitable real skew-adjoint Fredholm operators defines a real Pfaffian line over its parameter space. A local Pfaffian is a section of that line; a consistent global sign amounts to an orientation/trivialization compatible with gauge identifications.

\paragraph{Editorial reconstruction.} The global anomaly is the nontrivial sign holonomy around a loop in the quotient by gauge transformations. A local numerical square root of a determinant does not provide this compatibility.

\paragraph{Literature and fate.} Later determinant/Pfaffian-line and eta-invariant frameworks provide a rigorous formulation of the notebook's sign issue. They clarify this specific operator-family problem without claiming a construction of every quantum field theory. Locator and verification scope: Wang--Wen--Witten section 2.2; Dai--Freed line-valued formulation. \cite{wangwenwitten2019new,daifreed1994eta}.

\paragraph{Present status.} As of 12 November 2026: the stated qualified formulation is established within the cited and explicitly checked scope.

\paragraph{Audit.} Literal correctness: correct with missing hypotheses. Publication match: exact. Verification: read the result. Development: annotated. Confidence: high.

\noindent References: \citep{wangwenwitten2019new,daifreed1994eta}.\par

\subsubsection{1983-604.4 --- Pauli matrices and the missing factor two}

\label{item:1983-604.4}

\noindent\textbf{Kind:} formula. \textbf{Fate:} E. \textbf{Pages:} 1126.\par

\paragraph{Notebook content.} The matrix commutator is written in a Pauli-matrix calculation.

\paragraph{Precise statement.} For the standard Pauli matrices, \(\sigma_i\sigma_j=\delta_{ij}I+i\sum_k\epsilon_{ijk}\sigma_k\), hence \([\sigma_i,\sigma_j]=2i\sum_k\epsilon_{ijk}\sigma_k\). If generators are \(\sigma_i/2\), the structure constants change accordingly.

\paragraph{Editorial reconstruction.} A factor \(i\) without the 2 is inconsistent with the unscaled Pauli convention. The correction is checked by multiplying \(\sigma_1,\sigma_2\); it is a convention issue only if all displayed generators are rescaled consistently.

\paragraph{Literature and fate.} The corrected formulation belongs to established mathematics available before the talk. Locator and verification scope: Direct2-by-2 matrix multiplication.

\paragraph{Present status.} As of 12 November 2026: the stated qualified formulation is established within the cited and explicitly checked scope.

\paragraph{Audit.} Literal correctness: false as written. Publication match: none found. Verification: indirect. Development: annotated. Confidence: high.

\subsubsection{1983-604.5 --- The nontrivial gauge component, repeated}

\label{item:1983-604.5}

\noindent\textbf{Kind:} topological claim. \textbf{Fate:} E. \textbf{Pages:} 1126.\par

\paragraph{Notebook content.} The nontrivial class of a map \(S^4\to SU(2)\) reappears.

\paragraph{Precise statement.} For the trivial \(SU(2)\)-bundle over \(S^4\), \(\pi_0(\mathcal G)\simeq\pi_4(SU(2))=\mathbb Z/2\). A representative \(g\) of the nontrivial component is used to close the path of connections after gauge identification.

\paragraph{Editorial reconstruction.} This is a repeated occurrence of 1983-603.4, retained separately. It is not a formula for the gauge group of every nontrivial bundle over every4-manifold.

\paragraph{Literature and fate.} The corrected formulation belongs to established mathematics available before the talk. Locator and verification scope: Hatcher section 4.J; Wang--Wen--Witten section 2.2. \cite{hatcher2002algebraic,wangwenwitten2019new}.

\paragraph{Present status.} As of 12 November 2026: the stated qualified formulation is established within the cited and explicitly checked scope.

\paragraph{Audit.} Literal correctness: correct with missing hypotheses. Publication match: exact. Verification: read the result. Development: annotated. Confidence: high.

\noindent References: \citep{hatcher2002algebraic,wangwenwitten2019new}.\par

\subsubsection{1983-604.6 --- Mapping-torus clutching of the bundle}

\label{item:1983-604.6}

\noindent\textbf{Kind:} construction. \textbf{Fate:} E. \textbf{Pages:} 1126--1127.\par

\paragraph{Notebook content.} The path from a connection to its gauge transform is closed in one extra dimension.

\paragraph{Precise statement.} For a path \(A_t\) from \(A_0\) to \(A_0^g\) on a bundle over a closed spin4-manifold \(M\), glue the bundles over \(M\times\{0\}\) and \(M\times\{1\}\) using \(g\). The result is a gauge bundle over \(M\times S^1\) with a compatible connection; a spin structure around the circle must also be selected.

\paragraph{Editorial reconstruction.} For a pure gauge transformation the base is the product mapping torus, while the bundle can be nontrivially clutched. If a diffeomorphism is also used, the base itself can be a nontrivial mapping torus.

\paragraph{Literature and fate.} The corrected formulation belongs to established mathematics available before the talk. Locator and verification scope: Section 2.2, mapping-torus construction. \cite{wangwenwitten2019new}.

\paragraph{Present status.} As of 12 November 2026: the stated qualified formulation is established within the cited and explicitly checked scope.

\paragraph{Audit.} Literal correctness: correct with missing hypotheses. Publication match: exact. Verification: read the result. Development: annotated. Confidence: high.

\noindent References: \citep{wangwenwitten2019new}.\par

\subsubsection{1983-604.7 --- Parity of paired eigenvalue crossings}

\label{item:1983-604.7}

\noindent\textbf{Kind:} claim. \textbf{Fate:} E. \textbf{Pages:} 1127.\par

\paragraph{Notebook content.} An odd number of crossings gives the anomalous sign.

\paragraph{Precise statement.} For the real structure of the \(SU(2)\) Weyl doublet, transporting the Pfaffian along the path changes its sign by \((-1)^\nu\), where \(\nu\) is the parity of the appropriate eigenvalue-pair crossings through zero. This equals the mod2 index on the clutched five-manifold.

\paragraph{Editorial reconstruction.} Counting every eigenvalue of an ungraded full Dirac operator can double the count and erase the anomaly. The paired spectrum and real structure must be specified before reducing modulo two.

\paragraph{Literature and fate.} The corrected formulation belongs to established mathematics available before the talk. Locator and verification scope: Wang--Wen--Witten section 2.2; Atiyah--Singer section 2. \cite{wangwenwitten2019new,atiyahsinger1971index5}.

\paragraph{Present status.} As of 12 November 2026: the stated qualified formulation is established within the cited and explicitly checked scope.

\paragraph{Audit.} Literal correctness: correct with missing hypotheses. Publication match: exact. Verification: read the result. Development: annotated. Confidence: high.

\noindent References: \citep{wangwenwitten2019new,atiyahsinger1971index5}.\par

\subsubsection{1983-604.8 --- The local crossing model and the adiabatic argument}

\label{item:1983-604.8}

\noindent\textbf{Kind:} method. \textbf{Fate:} B. \textbf{Pages:} 1127.\par

\paragraph{Notebook content.} The notebook relates a crossing to a kernel of an operator with one extra derivative.

\paragraph{Precise statement.} Near a simple crossing, the one-dimensional model \((\partial_s+as)u=0\) has the solution \(u(s)=C\exp(-as^2/2)\), which is square-integrable on \(\mathbb R\) for \(a>0\). This explains locally how an eigenvalue crossing contributes to a suspended operator.

\paragraph{Editorial reconstruction.} The local Gaussian is a heuristic building block, not by itself a proof for a compact mapping torus with bundle gluing. Later adiabatic eta-invariant and determinant-line holonomy theorems supply precise global statements, with their spin and sign conventions.

\paragraph{Literature and fate.} The scalar crossing argument is elementary; the later1994 holonomy theorem gives a rigorous global framework. The notebook's heuristic is not asserted to be a complete proof of that theorem. Locator and verification scope: Dai--Freed Corollary5.11 and equation 5.12; mapping-torus review in section 2.2 of Wang--Wen--Witten. \cite{daifreed1994eta,wangwenwitten2019new}.

\paragraph{Present status.} As of 12 November 2026: the stated qualified formulation is established within the cited and explicitly checked scope.

\paragraph{Audit.} Literal correctness: correct with missing hypotheses. Publication match: partial. Verification: read the result. Development: annotated. Confidence: high.

\noindent References: \citep{daifreed1994eta,wangwenwitten2019new}.\par

\subsubsection{1983-604.9 --- Real skew-adjointness is essential for the mod-two index}

\label{item:1983-604.9}

\noindent\textbf{Kind:} claim. \textbf{Fate:} E. \textbf{Pages:} 1127.\par

\paragraph{Notebook content.} Skew-Hermitian symmetry is used to justify stability of kernel dimension modulo two.

\paragraph{Precise statement.} The parity \(\dim_{\mathbb R}\ker D\bmod2\) is stable in continuous families of real skew-adjoint Fredholm operators. Complex skew-Hermiticity alone does not suffice: on a one-dimensional complex space, \(D_t=it\) is skew-Hermitian but its kernel has complex dimension1 at \(t=0\) and 0 otherwise.

\paragraph{Editorial reconstruction.} The real symmetry forces nonzero finite-dimensional spectral blocks to have even real dimension. The \(SU(2)\) anomaly supplies the needed real structure through the spin and gauge representations, not through the word skew-Hermitian alone.

\paragraph{Literature and fate.} The corrected formulation belongs to established mathematics available before the talk. Locator and verification scope: Atiyah--Singer section 2; Wang--Wen--Witten section 2.2. \cite{atiyahsinger1971index5,wangwenwitten2019new}.

\paragraph{Present status.} As of 12 November 2026: the stated qualified formulation is established within the cited and explicitly checked scope.

\paragraph{Audit.} Literal correctness: false as written. Publication match: exact. Verification: read the result. Development: annotated. Confidence: high.

\noindent References: \citep{atiyahsinger1971index5,wangwenwitten2019new}.\par

\subsubsection{1983-604.10 --- Gelfand's request for a mathematical framework}

\label{item:1983-604.10}

\noindent\textbf{Kind:} question. \textbf{Fate:} B. \textbf{Pages:} 1127, closing question.\par

\paragraph{Notebook content.} Gelfand asks for a mathematical setting that makes the anomaly reasoning precise.

\paragraph{Precise statement.} A precise specialization is to construct determinant/Pfaffian lines for elliptic families, define their transport under gauge identifications, and compute the holonomy around a mapping-torus loop from spectral data.

\paragraph{Editorial reconstruction.} This specialization is editorial and narrower than the bare request for a general framework. Dai--Freed's exponentiated eta invariants, gluing theorem and holonomy formula realize substantial parts of it. They do not by themselves construct the interacting gauge-field path integral or resolve every possible interpretation of the original question.

\paragraph{Literature and fate.} The1994 determinant-line/eta-invariant framework and the established real Pfaffian interpretation provide a rigorous partial answer to the specifically reconstructed operator-family question. The broad original program remains too undefined for a binary solved/open judgment. Locator and verification scope: Dai--Freed Proposition 2.15, Theorem 2.20, Corollary5.11; real mod2 discussion in Wang--Wen--Witten. \cite{daifreed1994eta,wangwenwitten2019new}.

\paragraph{Present status.} As of 12 November 2026: the stated determinant/Pfaffian-line specialization has an established rigorous framework; no claim is made that the unspecified general field-theory program is solved.

\paragraph{Audit.} Literal correctness: correct with missing hypotheses. Publication match: partial. Verification: read the result. Development: annotated. Confidence: high.

\noindent Canonical problem: GP-1983-03.\par

\noindent References: \citep{daifreed1994eta,wangwenwitten2019new}.\par

\subsection{Record 605: Symmetric functions, standard bases, and tableaux}

\label{talk:605}

\noindent\textbf{Date:} 1983-10-17. \textbf{Speaker:} A. V. Zelevinsky. \textbf{Pages:} 1128--1130.\par

The lecture introduces the graded ring of symmetric functions, its elementary, complete, power-sum and Schur bases, and the Kostka-tableau rule. The first two lines of P001130 finish this lecture before the 24 October heading. Strict positivity of every Kostka number and unrestricted integral generation by power sums require correction.

\subsubsection{1983-605.1 --- The graded ring of symmetric functions}

\label{item:1983-605.1}

\noindent\textbf{Kind:} definition. \textbf{Fate:} E. \textbf{Pages:} 1128.\par

\paragraph{Notebook content.} Notebook occurrence: P001128, opening definition. The graded ring of symmetric functions.

\paragraph{Precise statement.} For each degree \(d\), put \(\Lambda_N^d=\mathbb Z[x_1,\ldots,x_N]^{S_N}_d\). Specialization \(x_{N+1}=0\) induces an isomorphism \(\Lambda_{N+1}^d\to\Lambda_N^d\) for \(N\ge d\). Define \(\Lambda=\bigoplus_{d\ge0}\varprojlim_N\Lambda_N^d\).

\paragraph{Editorial reconstruction.} The degree bound is essential. The unrestricted inverse limit without a bounded-degree condition is a completion, larger than the polynomial ring of symmetric functions.

\paragraph{Literature and fate.} Classical symmetric-function mathematics, already available by the 1979 first edition; the opened 1995 edition verifies the modern formulation. Locator and verification scope: Chapter I, section 2, definition of the graded inverse limit. \cite{macdonald1995symmetric}.

\paragraph{Present status.} As of 12 November2026: the stated corrected formulation is established. Verification is limited to the identified result and the explicit elementary deductions; no broader literature-completeness claim is made.

\paragraph{Audit.} Literal correctness: correct with missing hypotheses. Publication match: exact. Verification: read the result. Development: annotated. Confidence: high.

\noindent References: \citep{macdonald1995symmetric}.\par

\subsubsection{1983-605.2 --- Partitions and the monomial basis}

\label{item:1983-605.2}

\noindent\textbf{Kind:} definition. \textbf{Fate:} E. \textbf{Pages:} 1128.\par

\paragraph{Notebook content.} Notebook occurrence: P001128, partitions and m_lambda. Partitions and the monomial basis.

\paragraph{Precise statement.} For a partition \(\lambda=(\lambda_1\ge\cdots\ge\lambda_r>0)\), the monomial symmetric function \(m_\lambda(x_1,\ldots,x_N)\) is the sum of distinct monomials whose exponent sequence is a permutation of \((\lambda_1,\ldots,\lambda_r,0,\ldots,0)\). For \(|\lambda|=d\), these elements form a \(\mathbb Z\)-basis of \(\Lambda^d\).

\paragraph{Editorial reconstruction.} Repeated parts must not lead to repeated summands; summing over every permutation without dividing by the stabilizer gives the wrong normalization.

\paragraph{Literature and fate.} Classical symmetric-function mathematics, already available by the 1979 first edition; the opened 1995 edition verifies the modern formulation. Locator and verification scope: Chapter I, sections 1-2. \cite{macdonald1995symmetric}.

\paragraph{Present status.} As of 12 November2026: the stated corrected formulation is established. Verification is limited to the identified result and the explicit elementary deductions; no broader literature-completeness claim is made.

\paragraph{Audit.} Literal correctness: correct with missing hypotheses. Publication match: exact. Verification: read the result. Development: annotated. Confidence: high.

\noindent References: \citep{macdonald1995symmetric}.\par

\subsubsection{1983-605.3 --- Elementary symmetric functions and characteristic coefficients}

\label{item:1983-605.3}

\noindent\textbf{Kind:} definition. \textbf{Fate:} E. \textbf{Pages:} 1128.\par

\paragraph{Notebook content.} Notebook occurrence: P001128, e_n and product formula. Elementary symmetric functions and characteristic coefficients.

\paragraph{Precise statement.} Set \(e_r=\sum_{i_1<\cdots<i_r}x_{i_1}\cdots x_{i_r}\), with \(e_0=1\). Then \(E(t)=\sum_{r\ge0}e_rt^r=\prod_i(1+tx_i)\). For an \(N\times N\) matrix with eigenvalues \(x_i\), \(\det(zI-A)=\sum_{r=0}^N(-1)^re_rz^{N-r}\).

\paragraph{Editorial reconstruction.} The signs depend on whether the characteristic polynomial is \(\det(zI-A)\) or \(\det(A-zI)\); the displayed convention fixes them. Infinite products are interpreted degree by degree.

\paragraph{Literature and fate.} Classical symmetric-function mathematics, already available by the 1979 first edition; the opened 1995 edition verifies the modern formulation. Locator and verification scope: Chapter I, section 2, elementary functions. \cite{macdonald1995symmetric}.

\paragraph{Present status.} As of 12 November2026: the stated corrected formulation is established. Verification is limited to the identified result and the explicit elementary deductions; no broader literature-completeness claim is made.

\paragraph{Audit.} Literal correctness: correct. Publication match: exact. Verification: read the result. Development: annotated. Confidence: high.

\noindent References: \citep{macdonald1995symmetric}.\par

\subsubsection{1983-605.4 --- Complete symmetric functions and their generating function}

\label{item:1983-605.4}

\noindent\textbf{Kind:} definition and exercise. \textbf{Fate:} E. \textbf{Pages:} 1128.\par

\paragraph{Notebook content.} Notebook occurrence: P001128, h_n and exercise. Complete symmetric functions and their generating function.

\paragraph{Precise statement.} Set \(h_r=\sum_{i_1\le\cdots\le i_r}x_{i_1}\cdots x_{i_r}=\sum_{\lambda\vdash r}m_\lambda\), \(h_0=1\). Expanding geometric series proves \(H(t)=\sum_{r\ge0}h_rt^r=\prod_i(1-tx_i)^{-1}\).

\paragraph{Editorial reconstruction.} This is a formal identity; it requires no analytic convergence assumption. The requested verification follows by selecting a multiplicity for each variable.

\paragraph{Literature and fate.} Classical symmetric-function mathematics, already available by the 1979 first edition; the opened 1995 edition verifies the modern formulation. Locator and verification scope: Chapter I, section 2, complete functions. \cite{macdonald1995symmetric}.

\paragraph{Present status.} As of 12 November2026: the stated corrected formulation is established. Verification is limited to the identified result and the explicit elementary deductions; no broader literature-completeness claim is made.

\paragraph{Audit.} Literal correctness: correct. Publication match: exact. Verification: read the result. Development: annotated. Confidence: high.

\noindent Canonical problem: GP-1983-06.\par

\noindent References: \citep{macdonald1995symmetric}.\par

\subsubsection{1983-605.5 --- Power sums and the logarithmic derivative}

\label{item:1983-605.5}

\noindent\textbf{Kind:} definition and identity. \textbf{Fate:} E. \textbf{Pages:} 1128.\par

\paragraph{Notebook content.} Notebook occurrence: P001128, p_n and P(t). Power sums and the logarithmic derivative.

\paragraph{Precise statement.} For \(r\ge1\), set \(p_r=\sum_i x_i^r\). Then \(P(t)=\sum_{r\ge1}p_rt^r=tH'(t)/H(t)=\sum_i tx_i/(1-tx_i)\).

\paragraph{Editorial reconstruction.} The constant coefficient of \(P\) is zero; \(p_0\) is not a stable symmetric function under adjoining a zero variable.

\paragraph{Literature and fate.} Classical symmetric-function mathematics, already available by the 1979 first edition; the opened 1995 edition verifies the modern formulation. Locator and verification scope: Chapter I, section 2, equations for power sums. \cite{macdonald1995symmetric}.

\paragraph{Present status.} As of 12 November2026: the stated corrected formulation is established. Verification is limited to the identified result and the explicit elementary deductions; no broader literature-completeness claim is made.

\paragraph{Audit.} Literal correctness: correct. Publication match: exact. Verification: read the result. Development: annotated. Confidence: high.

\noindent References: \citep{macdonald1995symmetric}.\par

\subsubsection{1983-605.6 --- Integral bases and the rational power-sum basis}

\label{item:1983-605.6}

\noindent\textbf{Kind:} theorem. \textbf{Fate:} E. \textbf{Pages:} 1128--1129.\par

\paragraph{Notebook content.} Notebook occurrence: P001128-P001129, basis assertion. Integral bases and the rational power-sum basis.

\paragraph{Precise statement.} \(\Lambda=\mathbb Z[e_1,e_2,\ldots]=\mathbb Z[h_1,h_2,\ldots]\), with algebraically independent generators. Thus \(e_\lambda=\prod_j e_{\lambda_j}\) and \(h_\lambda=\prod_j h_{\lambda_j}\) are integral bases. The \(p_\lambda\) form a basis only after tensoring with \(\mathbb Q\): \(\Lambda_{\mathbb Q}=\mathbb Q[p_1,p_2,\ldots]\).

\paragraph{Editorial reconstruction.} For example \(e_2=(p_1^2-p_2)/2\), so the power sums do not freely generate \(\Lambda\) over \(\mathbb Z\).

\paragraph{Literature and fate.} Classical symmetric-function mathematics, already available by the 1979 first edition; the opened 1995 edition verifies the modern formulation. Locator and verification scope: Chapter I, section 2. \cite{macdonald1995symmetric}.

\paragraph{Present status.} As of 12 November2026: the stated corrected formulation is established. Verification is limited to the identified result and the explicit elementary deductions; no broader literature-completeness claim is made.

\paragraph{Audit.} Literal correctness: correct with missing hypotheses. Publication match: exact. Verification: read the result. Development: annotated. Confidence: high.

\noindent References: \citep{macdonald1995symmetric}.\par

\subsubsection{1983-605.7 --- Alternants and the Vandermonde denominator}

\label{item:1983-605.7}

\noindent\textbf{Kind:} construction. \textbf{Fate:} E. \textbf{Pages:} 1129.\par

\paragraph{Notebook content.} Notebook occurrence: P001129, a_alpha. Alternants and the Vandermonde denominator.

\paragraph{Precise statement.} For distinct nonnegative integers \(\alpha_1,\ldots,\alpha_N\), put \(a_\alpha(x)=\det(x_i^{\alpha_j})_{1\le i,j\le N}\). With \(\delta=(N-1,N-2,\ldots,0)\), \(a_\delta=\prod_{i<j}(x_i-x_j)\); every alternating polynomial is divisible by this Vandermonde polynomial, and the quotient is symmetric.

\paragraph{Editorial reconstruction.} Distinctness prevents the determinant from vanishing trivially; nonnegative exponents ensure polynomial rather than Laurent-polynomial quotients. The determinant convention fixes the Vandermonde sign.

\paragraph{Literature and fate.} Classical symmetric-function mathematics, already available by the 1979 first edition; the opened 1995 edition verifies the modern formulation. Locator and verification scope: Chapter I, section 3, alternant definition. \cite{macdonald1995symmetric}.

\paragraph{Present status.} As of 12 November2026: the stated corrected formulation is established. Verification is limited to the identified result and the explicit elementary deductions; no broader literature-completeness claim is made.

\paragraph{Audit.} Literal correctness: correct. Publication match: exact. Verification: read the result. Development: annotated. Confidence: high.

\noindent References: \citep{macdonald1995symmetric}.\par

\subsubsection{1983-605.8 --- Schur functions, stability, and the integral Schur basis}

\label{item:1983-605.8}

\noindent\textbf{Kind:} definition and theorem. \textbf{Fate:} E. \textbf{Pages:} 1129.\par

\paragraph{Notebook content.} Notebook occurrence: P001129, s_lambda. Schur functions, stability, and the integral Schur basis.

\paragraph{Precise statement.} For \(\ell(\lambda)\le N\), define \(s_\lambda(x_1,\ldots,x_N)=a_{\lambda+\delta}/a_\delta\). These are symmetric polynomials satisfying \(s_\lambda(x_1,\ldots,x_N,0)=s_\lambda(x_1,\ldots,x_N)\). Their stable values form a \(\mathbb Z\)-basis of \(\Lambda\).

\paragraph{Editorial reconstruction.} Stability means equality after specialization of the additional variable to zero, not equality as polynomials in different sets of independent variables.

\paragraph{Literature and fate.} Classical symmetric-function mathematics, already available by the 1979 first edition; the opened 1995 edition verifies the modern formulation. Locator and verification scope: Chapter I, section 3. \cite{macdonald1995symmetric}.

\paragraph{Present status.} As of 12 November2026: the stated corrected formulation is established. Verification is limited to the identified result and the explicit elementary deductions; no broader literature-completeness claim is made.

\paragraph{Audit.} Literal correctness: correct. Publication match: exact. Verification: read the result. Development: annotated. Confidence: high.

\noindent References: \citep{macdonald1995symmetric}.\par

\subsubsection{1983-605.9 --- The transition-matrix and multiplication questions}

\label{item:1983-605.9}

\noindent\textbf{Kind:} question. \textbf{Fate:} E. \textbf{Pages:} 1129.\par

\paragraph{Notebook content.} Notebook occurrence: P001129, purpose of the theory. The transition-matrix and multiplication questions.

\paragraph{Precise statement.} For the standard bases \(m_\lambda,e_\lambda,h_\lambda,p_\lambda,s_\lambda\), determine change-of-basis coefficients and multiplication constants. In particular, determine the coefficients \(c_{\lambda\mu}^{\nu}\) in \(s_\lambda s_\mu=\sum_\nu c_{\lambda\mu}^{\nu}s_\nu\).

\paragraph{Editorial reconstruction.} This is a programmatic teaching question about classical bases. The Littlewood--Richardson rule gives the Schur multiplication constants; it is a literature completion of the question, not a formula written in this notebook occurrence.

\paragraph{Literature and fate.} Classical transition rules and the Littlewood--Richardson rule were available before 1983. The1995 edition is a verified modern reference, not the date of resolution. Locator and verification scope: Chapter I, sections 6 and 9. \cite{macdonald1995symmetric}.

\paragraph{Present status.} As of 12 November2026: the stated corrected formulation is established. Verification is limited to the identified result and the explicit elementary deductions; no broader literature-completeness claim is made.

\paragraph{Audit.} Literal correctness: correct. Publication match: exact. Verification: read the result. Development: annotated. Confidence: high.

\noindent Canonical problem: GP-1983-04.\par

\noindent References: \citep{macdonald1995symmetric}.\par

\subsubsection{1983-605.10 --- Kostka coefficients are nonnegative, not everywhere positive}

\label{item:1983-605.10}

\noindent\textbf{Kind:} claim. \textbf{Fate:} E. \textbf{Pages:} 1129.\par

\paragraph{Notebook content.} Notebook occurrence: P001129, s_lambda expansion. Kostka coefficients are nonnegative, not everywhere positive.

\paragraph{Precise statement.} The expansion \(s_\lambda=\sum_{\mu\vdash|\lambda|}K_{\lambda\mu}m_\mu\) has \(K_{\lambda\mu}\in\mathbb Z_{\ge0}\), \(K_{\lambda\lambda}=1\), and \(K_{\lambda\mu}=0\) unless \(\lambda\) dominates \(\mu\).

\paragraph{Editorial reconstruction.} A literal strict positivity assertion is false: \(s_{(1,1)}=m_{(1,1)}\), so \(K_{(1,1),(2)}=0\). The nonnegative assertion is the correct theorem.

\paragraph{Literature and fate.} Classical symmetric-function mathematics, already available by the 1979 first edition; the opened 1995 edition verifies the modern formulation. Locator and verification scope: Chapter I, section 6, equations(6.4)-(6.5). \cite{macdonald1995symmetric}.

\paragraph{Present status.} As of 12 November2026: the stated corrected formulation is established. Verification is limited to the identified result and the explicit elementary deductions; no broader literature-completeness claim is made.

\paragraph{Audit.} Literal correctness: false as written. Publication match: exact. Verification: read the result. Development: annotated. Confidence: high.

\noindent References: \citep{macdonald1995symmetric}.\par

\subsubsection{1983-605.11 --- The semistandard-tableau rule}

\label{item:1983-605.11}

\noindent\textbf{Kind:} theorem. \textbf{Fate:} E. \textbf{Pages:} 1129.\par

\paragraph{Notebook content.} Notebook occurrence: P001129, filled Young diagram. The semistandard-tableau rule.

\paragraph{Precise statement.} For partitions \(\lambda,\mu\) of equal size, \(K_{\lambda\mu}\) equals the number of semistandard Young tableaux of shape \(\lambda\) and content \(\mu\): entries weakly increase along each row and strictly increase down each column in English diagram convention.

\paragraph{Editorial reconstruction.} An upward column arrow in a diagram using the opposite orientation encodes the same condition after changing convention. The row/column convention must be fixed before interpreting the sketch.

\paragraph{Literature and fate.} Classical symmetric-function mathematics, already available by the 1979 first edition; the opened 1995 edition verifies the modern formulation. Locator and verification scope: Chapter I, sections 5-6, tableau formula. \cite{macdonald1995symmetric}.

\paragraph{Present status.} As of 12 November2026: the stated corrected formulation is established. Verification is limited to the identified result and the explicit elementary deductions; no broader literature-completeness claim is made.

\paragraph{Audit.} Literal correctness: correct. Publication match: exact. Verification: read the result. Development: annotated. Confidence: high.

\noindent References: \citep{macdonald1995symmetric}.\par

\subsubsection{1983-605.12 --- Exercise on the other transition matrices}

\label{item:1983-605.12}

\noindent\textbf{Kind:} exercise. \textbf{Fate:} E. \textbf{Pages:} 1130.\par

\paragraph{Notebook content.} Notebook occurrence: P001130, first two lines before 24 October heading. Exercise on the other transition matrices.

\paragraph{Precise statement.} Compute the remaining transition matrices between the standard symmetric-function bases in each fixed degree. These are finite matrices indexed by partitions of that degree; inverse matrices and the involution interchanging \(e_r,h_r\) recover the remaining changes of basis.

\paragraph{Editorial reconstruction.} This occurrence is distinct from the general question on P001129. It is the end of the 17 October lecture, before the 24 October date line. Concrete matrix-counting solutions recur in 1983-608.

\paragraph{Literature and fate.} Classical symmetric-function mathematics, already available by the 1979 first edition; the opened 1995 edition verifies the modern formulation. Locator and verification scope: Chapter I, section 6, transition-matrix diagram. \cite{macdonald1995symmetric}.

\paragraph{Present status.} As of 12 November2026: the stated corrected formulation is established. Verification is limited to the identified result and the explicit elementary deductions; no broader literature-completeness claim is made.

\paragraph{Audit.} Literal correctness: correct. Publication match: exact. Verification: read the result. Development: annotated. Confidence: high.

\noindent Canonical problem: GP-1983-04.\par

\noindent References: \citep{macdonald1995symmetric}.\par

\subsubsection{1983-605.13 --- The forgotten symmetric-function basis}

\label{item:1983-605.13}

\noindent\textbf{Kind:} remark and definition. \textbf{Fate:} E. \textbf{Pages:} 1130.\par

\paragraph{Notebook content.} Notebook occurrence: P001130, note on forgotten functions. The forgotten symmetric-function basis.

\paragraph{Precise statement.} The forgotten functions are \(f_\lambda=\omega(m_\lambda)\), where the involution \(\omega\) is defined by \(\omega(e_r)=h_r\). With the Hall inner product satisfying \(\langle h_\lambda,m_\mu\rangle=\delta_{\lambda\mu}\), one has \(\langle e_\lambda,f_\mu\rangle=\delta_{\lambda\mu}\).

\paragraph{Editorial reconstruction.} The notebook names the forgotten basis; its precise definition and duality are supplied editorially from the cited reference.

\paragraph{Literature and fate.} Classical symmetric-function mathematics, already available by the 1979 first edition; the opened 1995 edition verifies the modern formulation. Locator and verification scope: Chapter I, section 2, definition of f_lambda. \cite{macdonald1995symmetric}.

\paragraph{Present status.} As of 12 November2026: the stated corrected formulation is established. Verification is limited to the identified result and the explicit elementary deductions; no broader literature-completeness claim is made.

\paragraph{Audit.} Literal correctness: correct. Publication match: exact. Verification: read the result. Development: annotated. Confidence: high.

\noindent References: \citep{macdonald1995symmetric}.\par

\subsection{Record 606: Formal wave functions, Cherednik's identity, and tau-functions}

\label{talk:606}

\noindent\textbf{Date:} 1983-10-24. \textbf{Speaker:} I. M. Gelfand. \textbf{Pages:} 1130--1131.\par

The24 October heading on P001130 begins a formal differential-operator calculation continuing through equation 17 on P001131. The logarithmic derivative, fractional powers and tau quotient have identifiable classical sources. The Lax sign, quotient derivative, two different potentials called tau, and omitted spectral powers require explicit normalization. The theorem number 3.1 does not match the located 1978 Cherednik paper, where the identity is Theorem 1.4.

\subsubsection{1983-606.1 --- Normalized scalar differential operator and formal eigenfunction}

\label{item:1983-606.1}

\noindent\textbf{Kind:} construction. \textbf{Fate:} E. \textbf{Pages:} 1130.\par

\paragraph{Notebook content.} Notebook occurrence: P001130, L_n and Psi. Normalized scalar differential operator and formal eigenfunction.

\paragraph{Precise statement.} Let \(R\) be a commutative differential \(\mathbb Q\)-algebra with derivation \(\partial=\partial_{x_1}\), and \(L_n=\partial^n+\sum_{j=0}^{n-2}u_j\partial^j\), \(n\ge2\). After adjoining the necessary antiderivatives, solve \(L_n\Psi=k^n\Psi\) with \(\Psi=e^{kx_1}(1+\sum_{r\ge1}c_rk^{-r})\).

\paragraph{Editorial reconstruction.} The coefficient of \(\partial^{n-1}\) must vanish for the displayed leading factor. Otherwise it is absorbed in a nonconstant leading amplitude or removed by a gauge transformation. The coefficients \(c_r\) need not belong to the original differential-polynomial ring. Related earlier or repeated occurrences: \texttt{1975-379.3}.

\paragraph{Literature and fate.} The formal-operator and wave-function framework was published or circulating before the talk. The cited result locates the corrected identity; a later exposition is not treated as its discovery date. Locator and verification scope: Section 1, Proposition 1.6 and Remark1.8. \cite{cherednik1978differential}.

\paragraph{Present status.} As of 12 November2026: the stated corrected formulation is established. Verification is limited to the identified result and the explicit elementary deductions; no broader literature-completeness claim is made.

\paragraph{Audit.} Literal correctness: correct. Publication match: exact. Verification: read the result. Development: annotated. Confidence: high.

\noindent References: \citep{cherednik1978differential}.\par

\subsubsection{1983-606.2 --- Recursive computation of the wave coefficients}

\label{item:1983-606.2}

\noindent\textbf{Kind:} exercise and method. \textbf{Fate:} E. \textbf{Pages:} 1130.\par

\paragraph{Notebook content.} Notebook occurrence: P001130, exercise following Psi expansion. Recursive computation of the wave coefficients.

\paragraph{Precise statement.} Substitution into \(L_n\Psi=k^n\Psi\) determines the derivatives of successive \(c_r\), hence the coefficients after choices of integration constants. For \(L_2=\partial^2+u\), \(2c_1'=-u\) and \(2c_{r+1}'=-c_r''-uc_r\), with \(c_0=1\).

\paragraph{Editorial reconstruction.} Multiplication of \(\Psi\) by a formal series \(1+\sum a_rk^{-r}\) with \(\partial a_r=0\) preserves the eigenvalue equation. Thus uniqueness of \(\Psi\) requires normalization of these constants; the logarithmic derivative is independent of this freedom. Related earlier or repeated occurrences: \texttt{1975-379.3}.

\paragraph{Literature and fate.} The formal-operator and wave-function framework was published or circulating before the talk. The cited result locates the corrected identity; a later exposition is not treated as its discovery date. Locator and verification scope: Section 1, Proposition 1.6 and Remark1.8. \cite{cherednik1978differential}.

\paragraph{Present status.} As of 12 November2026: the stated corrected formulation is established. Verification is limited to the identified result and the explicit elementary deductions; no broader literature-completeness claim is made.

\paragraph{Audit.} Literal correctness: correct. Publication match: exact. Verification: read the result. Development: annotated. Confidence: high.

\noindent Canonical problem: GP-1983-05.\par

\noindent References: \citep{cherednik1978differential}.\par

\subsubsection{1983-606.3 --- The logarithmic derivative is a local differential polynomial}

\label{item:1983-606.3}

\noindent\textbf{Kind:} definition and theorem. \textbf{Fate:} E. \textbf{Pages:} 1130.\par

\paragraph{Notebook content.} Notebook occurrence: P001130, chi expansion. The logarithmic derivative is a local differential polynomial.

\paragraph{Precise statement.} Set \(\chi=\partial\Psi/\Psi=k+\sum_{r\ge1}\chi_rk^{-r}\). Its coefficients are differential polynomials in the \(u_j\) with rational coefficients, even though the \(c_r\) can require integration. For \(L_2=\partial^2+u\), \(\chi_1=-u/2\), \(\chi_2=u'/4\).

\paragraph{Editorial reconstruction.} The Riccati equation \(\chi'+\chi^2+u=k^2\) verifies the two coefficients and fixes the sign convention. The derivative is with respect to \(x_1\), not an unspecified indexed coordinate.

\paragraph{Literature and fate.} The formal-operator and wave-function framework was published or circulating before the talk. The cited result locates the corrected identity; a later exposition is not treated as its discovery date. Locator and verification scope: Section 1, Proposition 1.6. \cite{cherednik1978differential}.

\paragraph{Present status.} As of 12 November2026: the stated corrected formulation is established. Verification is limited to the identified result and the explicit elementary deductions; no broader literature-completeness claim is made.

\paragraph{Audit.} Literal correctness: correct. Publication match: exact. Verification: read the result. Development: annotated. Confidence: high.

\noindent References: \citep{cherednik1978differential}.\par

\subsubsection{1983-606.4 --- Differential parts of fractional powers}

\label{item:1983-606.4}

\noindent\textbf{Kind:} construction. \textbf{Fate:} E. \textbf{Pages:} 1130.\par

\paragraph{Notebook content.} Notebook occurrence: P001130, the operators L_m. Differential parts of fractional powers.

\paragraph{Precise statement.} There is a unique monic formal pseudodifferential root \(L_n^{1/n}=\partial+\sum_{r\ge1}a_r\partial^{-r}\). For integers \(m\ge0\), define \(L_m=(L_n^{m/n})_+\), where the plus sign keeps nonnegative powers of \(\partial\). Then \(L_m\Psi=(k^m+O(k^{-1}))\Psi\).

\paragraph{Editorial reconstruction.} The construction gives differential operators for \(m\ge0\). A statement that their differential order equals an arbitrary negative integer is not meaningful; negative powers belong to the pseudodifferential algebra and have zero differential part. Related earlier or repeated occurrences: \texttt{1973-327.1}, \texttt{1982-592.8}. The reused problem identifier refers only to the (n=2), odd-flow KdV specialization; the general-(n) construction is broader.

\paragraph{Literature and fate.} The formal-operator and wave-function framework was published or circulating before the talk. The cited result locates the corrected identity; a later exposition is not treated as its discovery date. Locator and verification scope: Cherednik, Proposition 1.3; Gelfand--Dikii, publication metadata. \cite{gelfanddikii1976fractional,cherednik1978differential}.

\paragraph{Present status.} As of 12 November2026: the stated corrected formulation is established. Verification is limited to the identified result and the explicit elementary deductions; no broader literature-completeness claim is made.

\paragraph{Audit.} Literal correctness: correct. Publication match: exact. Verification: read the result. Development: annotated. Confidence: high.

\noindent Canonical problem: GP-1973-53.\par

\noindent References: \citep{gelfanddikii1976fractional,cherednik1978differential}.\par

\subsubsection{1983-606.5 --- Compatibility fixes the Lax commutator sign}

\label{item:1983-606.5}

\noindent\textbf{Kind:} claim. \textbf{Fate:} E. \textbf{Pages:} 1130.\par

\paragraph{Notebook content.} Notebook occurrence: P001130, Lax equation. Compatibility fixes the Lax commutator sign.

\paragraph{Precise statement.} With the auxiliary equations \(\partial_{x_m}\Psi=L_m\Psi\), the eigenvalue equation implies \(\partial_{x_m}L_n=[L_m,L_n]\). Equivalently, \(\partial_{x_m}L_n+[L_n,L_m]=0\).

\paragraph{Editorial reconstruction.} Differentiating \(L_n\Psi=k^n\Psi\) gives \((\partial_{x_m}L_n+[L_n,L_m])\Psi=0\). The notebook's opposite commutator sign is incompatible with the stated auxiliary equation unless every corresponding time is reversed. Related earlier or repeated occurrences: \texttt{1982-592.8}.

\paragraph{Literature and fate.} The formal-operator and wave-function framework was published or circulating before the talk. The cited result locates the corrected identity; a later exposition is not treated as its discovery date. Locator and verification scope: Cherednik section 1; Segal--Wilson section 4. \cite{cherednik1978differential,segalwilson1985loop}.

\paragraph{Present status.} As of 12 November2026: the stated corrected formulation is established. Verification is limited to the identified result and the explicit elementary deductions; no broader literature-completeness claim is made.

\paragraph{Audit.} Literal correctness: false as written. Publication match: exact. Verification: read the result. Development: annotated. Confidence: high.

\noindent References: \citep{cherednik1978differential,segalwilson1985loop}.\par

\subsubsection{1983-606.6 --- The formal resolvent and its scalar remainder}

\label{item:1983-606.6}

\noindent\textbf{Kind:} definition. \textbf{Fate:} E. \textbf{Pages:} 1130--1131.\par

\paragraph{Notebook content.} Notebook occurrence: P001130-P001131, Lambda and epsilon. The formal resolvent and its scalar remainder.

\paragraph{Precise statement.} Put \(\Lambda(k)=\sum_{m\ge0}L_mk^{-m-1}\). There is a unique scalar series \(\varepsilon(k)=-1+\sum_{j\ge1}\varepsilon_jk^{-j-1}\) such that \((\partial-\chi(k))\Lambda(k)=\varepsilon(k)\). The coefficients \(\varepsilon_j\) are differential polynomials in those of \(L_n\).

\paragraph{Editorial reconstruction.} All products are in the algebra of formal series with differential-operator coefficients. The leading term \(-1\) follows from multiplying \(-k\) by \(k^{-1}L_0\), where \(L_0=1\).

\paragraph{Literature and fate.} The formal-operator and wave-function framework was published or circulating before the talk. The cited result locates the corrected identity; a later exposition is not treated as its discovery date. Locator and verification scope: Section 1, Proposition 1.3 and Theorem 1.4. \cite{cherednik1978differential}.

\paragraph{Present status.} As of 12 November2026: the stated corrected formulation is established. Verification is limited to the identified result and the explicit elementary deductions; no broader literature-completeness claim is made.

\paragraph{Audit.} Literal correctness: correct. Publication match: exact. Verification: read the result. Development: annotated. Confidence: high.

\noindent References: \citep{cherednik1978differential}.\par

\subsubsection{1983-606.7 --- Cherednik's operator identity and the unmatched theorem number}

\label{item:1983-606.7}

\noindent\textbf{Kind:} theorem. \textbf{Fate:} E. \textbf{Pages:} 1131.\par

\paragraph{Notebook content.} Notebook occurrence: P001131, Theorem 3.1 in the notes. Cherednik's operator identity and the incorrect theorem number.

\paragraph{Precise statement.} In the preceding formal setup, \((\partial-\chi)\Lambda=\varepsilon\). This is Theorem 1.4 in Cherednik's1978 paper, following Proposition 1.3.

\paragraph{Editorial reconstruction.} The scan explicitly reads \ru{Теор. Чередника3.1}. This number does not match the located 1978 paper, in which the same identity is Theorem 1.4 and Definition 3.1 concerns a nonlinear Schrodinger reduction. Another source or numbering convention is not excluded; the notebook supplies no title sufficient to settle that bibliographic question.

\paragraph{Literature and fate.} The formal-operator and wave-function framework was published or circulating before the talk. The cited result locates the corrected identity; a later exposition is not treated as its discovery date. Locator and verification scope: Theorem 1.4, Russian p. 47; contrast Definition 3.1. \cite{cherednik1978differential}.

\paragraph{Present status.} As of 12 November2026: the stated corrected formulation is established. Verification is limited to the identified result and the explicit elementary deductions; no broader literature-completeness claim is made.

\paragraph{Audit.} Literal correctness: correct with missing hypotheses. Publication match: partial. Verification: read the result. Development: annotated. Confidence: high.

\noindent References: \citep{cherednik1978differential}.\par

\subsubsection{1983-606.8 --- The derivative in equation 12 must act on the whole quotient}

\label{item:1983-606.8}

\noindent\textbf{Kind:} claim. \textbf{Fate:} E. \textbf{Pages:} 1131.\par

\paragraph{Notebook content.} Notebook occurrence: P001131, equation 12. The derivative in equation 12 must act on the whole quotient.

\paragraph{Precise statement.} The well-defined formal identity is \(\sum_{m\ge0}k^{-m-1}\partial_{x_1}\bigl((L_m\Psi)/\Psi\bigr)=\partial_k\chi+\varepsilon\). Each differentiated quotient has no nonnegative powers of \(k\), so each coefficient receives only finitely many contributions.

\paragraph{Editorial reconstruction.} The scan places the derivative on \(L_m\Psi\) and multiplies by \(\Psi^{-1}\) afterwards. That leaves infinitely many identical leading contributions and omits the quotient-rule term. Moving \(\Psi^{-1}\) inside the derivative repairs the formula consistently with equations14--15.

\paragraph{Literature and fate.} The formal-operator and wave-function framework was published or circulating before the talk. The cited result locates the corrected identity; a later exposition is not treated as its discovery date. Locator and verification scope: Theorem 1.4 supplies the underlying resolvent identity; the displayed repair is checked coefficientwise. \cite{cherednik1978differential}.

\paragraph{Present status.} As of 12 November2026: the stated corrected formulation is established. Verification is limited to the identified result and the explicit elementary deductions; no broader literature-completeness claim is made.

\paragraph{Audit.} Literal correctness: false as written. Publication match: exact. Verification: read the result. Development: annotated. Confidence: high.

\noindent References: \citep{cherednik1978differential}.\par

\subsubsection{1983-606.9 --- Densities epsilon\_j and j chi\_j differ by a total derivative}

\label{item:1983-606.9}

\noindent\textbf{Kind:} claim. \textbf{Fate:} E. \textbf{Pages:} 1131.\par

\paragraph{Notebook content.} Notebook occurrence: P001131, equation 13. Densities epsilon_j and j chi_j differ by a total derivative.

\paragraph{Precise statement.} For each \(j\ge1\), \(\varepsilon_j-j\chi_j\) is a total \(x_1\)-derivative of a differential polynomial. In particular \(\varepsilon_1=\chi_1\).

\paragraph{Editorial reconstruction.} The assertion concerns equivalence of local densities modulo total derivatives. It does not say the two differential polynomials are literally identical for all \(j\).

\paragraph{Literature and fate.} The formal-operator and wave-function framework was published or circulating before the talk. The cited result locates the corrected identity; a later exposition is not treated as its discovery date. Locator and verification scope: Section 1 formal resolvent identity and its coefficient consequences. \cite{cherednik1978differential}.

\paragraph{Present status.} As of 12 November2026: the stated corrected formulation is established. Verification is limited to the identified result and the explicit elementary deductions; no broader literature-completeness claim is made.

\paragraph{Audit.} Literal correctness: correct. Publication match: exact. Verification: read the result. Development: annotated. Confidence: high.

\noindent References: \citep{cherednik1978differential}.\par

\subsubsection{1983-606.10 --- Mixed derivatives of the wave function}

\label{item:1983-606.10}

\noindent\textbf{Kind:} identity. \textbf{Fate:} E. \textbf{Pages:} 1131.\par

\paragraph{Notebook content.} Notebook occurrence: P001131, equation 14. Mixed derivatives of the wave function.

\paragraph{Precise statement.} If \(\partial_{x_m}\Psi=L_m\Psi\), then \(\partial_{x_1}((L_m\Psi)/\Psi)=\partial_{x_m}(\partial_{x_1}\Psi/\Psi)=\partial_{x_m}\chi\).

\paragraph{Editorial reconstruction.} This is the commutation of mixed derivatives of \(\log\Psi\); it is formal and independent of analytic convergence. It also identifies the missing quotient-rule term in equation 12.

\paragraph{Literature and fate.} The formal-operator and wave-function framework was published or circulating before the talk. The cited result locates the corrected identity; a later exposition is not treated as its discovery date. Locator and verification scope: Section 4 auxiliary linear equations; direct differentiation. \cite{segalwilson1985loop}.

\paragraph{Present status.} As of 12 November2026: the stated corrected formulation is established. Verification is limited to the identified result and the explicit elementary deductions; no broader literature-completeness claim is made.

\paragraph{Audit.} Literal correctness: correct. Publication match: exact. Verification: read the result. Development: annotated. Confidence: high.

\noindent References: \citep{segalwilson1985loop}.\par

\subsubsection{1983-606.11 --- Coefficient relation among the commuting KP flows}

\label{item:1983-606.11}

\noindent\textbf{Kind:} identity. \textbf{Fate:} E. \textbf{Pages:} 1131.\par

\paragraph{Notebook content.} Notebook occurrence: P001131, equation 15. Coefficient relation among the commuting KP flows.

\paragraph{Precise statement.} Comparing coefficients in the corrected equation 12 gives \(\varepsilon_\ell-\ell\chi_\ell=\sum_{m=1}^{\ell-1}\partial_{x_m}\chi_{\ell-m}\), for \(\ell\ge1\), with the empty sum zero.

\paragraph{Editorial reconstruction.} The powers are fixed by \(k^{-m-1}k^{-(\ell-m)}=k^{-\ell-1}\). The upper endpoint \(\ell-1\) excludes the constant leading spectral term, whose derivative vanishes.

\paragraph{Literature and fate.} The formal-operator and wave-function framework was published or circulating before the talk. The cited result locates the corrected identity; a later exposition is not treated as its discovery date. Locator and verification scope: Formal coefficient deduction from the resolvent and auxiliary equations. \cite{cherednik1978differential,segalwilson1985loop}.

\paragraph{Present status.} As of 12 November2026: the stated corrected formulation is established. Verification is limited to the identified result and the explicit elementary deductions; no broader literature-completeness claim is made.

\paragraph{Audit.} Literal correctness: correct. Publication match: exact. Verification: read the result. Development: annotated. Confidence: high.

\noindent References: \citep{cherednik1978differential,segalwilson1985loop}.\par

\subsubsection{1983-606.12 --- A potential for epsilon and its distinction from the tau-function}

\label{item:1983-606.12}

\noindent\textbf{Kind:} claim. \textbf{Fate:} E. \textbf{Pages:} 1131.\par

\paragraph{Notebook content.} Notebook occurrence: P001131, equation 16. A potential for epsilon and its distinction from the tau-function.

\paragraph{Precise statement.} For a compatible KP solution, the densities admit locally a potential \(U\) with \(\varepsilon_j=\partial_{x_j}U\). In the standard tau normalization one may take \(U=-\partial_{x_1}\log\tau_{\mathrm{KP}}\), hence \(\varepsilon_j=-\partial_{x_1}\partial_{x_j}\log\tau_{\mathrm{KP}}\).

\paragraph{Editorial reconstruction.} Equation 16 calls the first potential \(\tau\). It cannot simultaneously be the multiplicative Sato tau-function in equation 17 without a change of notation and a logarithmic derivative. Existence is local or formal; a globally single-valued potential needs a period condition.

\paragraph{Literature and fate.} The formal-operator and wave-function framework was published or circulating before the talk. The cited result locates the corrected identity; a later exposition is not treated as its discovery date. Locator and verification scope: Proposition 5.14, equation 5.16, and the logarithmic-derivative normalization. \cite{segalwilson1985loop}.

\paragraph{Present status.} As of 12 November2026: the stated corrected formulation is established. Verification is limited to the identified result and the explicit elementary deductions; no broader literature-completeness claim is made.

\paragraph{Audit.} Literal correctness: correct with missing hypotheses. Publication match: exact. Verification: read the result. Development: annotated. Confidence: high.

\noindent References: \citep{segalwilson1985loop}.\par

\subsubsection{1983-606.13 --- The wave-function tau quotient requires spectral powers}

\label{item:1983-606.13}

\noindent\textbf{Kind:} formula. \textbf{Fate:} E. \textbf{Pages:} 1131.\par

\paragraph{Notebook content.} Notebook occurrence: P001131, equation 17. The wave-function tau quotient requires spectral powers.

\paragraph{Precise statement.} \(\Psi(x,k)=\exp(\sum_{j\ge1}x_jk^j)\,\tau_{\mathrm{KP}}(x-[k^{-1}])/\tau_{\mathrm{KP}}(x)\), where \([k^{-1}]_j=k^{-j}/j\). Equivalently the numerator is \(\exp(-\sum_{j\ge1}k^{-j}\partial_{x_j}/j)\tau_{\mathrm{KP}}(x)\).

\paragraph{Editorial reconstruction.} The scan omits \(k^{-j}\) in the shift operator. Without these powers the numerator is independent of the spectral parameter and cannot yield the normalized inverse-power wave expansion. The repaired formula matches Segal--Wilson equation 5.16 exactly, and the same construction was already recorded in 1982. Related earlier or repeated occurrences: \texttt{1982-592.7}.

\paragraph{Literature and fate.} The construction was already circulating before this talk, as witnessed by1982-592.7; the 1985 publication is a later precise exposition, not evidence of a new post1983 solution. Locator and verification scope: Proposition 5.14 and equation 5.16, pp. 32--33. \cite{segalwilson1985loop}.

\paragraph{Present status.} As of 12 November2026: the stated corrected formulation is established. Verification is limited to the identified result and the explicit elementary deductions; no broader literature-completeness claim is made.

\paragraph{Audit.} Literal correctness: false as written. Publication match: exact. Verification: read the result. Development: annotated. Confidence: high.

\noindent References: \citep{segalwilson1985loop}.\par

\subsubsection{1983-606.14 --- Attribution to Flaschka}

\label{item:1983-606.14}

\noindent\textbf{Kind:} remark. \textbf{Fate:} F. \textbf{Pages:} 1131.\par

\paragraph{Notebook content.} Notebook occurrence: P001131, closing attribution before Olshanski. Attribution to Flaschka.

\paragraph{Precise statement.} The notes attach Flaschka's name to this circle of tau-function identities without a title, date, or uniquely identifying theorem.

\paragraph{Editorial reconstruction.} The mathematical tau identity is identified in the previous occurrence. This separate bibliographic attribution cannot be resolved to one publication from the surviving wording.

\paragraph{Literature and fate.} No unique publication match has been established for this attribution or incomplete remark. Locator and verification scope: Introduction and section 5 identify the Sato/Date--Jimbo--Kashiwara--Miwa context, but do not identify the notebook's particular Flaschka reference. \cite{segalwilson1985loop}.

\paragraph{Present status.} As of 12 November2026: The attributed Flaschka publication remains unidentified within the documented search scope; the mathematical formula itself is established.

\paragraph{Audit.} Literal correctness: correct. Publication match: thematic only. Verification: unresolved lead. Development: annotated. Confidence: low.

\noindent References: \citep{segalwilson1985loop}.\par

\subsection{Record 607: Representations of the stable unitary group from oscillator and spin models}

\label{talk:607}

\noindent\textbf{Date:} 1983-10-24. \textbf{Speaker:} G. I. Olshanskii. \textbf{Pages:} 1131--1132.\par

The named Olshanskii heading follows the tau-function discussion on P001131. Conjugated embeddings into stable symplectic groups produce representation families, whose infinitesimal oscillator formulas continue on P001132. The unqualified full-Fock irreducibility assertion omits the component selection required in the published dual-pair construction.

\subsubsection{1983-607.1 --- Stable unitary groups and the bosonic oscillator representation}

\label{item:1983-607.1}

\noindent\textbf{Kind:} construction. \textbf{Fate:} E. \textbf{Pages:} 1131.\par

\paragraph{Notebook content.} Notebook occurrence: P001131, Olshanski opening. Stable unitary groups and the bosonic oscillator representation.

\paragraph{Precise statement.} Use finite-support block embeddings \(SU(N)\hookrightarrow SU(N+1)\) to form \(SU(\infty)\). On the bosonic Fock completion of \(\mathbb C[z_1,z_2,\ldots]\), finite-dimensional metaplectic representations are compatible and give an oscillator representation of the corresponding direct-limit metaplectic group.

\paragraph{Editorial reconstruction.} The metaplectic double cover, the Fock Hilbert completion, and the topology used for continuity must be specified; a purely algebraic polynomial model is a dense invariant domain, not the full Hilbert representation. Related earlier or repeated occurrences: \texttt{1974-331.3}, \texttt{1975-371.1}.

\paragraph{Literature and fate.} Olshanskii's1983 paper predates the October talk and gives the relevant dual-pair construction, with the qualifications stated here. Locator and verification scope: Sections5--6, oscillator representation and dual-pair decomposition. \cite{olshanskii1983pairs}.

\paragraph{Present status.} As of 12 November2026: the stated corrected formulation is established. Verification is limited to the identified result and the explicit elementary deductions; no broader literature-completeness claim is made.

\paragraph{Audit.} Literal correctness: correct. Publication match: exact. Verification: read the result. Development: annotated. Confidence: high.

\noindent References: \citep{olshanskii1983pairs}.\par

\subsubsection{1983-607.2 --- Complex block coordinates for the real symplectic Lie algebra}

\label{item:1983-607.2}

\noindent\textbf{Kind:} definition. \textbf{Fate:} E. \textbf{Pages:} 1131.\par

\paragraph{Notebook content.} Notebook occurrence: P001131, displayed matrix conditions. Complex block coordinates for the real symplectic Lie algebra.

\paragraph{Precise statement.} In complex coordinates, \(\operatorname{Sp}(2N,\mathbb R)\) is the subgroup preserving both the standard indefinite Hermitian form and the complex symplectic form. Its Lie algebra consists of matrices \(\begin{pmatrix}A&B\\ \overline B&\overline A\end{pmatrix}\) with \(A^*=-A\), \(B^t=B\). Pass to finite-support blocks for the stable algebra.

\paragraph{Editorial reconstruction.} The signs in a block realization depend on the chosen symplectic matrix; the displayed convention is equivalent to the usual real symplectic algebra by a fixed change of basis.

\paragraph{Literature and fate.} Olshanskii's1983 paper predates the October talk and gives the relevant dual-pair construction, with the qualifications stated here. Locator and verification scope: Sections5--6 supply the stable symplectic setting; block conditions checked directly. \cite{olshanskii1983pairs}.

\paragraph{Present status.} As of 12 November2026: the stated corrected formulation is established. Verification is limited to the identified result and the explicit elementary deductions; no broader literature-completeness claim is made.

\paragraph{Audit.} Literal correctness: correct. Publication match: exact. Verification: read the result. Development: annotated. Confidence: high.

\noindent References: \citep{olshanskii1983pairs}.\par

\subsubsection{1983-607.3 --- Conjugated embeddings and the parameter family}

\label{item:1983-607.3}

\noindent\textbf{Kind:} construction. \textbf{Fate:} E. \textbf{Pages:} 1131.\par

\paragraph{Notebook content.} Notebook occurrence: P001131, tau_s formula. Conjugated embeddings and the parameter family.

\paragraph{Precise statement.} For \(g\in SU(\infty)\), set \(\tau_0(g)=\operatorname{diag}(\overline g,g)\), \(S_s=\begin{pmatrix}\cosh s\,I&\sinh s\,I\\ \sinh s\,I&\cosh s\,I\end{pmatrix}\), and \(\tau_s(g)=S_s\tau_0(g)S_{-s}\). Every \(\tau_s(g)-I\) has finite support, so the formula defines a homomorphism into the stable symplectic group.

\paragraph{Editorial reconstruction.} Although the uniform infinite matrix \(S_s\) itself is not a finite-support element and need not be implementable on the original Fock representation, conjugating each finite-support \(g\) is well-defined. Composition with the oscillator representation requires the compatible lift to the metaplectic cover.

\paragraph{Literature and fate.} Olshanskii's1983 paper predates the October talk and gives the relevant dual-pair construction, with the qualifications stated here. Locator and verification scope: Section 6, external conjugations used to construct representation families. \cite{olshanskii1983pairs}.

\paragraph{Present status.} As of 12 November2026: the stated corrected formulation is established. Verification is limited to the identified result and the explicit elementary deductions; no broader literature-completeness claim is made.

\paragraph{Audit.} Literal correctness: correct. Publication match: partial. Verification: read the result. Development: annotated. Confidence: high.

\noindent References: \citep{olshanskii1983pairs}.\par

\subsubsection{1983-607.4 --- The full bosonic Fock space cannot give an irreducible representation}

\label{item:1983-607.4}

\noindent\textbf{Kind:} claim. \textbf{Fate:} E. \textbf{Pages:} 1131.\par

\paragraph{Notebook content.} Notebook occurrence: P001131, irreducibility of T_s^+. The full bosonic Fock space cannot give an irreducible representation.

\paragraph{Precise statement.} The composite \(T_s^+=W\circ\widetilde\tau_s\) is unitary, but on full bosonic Fock space it preserves both the even- and odd-degree subspaces. Therefore the blanket assertion of irreducibility for \(s\ne0\) is false without selecting an appropriate component.

\paragraph{Editorial reconstruction.} Every infinitesimal oscillator generator changes degree by \(0\) or \(2\), so parity commutes with the representation and both parity spaces are nonzero. Olshanskii's published construction includes decomposition under a commuting compact group before its irreducibility theorem; for a one-copy orthogonal commutant this includes the parity selection omitted in the notes.

\paragraph{Literature and fate.} Olshanskii's1983 paper predates the October talk and gives the relevant dual-pair construction, with the qualifications stated here. Locator and verification scope: Section 5 decomposition under O(k); section 6, Theorem 2. \cite{olshanskii1983pairs}.

\paragraph{Present status.} As of 12 November2026: the stated corrected formulation is established. Verification is limited to the identified result and the explicit elementary deductions; no broader literature-completeness claim is made.

\paragraph{Audit.} Literal correctness: false as written. Publication match: partial. Verification: read the result. Development: annotated. Confidence: high.

\noindent References: \citep{olshanskii1983pairs}.\par

\subsubsection{1983-607.5 --- The fermionic analogue}

\label{item:1983-607.5}

\noindent\textbf{Kind:} construction and claim. \textbf{Fate:} E. \textbf{Pages:} 1131.\par

\paragraph{Notebook content.} Notebook occurrence: P001131, T_t^-. The fermionic analogue.

\paragraph{Precise statement.} The corresponding spin representation on a fermionic Fock space can be restricted along a suitable family of embeddings to produce unitary representations of the stable unitary group. Any irreducibility assertion refers to the specified summands in the dual-pair construction.

\paragraph{Editorial reconstruction.} The notebook writes \(T_t^-=S\circ\tau_t^-\) but does not give the map \(\tau_t^-\), its parameter range, or the chosen summand. Olshanskii's sections5--6 contain a spin construction with a symplectic commutant and parameters \(0<b<1\); an exact conversion between that parameter and the notebook's \(t\) has not been established.

\paragraph{Literature and fate.} Olshanskii's1983 paper predates the October talk and gives the relevant dual-pair construction, with the qualifications stated here. Locator and verification scope: Sections5--6, fermionic/spin analogue and Theorem 2. \cite{olshanskii1983pairs}.

\paragraph{Present status.} As of 12 November2026: The dual-pair construction is published, while the precise parameter convention and component of this notebook formula remain unverified.

\paragraph{Audit.} Literal correctness: not determinable. Publication match: partial. Verification: read the result. Development: annotated. Confidence: medium.

\noindent References: \citep{olshanskii1983pairs}.\par

\subsubsection{1983-607.6 --- Quadratic differential operators realize the oscillator Lie algebra}

\label{item:1983-607.6}

\noindent\textbf{Kind:} construction. \textbf{Fate:} E. \textbf{Pages:} 1132.\par

\paragraph{Notebook content.} Notebook occurrence: P001132, three displayed operator assignments. Quadratic differential operators realize the oscillator Lie algebra.

\paragraph{Precise statement.} On polynomials, the three block types act by \(z_i\partial_j+\delta_{ij}/2\), \(\partial_i\partial_j\), and \(-z_iz_j\), respectively, for blocks \(\operatorname{diag}(-E_{ji},E_{ij})\), upper-right \(E_{ij}+E_{ji}\), and lower-left \(E_{ij}+E_{ji}\), in the notebook's polarization. They satisfy the corresponding symplectic commutators.

\paragraph{Editorial reconstruction.} For \(i=j\), the off-diagonal matrix block is \(2E_{ii}\); no extra factor \(1/2\) should be inserted in the resulting operators. For example \([\partial_i^2,-z_i^2]=-4z_i\partial_i-2\), consistent with the diagonal generator and its vacuum shift. Related earlier or repeated occurrences: \texttt{1975-371.1}.

\paragraph{Literature and fate.} Olshanskii's1983 paper predates the October talk and gives the relevant dual-pair construction, with the qualifications stated here. Locator and verification scope: Section 5 oscillator construction; displayed commutator independently checked. \cite{olshanskii1983pairs}.

\paragraph{Present status.} As of 12 November2026: the stated corrected formulation is established. Verification is limited to the identified result and the explicit elementary deductions; no broader literature-completeness claim is made.

\paragraph{Audit.} Literal correctness: correct. Publication match: exact. Verification: read the result. Development: annotated. Confidence: high.

\noindent References: \citep{olshanskii1983pairs}.\par

\subsection{Record 608: Symmetric-function transition matrices}

\label{talk:608}

\noindent\textbf{Date:} 1983-10-31. \textbf{Speaker:} A. V. Zelevinsky. \textbf{Pages:} 1133--1134.\par

The lecture repeats the graded ring and standard bases, then interprets transition coefficients by nonnegative, zero-one and weighted-assignment matrices. On P001134 the explicit I.M.G. partition-number aside interrupts the exposition; the Schur alternant formula below it resumes the symmetric-function discussion. These repeated occurrences remain indexed separately.

\subsubsection{1983-608.1 --- Stable homogeneous symmetric functions, repeated}

\label{item:1983-608.1}

\noindent\textbf{Kind:} definition. \textbf{Fate:} E. \textbf{Pages:} 1133.\par

\paragraph{Notebook content.} Notebook occurrence: P001133, opening Lambda definition. Stable homogeneous symmetric functions, repeated.

\paragraph{Precise statement.} For degree \(d\), the specialization maps \(\mathbb Z[x_1,\ldots,x_{N+1}]^{S_{N+1}}_d\to\mathbb Z[x_1,\ldots,x_N]^{S_N}_d\) are isomorphisms for \(N\ge d\); their stable values define \(\Lambda^d\), and \(\Lambda=\bigoplus_d\Lambda^d\).

\paragraph{Editorial reconstruction.} This repeats the construction of 1983-605.1 and is retained as a new occurrence. The graded limit prevents unbounded-degree completion. Related earlier or repeated occurrences: \texttt{1983-605.1}.

\paragraph{Literature and fate.} Classical symmetric-function mathematics, already available by the 1979 first edition; the opened 1995 edition verifies the modern formulation. Locator and verification scope: Chapter I, section 2. \cite{macdonald1995symmetric}.

\paragraph{Present status.} As of 12 November2026: the stated corrected formulation is established. Verification is limited to the identified result and the explicit elementary deductions; no broader literature-completeness claim is made.

\paragraph{Audit.} Literal correctness: correct. Publication match: exact. Verification: read the result. Development: annotated. Confidence: high.

\noindent References: \citep{macdonald1995symmetric}.\par

\subsubsection{1983-608.2 --- Two notations for partitions}

\label{item:1983-608.2}

\noindent\textbf{Kind:} definition. \textbf{Fate:} E. \textbf{Pages:} 1133.\par

\paragraph{Notebook content.} Notebook occurrence: P001133, lambda=(lambda_i)=(1^m1...). Two notations for partitions.

\paragraph{Precise statement.} A partition \(\lambda\) can be written as a nonincreasing finite list of positive parts or in frequency notation \(\lambda=(1^{m_1}2^{m_2}\cdots)\), where \(m_r=\#\{i:\lambda_i=r\}\). Its size is \(|\lambda|=\sum_i\lambda_i=\sum_r rm_r\).

\paragraph{Editorial reconstruction.} Zeros may be appended for finite-variable formulas but do not count as positive parts. Related earlier or repeated occurrences: \texttt{1976-415.1}, \texttt{1983-605.2}.

\paragraph{Literature and fate.} Classical symmetric-function mathematics, already available by the 1979 first edition; the opened 1995 edition verifies the modern formulation. Locator and verification scope: Chapter I, section 1. \cite{macdonald1995symmetric}.

\paragraph{Present status.} As of 12 November2026: the stated corrected formulation is established. Verification is limited to the identified result and the explicit elementary deductions; no broader literature-completeness claim is made.

\paragraph{Audit.} Literal correctness: correct. Publication match: exact. Verification: read the result. Development: annotated. Confidence: high.

\noindent References: \citep{macdonald1995symmetric}.\par

\subsubsection{1983-608.3 --- Monomial symmetric functions, repeated}

\label{item:1983-608.3}

\noindent\textbf{Kind:} definition. \textbf{Fate:} E. \textbf{Pages:} 1133.\par

\paragraph{Notebook content.} Notebook occurrence: P001133, m_lambda. Monomial symmetric functions, repeated.

\paragraph{Precise statement.} For each partition \(\lambda\), \(m_\lambda\) is the sum of all distinct monomials with exponent multiset \(\lambda\). The \(m_\lambda\) with \(|\lambda|=d\) form the integral basis of \(\Lambda^d\).

\paragraph{Editorial reconstruction.} Terms with repeated exponent sequences are counted once. This repetition supplies the basis used in the subsequent coefficient-counting exercises. Related earlier or repeated occurrences: \texttt{1983-605.2}.

\paragraph{Literature and fate.} Classical symmetric-function mathematics, already available by the 1979 first edition; the opened 1995 edition verifies the modern formulation. Locator and verification scope: Chapter I, section 2. \cite{macdonald1995symmetric}.

\paragraph{Present status.} As of 12 November2026: the stated corrected formulation is established. Verification is limited to the identified result and the explicit elementary deductions; no broader literature-completeness claim is made.

\paragraph{Audit.} Literal correctness: correct. Publication match: exact. Verification: read the result. Development: annotated. Confidence: high.

\noindent References: \citep{macdonald1995symmetric}.\par

\subsubsection{1983-608.4 --- Complete symmetric functions, repeated}

\label{item:1983-608.4}

\noindent\textbf{Kind:} definition. \textbf{Fate:} E. \textbf{Pages:} 1133.\par

\paragraph{Notebook content.} Notebook occurrence: P001133, h_n. Complete symmetric functions, repeated.

\paragraph{Precise statement.} For \(n\ge0\), \(h_n=\sum_{\lambda\vdash n}m_\lambda\), with \(h_0=1\), and \(h_\lambda=\prod_i h_{\lambda_i}\).

\paragraph{Editorial reconstruction.} The following matrix formula concerns products \(h_\lambda\), not only the single homogeneous generator \(h_n\). Related earlier or repeated occurrences: \texttt{1983-605.4}.

\paragraph{Literature and fate.} Classical symmetric-function mathematics, already available by the 1979 first edition; the opened 1995 edition verifies the modern formulation. Locator and verification scope: Chapter I, section 2. \cite{macdonald1995symmetric}.

\paragraph{Present status.} As of 12 November2026: the stated corrected formulation is established. Verification is limited to the identified result and the explicit elementary deductions; no broader literature-completeness claim is made.

\paragraph{Audit.} Literal correctness: correct. Publication match: exact. Verification: read the result. Development: annotated. Confidence: high.

\noindent References: \citep{macdonald1995symmetric}.\par

\subsubsection{1983-608.5 --- Elementary symmetric functions, repeated}

\label{item:1983-608.5}

\noindent\textbf{Kind:} definition. \textbf{Fate:} E. \textbf{Pages:} 1133.\par

\paragraph{Notebook content.} Notebook occurrence: P001133, e_n. Elementary symmetric functions, repeated.

\paragraph{Precise statement.} For \(n\ge0\), \(e_n=m_{(1^n)}\), with \(e_0=1\), and \(e_\lambda=\prod_i e_{\lambda_i}\). In \(N\) variables \(e_n=0\) for \(n>N\).

\paragraph{Editorial reconstruction.} The stable ring has generators of all positive degrees even though finite-variable specializations eventually vanish. Related earlier or repeated occurrences: \texttt{1983-605.3}.

\paragraph{Literature and fate.} Classical symmetric-function mathematics, already available by the 1979 first edition; the opened 1995 edition verifies the modern formulation. Locator and verification scope: Chapter I, section 2. \cite{macdonald1995symmetric}.

\paragraph{Present status.} As of 12 November2026: the stated corrected formulation is established. Verification is limited to the identified result and the explicit elementary deductions; no broader literature-completeness claim is made.

\paragraph{Audit.} Literal correctness: correct. Publication match: exact. Verification: read the result. Development: annotated. Confidence: high.

\noindent References: \citep{macdonald1995symmetric}.\par

\subsubsection{1983-608.6 --- Power sums, repeated}

\label{item:1983-608.6}

\noindent\textbf{Kind:} definition. \textbf{Fate:} E. \textbf{Pages:} 1133.\par

\paragraph{Notebook content.} Notebook occurrence: P001133, p_n. Power sums, repeated.

\paragraph{Precise statement.} For \(n\ge1\), \(p_n=m_{(n)}=\sum_i x_i^n\), and \(p_\lambda=\prod_jp_{\lambda_j}\).

\paragraph{Editorial reconstruction.} When two parts of \(\lambda\) coincide, their positions in this product remain distinct choices in a coefficient-counting expansion. Related earlier or repeated occurrences: \texttt{1983-605.5}.

\paragraph{Literature and fate.} Classical symmetric-function mathematics, already available by the 1979 first edition; the opened 1995 edition verifies the modern formulation. Locator and verification scope: Chapter I, section 2. \cite{macdonald1995symmetric}.

\paragraph{Present status.} As of 12 November2026: the stated corrected formulation is established. Verification is limited to the identified result and the explicit elementary deductions; no broader literature-completeness claim is made.

\paragraph{Audit.} Literal correctness: correct. Publication match: exact. Verification: read the result. Development: annotated. Confidence: high.

\noindent References: \citep{macdonald1995symmetric}.\par

\subsubsection{1983-608.7 --- Coefficient rings in the basis theorem, repeated}

\label{item:1983-608.7}

\noindent\textbf{Kind:} theorem. \textbf{Fate:} E. \textbf{Pages:} 1133.\par

\paragraph{Notebook content.} Notebook occurrence: P001133, bases of Lambda. Coefficient rings in the basis theorem, repeated.

\paragraph{Precise statement.} The families \(m_\lambda,h_\lambda,e_\lambda\) are bases over \(\mathbb Z\), whereas \(p_\lambda\) are a basis over \(\mathbb Q\). The elementary and complete generators freely generate the integral polynomial algebra \(\Lambda\).

\paragraph{Editorial reconstruction.} One cannot omit the rationalization for power sums: \(p_1^2-p_2=2e_2\). Related earlier or repeated occurrences: \texttt{1983-605.6}.

\paragraph{Literature and fate.} Classical symmetric-function mathematics, already available by the 1979 first edition; the opened 1995 edition verifies the modern formulation. Locator and verification scope: Chapter I, section 2. \cite{macdonald1995symmetric}.

\paragraph{Present status.} As of 12 November2026: the stated corrected formulation is established. Verification is limited to the identified result and the explicit elementary deductions; no broader literature-completeness claim is made.

\paragraph{Audit.} Literal correctness: correct. Publication match: exact. Verification: read the result. Development: annotated. Confidence: high.

\noindent References: \citep{macdonald1995symmetric}.\par

\subsubsection{1983-608.8 --- Nonnegative integer matrices give the h-to-m coefficients}

\label{item:1983-608.8}

\noindent\textbf{Kind:} theorem. \textbf{Fate:} E. \textbf{Pages:} 1133--1134.\par

\paragraph{Notebook content.} Notebook occurrence: P001133-P001134, matrix-counting rule. Nonnegative integer matrices give the h-to-m coefficients.

\paragraph{Precise statement.} In \(h_\lambda=\sum_{\mu}N_{\lambda\mu}m_\mu\), the coefficient \(N_{\lambda\mu}\) is the number of matrices \(A=(a_{ij})\) of nonnegative integers with row sums \(\lambda_i\) and column sums \(\mu_j\). Rows and columns are labelled by their positions, including repeated parts.

\paragraph{Editorial reconstruction.} Expand the \(i\)-th factor \(h_{\lambda_i}\) by recording how many times each variable occurs. The resulting row is a weak composition, and multiplication forces the specified column totals. Related earlier or repeated occurrences: \texttt{1983-605.12}.

\paragraph{Literature and fate.} Classical symmetric-function mathematics, already available by the 1979 first edition; the opened 1995 edition verifies the modern formulation. Locator and verification scope: Chapter I, section 6, equation 6.7(ii). \cite{macdonald1995symmetric}.

\paragraph{Present status.} As of 12 November2026: the stated corrected formulation is established. Verification is limited to the identified result and the explicit elementary deductions; no broader literature-completeness claim is made.

\paragraph{Audit.} Literal correctness: correct. Publication match: exact. Verification: read the result. Development: annotated. Confidence: high.

\noindent Canonical problem: GP-1983-04.\par

\noindent References: \citep{macdonald1995symmetric}.\par

\subsubsection{1983-608.9 --- Zero-one matrices give the e-to-m coefficients}

\label{item:1983-608.9}

\noindent\textbf{Kind:} theorem. \textbf{Fate:} E. \textbf{Pages:} 1134.\par

\paragraph{Notebook content.} Notebook occurrence: P001134, zero-one matrix rule. Zero-one matrices give the e-to-m coefficients.

\paragraph{Precise statement.} In \(e_\lambda=\sum_{\mu}M_{\lambda\mu}m_\mu\), \(M_{\lambda\mu}\) is the number of zero-one matrices with row sums \(\lambda_i\) and column sums \(\mu_j\).

\paragraph{Editorial reconstruction.} Every elementary factor chooses each variable at most once, so every entry in its row is zero or one. The matrix rule remains valid when a row sum exceeds the number of columns, in which case the coefficient is zero.

\paragraph{Literature and fate.} Classical symmetric-function mathematics, already available by the 1979 first edition; the opened 1995 edition verifies the modern formulation. Locator and verification scope: Chapter I, section 6, equation 6.7(i). \cite{macdonald1995symmetric}.

\paragraph{Present status.} As of 12 November2026: the stated corrected formulation is established. Verification is limited to the identified result and the explicit elementary deductions; no broader literature-completeness claim is made.

\paragraph{Audit.} Literal correctness: correct. Publication match: exact. Verification: read the result. Development: annotated. Confidence: high.

\noindent Canonical problem: GP-1983-04.\par

\noindent References: \citep{macdonald1995symmetric}.\par

\subsubsection{1983-608.10 --- The p-to-m transition exercise}

\label{item:1983-608.10}

\noindent\textbf{Kind:} exercise. \textbf{Fate:} E. \textbf{Pages:} 1134.\par

\paragraph{Notebook content.} Notebook occurrence: P001134, remaining p_lambda exercise. The p-to-m transition exercise.

\paragraph{Precise statement.} The coefficient of \(m_\mu\) in \(p_\lambda\), with \(\ell(\lambda)=r\) and \(\ell(\mu)=s\), equals the number of functions \(f:\{1,\ldots,r\}\to\{1,\ldots,s\}\) satisfying \(\sum_{j:f(j)=i}\lambda_j=\mu_i\) for every \(i\).

\paragraph{Editorial reconstruction.} Each labelled factor \(p_{\lambda_j}\) chooses one variable. The formula counts those choices yielding the fixed monomial \(x_1^{\mu_1}\cdots x_s^{\mu_s}\); equal parts of \(\lambda\) are still separate factors. This is an elementary solution of the stated exercise.

\paragraph{Literature and fate.} Classical symmetric-function mathematics, already available by the 1979 first edition; the opened 1995 edition verifies the modern formulation. Locator and verification scope: Chapter I, section 6, equations6.8--6.9. \cite{macdonald1995symmetric}.

\paragraph{Present status.} As of 12 November2026: the stated corrected formulation is established. Verification is limited to the identified result and the explicit elementary deductions; no broader literature-completeness claim is made.

\paragraph{Audit.} Literal correctness: correct. Publication match: exact. Verification: read the result. Development: annotated. Confidence: high.

\noindent Canonical problem: GP-1983-04.\par

\noindent References: \citep{macdonald1995symmetric}.\par

\subsubsection{1983-608.11 --- Generating-series relations and Newton identities}

\label{item:1983-608.11}

\noindent\textbf{Kind:} method and identities. \textbf{Fate:} E. \textbf{Pages:} 1134.\par

\paragraph{Notebook content.} Notebook occurrence: P001134, E,H,P displayed together. Generating-series relations and Newton identities.

\paragraph{Precise statement.} The series \(E(t)=\prod_i(1+tx_i)\), \(H(t)=\prod_i(1-tx_i)^{-1}\), and \(P(t)=\sum_{r\ge1}p_rt^r\) satisfy \(E(-t)H(t)=1\) and \(P(t)=tH'(t)/H(t)\). Consequently \(nh_n=\sum_{r=1}^np_rh_{n-r}\) for \(n\ge1\).

\paragraph{Editorial reconstruction.} These relations operate formally in \(\Lambda[[t]]\); coefficient extraction proves the recurrence without analytic convergence. The Newton recurrence is an explicit consequence of the displayed notebook series. Related earlier or repeated occurrences: \texttt{1983-605.3}, \texttt{1983-605.4}, \texttt{1983-605.5}.

\paragraph{Literature and fate.} Classical symmetric-function mathematics, already available by the 1979 first edition; the opened 1995 edition verifies the modern formulation. Locator and verification scope: Chapter I, section 2. \cite{macdonald1995symmetric}.

\paragraph{Present status.} As of 12 November2026: the stated corrected formulation is established. Verification is limited to the identified result and the explicit elementary deductions; no broader literature-completeness claim is made.

\paragraph{Audit.} Literal correctness: correct. Publication match: exact. Verification: read the result. Development: annotated. Confidence: high.

\noindent References: \citep{macdonald1995symmetric}.\par

\subsubsection{1983-608.12 --- Schur alternants resumed after Gelfand's aside}

\label{item:1983-608.12}

\noindent\textbf{Kind:} definition. \textbf{Fate:} E. \textbf{Pages:} 1134.\par

\paragraph{Notebook content.} Notebook occurrence: P001134, below I.M.G. partition discussion. Schur alternants resumed after Gelfand's aside.

\paragraph{Precise statement.} With \(\delta=(N-1,\ldots,0)\), the resumed Schur formula is \(s_\lambda(x)=\det(x_i^{\lambda_j+N-j})/\det(x_i^{N-j})\). The denominator is the Vandermonde determinant, and the quotient is a symmetric polynomial stable under adjoining a zero variable.

\paragraph{Editorial reconstruction.} Physical page order interleaves Gelfand's partition-number aside before this return to Zelevinsky's symmetric-function exposition. The occurrence is assigned to608 by continuity of the Schur topic, with the interruption retained explicitly in coverage notes. Related earlier or repeated occurrences: \texttt{1983-605.7}, \texttt{1983-605.8}.

\paragraph{Literature and fate.} Classical symmetric-function mathematics, already available by the 1979 first edition; the opened 1995 edition verifies the modern formulation. Locator and verification scope: Chapter I, section 3. \cite{macdonald1995symmetric}.

\paragraph{Present status.} As of 12 November2026: the stated corrected formulation is established. Verification is limited to the identified result and the explicit elementary deductions; no broader literature-completeness claim is made.

\paragraph{Audit.} Literal correctness: correct. Publication match: exact. Verification: read the result. Development: annotated. Confidence: high.

\noindent References: \citep{macdonald1995symmetric}.\par

\subsection{Record 609: Partition numbers and an unattributed anomaly recap}

\label{talk:609}

\noindent\textbf{Date:} 1983-10-31. \textbf{Speaker:} I. M. Gelfand; an unattributed continuation. \textbf{Pages:} 1134--1135.\par

The explicit \ru{И.М.Г.} passage on P001134 concerns partition asymptotics and Euler's product. The Witten-anomaly recap at the beginning of P001135 has no speaker label and precedes the named Perelomov heading. It is included as an unattributed continuation of the same session, with no claim that Gelfand authored it. The multi-term partition formula is not written, so its exact version cannot be uniquely identified.

\subsubsection{1983-609.1 --- The partition-number function}

\label{item:1983-609.1}

\noindent\textbf{Kind:} definition. \textbf{Fate:} E. \textbf{Pages:} 1134.\par

\paragraph{Notebook content.} Notebook occurrence: P001134, I.M.G. interruption. The partition-number function.

\paragraph{Precise statement.} Let \(p(n)\) be the number of unordered partitions of the nonnegative integer \(n\) into positive integers, with \(p(0)=1\).

\paragraph{Editorial reconstruction.} Order is ignored; compositions would have a different counting function. The explicit \ru{И.М.Г.} heading supports Gelfand's authorship of this aside.

\paragraph{Literature and fate.} The identified result was available before 1983; the notes do not constitute a new open problem. Locator and verification scope: Section 1.1, equation 1.11. \cite{hardyramanujan1918asymptotic}.

\paragraph{Present status.} As of 12 November2026: the stated corrected formulation is established. Verification is limited to the identified result and the explicit elementary deductions; no broader literature-completeness claim is made.

\paragraph{Audit.} Literal correctness: correct. Publication match: exact. Verification: read the result. Development: annotated. Confidence: high.

\noindent References: \citep{hardyramanujan1918asymptotic}.\par

\subsubsection{1983-609.2 --- Euler's product for partition numbers}

\label{item:1983-609.2}

\noindent\textbf{Kind:} identity. \textbf{Fate:} E. \textbf{Pages:} 1134.\par

\paragraph{Notebook content.} Notebook occurrence: P001134, generating product. Euler's product for partition numbers.

\paragraph{Precise statement.} \(\sum_{n\ge0}p(n)t^n=\prod_{j\ge1}(1-t^j)^{-1}\) in \(\mathbb Z[[t]]\). It also converges analytically for \(|t|<1\).

\paragraph{Editorial reconstruction.} Expanding the \(j\)-th factor chooses the multiplicity of the part \(j\). For each fixed coefficient only finitely many factors contribute, which proves the formal identity. Related earlier or repeated occurrences: \texttt{1975-366.1}, \texttt{1976-408.3}.

\paragraph{Literature and fate.} The identified result was available before 1983; the notes do not constitute a new open problem. Locator and verification scope: Equation 1.11. \cite{hardyramanujan1918asymptotic}.

\paragraph{Present status.} As of 12 November2026: the stated corrected formulation is established. Verification is limited to the identified result and the explicit elementary deductions; no broader literature-completeness claim is made.

\paragraph{Audit.} Literal correctness: correct. Publication match: exact. Verification: read the result. Development: annotated. Confidence: high.

\noindent References: \citep{hardyramanujan1918asymptotic}.\par

\subsubsection{1983-609.3 --- Ramanujan's partition asymptotics and the distinction from a convergent exact series}

\label{item:1983-609.3}

\noindent\textbf{Kind:} remark and theorem. \textbf{Fate:} E. \textbf{Pages:} 1134.\par

\paragraph{Notebook content.} Notebook occurrence: P001134, comment on an explicit formula with further terms. Ramanujan's partition asymptotics and the distinction from a convergent exact series.

\paragraph{Precise statement.} A precise classical form of the partition asymptotic is \(p(n)\sim(4n\sqrt3)^{-1}\exp(\pi\sqrt{2n/3})\) as \(n\to\infty\). Hardy--Ramanujan supplied more accurate finite expansions; Rademacher later obtained a convergent exact series.

\paragraph{Editorial reconstruction.} The notebook does not write the multi-term formula or identify a truncation. Thus the broad reference to Ramanujan is compatible with Hardy--Ramanujan asymptotics, but cannot uniquely identify the exact version. Rademacher's convergent series is a related1937 result and must not be silently attributed to Ramanujan.

\paragraph{Literature and fate.} The identified result was available before 1983; the notes do not constitute a new open problem. Locator and verification scope: Hardy--Ramanujan section 1.4, equation 1.41, and sections1.5--1.6; Rademacher publisher metadata. \cite{hardyramanujan1918asymptotic,rademacher1937partition}.

\paragraph{Present status.} As of 12 November2026: the stated corrected formulation is established. Verification is limited to the identified result and the explicit elementary deductions; no broader literature-completeness claim is made.

\paragraph{Audit.} Literal correctness: correct. Publication match: partial. Verification: read the result. Development: annotated. Confidence: medium.

\noindent References: \citep{hardyramanujan1918asymptotic,rademacher1937partition}.\par

\subsubsection{1983-609.4 --- The reference to Andrews's book}

\label{item:1983-609.4}

\noindent\textbf{Kind:} remark. \textbf{Fate:} E. \textbf{Pages:} 1134.\par

\paragraph{Notebook content.} Notebook occurrence: P001134, bibliographic aside. The reference to Andrews's book.

\paragraph{Precise statement.} The speaker recommends Andrews's treatment of partition theory, without a complete title or theorem locator in the notes.

\paragraph{Editorial reconstruction.} \emph{The Theory of Partitions} first appeared in 1976 and is a plausible identification, supported by the topic and author. The note's short appreciative wording does not identify an additional mathematical assertion.

\paragraph{Literature and fate.} The identified result was available before 1983; the notes do not constitute a new open problem. Locator and verification scope: Author bibliography, book list; Cambridge publisher record. \cite{andrews1976partitions}.

\paragraph{Present status.} As of 12 November2026: the stated corrected formulation is established. Verification is limited to the identified result and the explicit elementary deductions; no broader literature-completeness claim is made.

\paragraph{Audit.} Literal correctness: correct. Publication match: partial. Verification: publisher metadata checked. Development: annotated. Confidence: medium.

\noindent References: \citep{andrews1976partitions}.\par

\subsubsection{1983-609.5 --- Unattributed reprise of the odd-doublet anomaly}

\label{item:1983-609.5}

\noindent\textbf{Kind:} claim. \textbf{Fate:} E. \textbf{Pages:} 1135.\par

\paragraph{Notebook content.} Notebook occurrence: P001135, opening Witten recap before the Perelomov heading. Unattributed reprise of the odd-doublet anomaly.

\paragraph{Precise statement.} For a four-dimensional spin \(SU(2)\) gauge theory containing an odd number of left-handed Weyl doublets in the fundamental representation, the fermion Pfaffian changes sign under a gauge transformation representing the nonzero element of \(\pi_4(SU(2))=\mathbb Z/2\). This obstructs a gauge-invariant choice of its sign.

\paragraph{Editorial reconstruction.} This recap has no new speaker label. It is indexed here as an unattributed continuation of the 31 October session, not as evidence that Gelfand or Perelomov made the assertion. The odd-doublet and spin hypotheses exclude the misleading claim that every \(SU(2)\) theory is anomalous. Related earlier or repeated occurrences: \texttt{1983-603.1}, \texttt{1983-604.7}.

\paragraph{Literature and fate.} The identified result was available before 1983; the notes do not constitute a new open problem. Locator and verification scope: Witten 1982 abstract; Wang--Wen--Witten sections2.1--2.3. \cite{witten1982anomaly,wangwenwitten2019new}.

\paragraph{Present status.} As of 12 November2026: the stated corrected formulation is established. Verification is limited to the identified result and the explicit elementary deductions; no broader literature-completeness claim is made.

\paragraph{Audit.} Literal correctness: correct. Publication match: exact. Verification: read the result. Development: annotated. Confidence: high.

\noindent References: \citep{witten1982anomaly,wangwenwitten2019new}.\par

\subsubsection{1983-609.6 --- The repeated Grassmann Gaussian factor}

\label{item:1983-609.6}

\noindent\textbf{Kind:} claim. \textbf{Fate:} E. \textbf{Pages:} 1135.\par

\paragraph{Notebook content.} Notebook occurrence: P001135, opening determinant expression. The repeated Grassmann Gaussian factor.

\paragraph{Precise statement.} For \(2r\) Grassmann variables and an antisymmetric matrix \(A\), the Berezin Gaussian integral is \(\operatorname{Pf}(A)\), up to the fixed ordering convention for the measure, and its square is \(\det A\). A bosonic Gaussian has a reciprocal square root instead.

\paragraph{Editorial reconstruction.} The half-power must therefore be interpreted with the fermionic sign and Pfaffian choice. This finite-dimensional identity explains the anomaly discussion but is not a definition of the divergent infinite-dimensional gauge-field integral. The recap is unattributed. Related earlier or repeated occurrences: \texttt{1983-603.8}.

\paragraph{Literature and fate.} The identified result was available before 1983; the notes do not constitute a new open problem. Locator and verification scope: Wang--Wen--Witten section 2.2, real fermion Pfaffian. \cite{witten1982anomaly,wangwenwitten2019new}.

\paragraph{Present status.} As of 12 November2026: the stated corrected formulation is established. Verification is limited to the identified result and the explicit elementary deductions; no broader literature-completeness claim is made.

\paragraph{Audit.} Literal correctness: correct. Publication match: exact. Verification: read the result. Development: annotated. Confidence: high.

\noindent References: \citep{witten1982anomaly,wangwenwitten2019new}.\par

\subsection{Record 610: Chiral fields, Skyrme terms, and multivalued actions}

\label{talk:610}

\noindent\textbf{Date:} 1983-10-31. \textbf{Speaker:} A. M. Perelomov. \textbf{Pages:} 1135--1136.\par

The named Perelomov section begins below the anomaly recap on P001135 and ends on P001136; the 14 November material starts on P001137. It revisits two-dimensional instantons and the Skyrme correction, mentions Krichever's nonlinear d'Alembert construction, and compares Witten's topological term with Novikov's multivalued variational framework.

\subsubsection{1983-610.1 --- The chiral sigma-model action, repeated}

\label{item:1983-610.1}

\noindent\textbf{Kind:} definition. \textbf{Fate:} E. \textbf{Pages:} 1135.\par

\paragraph{Notebook content.} Notebook occurrence: P001135, Perelomov heading and Lagrangian. The chiral sigma-model action, repeated.

\paragraph{Precise statement.} For a map \(g:M\to G\) from a Riemannian or Lorentzian spacetime to a compact Lie group with a fixed bi-invariant metric, the principal chiral action is \(S_2(g)=\frac12\int_M\langle (\partial_\mu g)g^{-1},(\partial^\mu g)g^{-1}\rangle\,d\operatorname{vol}\). More generally the sigma model is the Dirichlet action of a map into a specified Riemannian homogeneous space.

\paragraph{Editorial reconstruction.} Left and right Maurer--Cartan conventions give the same action because the target metric is invariant under conjugation. Compact-matrix trace formulas need the sign that makes the Euclidean target metric positive. Related earlier or repeated occurrences: \texttt{1983-601.4}.

\paragraph{Literature and fate.} The identified construction was published before the October 1983 talk; the precise scope is stated separately from the notebook's shorthand. Locator and verification scope: Perelomov section 1; Novikov section 5, equation 19. \cite{perelomov1981instantons,novikov1982hamiltonian}.

\paragraph{Present status.} As of 12 November2026: the stated corrected formulation is established. Verification is limited to the identified result and the explicit elementary deductions; no broader literature-completeness claim is made.

\paragraph{Audit.} Literal correctness: correct. Publication match: exact. Verification: read the result. Development: annotated. Confidence: high.

\noindent References: \citep{perelomov1981instantons,novikov1982hamiltonian}.\par

\subsubsection{1983-610.2 --- Two-dimensional instantons on compact coadjoint orbits}

\label{item:1983-610.2}

\noindent\textbf{Kind:} claim. \textbf{Fate:} E. \textbf{Pages:} 1135.\par

\paragraph{Notebook content.} Notebook occurrence: P001135, classification of instantons. Two-dimensional instantons on compact coadjoint orbits.

\paragraph{Precise statement.} For a compact coadjoint orbit equipped with a compatible invariant Kahler structure, a holomorphic or antiholomorphic map from a Riemann surface satisfies the sigma-model first-order self-duality equations and saturates the corresponding topological energy bound.

\paragraph{Editorial reconstruction.} The statement that all instantons are found is interpreted within this self-dual sector. It does not classify every harmonic map or every stationary point of the Dirichlet energy; Perelomov's review separately discusses non-self-dual solutions. Related earlier or repeated occurrences: \texttt{1983-601.9}, \texttt{1983-601.13}.

\paragraph{Literature and fate.} The identified construction was published before the October 1983 talk; the precise scope is stated separately from the notebook's shorthand. Locator and verification scope: Section 2(b), equations2.39--2.48; section 3(c) for other harmonic solutions. \cite{perelomov1981instantons}.

\paragraph{Present status.} As of 12 November2026: the stated corrected formulation is established. Verification is limited to the identified result and the explicit elementary deductions; no broader literature-completeness claim is made.

\paragraph{Audit.} Literal correctness: correct with missing hypotheses. Publication match: exact. Verification: read the result. Development: annotated. Confidence: high.

\noindent References: \citep{perelomov1981instantons}.\par

\subsubsection{1983-610.3 --- The Belavin--Polyakov rational-map example, repeated}

\label{item:1983-610.3}

\noindent\textbf{Kind:} example. \textbf{Fate:} E. \textbf{Pages:} 1135.\par

\paragraph{Notebook content.} Notebook occurrence: P001135, CP1 example. The Belavin--Polyakov rational-map example, repeated.

\paragraph{Precise statement.} For the \(S^2\simeq\mathbb{CP}^1\) target, finite-energy self-dual maps on the compactified plane \(S^2\) are rational functions in a complex coordinate; anti-self-dual maps are antiholomorphic rational functions. Their topological degree controls the saturated energy bound.

\paragraph{Editorial reconstruction.} The rational-function statement assumes finite energy and the constant limit at infinity needed for compactification. It describes the self-dual solutions in this model. Related earlier or repeated occurrences: \texttt{1983-601.14}.

\paragraph{Literature and fate.} The identified construction was published before the October 1983 talk; the precise scope is stated separately from the notebook's shorthand. Locator and verification scope: Section 1, the CP1 model and rational instantons. \cite{perelomov1981instantons}.

\paragraph{Present status.} As of 12 November2026: the stated corrected formulation is established. Verification is limited to the identified result and the explicit elementary deductions; no broader literature-completeness claim is made.

\paragraph{Audit.} Literal correctness: correct. Publication match: exact. Verification: read the result. Development: annotated. Confidence: high.

\noindent References: \citep{perelomov1981instantons}.\par

\subsubsection{1983-610.4 --- Krichever's nonlinear analogue of d'Alembert's formula}

\label{item:1983-610.4}

\noindent\textbf{Kind:} method and theorem. \textbf{Fate:} E. \textbf{Pages:} 1135.\par

\paragraph{Notebook content.} Notebook occurrence: P001135, reference to Krichever. Krichever's nonlinear analogue of d'Alembert's formula.

\paragraph{Precise statement.} For the \(1+1\)-dimensional principal chiral-field equation, the local solution can be reconstructed from its data on the two intersecting characteristics by solving two matrix ordinary differential equations in a spectral parameter and a normalized matrix Riemann--Hilbert factorization with singularities at the two spectral points. In the scalar rank-one case the construction reduces to superposition of left- and right-moving waves.

\paragraph{Editorial reconstruction.} Krichever's Theorem 1 and its corollary identify the two input matrix functions with the characteristic data. The reconstruction is asserted where the normalized factorization exists, in particular locally near compatible regular initial data; an unrestricted global factorization for arbitrary data is not supplied by the brief notebook remark.

\paragraph{Literature and fate.} The identified construction was published before the October 1983 talk; the precise scope is stated separately from the notebook's shorthand. Locator and verification scope: Theorem 1 and corollary, pp. 288--289; rank-one comparison p. 290. \cite{krichever1980dalembert}.

\paragraph{Present status.} As of 12 November2026: the stated corrected formulation is established. Verification is limited to the identified result and the explicit elementary deductions; no broader literature-completeness claim is made.

\paragraph{Audit.} Literal correctness: correct. Publication match: exact. Verification: read the result. Development: annotated. Confidence: high.

\noindent References: \citep{krichever1980dalembert}.\par

\subsubsection{1983-610.5 --- The brief reference to solitons}

\label{item:1983-610.5}

\noindent\textbf{Kind:} remark. \textbf{Fate:} F. \textbf{Pages:} 1135.\par

\paragraph{Notebook content.} Notebook occurrence: P001135, soliton remark. The brief reference to solitons.

\paragraph{Precise statement.} The notes mention solitons without specifying a field equation, target, boundary conditions, conserved charge, or existence theorem at this point.

\paragraph{Editorial reconstruction.} The neighboring principal chiral and Skyrme models give possible contexts, but they lead to different questions. No unique theorem or new open problem can be reconstructed from this isolated remark.

\paragraph{Literature and fate.} No unique publication match has been established for this attribution or incomplete remark. Locator and verification scope: Review used only as context, not an exact identification. \cite{perelomov1981instantons}.

\paragraph{Present status.} As of 12 November2026: The precise soliton claim remains unidentified; no assertion of unresolved mathematical status is made.

\paragraph{Audit.} Literal correctness: not determinable. Publication match: thematic only. Verification: unresolved lead. Development: annotated. Confidence: low.

\noindent References: \citep{perelomov1981instantons}.\par

\subsubsection{1983-610.6 --- The quartic Skyrme correction, repeated}

\label{item:1983-610.6}

\noindent\textbf{Kind:} construction. \textbf{Fate:} E. \textbf{Pages:} 1135.\par

\paragraph{Notebook content.} Notebook occurrence: P001135, higher-order Lagrangian term. The quartic Skyrme correction, repeated.

\paragraph{Precise statement.} Writing \(R_\mu=(\partial_\mu g)g^{-1}\), a Skyrme term is a signed multiple of \(\int\operatorname{tr}([R_\mu,R_\nu][R^\mu,R^\nu])\). In the static three-dimensional positive energy, the quadratic and quartic contributions scale oppositely under \(g_\lambda(x)=g(x/\lambda)\): \(E_2(g_\lambda)=\lambda E_2(g)\), \(E_4(g_\lambda)=\lambda^{-1}E_4(g)\).

\paragraph{Editorial reconstruction.} These opposite scalings remove the immediate shrinking argument for the quadratic model. They do not by themselves prove the existence, regularity, or uniqueness of a minimizing soliton. Related earlier or repeated occurrences: \texttt{1983-601.16}, \texttt{1983-601.17}.

\paragraph{Literature and fate.} The identified construction was published before the October 1983 talk; the precise scope is stated separately from the notebook's shorthand. Locator and verification scope: pp. 433--435, Skyrme-model discussion. \cite{witten1983current}.

\paragraph{Present status.} As of 12 November2026: the stated corrected formulation is established. Verification is limited to the identified result and the explicit elementary deductions; no broader literature-completeness claim is made.

\paragraph{Audit.} Literal correctness: correct. Publication match: exact. Verification: read the result. Development: annotated. Confidence: high.

\noindent References: \citep{witten1983current}.\par

\subsubsection{1983-610.7 --- Compactification and extension to a five-dimensional filling}

\label{item:1983-610.7}

\noindent\textbf{Kind:} construction. \textbf{Fate:} E. \textbf{Pages:} 1136.\par

\paragraph{Notebook content.} Notebook occurrence: P001136, S3 times S1 to S3 times D2. Compactification and extension to a five-dimensional filling.

\paragraph{Precise statement.} A field history constant at spatial infinity gives a map \(S^3\times S^1\to SU(3)\). To extend it over \(S^3\times D^2\), the boundary loop in the relevant component of \(\operatorname{Map}(S^3,SU(3))\) must be null-homotopic. For the group target, \(\pi_1(SU(3))=0\) and \(\pi_4(SU(3))=0\) remove the evaluation and based-mapping obstructions.

\paragraph{Editorial reconstruction.} Simply connectedness of \(SU(3)\) alone does not suffice. The missing \(\pi_4\) condition distinguishes \(SU(3)\) from \(SU(2)\), whose \(\pi_4\) is \(\mathbb Z/2\). Related earlier or repeated occurrences: \texttt{1983-601.18}.

\paragraph{Literature and fate.} The identified construction was published before the October 1983 talk; the precise scope is stated separately from the notebook's shorthand. Locator and verification scope: Witten pp. 434--436; Hatcher section 4.J for the low sphere homotopy group. \cite{witten1983current,hatcher2002algebraic}.

\paragraph{Present status.} As of 12 November2026: the stated corrected formulation is established. Verification is limited to the identified result and the explicit elementary deductions; no broader literature-completeness claim is made.

\paragraph{Audit.} Literal correctness: correct with missing hypotheses. Publication match: exact. Verification: read the result. Development: annotated. Confidence: high.

\noindent References: \citep{witten1983current,hatcher2002algebraic}.\par

\subsubsection{1983-610.8 --- The degree-three and degree-five cohomology generators}

\label{item:1983-610.8}

\noindent\textbf{Kind:} theorem. \textbf{Fate:} E. \textbf{Pages:} 1136.\par

\paragraph{Notebook content.} Notebook occurrence: P001136, SU3 generators. The degree-three and degree-five cohomology generators.

\paragraph{Precise statement.} The integral cohomology ring is \(H^*(SU(3);\mathbb Z)\simeq\bigwedge_{\mathbb Z}(x_3,x_5)\), with \(\deg x_j=j\). Invariant closed differential forms represent the corresponding real cohomology classes; an integral normalization fixes their periods.

\paragraph{Editorial reconstruction.} Differential forms represent cohomology classes, not homology classes. A chosen bi-invariant metric permits the invariant representatives to be taken harmonic; uniqueness refers to the harmonic representative of a fixed cohomology class. Related earlier or repeated occurrences: \texttt{1983-601.19}. An integral-cohomology representative is distinct from Witten's physical five-form, which can have half-integral periods after division by2pi on unrestricted cycles.

\paragraph{Literature and fate.} The identified construction was published before the October 1983 talk; the precise scope is stated separately from the notebook's shorthand. Locator and verification scope: Witten current-algebra discussion of the invariant5-form; global paper abstract for normalization. \cite{witten1983current,witten1983global}.

\paragraph{Present status.} As of 12 November2026: the stated corrected formulation is established. Verification is limited to the identified result and the explicit elementary deductions; no broader literature-completeness claim is made.

\paragraph{Audit.} Literal correctness: correct with missing hypotheses. Publication match: exact. Verification: read the result. Development: annotated. Confidence: high.

\noindent References: \citep{witten1983current,witten1983global}.\par

\subsubsection{1983-610.9 --- Quantization of the Wess--Zumino coefficient}

\label{item:1983-610.9}

\noindent\textbf{Kind:} claim. \textbf{Fate:} E. \textbf{Pages:} 1136.\par

\paragraph{Notebook content.} Notebook occurrence: P001136, integral coefficient N. Quantization of the Wess--Zumino coefficient.

\paragraph{Precise statement.} In Witten's normalization, let \(\Gamma(U)=\int_Q\widetilde U^*\omega_5\), with spacetime \(S^3\times S^1\) and filling \(Q=S^3\times D^2\). The allowed differences of fillings have periods in \(2\pi\mathbb Z\), so \(\exp(iN\Gamma(U))\) is well defined for \(N\in\mathbb Z\). The action is defined modulo \(2\pi\), not integer-valued.

\paragraph{Editorial reconstruction.} The allowed cycle class is essential. Witten notes that some closed five-cycles in \(SU(3)\) have odd multiples of \(\pi\) as periods of his form. They do not arise from the stated fillings: classes represented by \(S^3\times S^2\) are spherical in this setting. Thus his physical normalization must not be identified with a primitive integral-cohomology normalization. A form with all periods in \(2\pi\mathbb Z\) gives a more general sufficient construction, but changes this normalization. The rotation phase \((-1)^N\) in 1983-601.24 uses Witten's \(N\).

\paragraph{Literature and fate.} This precise restriction is stated in Witten's1983 paper, p. 434, the technical paragraph immediately before baryon number, and predates the talk. The general multivalued-action comparison is Novikov 1982 section 5. \cite{witten1983current,novikov1982hamiltonian}

\paragraph{Present status.} As of 12 November2026: the stated corrected formulation is established. Verification is limited to the identified result and the explicit elementary deductions; no broader literature-completeness claim is made.

\paragraph{Audit.} Literal correctness: correct with missing hypotheses. Publication match: exact. Verification: read the result. Development: annotated. Confidence: high.

\noindent References: \citep{witten1983current,novikov1982hamiltonian}.\par

\subsubsection{1983-610.10 --- The novelty claim concerns coefficient quantization}

\label{item:1983-610.10}

\noindent\textbf{Kind:} historical claim. \textbf{Fate:} E. \textbf{Pages:} 1136.\par

\paragraph{Notebook content.} Notebook occurrence: P001136, remark that the integral restriction was not known before. The novelty claim concerns coefficient quantization.

\paragraph{Precise statement.} The surrounding argument concerns the quantization of the coefficient of the topological Wess--Zumino functional, not the discovery of \(H^5(SU(3);\mathbb Z)\).

\paragraph{Editorial reconstruction.} A broad assertion that integral periods or the cohomology generator were previously unknown would be false. Witten's1983 abstract explicitly announces a quantization law, while Novikov's1982 paper already discusses single-valued amplitudes and integral periods in a general multivalued-action framework. The narrower historical priority claim is therefore not settled by the notebook wording alone.

\paragraph{Literature and fate.} The identified construction was published before the October 1983 talk; the precise scope is stated separately from the notebook's shorthand. Locator and verification scope: Witten abstract; Novikov section 5, pp. 35--36. \cite{witten1983global,novikov1982hamiltonian}.

\paragraph{Present status.} As of 12 November2026: the stated corrected formulation is established. Verification is limited to the identified result and the explicit elementary deductions; no broader literature-completeness claim is made.

\paragraph{Audit.} Literal correctness: not determinable. Publication match: partial. Verification: read the result. Development: annotated. Confidence: medium.

\noindent References: \citep{witten1983global,novikov1982hamiltonian}.\par

\subsubsection{1983-610.11 --- Novikov's multivalued variational functionals}

\label{item:1983-610.11}

\noindent\textbf{Kind:} method. \textbf{Fate:} E. \textbf{Pages:} 1136.\par

\paragraph{Notebook content.} Notebook occurrence: P001136, final Novikov comparison. Novikov's multivalued variational functionals.

\paragraph{Precise statement.} For a closed \((q+1)\)-form \(\Omega\) on a target \(X\) and a closed oriented \(q\)-manifold \(M\), transgression gives a closed1-form on the mapping space: \(\alpha_\phi(v)=\int_M\phi^*(\iota_v\Omega)\). Its local primitives form a multivalued action; it becomes single-valued on the cover killing its period homomorphism. An exponentiated action is single-valued when the relevant periods lie in \(2\pi\mathbb Z\).

\paragraph{Editorial reconstruction.} The displayed formula is a precise reconstruction of the comparison, and applies to the admissible smooth mapping component with the stated boundary conditions. Integral periods of \(\Omega/(2\pi)\) suffice; necessity only concerns the cycles actually realized by loops in that mapping space, unless further realization hypotheses are imposed.

\paragraph{Literature and fate.} The identified construction was published before the October 1983 talk; the precise scope is stated separately from the notebook's shorthand. Locator and verification scope: Section 5, equations19--22 and lemma on p. 36 of the author PDF. \cite{novikov1982hamiltonian}.

\paragraph{Present status.} As of 12 November2026: the stated corrected formulation is established. Verification is limited to the identified result and the explicit elementary deductions; no broader literature-completeness claim is made.

\paragraph{Audit.} Literal correctness: correct. Publication match: exact. Verification: read the result. Development: annotated. Confidence: high.

\noindent References: \citep{novikov1982hamiltonian}.\par

\subsection{Record 613: The one-dimensional Ising model}

\label{talk:613}

\noindent\textbf{Date:} 1983-11-14. \textbf{Speaker:} I. M. Gelfand. \textbf{Pages:} 1137.\par

Periodic spins, action normalizations and local statistical weights. First physical heading on this date.

\subsubsection{1983-613.1 --- Periodic Ising action}

\label{item:1983-613.1}

\noindent\textbf{Kind:} definition. \textbf{Fate:} E. \textbf{Pages:} 1137.\par

\paragraph{Notebook content.} The spins take values $s_i=\pm1$ and the last site is identified with the first.

\paragraph{Precise statement.} For $N\ge2$, $s_{N+1}=s_1$, coupling $\beta\in\mathbb R$ and site fields $h_i\in\mathbb R$, set $S(s)=-\beta\sum_{i=1}^N s_i s_{i+1}-\sum_i h_i s_i$. The finite partition function is $Z=\sum_{s\in\{-1,1\}^N}e^{-S(s)}$. This convention counts the two oriented bonds separately when $N=2$; alternatively take $N\ge3$.

\paragraph{Editorial reconstruction.} The sign and whether the fields already absorb inverse temperature must be fixed before differentiating the partition function. No thermodynamic or infinite-volume claim is implicit.

\paragraph{Literature and fate.} This is pre-existing finite-volume Ising notation; Ising1925 is an historical anchor, not an exact match to every field convention. Publisher metadata, not the original proof, was checked.

\paragraph{Present status.} Status as of 12 November 2026: the reconstructed result is established; the cited verification has the stated limits.

\paragraph{Assessment.} Established mathematics; useful as a precise critical annotation.

\paragraph{Audit.} Literal correctness: correct with missing hypotheses. Publication match: partial. Verification: publisher metadata checked. Development: annotated. Confidence: high.

\noindent References: \citep{ising1925ferromagnetismus}.\par

\subsubsection{1983-613.2 --- Squared-difference action and additive constant}

\label{item:1983-613.2}

\noindent\textbf{Kind:} identity. \textbf{Fate:} E. \textbf{Pages:} 1137.\par

\paragraph{Notebook content.} The interaction is rewritten as a sum of squared differences.

\paragraph{Precise statement.} Since $s_i^2=1$, $\frac\beta2\sum_i(s_i-s_{i+1})^2=\beta N-\beta\sum_i s_i s_{i+1}$. Thus replacing the interaction action by the squared form adds $\beta N$ and multiplies $Z$ by $e^{-\beta N}$; normalized Gibbs expectations are unchanged.

\paragraph{Editorial reconstruction.} The notebook equality omits the additive constant. It is harmless for normalized correlations but changes the partition function and free energy.

\paragraph{Literature and fate.} The correction follows by expanding each square. It does not require a new theorem or a speculative publication identification.

\paragraph{Present status.} Status as of 12 November 2026: the reconstructed result is established; the cited verification has the stated limits.

\paragraph{Assessment.} Established mathematics; useful as a precise critical annotation.

\paragraph{Audit.} Literal correctness: false as written. Publication match: none found. Verification: read the result. Development: annotated. Confidence: high.

\noindent References: \citep{ising1925ferromagnetismus}.\par

\subsubsection{1983-613.3 --- Local weights and transfer construction}

\label{item:1983-613.3}

\noindent\textbf{Kind:} construction. \textbf{Fate:} E. \textbf{Pages:} 1137.\par

\paragraph{Notebook content.} Local interaction weights are multiplied and summed over spin configurations.

\paragraph{Precise statement.} For a homogeneous dimensionless field $h$, define $T_{ab}=\exp(\beta ab+h(a+b)/2)$ for $a,b\in\{-1,1\}$. Then $Z_N=\operatorname{Tr}(T^N)$ for a periodic chain. More generally, allocating each site field between its adjacent bonds produces $Z_N=\operatorname{Tr}(T_1\cdots T_N)$.

\paragraph{Editorial reconstruction.} This finite matrix formula makes the source local-weight prescription explicit. It requires consistent field allocation; it is not a claim that arbitrary higher-dimensional models reduce to a fixed-size transfer matrix.

\paragraph{Literature and fate.} An elementary, pre-existing transfer construction. The notebook bibliography fragment Phys.Rev.D17,1978 was not identified and is not silently substituted for a verified citation.

\paragraph{Present status.} Status as of 12 November 2026: the reconstructed result is established; the cited verification has the stated limits.

\paragraph{Assessment.} Established mathematics; useful as a precise critical annotation.

\paragraph{Audit.} Literal correctness: correct with missing hypotheses. Publication match: thematic only. Verification: publisher metadata checked. Development: annotated. Confidence: high.

\noindent References: \citep{ising1925ferromagnetismus}.\par

\subsection{Record 612: Gibbs measures and correlations}

\label{talk:612}

\noindent\textbf{Date:} 1983-11-14. \textbf{Speaker:} V. A. Chulaevskii. \textbf{Pages:} 1137.\par

Marginal speaker heading attached to the correlation discussion following Gelfand.

\subsubsection{1983-612.1 --- Finite Gibbs measure}

\label{item:1983-612.1}

\noindent\textbf{Kind:} definition. \textbf{Fate:} E. \textbf{Pages:} 1137.\par

\paragraph{Notebook content.} The probability of a spin configuration is its exponential weight divided by the partition function.

\paragraph{Precise statement.} For a finite configuration set $\Omega$, a real Hamiltonian $H$ and inverse temperature $\beta>0$, define $\mu_{\beta,h}(s)=Z(\beta,h)^{-1}\exp[-\beta H(s)+\beta\sum_i h_i s_i]$. Here $Z$ is the finite positive normalizing sum. No DLR or infinite-volume existence assertion is needed.

\paragraph{Editorial reconstruction.} This field convention differs from613.1 by absorbing $\beta$ into the field. Both are legitimate, but derivative factors must follow the chosen convention.

\paragraph{Literature and fate.} The finite Gibbs construction is pre-existing. The source gives no independent infinite-volume theorem to classify.

\paragraph{Present status.} Status as of 12 November 2026: the reconstructed result is established; the cited verification has the stated limits.

\paragraph{Assessment.} Established mathematics; useful as a precise critical annotation.

\paragraph{Audit.} Literal correctness: correct with missing hypotheses. Publication match: partial. Verification: publisher metadata checked. Development: annotated. Confidence: high.

\noindent References: \citep{ising1925ferromagnetismus}.\par

\subsubsection{1983-612.2 --- Correlations as source derivatives}

\label{item:1983-612.2}

\noindent\textbf{Kind:} identity. \textbf{Fate:} E. \textbf{Pages:} 1137.\par

\paragraph{Notebook content.} The notes connect a two-spin correlation with derivatives of $Z$, using spin labels where source variables are needed.

\paragraph{Precise statement.} With the fields in612.1, $\langle s_i s_j\rangle=\beta^{-2}Z^{-1}\partial_{h_i}\partial_{h_j}Z$ and $\operatorname{Cov}(s_i,s_j)=\beta^{-2}\partial_{h_i}\partial_{h_j}\log Z$. These identities include $i=j$. For dimensionless fields the factors $\beta^{-2}$ are absent.

\paragraph{Editorial reconstruction.} Differentiate the finite sum with respect to external fields, not the discrete variables $s_i$. Distinguish the full and connected correlation functions.

\paragraph{Literature and fate.} The identities follow by direct differentiation. The source formula needs a variable correction and a convention-dependent factor.

\paragraph{Present status.} Status as of 12 November 2026: the reconstructed result is established; the cited verification has the stated limits.

\paragraph{Assessment.} Established mathematics; useful as a precise critical annotation.

\paragraph{Audit.} Literal correctness: false as written. Publication match: none found. Verification: read the result. Development: annotated. Confidence: high.

\noindent References: \citep{ising1925ferromagnetismus}.\par

\subsection{Record 611: Partitions and Dyson ranks}

\label{talk:611}

\noindent\textbf{Date:} 1983-11-14. \textbf{Speaker:} A. V. Zelevinsky. \textbf{Pages:} 1138.\par

Ramanujan congruences, rank equidistribution and direct bijection questions.

\subsubsection{1983-611.1 --- Partitions, generating function and symmetric groups}

\label{item:1983-611.1}

\noindent\textbf{Kind:} definition. \textbf{Fate:} E. \textbf{Pages:} 1138.\par

\paragraph{Notebook content.} Integer partitions enumerate conjugacy classes and irreducible representations of the symmetric group.

\paragraph{Precise statement.} A partition $\lambda\vdash n$ is a weakly decreasing finite sequence of positive integers with sum $n$. Writing $p(n)$ for their number, $\sum_{n\ge0}p(n)q^n=\prod_{m\ge1}(1-q^m)^{-1}$ as a formal power series, and analytically for $|q|<1$. Cycle types identify conjugacy classes of $S_n$ with partitions; over $\mathbb C$, irreducible representations are indexed by the same set.

\paragraph{Editorial reconstruction.} The equality of the two counts alone is not a canonical bijection between conjugacy classes and representations; the partition labels provide a conventional common indexing.

\paragraph{Literature and fate.} Pre-existing representation theory. Earlier occurrences1976-415.1,.3 cover the basic indexing; this is a recurrence, not a new problem.

\paragraph{Present status.} Status as of 12 November 2026: the reconstructed result is established; the cited verification has the stated limits.

\paragraph{Assessment.} Established mathematics; useful as a precise critical annotation.

\paragraph{Audit.} Literal correctness: correct with missing hypotheses. Publication match: partial. Verification: publisher metadata checked. Development: annotated. Confidence: high.

\noindent References: \citep{framerobinsonthrall1954hook}.\par

\subsubsection{1983-611.2 --- Ramanujan congruences modulo five and seven}

\label{item:1983-611.2}

\noindent\textbf{Kind:} result. \textbf{Fate:} E. \textbf{Pages:} 1138.\par

\paragraph{Notebook content.} The partition numbers at $5n+4$ and $7n+5$ are divisible by five and seven.

\paragraph{Precise statement.} For every integer $n\ge0$, $p(5n+4)\equiv0\pmod5$ and $p(7n+5)\equiv0\pmod7$. The asserted congruences concern all partitions of those integers.

\paragraph{Editorial reconstruction.} These are arithmetic statements. A rank equidistribution theorem gives a refinement; it does not itself supply the requested direct combinatorial maps.

\paragraph{Literature and fate.} The congruences and their rank refinements predate1983. Atkin--Swinnerton-Dyer1954 is the verified bibliographic anchor for the refinement; its original full proof was not accessible.

\paragraph{Present status.} Status as of 12 November 2026: the reconstructed result is established; the cited verification has the stated limits.

\paragraph{Assessment.} Established mathematics; useful as a precise critical annotation.

\paragraph{Audit.} Literal correctness: correct with missing hypotheses. Publication match: partial. Verification: publisher metadata checked. Development: annotated. Confidence: high.

\noindent References: \citep{atkinswinnertondyer1954partitions}.\par

\subsubsection{1983-611.3 --- Dyson rank and equidistribution}

\label{item:1983-611.3}

\noindent\textbf{Kind:} result. \textbf{Fate:} E. \textbf{Pages:} 1138.\par

\paragraph{Notebook content.} The rank separates the partition sets into equally numerous residue classes.

\paragraph{Precise statement.} For a nonempty partition define $r(\lambda)=\lambda_1-\ell(\lambda)$. If $N(a,m,n)$ counts partitions of $n$ whose rank is $a\pmod m$, then $N(a,5,5n+4)=p(5n+4)/5$ for each $a\in\mathbb Z/5\mathbb Z$ and $N(a,7,7n+5)=p(7n+5)/7$ for each $a\in\mathbb Z/7\mathbb Z$.

\paragraph{Editorial reconstruction.} The two moduli and two progressions must be kept paired. The empty partition is irrelevant to the stated progressions.

\paragraph{Literature and fate.} Atkin--Swinnerton-Dyer1954, Theorems4--5, is the historical result. The numbering was checked through indexed original text and Stanton2019; full original proof not read.

\paragraph{Present status.} Status as of 12 November 2026: the reconstructed result is established; the cited verification has the stated limits.

\paragraph{Assessment.} Established mathematics; useful as a precise critical annotation.

\paragraph{Audit.} Literal correctness: correct with missing hypotheses. Publication match: exact. Verification: publisher metadata checked. Development: annotated. Confidence: high.

\noindent References: \citep{atkinswinnertondyer1954partitions,stanton2019problems}.\par

\subsubsection{1983-611.4 --- Direct bijections between rank classes}

\label{item:1983-611.4}

\noindent\textbf{Kind:} question. \textbf{Fate:} C. \textbf{Pages:} 1138.\par

\paragraph{Notebook content.} A bijective explanation of the equally sized rank classes is requested.

\paragraph{Precise statement.} For each $n$, seek a directly specified combinatorial cyclic action of order5 on partitions of $5n+4$, respectively order7 on partitions of $7n+5$, whose nontrivial orbits contain exactly one partition of each rank residue. The intended construction must explain the statistic through operations on partitions rather than pairing independently enumerated rank classes after invoking their equal cardinalities.

\paragraph{Editorial reconstruction.} As a bare set-theoretic existence question this is already answered by equinumerosity. The meaningful unresolved request is for an explicit explanatory combinatorial construction; the word natural is not a formal complexity constraint.

\paragraph{Literature and fate.} Stanton2019, Some Problems section3,p4 and OPAC-001, explicitly asks for such rank-class bijections and gives an obstruction to one polynomial approach. No later verified resolution was located in the bounded search. This is C, not a claim that literal finite-set bijections do not exist.

\paragraph{Present status.} As of12November2026: last explicitly opened open-problem statement2019; later universal openness is not independently certified.

\paragraph{Assessment.} A focused combinatorial problem; no new construction is claimed.

\paragraph{Audit.} Literal correctness: correct with missing hypotheses. Publication match: exact. Verification: read the result. Development: annotated. Confidence: high.

\noindent Canonical problem: GP-1983-30.\par

\noindent References: \citep{stanton2019problems,stanton2019rankcrank}.\par

\subsection{Record 614: Asymptotic representation theory and the K-functor}

\label{talk:614}

\noindent\textbf{Date:} 1983-11-21. \textbf{Speaker:} A. M. Vershik. \textbf{Pages:} 1139--1142.\par

Joint work with S. V. Kerov: dimensions, random diagrams, characters and the ordered Grothendieck ring.

\subsubsection{1983-614.1 --- Largest irreducible dimension and its exponential scale}

\label{item:1983-614.1}

\noindent\textbf{Kind:} question. \textbf{Fate:} C. \textbf{Pages:} 1139--1142.\par

\paragraph{Notebook content.} The maximum dimension is compared with $\sqrt{n!}$; a limiting constant and an approximate numerical value are discussed.

\paragraph{Precise statement.} Let $d_n=\max_{\lambda\vdash n}f^\lambda$ and $a_n=-n^{-1/2}\log(d_n/\sqrt{n!})$. The precise asymptotic questions are existence and evaluation of $\lim a_n$, and control of the remainder. Aggarwal--Elboim prove that there are $d>0,C<\infty$ such that $|a_n-d|\le C/\log n$ for all sufficiently large $n$. This does not evaluate $d$.

\paragraph{Editorial reconstruction.} The notes compress a limsup, a numerical estimate and a full asymptotic into one assertion. Different conventions use $2a_n$, so a quoted decimal cannot be transferred without its normalization. The source does not justify an exact value near1.2.

\paragraph{Literature and fate.} Vershik--Kerov1985, section1 Theorem1, gives positive exponential upper and lower bounds and explicitly leaves convergence unsettled. Aggarwal--Elboim, Theorem1,p2, first posted25May2026, proves convergence with the stated error. No journal publication was verified. Evaluation and sharper asymptotics remain outside the checked theorem.

\paragraph{Present status.} 12November2026: existence resolved in the opened2026preprint; explicit constant not evaluated there; no peer-reviewed status verified.

\paragraph{Assessment.} A precise asymptotic programme; convergence is now a verified preprint result, so development must target an explicit constant or a genuinely sharper bound.

\paragraph{Audit.} Literal correctness: false as written. Publication match: partial. Verification: read the result. Development: annotated. Confidence: high.

\noindent Canonical problem: GP-1983-31.\par

\noindent References: \citep{vershikkerov1985dimensions,aggarwalelboim2026maximal}.\par

\subsubsection{1983-614.2 --- The giant-representation conjecture and the computation at81}

\label{item:1983-614.2}

\noindent\textbf{Kind:} remark. \textbf{Fate:} A. \textbf{Pages:} 1139--1142.\par

\paragraph{Notebook content.} A numerical conjecture attributed to McKay is reported as having failed.

\paragraph{Precise statement.} The conjectural lower bound $d_n/\sqrt{n!}>1/n$ is false; the historical counterexample reported by Vershik--Kerov occurs at $n=81$. Their exponential upper bound also implies failure for every sufficiently large $n$, since $e^{-c\sqrt n}<1/n$ eventually for every $c>0$.

\paragraph{Editorial reconstruction.} The published account says the motivating tabulation covered $n\le75$, followed by the counterexample81. This is different from claiming that the original table already extended through81.

\paragraph{Literature and fate.} Vershik--Kerov1985, first page and Theorem1, verifies both the chronology and the eventual exponential obstruction. The seminar is an early announcement relative to that paper, although the finite counterexample itself is not claimed to originate with these speakers.

\paragraph{Present status.} Status as of 12 November 2026: see the dated fate assessment; no stronger literature-completeness claim is made.

\paragraph{Assessment.} Established mathematics; useful as a precise critical annotation.

\paragraph{Audit.} Literal correctness: false as written. Publication match: exact. Verification: read the result. Development: annotated. Confidence: high.

\noindent References: \citep{vershikkerov1985dimensions}.\par

\subsubsection{1983-614.3 --- Hook lengths, dimensions and standard tableaux}

\label{item:1983-614.3}

\noindent\textbf{Kind:} result. \textbf{Fate:} E. \textbf{Pages:} 1139--1142.\par

\paragraph{Notebook content.} Dimensions are calculated by hooks and interpreted as numbers of standard tableaux; two staircase examples occur.

\paragraph{Precise statement.} For $\lambda\vdash n$, with conjugate partition $\lambda^{\prime}$, define $h_{ij}=\lambda_i-j+\lambda_j^{\prime}-i+1$. Then $f^\lambda=n!/\prod_{(i,j)\in\lambda}h_{ij}$, and $f^\lambda$ counts fillings by $1,\ldots,n$ increasing along rows and columns. In particular $f^{(3,2,1)}=16$ and $f^{(5,4,3,2,1)}=292864=13\cdot11\cdot2^{11}$.

\paragraph{Editorial reconstruction.} The row and column contributions both enter each hook. The examples were checked with exact integer arithmetic; no numerical fitting is involved.

\paragraph{Literature and fate.} The hook theorem predates1983; Frame--Robinson--Thrall1954 is the historical anchor, and the formula is reproduced in the opened Vershik--Kerov1985 article. Earlier1976-415.5,.9 are recurrences.

\paragraph{Present status.} Status as of 12 November 2026: the reconstructed result is established; the cited verification has the stated limits.

\paragraph{Assessment.} Established mathematics; useful as a precise critical annotation.

\paragraph{Audit.} Literal correctness: correct with missing hypotheses. Publication match: exact. Verification: read the result. Development: annotated. Confidence: high.

\noindent References: \citep{framerobinsonthrall1954hook,vershikkerov1985dimensions}.\par

\subsubsection{1983-614.4 --- Plancherel probability on Young diagrams}

\label{item:1983-614.4}

\noindent\textbf{Kind:} definition. \textbf{Fate:} E. \textbf{Pages:} 1139--1142.\par

\paragraph{Notebook content.} The square of dimension divided by $n!$ defines a distribution.

\paragraph{Precise statement.} For $\lambda\vdash n$, let $\mathbb P_n(\lambda)=(f^\lambda)^2/n!$. The identity $\sum_\lambda(f^\lambda)^2=n!$ makes this a probability measure. It is the distribution of the shape under the Robinson--Schensted correspondence applied to a uniform random permutation in $S_n$. It is different from the uniform measure on partitions.

\paragraph{Editorial reconstruction.} The regular representation contains each irreducible with multiplicity its dimension, explaining the squared weight. Conflating this measure with equal weights changes the limit curve.

\paragraph{Literature and fate.} Pre-existing Plancherel and Robinson--Schensted theory, already present in1976-415.12. The opened1977 and1985 primary texts use this measure explicitly.

\paragraph{Present status.} Status as of 12 November 2026: the reconstructed result is established; the cited verification has the stated limits.

\paragraph{Assessment.} Established mathematics; useful as a precise critical annotation.

\paragraph{Audit.} Literal correctness: correct with missing hypotheses. Publication match: exact. Verification: read the result. Development: annotated. Confidence: high.

\noindent References: \citep{vershikkerov1977asymptotics,vershikkerov1985dimensions}.\par

\subsubsection{1983-614.5 --- Plancherel limit shape and coordinate normalization}

\label{item:1983-614.5}

\noindent\textbf{Kind:} result. \textbf{Fate:} E. \textbf{Pages:} 1139--1142.\par

\paragraph{Notebook content.} A rotated deterministic curve is displayed for large Young diagrams.

\paragraph{Precise statement.} Use Russian coordinates $u=x-y,v=x+y$, scale both by $n^{-1/2}$, and let $\omega_n$ be the boundary profile. Under $\mathbb P_n$, $\sup_u|\omega_n(u)-\Omega(u)|\to0$ in probability, where $\Omega(u)=\frac2\pi[u\arcsin(u/2)+\sqrt{4-u^2}]$ for $|u|\le2$ and $\Omega(u)=|u|$ outside. In coordinates scaled by another factor $1/2$, the same curve is $\Omega(2u)/2$ on $[-1,1]$.

\paragraph{Editorial reconstruction.} The apparent factor-of-two discrepancy is a coordinate convention, not a different shape. The probability measure and scale are essential hypotheses. The deterministic maximizers also approach the same macroscopic curve in the later1985 theorem.

\paragraph{Literature and fate.} Vershik--Kerov1977 establishes the probabilistic limit; the1985 paper restates and proves the shape theorem and treats maximizing diagrams. This is pre-existing for the Plancherel claim, with the later strengthening identified separately.

\paragraph{Present status.} Status as of 12 November 2026: the reconstructed result is established; the cited verification has the stated limits.

\paragraph{Assessment.} Established mathematics; useful as a precise critical annotation.

\paragraph{Audit.} Literal correctness: correct with missing hypotheses. Publication match: exact. Verification: read the result. Development: annotated. Confidence: high.

\noindent References: \citep{vershikkerov1977asymptotics,vershikkerov1985dimensions}.\par

\subsubsection{1983-614.6 --- Searching numerically for maximizing diagrams}

\label{item:1983-614.6}

\noindent\textbf{Kind:} method. \textbf{Fate:} F. \textbf{Pages:} 1139--1142.\par

\paragraph{Notebook content.} Computer exploration and the shape of diagrams with large dimensions are mentioned.

\paragraph{Precise statement.} For a specified $n$, one can enumerate the finite set of partitions, evaluate every hook product exactly and select the maximum. Such a computation certifies $d_n$ only for that $n$ if enumeration is complete; selected candidate shapes give lower bounds only.

\paragraph{Editorial reconstruction.} The notes do not specify an algorithm, search range, certificate or complexity claim. No efficient-search theorem or new open complexity problem can be recovered from the remark.

\paragraph{Literature and fate.} The1985 primary account confirms numerical exploration but does not identify every notebook computational detail. The unspecified part is F.

\paragraph{Present status.} Status as of 12 November 2026: see the dated fate assessment; no stronger literature-completeness claim is made.

\paragraph{Assessment.} Useful provenance warning; no independent research problem is assigned.

\paragraph{Audit.} Literal correctness: not determinable. Publication match: thematic only. Verification: read the result. Development: annotated. Confidence: medium.

\noindent References: \citep{vershikkerov1985dimensions}.\par

\subsubsection{1983-614.7 --- The partition asymptotic prefactor}

\label{item:1983-614.7}

\noindent\textbf{Kind:} result. \textbf{Fate:} E. \textbf{Pages:} 1139--1142.\par

\paragraph{Notebook content.} An exponential estimate for $p(n)$ is written with an incorrect power of $n$ in the prefactor.

\paragraph{Precise statement.} As $n\to\infty$, $p(n)\sim(4\sqrt3\,n)^{-1}\exp(\pi\sqrt{2n/3})$. In particular the prefactor is proportional to $n^{-1}$, not $n^{-1/2}$.

\paragraph{Editorial reconstruction.} The exponent alone gives a logarithmic asymptotic; retaining a prefactor requires correcting its power.

\paragraph{Literature and fate.} Hardy--Ramanujan1918, with an earlier publisher-recorded abstract2January1917, predates the seminar. Publisher metadata was checked, not the original full proof.

\paragraph{Present status.} Status as of 12 November 2026: the reconstructed result is established; the cited verification has the stated limits.

\paragraph{Assessment.} Established mathematics; useful as a precise critical annotation.

\paragraph{Audit.} Literal correctness: false as written. Publication match: exact. Verification: publisher metadata checked. Development: annotated. Confidence: high.

\noindent References: \citep{hardyramanujan1918asymptotic}.\par

\subsubsection{1983-614.8 --- Uniform partitions and the independent-geometric model}

\label{item:1983-614.8}

\noindent\textbf{Kind:} result. \textbf{Fate:} A. \textbf{Pages:} 1139--1142.\par

\paragraph{Notebook content.} A different limit curve arises when partitions have equal probabilities, and multiplicities are related to geometric variables.

\paragraph{Precise statement.} For a uniform partition of $n$, put $Y_n(x)=n^{-1/2}\#\{i:\lambda_i\ge x\sqrt n\}$. On compact subsets of $(0,\infty)$, $Y_n$ converges in probability to $y(x)=-\frac{\sqrt6}{\pi}\log(1-e^{-\pi x/\sqrt6})$; equivalently $e^{-\pi x/\sqrt6}+e^{-\pi y/\sqrt6}=1$. Under the grand-canonical law with $0<q<1$, multiplicities $m_j$ are independent with $\Pr(m_j=k)=(1-q^j)q^{jk}$; conditioning on $\sum j m_j=n$ makes partitions uniform and removes independence.

\paragraph{Editorial reconstruction.} The limiting curve uses unrotated coordinates scaled by $\sqrt n$. Independence belongs to the unconditioned model, not the fixed-size partition measure.

\paragraph{Literature and fate.} The source announces this distinct uniform-partition shape; the opened Vershik--Kerov1985 final section records it and credits Vershik--Szalay. It is classified A as a near-term documented development, without claiming the1985 paper is the first circulation of every component.

\paragraph{Present status.} Status as of 12 November 2026: see the dated fate assessment; no stronger literature-completeness claim is made.

\paragraph{Assessment.} Established mathematics; useful as a precise critical annotation.

\paragraph{Audit.} Literal correctness: correct with missing hypotheses. Publication match: exact. Verification: read the result. Development: annotated. Confidence: high.

\noindent References: \citep{vershikkerov1985dimensions}.\par

\subsubsection{1983-614.9 --- Random-permutation cycle lengths have a random limit}

\label{item:1983-614.9}

\noindent\textbf{Kind:} remark. \textbf{Fate:} E. \textbf{Pages:} 1139--1142.\par

\paragraph{Notebook content.} The contrast with random permutations is expressed as an absence of a limit.

\paragraph{Precise statement.} If $L_1^{(n)}\ge L_2^{(n)}\ge\cdots$ are the ordered cycle lengths of a uniform random permutation, their normalized ordered sequence has a nondegenerate limiting distribution. Thus there is no deterministic partition-shaped profile of the preceding type, but it is false that normalized cycle lengths have no limiting law.

\paragraph{Editorial reconstruction.} A random limiting distribution must be distinguished from convergence to a deterministic curve. The statement does not concern the Robinson--Schensted shape of the same random permutation.

\paragraph{Literature and fate.} Shepp--Lloyd1966 establishes limiting distributions of ordered cycle lengths; its introductory and limiting-distribution discussion was read. The literal no-limit formulation therefore needs qualification.

\paragraph{Present status.} Status as of 12 November 2026: the reconstructed result is established; the cited verification has the stated limits.

\paragraph{Assessment.} Established mathematics; useful as a precise critical annotation.

\paragraph{Audit.} Literal correctness: false as written. Publication match: exact. Verification: read the result. Development: annotated. Confidence: high.

\noindent References: \citep{shepplloyd1966cycles}.\par

\subsubsection{1983-614.10 --- Set partitions and Bell numbers}

\label{item:1983-614.10}

\noindent\textbf{Kind:} remark. \textbf{Fate:} E. \textbf{Pages:} 1139--1142.\par

\paragraph{Notebook content.} Partitions of a finite labelled set are distinguished from partitions of an integer.

\paragraph{Precise statement.} The Bell number $B_n$ counts decompositions of $\{1,\ldots,n\}$ into unordered nonempty blocks, with exponential generating function $\sum_{n\ge0}B_nz^n/n!=\exp(e^z-1)$. A block-size multiplicity vector $(m_j)$ with $\sum j m_j=n$ has $n!/\prod_j[(j!)^{m_j}m_j!]$ realizations. Consequently a uniform set partition does not induce the uniform distribution on integer partitions.

\paragraph{Editorial reconstruction.} This identifies the third probabilistic object in the comparison without inventing a limit-shape theorem that the notes do not state.

\paragraph{Literature and fate.} The counting identities are elementary pre-existing combinatorics; the exponential formula gives a direct derivation. No exact external publication match is asserted.

\paragraph{Present status.} Status as of 12 November 2026: the reconstructed result is established; the cited verification has the stated limits.

\paragraph{Assessment.} Established mathematics; useful as a precise critical annotation.

\paragraph{Audit.} Literal correctness: correct with missing hypotheses. Publication match: none found. Verification: read the result. Development: annotated. Confidence: high.

\subsubsection{1983-614.11 --- Extremal characters of the infinite symmetric group}

\label{item:1983-614.11}

\noindent\textbf{Kind:} result. \textbf{Fate:} E. \textbf{Pages:} 1139--1142.\par

\paragraph{Notebook content.} Thoma-type parameters describe indecomposable characters of $S_\infty$.

\paragraph{Precise statement.} Let $S_\infty$ be the finite-support permutations of a countable set. Its extreme normalized central positive-definite functions are indexed by decreasing nonnegative sequences $\alpha,\beta$ with $\sum_i(\alpha_i+\beta_i)\le1$. If $m_k(\sigma)$ counts $k$-cycles for $k\ge2$, then $\chi(\sigma)=\prod_{k\ge2}(\sum_i\alpha_i^k+(-1)^{k-1}\sum_i\beta_i^k)^{m_k(\sigma)}$. Fixed points contribute the normalization factor1.

\paragraph{Editorial reconstruction.} Extremality and positive definiteness are essential. Arbitrary central functions are not classified by a single parameter pair; nonextreme characters are mixtures.

\paragraph{Literature and fate.} This is a recurrence of the pre-existing Thoma classification, already indexed at1964-11.1 and1964-13.7. The opened Vershik--Kerov character/K-functor papers provide the relevant seminar context, not a new1983priority claim.

\paragraph{Present status.} Status as of 12 November 2026: the reconstructed result is established; the cited verification has the stated limits.

\paragraph{Assessment.} Established mathematics; useful as a precise critical annotation.

\paragraph{Audit.} Literal correctness: correct with missing hypotheses. Publication match: partial. Verification: read the result. Development: annotated. Confidence: high.

\noindent References: \citep{vershikkerov1982unitary,vershikkerov1983kfunctor}.\par

\subsubsection{1983-614.12 --- Extreme characters of the infinite unitary group}

\label{item:1983-614.12}

\noindent\textbf{Kind:} result. \textbf{Fate:} E. \textbf{Pages:} 1139--1142.\par

\paragraph{Notebook content.} A product formula in positive and negative powers of an eigenvalue is displayed.

\paragraph{Precise statement.} Let $U(\infty)=\bigcup_nU(n)$ and require central positive-definite normalized characters continuous on each $U(n)$. An extreme character has $\chi(U)=\prod_{z\in\operatorname{spec}(U)}\Phi(z)$, where $\Phi(z)=e^{\gamma^+(z-1)+\gamma^-(z^{-1}-1)}\prod_{i\ge1}\frac{1+\beta_i^+(z-1)}{1-\alpha_i^+(z-1)}\frac{1+\beta_i^-(z^{-1}-1)}{1-\alpha_i^-(z^{-1}-1)}$. Here $\alpha_i^\pm,\beta_i^\pm$ are decreasing nonnegative sequences, $\gamma^\pm\ge0$, $\sum_i(\alpha_i^++\alpha_i^-+\beta_i^++\beta_i^-)<\infty$, and $\beta_1^++\beta_1^-\le1$. These parameters give all extreme characters.

\paragraph{Editorial reconstruction.} The unitary group is an inductive union, not a Haar-probability group. The finite-level continuity, normalization and extremality clauses, and the nonredundancy restriction on the leading beta parameters, complete the displayed product.

\paragraph{Literature and fate.} Vershik--Kerov1982, Theorems3--4 and its parameter product, was opened and read. It predates the seminar; the received date19April1982 is recorded separately from publication.

\paragraph{Present status.} Status as of 12 November 2026: the reconstructed result is established; the cited verification has the stated limits.

\paragraph{Assessment.} Established mathematics; useful as a precise critical annotation.

\paragraph{Audit.} Literal correctness: correct with missing hypotheses. Publication match: exact. Verification: read the result. Development: annotated. Confidence: high.

\noindent References: \citep{vershikkerov1982unitary}.\par

\subsubsection{1983-614.13 --- Total positivity behind the character product}

\label{item:1983-614.13}

\noindent\textbf{Kind:} method. \textbf{Fate:} E. \textbf{Pages:} 1139--1142.\par

\paragraph{Notebook content.} The character formula is related to completely positive minors of a Toeplitz matrix.

\paragraph{Precise statement.} For $\Phi(z)=\sum_{k\in\mathbb Z}p_kz^k$ from614.12, the doubly infinite Toeplitz matrix $(p_{i-j})_{i,j\in\mathbb Z}$ is totally nonnegative: every finite minor is nonnegative. Conversely, normalized summable doubly infinite totally nonnegative sequences are described by the corresponding Edrei-type product, with a consistent choice of winding/nonredundancy convention.

\paragraph{Editorial reconstruction.} Total positivity here means minors, not merely $p_k\ge0$. The normalization is $\sum_kp_k=\Phi(1)=1$. The link is a structural ingredient of the character classification, not an unrestricted positivity theorem for arbitrary class functions.

\paragraph{Literature and fate.} The opened1982announcement connects the product and classification; the older total-positivity theorem is pre-existing. No independent proof of its full classification was read in this run.

\paragraph{Present status.} Status as of 12 November 2026: the reconstructed result is established; the cited verification has the stated limits.

\paragraph{Assessment.} Established mathematics; useful as a precise critical annotation.

\paragraph{Audit.} Literal correctness: correct with missing hypotheses. Publication match: partial. Verification: indirect. Development: annotated. Confidence: high.

\noindent References: \citep{vershikkerov1982unitary}.\par

\subsubsection{1983-614.14 --- Weyl alternants and normalized character limits}

\label{item:1983-614.14}

\noindent\textbf{Kind:} formula. \textbf{Fate:} E. \textbf{Pages:} 1139--1142.\par

\paragraph{Notebook content.} A determinant quotient and an asymptotic passage to infinite characters are written in abbreviated form.

\paragraph{Precise statement.} For an integral dominant weight $\lambda_1\ge\cdots\ge\lambda_n$ and pairwise distinct nonzero eigenvalues $z_i$, $\chi_\lambda(z)=\det[z_i^{\lambda_j+n-j}]_{i,j=1}^n/\det[z_i^{n-j}]_{i,j=1}^n$, extending continuously to repeated eigenvalues. Its value at the identity is $\prod_{i<j}(\lambda_i-\lambda_j+j-i)/(j-i)$; normalize by this dimension before taking character limits. If such normalized characters converge to an extreme $U(\infty)$ character, its limit has the product in614.12.

\paragraph{Editorial reconstruction.} The malformed source determinant needs consistent row and column indices and the Weyl shifts. Limits cannot be inferred from a few fixed row ratios alone: a regular asymptotic sequence must control the positive and negative Frobenius data and their total mass. Those hypotheses are not supplied in the notebook.

\paragraph{Literature and fate.} The finite formula is pre-existing; the precise asymptotic framework is in Vershik--Kerov1982. The notebook compression does not justify an unconditional limit theorem, and no stronger one is inserted.

\paragraph{Present status.} Status as of 12 November 2026: the reconstructed result is established; the cited verification has the stated limits.

\paragraph{Assessment.} Established mathematics; useful as a precise critical annotation.

\paragraph{Audit.} Literal correctness: false as written. Publication match: partial. Verification: read the result. Development: annotated. Confidence: high.

\noindent References: \citep{vershikkerov1982unitary}.\par

\subsubsection{1983-614.15 --- The unidentified Miwa--Jimbo operator remark}

\label{item:1983-614.15}

\noindent\textbf{Kind:} remark. \textbf{Fate:} F. \textbf{Pages:} 1139--1142.\par

\paragraph{Notebook content.} An operator associated with Miwa and Jimbo is mentioned without enough defining data.

\paragraph{Precise statement.} The source supplies no unambiguous domain, kernel, spectral parameter or operator identity. It therefore does not determine a mathematical statement whose truth or later solution can be assessed uniquely.

\paragraph{Editorial reconstruction.} Possible links to infinite characters, tau functions and integrable hierarchies are thematic possibilities, not verified identifications. They are not promoted to new problems.

\paragraph{Literature and fate.} The exact publication and operator remain unverified after the targeted searches. Class F records this underdetermination rather than claiming nonpublication or openness.

\paragraph{Present status.} Status as of 12 November 2026: see the dated fate assessment; no stronger literature-completeness claim is made.

\paragraph{Assessment.} Do not develop until an independent source identifies the operator.

\paragraph{Audit.} Literal correctness: not determinable. Publication match: none found. Verification: unresolved lead. Development: annotated. Confidence: low.

\subsubsection{1983-614.16 --- The ordered Grothendieck ring of the infinite symmetric group}

\label{item:1983-614.16}

\noindent\textbf{Kind:} construction. \textbf{Fate:} E. \textbf{Pages:} 1139--1142.\par

\paragraph{Notebook content.} The K-functor is expressed by symmetric functions with the degree-one generator set equal to1.

\paragraph{Precise statement.} For the inductive system $\mathbb CS_n\subset\mathbb CS_{n+1}$, the Grothendieck group is the direct limit of representation groups under the branching inclusion. Via Frobenius characteristics, the connecting map is multiplication by $s_{(1)}$. The induction product yields the ring $K_0(\mathbb CS_\infty)\cong\Lambda_{\mathbb Z}/(s_{(1)}-1)$, where $\Lambda_{\mathbb Z}$ is the integral symmetric-function ring.

\paragraph{Editorial reconstruction.} The ring product is the one induced by external tensor product and induction, not pointwise multiplication of characters at a fixed finite group. The quotient forgets homogeneous degree by the specified connecting maps.

\paragraph{Literature and fate.} Vershik--Kerov1983, section1 equation(9),p132, is an exact near-contemporary match explicitly cited by the notebook as volume123; it was already published in the seminar year.

\paragraph{Present status.} Status as of 12 November 2026: the reconstructed result is established; the cited verification has the stated limits.

\paragraph{Assessment.} Established mathematics; useful as a precise critical annotation.

\paragraph{Audit.} Literal correctness: correct with missing hypotheses. Publication match: exact. Verification: read the result. Development: annotated. Confidence: high.

\noindent References: \citep{vershikkerov1983kfunctor}.\par

\subsubsection{1983-614.17 --- The positive cone in the K-functor}

\label{item:1983-614.17}

\noindent\textbf{Kind:} question. \textbf{Fate:} E. \textbf{Pages:} 1139--1142.\par

\paragraph{Notebook content.} The positivity structure is asked for after the ring identification.

\paragraph{Precise statement.} Under614.16 the positive cone is the union of the images of nonnegative integral combinations of Schur functions at finite levels. Equivalently, after homogenizing a representative $f$ to one degree using $s_{(1)}=1$, its class is positive exactly when $s_{(1)}^mf$ is Schur-positive for some $m\ge0$. This is an exact ordered-direct-limit criterion; mere nonnegativity of arbitrary polynomial coefficients is unrelated.

\paragraph{Editorial reconstruction.} The source question must be read with the positive representation cone, not an arbitrarily chosen order on the quotient ring. The cited paper develops the intrinsic positivity description further.

\paragraph{Literature and fate.} Vershik--Kerov1983, sections6.1 and6.7,pp145--147, was read for this point. The relevant positivity analysis is already in the paper cited by the notes, so this is E, not a newly open1983question.

\paragraph{Present status.} Status as of 12 November 2026: the reconstructed result is established; the cited verification has the stated limits.

\paragraph{Assessment.} Resolved question useful for explaining the ordered ring; no new development is warranted.

\paragraph{Audit.} Literal correctness: correct with missing hypotheses. Publication match: exact. Verification: read the result. Development: annotated. Confidence: high.

\noindent Canonical problem: GP-1983-32.\par

\noindent References: \citep{vershikkerov1983kfunctor}.\par

\subsection{Record 615: Large-dimensional integrals and loop equations}

\label{talk:615}

\noindent\textbf{Date:} 1983-11-28. \textbf{Speaker:} A. A. Migdal; opening remark by I. M. Gelfand. \textbf{Pages:} 1143--1147.\par

Spherical and matrix examples, entropy, factorization, resolvents and loop equations; ends before the Wiegmann heading.

\subsubsection{1983-615.1 --- Opening remark on infinite-dimensional integration}

\label{item:1983-615.1}

\noindent\textbf{Kind:} remark. \textbf{Fate:} F. \textbf{Pages:} 1143--1147.\par

\paragraph{Notebook content.} Gelfand introduces the broad problem of extracting meaning from infinite-dimensional integrals before Migdal begins.

\paragraph{Precise statement.} No space of fields, measure, regularization, integrand or limit is specified in this opening remark. It is a programme for the following examples, not a determined integration theorem.

\paragraph{Editorial reconstruction.} The remark belongs to Gelfand within the physical span of record615. It must not be merged into a purported theorem attributed to Migdal.

\paragraph{Literature and fate.} No unique publication match follows from the general subject. The mathematical content remains underdetermined, class F.

\paragraph{Present status.} Status as of 12 November 2026: see the dated fate assessment; no stronger literature-completeness claim is made.

\paragraph{Assessment.} Preserve as context; no invented open problem.

\paragraph{Audit.} Literal correctness: Programmatic rather than a stated theorem.. Publication match: none. Verification: unresolvedlead. Development: annotated. Confidence: medium.

\subsubsection{1983-615.2 --- Fourier transform of spherical measure}

\label{item:1983-615.2}

\noindent\textbf{Kind:} formula. \textbf{Fate:} E. \textbf{Pages:} 1143--1147.\par

\paragraph{Notebook content.} An integral constrained by a delta function to a sphere is evaluated through a Bessel function.

\paragraph{Precise statement.} For $n\ge2$ and uniform probability measure $\sigma$ on $S^{n-1}$, $\int e^{ik\cdot x}\,d\sigma(x)=\Gamma(n/2)(2/|k|)^{n/2-1}J_{n/2-1}(|k|)$, with value1 at $k=0$ by continuity. For Lebesgue measure with $\delta(|x|^2-1)$, multiply this expression by $\pi^{n/2}/\Gamma(n/2)$.

\paragraph{Editorial reconstruction.} The delta integral is half the surface-area integral because the radial derivative of $r^2-1$ at1 is2. Normalized and unnormalized transforms must not share the same prefactor.

\paragraph{Literature and fate.} This is a classical Bessel integral, checked against the official DLMF integral representations. It predates the seminar. \citep{nist2026bessel}.

\paragraph{Present status.} Status as of 12 November 2026: the reconstructed result is established; the cited verification has the limits stated below.

\paragraph{Assessment.} Established mathematics; useful as a precise critical annotation.

\paragraph{Audit.} Literal correctness: Correct after the stated hypotheses and normalization are supplied.. Publication match: exact. Verification: readtheresult. Development: annotated. Confidence: high.

\noindent References: \citep{nist2026bessel}.\par

\subsubsection{1983-615.3 --- High-dimensional spherical Gaussian scaling}

\label{item:1983-615.3}

\noindent\textbf{Kind:} result. \textbf{Fate:} E. \textbf{Pages:} 1143--1147.\par

\paragraph{Notebook content.} The high-dimensional limit is compared with a Gaussian integral.

\paragraph{Precise statement.} For fixed $m$ and $X^{(n)}$ uniform on $S^{n-1}$, $\sqrt n(X_1^{(n)},\ldots,X_m^{(n)})$ converges in distribution to $N(0,I_m)$. Hence the characteristic function at fixed $\xi\in\mathbb R^m$ converges to $e^{-|\xi|^2/2}$. Without the $\sqrt n$ rescaling, every fixed set of coordinates converges to0 and the transform tends to1.

\paragraph{Editorial reconstruction.} Write $X^{(n)}=G/|G|$ for independent standard normal coordinates and use $|G|^2/n\to1$ in probability. The result concerns fixed m, not an unspecified fraction of all coordinates.

\paragraph{Literature and fate.} The elementary proof reconstructs the normalization. Stam 1982 is an opened pre-seminar institutional bibliographic anchor; its full proof was not read. \citep{stam1982spheres}.

\paragraph{Present status.} Status as of 12 November 2026: the reconstructed result is established; the cited verification has the limits stated below.

\paragraph{Assessment.} Established mathematics; useful as a precise critical annotation.

\paragraph{Audit.} Literal correctness: Correct after the stated hypotheses and normalization are supplied.. Publication match: partial. Verification: publishermetadatachecked. Development: annotated. Confidence: high.

\noindent References: \citep{stam1982spheres}.\par

\subsubsection{1983-615.4 --- Auxiliary multiplier and saddle regimes}

\label{item:1983-615.4}

\noindent\textbf{Kind:} method. \textbf{Fate:} E. \textbf{Pages:} 1143--1147.\par

\paragraph{Notebook content.} The spherical constraint is represented by an auxiliary integral and then treated by a saddle point.

\paragraph{Precise statement.} For the delta integral in615.2, a Bromwich-contour representation gives $I_n(k)=\frac{\pi^{n/2}}{2\pi i}\int_{c-i\infty}^{c+i\infty}s^{-n/2}e^{s-|k|^2/(4s)}\,ds$, interpreted by inverse Laplace continuation with $c>0$. With $s=na$ and $|k|=n\kappa$, the phase per n, apart from constants, is $a-\frac12\log a-\kappa^2/(4a)$. Saddles solve $4a^2-2a+\kappa^2=0$ and coalesce at $\kappa=1/2$.

\paragraph{Editorial reconstruction.} The Gaussian integration denominator contains $4s$, hence $4na$ after scaling. Fixed k, k of order $\sqrt n$, and k of order n are different asymptotic regimes; ordinary isolated-saddle expansion is not uniform at coalescence.

\paragraph{Literature and fate.} The exact integral follows from the Gaussian transform and the Bessel representation. DLMF10.19 confirms distinct large-order regimes; the displayed saddle equation is a direct calculation, not a claim that the notes supply a uniform asymptotic theorem. \citep{nist2026bessel}.

\paragraph{Present status.} Status as of 12 November 2026: the reconstructed result is established; the cited verification has the limits stated below.

\paragraph{Assessment.} Established mathematics; useful as a precise critical annotation.

\paragraph{Audit.} Literal correctness: The scaled Gaussian denominator and regime qualifications require correction.. Publication match: partial. Verification: readtheresult. Development: annotated. Confidence: high.

\noindent References: \citep{nist2026bessel}.\par

\subsubsection{1983-615.5 --- Unitary external-field integral}

\label{item:1983-615.5}

\noindent\textbf{Kind:} construction. \textbf{Fate:} E. \textbf{Pages:} 1143--1147.\par

\paragraph{Notebook content.} Integration over a unitary matrix is converted formally to a matrix multiplier problem.

\paragraph{Precise statement.} For $A\in M_N(\mathbb C)$ and normalized Haar measure, $Z_N(A)=\int_{U(N)}\exp[N\operatorname{Tr}(AU+A^*U^*)]\,dU$ is finite and depends only on the singular values of A. A multiplier representation enforces $UU^*=I$ before integrating over unconstrained complex matrices; its contour and Gaussian convergence must be specified before using a saddle-point approximation.

\paragraph{Editorial reconstruction.} The compact integral is well defined. The multiplier expression is a method, not a probability measure on all complex matrices with an arbitrary-sign Gaussian. A large-N formula additionally needs a limiting source spectrum and a choice of phase.

\paragraph{Literature and fate.} Brezin--Gross 1980 is the direct thematic primary match, already present at 1982-565.7. Author bibliography and publisher metadata/abstract were checked, but its exact coefficient formula was not read; no stronger equation-level identification is claimed. \citep{brezingross1980external}.

\paragraph{Present status.} Status as of 12 November 2026: the reconstructed result is established; the cited verification has the limits stated below.

\paragraph{Assessment.} Established mathematics; useful as a precise critical annotation.

\paragraph{Audit.} Literal correctness: Correct after the stated hypotheses and normalization are supplied.. Publication match: partial. Verification: publishermetadatachecked. Development: annotated. Confidence: high.

\noindent References: \citep{brezingross1980external}.\par

\subsubsection{1983-615.6 --- Eigenvalue repulsion and matrix entropy}

\label{item:1983-615.6}

\noindent\textbf{Kind:} method. \textbf{Fate:} E. \textbf{Pages:} 1143--1147.\par

\paragraph{Notebook content.} The large number of matrix degrees of freedom contributes an entropy term in addition to the potential.

\paragraph{Precise statement.} For Hermitian matrices with weight $e^{-N\beta\operatorname{Tr}V(M)}dM$, diagonalization yields eigenvalue density proportional to $e^{-N\beta\sum_iV(\lambda_i)}\prod_{i<j}(\lambda_i-\lambda_j)^2$. For an empirical probability density rho, the leading energy functional is $\beta\int V\,d\rho-\iint\log|x-y|\,d\rho(x)d\rho(y)$, with the diagonal removed in the finite sum. Its stationary equation on regular support is $\beta V^{\prime}(x)=2\operatorname{PV}\int d\rho(y)/(x-y)$.

\paragraph{Editorial reconstruction.} The logarithmic term comes from the Vandermonde Jacobian. Replacing the matrix by a single scalar saddle and omitting this term loses an order-$N^2$ contribution. The continuum stationary equation presupposes a limiting equilibrium measure.

\paragraph{Literature and fate.} Brezin--Itzykson--Parisi--Zuber 1978, section3 equations6--18, was opened and read as the precise primary source for this construction. It is pre-existing. \citep{brezinitzyksonparisizuber1978planar}.

\paragraph{Present status.} Status as of 12 November 2026: the reconstructed result is established; the cited verification has the limits stated below.

\paragraph{Assessment.} Established mathematics; useful as a precise critical annotation.

\paragraph{Audit.} Literal correctness: Correct after the stated hypotheses and normalization are supplied.. Publication match: exact. Verification: readtheresult. Development: annotated. Confidence: high.

\noindent References: \citep{brezinitzyksonparisizuber1978planar}.\par

\subsubsection{1983-615.7 --- A convergent normalized matrix model}

\label{item:1983-615.7}

\noindent\textbf{Kind:} definition. \textbf{Fate:} E. \textbf{Pages:} 1143--1147.\par

\paragraph{Notebook content.} The notes switch signs in matrix weights while defining a trace observable.

\paragraph{Precise statement.} Fix $\beta>0$ and a real polynomial V of positive even degree with positive leading coefficient. On Hermitian $N\times N$ matrices let $d\mu_N=Z_N^{-1}e^{-N\beta\operatorname{Tr}V(M)}dM$, and define $O_t(M)=N^{-1}\operatorname{Tr}e^{tM}$, $\Psi_N(t)=\mathbb E_N O_t$. Then $Z_N<\infty$, all real-t observables exist, and $\Psi_N(0)=1$.

\paragraph{Editorial reconstruction.} This single consistent convention is an explicit reconstruction. With the opposite exponential sign and the same confining V the integral diverges. Dividing by Z and N is necessary for normalized factorization.

\paragraph{Literature and fate.} The Hermitian one-matrix model and its planar solution are pre-existing in BIPZ 1978. The notebook sign changes are not presented as a different convergent model. \citep{brezinitzyksonparisizuber1978planar}.

\paragraph{Present status.} Status as of 12 November 2026: the reconstructed result is established; the cited verification has the limits stated below.

\paragraph{Assessment.} Established mathematics; useful as a precise critical annotation.

\paragraph{Audit.} Literal correctness: Inconsistent signs and omitted normalizations in the source must be fixed.. Publication match: partial. Verification: readtheresult. Development: annotated. Confidence: high.

\noindent References: \citep{brezinitzyksonparisizuber1978planar}.\par

\subsubsection{1983-615.8 --- Wilson observables and ordering}

\label{item:1983-615.8}

\noindent\textbf{Kind:} construction. \textbf{Fate:} E. \textbf{Pages:} 1143--1147.\par

\paragraph{Notebook content.} Matrix exponentials are compared with an ordered exponential around a loop.

\paragraph{Precise statement.} For a smooth matrix-valued connection A along a parametrized closed curve C, define parallel transport by its first-order matrix ODE and write $U(C)=\mathcal P\exp(i\oint_C A)$ in the Hermitian convention. The normalized Wilson observable is $W(C)=N^{-1}\mathbb E\operatorname{Tr}U(C)$. Path ordering is essential unless the matrices at distinct points commute.

\paragraph{Editorial reconstruction.} The finite-dimensional exponential is a useful model for the loop variable, but the underlying expectation in quantum gauge theory requires a regulator and a defined measure. Gauge invariance follows from conjugation of the closed-loop holonomy.

\paragraph{Literature and fate.} Makeenko--Migdal 1979 and the opened Makeenko 2000 exposition give the loop formulation. Earlier 1979-482.1,.9,.11,.13 are recurrences; the 1983 lecture does not originate Wilson-loop methods. \citep{makeenkomigdal1979exact,makeenko2000large}.

\paragraph{Present status.} Status as of 12 November 2026: the reconstructed result is established; the cited verification has the limits stated below.

\paragraph{Assessment.} Established mathematics; useful as a precise critical annotation.

\paragraph{Audit.} Literal correctness: Correct after the stated hypotheses and normalization are supplied.. Publication match: exact. Verification: readtheresult. Development: annotated. Confidence: high.

\noindent References: \citep{makeenkomigdal1979exact,makeenko2000large}.\par

\subsubsection{1983-615.9 --- Duhamel differentiation of a matrix exponential}

\label{item:1983-615.9}

\noindent\textbf{Kind:} identity. \textbf{Fate:} E. \textbf{Pages:} 1143--1147.\par

\paragraph{Notebook content.} Differentiation of matrix exponentials is used to obtain loop equations.

\paragraph{Precise statement.} For finite matrices M and H and real t, $\left.\frac{d}{d\epsilon}e^{t(M+\epsilon H)}\right|_{\epsilon=0}=\int_0^t e^{sM}He^{(t-s)M}\,ds$. Thus $\partial(e^{tM})_{kl}/\partial M_{ij}=\int_0^t(e^{sM})_{ki}(e^{(t-s)M})_{jl}\,ds$, and $\partial\operatorname{Tr}V(M)/\partial M_{ij}=(V^{\prime}(M))_{ji}$ in the corresponding matrix-coordinate convention.

\paragraph{Editorial reconstruction.} Matrix factors cannot be commuted before taking a trace. The formula follows by differentiating the defining ODE, so it is valid without diagonalizing M or assuming distinct eigenvalues.

\paragraph{Literature and fate.} An established matrix identity. The later loop-equation exposition supplies context, while the finite-dimensional derivation is direct. \citep{makeenko2000large}.

\paragraph{Present status.} Status as of 12 November 2026: the reconstructed result is established; the cited verification has the limits stated below.

\paragraph{Assessment.} Established mathematics; useful as a precise critical annotation.

\paragraph{Audit.} Literal correctness: Correct after the stated hypotheses and normalization are supplied.. Publication match: partial. Verification: readtheresult. Development: annotated. Confidence: high.

\noindent References: \citep{makeenko2000large}.\par

\subsubsection{1983-615.10 --- Exact finite-N Schwinger--Dyson identity}

\label{item:1983-615.10}

\noindent\textbf{Kind:} identity. \textbf{Fate:} E. \textbf{Pages:} 1143--1147.\par

\paragraph{Notebook content.} Integration by parts relates a trace with an insertion of the potential derivative to a product of traces.

\paragraph{Precise statement.} For the convergent model615.7 and $t\ge0$, integration by parts on the real vector space of Hermitian matrices gives $\beta\mathbb E_N[N^{-1}\operatorname{Tr}(V^{\prime}(M)e^{tM})]=\int_0^t\mathbb E_N[O_sO_{t-s}]\,ds$. The confining potential makes the boundary terms vanish; the Hermitian-coordinate divergence of $e^{tM}$ is $\int_0^t\operatorname{Tr}e^{sM}\operatorname{Tr}e^{(t-s)M}\,ds$.

\paragraph{Editorial reconstruction.} This is an exact finite-N identity, before factorization. The use of a negative confining weight fixes the sign on the potential term.

\paragraph{Literature and fate.} The derivation is ordinary finite-dimensional integration by parts and Duhamel differentiation. It is the matrix-model form underlying the opened BIPZ construction, not a novel 1983 result. \citep{brezinitzyksonparisizuber1978planar}.

\paragraph{Present status.} Status as of 12 November 2026: the reconstructed result is established; the cited verification has the limits stated below.

\paragraph{Assessment.} Established mathematics; useful as a precise critical annotation.

\paragraph{Audit.} Literal correctness: Correct after the stated hypotheses and normalization are supplied.. Publication match: partial. Verification: readtheresult. Development: annotated. Confidence: high.

\noindent References: \citep{brezinitzyksonparisizuber1978planar}.\par

\subsubsection{1983-615.11 --- Normalized factorization and the planar equation}

\label{item:1983-615.11}

\noindent\textbf{Kind:} method. \textbf{Fate:} E. \textbf{Pages:} 1143--1147.\par

\paragraph{Notebook content.} The large-N step factorizes a two-trace expectation, but the notes blur it with a trace of a matrix product.

\paragraph{Precise statement.} Assume $\Psi_N\to\Psi$ with the derivatives required by the polynomial $V^{\prime}$ locally uniformly, and $\operatorname{Cov}(O_s,O_{t-s})\to0$ uniformly on each $0\le s\le t\le T$. Then615.10 implies $\beta V^{\prime}(\partial_t)\Psi(t)=\int_0^t\Psi(s)\Psi(t-s)\,ds$, with $\Psi(0)=1$. Factorization concerns $\mathbb E[O_sO_t]$, not $N^{-1}\mathbb E\operatorname{Tr}(e^{sM}e^{tM})=\Psi_N(s+t)$.

\paragraph{Editorial reconstruction.} The limit passage is conditional on the stated convergence; no universal factorization assertion is made for arbitrary scaling or phases. A product of normalized traces and one normalized trace of a product are distinct already for deterministic matrices.

\paragraph{Literature and fate.} Normalized large-N factorization and planar loop closure predate the seminar; Makeenko--Migdal 1979 and Makeenko 2000 explicitly state the distinction and the large-N counting. \citep{makeenkomigdal1979exact,makeenko2000large}.

\paragraph{Present status.} Status as of 12 November 2026: the reconstructed result is established; the cited verification has the limits stated below.

\paragraph{Assessment.} Established mathematics; useful as a precise critical annotation.

\paragraph{Audit.} Literal correctness: False if factorization is applied to one trace of a product.. Publication match: exact. Verification: readtheresult. Development: annotated. Confidence: high.

\noindent References: \citep{makeenkomigdal1979exact,makeenko2000large}.\par

\subsubsection{1983-615.12 --- Resolvent equation, radical sign and physical branch}

\label{item:1983-615.12}

\noindent\textbf{Kind:} construction. \textbf{Fate:} E. \textbf{Pages:} 1143--1147.\par

\paragraph{Notebook content.} The loop equation is transformed into a quadratic relation and solved by a square root.

\paragraph{Precise statement.} For the Hermitian convention, let $W(z)=\int(z-x)^{-1}d\rho(x)$ and $Q(z)=\int[V^{\prime}(z)-V^{\prime}(x)]/(z-x)\,d\rho(x)$. A planar equilibrium solution obeys $W^2-\beta V^{\prime}W+\beta Q=0$. Choose $W=(\beta V^{\prime}-\sqrt{\beta^2(V^{\prime})^2-4\beta Q})/2$ with $W(z)\sim1/z$ at infinity, and impose nonnegative density and the required support/filling conditions. Separately, the literal notebook equation $X^2+\beta V^{\prime}X=K$ has radical $\sqrt{\beta^2(V^{\prime})^2/4+K}$, not a minus K.

\paragraph{Editorial reconstruction.} The quadratic formula alone does not fix the polynomial Q, support or branch. The sign correction to the notebook algebra is distinct from translating its inconsistent measure to the convergent Hermitian convention.

\paragraph{Literature and fate.} BIPZ 1978, section3 equations6--18, gives the eigenvalue and resolvent construction and the branch condition. This is a precise pre-existing source; the algebraic radical correction is checked directly. \citep{brezinitzyksonparisizuber1978planar}.

\paragraph{Present status.} Status as of 12 November 2026: the reconstructed result is established; the cited verification has the limits stated below.

\paragraph{Assessment.} Established mathematics; useful as a precise critical annotation.

\paragraph{Audit.} Literal correctness: Wrong radical sign for the written equation; branch data omitted.. Publication match: exact. Verification: readtheresult. Development: annotated. Confidence: high.

\noindent References: \citep{brezinitzyksonparisizuber1978planar}.\par

\subsubsection{1983-615.13 --- Haar identities and topological expansion}

\label{item:1983-615.13}

\noindent\textbf{Kind:} method. \textbf{Fate:} E. \textbf{Pages:} 1143--1147.\par

\paragraph{Notebook content.} An infinitesimal change of a unitary matrix and a hierarchy in powers of N are proposed.

\paragraph{Precise statement.} For $X^*=-X$, Haar invariance gives $\int_{U(N)}\left.\frac{d}{dt}F(Ue^{tX})\right|_{t=0}dU=0$ for smooth F. In a formal Hermitian matrix perturbation expansion with the standard N-scaled action, a connected orientable vacuum ribbon graph of genus g has weight $N^{2-2g}$; normalized one-trace observables have powers $N^{-2g}$.

\paragraph{Editorial reconstruction.} A general finite change $U\mapsto U(1+\delta)$ does not remain unitary. The graph expansion is formal and does not establish convergence of a genus series or a uniform expansion at critical parameters.

\paragraph{Literature and fate.} BIPZ 1978 section2 supplies the genus organization; Haar integration by parts is elementary. Neither ingredient is a new seminar theorem. \citep{brezinitzyksonparisizuber1978planar}.

\paragraph{Present status.} Status as of 12 November 2026: the reconstructed result is established; the cited verification has the limits stated below.

\paragraph{Assessment.} Established mathematics; useful as a precise critical annotation.

\paragraph{Audit.} Literal correctness: The unitary variation needs a skew-Hermitian tangent direction; asymptotic claims need regime qualifications.. Publication match: exact. Verification: readtheresult. Development: annotated. Confidence: high.

\noindent References: \citep{brezinitzyksonparisizuber1978planar}.\par

\subsubsection{1983-615.14 --- Sources and the connected correction}

\label{item:1983-615.14}

\noindent\textbf{Kind:} construction. \textbf{Fate:} E. \textbf{Pages:} 1143--1147.\par

\paragraph{Notebook content.} A source-dependent functional is introduced to retain corrections beyond factorization.

\paragraph{Precise statement.} Treat $Z_N[J]=\int\exp[-N\beta\operatorname{Tr}V+N^2\int J(s)O_s\,ds]dM$ as a formal source series at J=0, or restrict to sources for which the integral exists. Let $\Psi_J(t)=N^{-2}\delta\log Z_N[J]/\delta J(t)$. Then $\delta\Psi_J(t)/\delta J(s)=N^2\operatorname{Cov}_J(O_t,O_s)$, and the exact source equation is $\beta V^{\prime}(\partial_t)\Psi_J(t)=\int_0^t[\Psi_J(u)\Psi_J(t-u)+N^{-2}\delta\Psi_J(u)/\delta J(t-u)]du+\int sJ(s)\Psi_J(t+s)ds$.

\paragraph{Editorial reconstruction.} The source normalization determines both powers of N. Differentiating $e^{sM}$ produces the factor s. General positive real sources can destroy convergence because an exponential of $e^{sM}$ outruns a polynomial potential; the formal-source interpretation avoids an unjustified analytic assertion.

\paragraph{Literature and fate.} The finite-dimensional formula is a direct integration-by-parts reconstruction. The notebook version has incompatible source normalization, so an exact publication match to that literal formula is not claimed. The opened loop-equation literature supplies the connected hierarchy context. \citep{makeenko2000large,brezinitzyksonparisizuber1978planar}.

\paragraph{Present status.} Status as of 12 November 2026: the reconstructed result is established; the cited verification has the limits stated below.

\paragraph{Assessment.} Established mathematics; useful as a precise critical annotation.

\paragraph{Audit.} Literal correctness: The displayed source coupling and connected term cannot be retained simultaneously without rescaling.. Publication match: partial. Verification: readtheresult. Development: annotated. Confidence: high.

\noindent References: \citep{makeenko2000large,brezinitzyksonparisizuber1978planar}.\par

\subsubsection{1983-615.15 --- From matrix identities to gauge-field loop equations}

\label{item:1983-615.15}

\noindent\textbf{Kind:} method. \textbf{Fate:} E. \textbf{Pages:} 1143--1147.\par

\paragraph{Notebook content.} The same integration-by-parts mechanism is applied formally to gauge fields and loops.

\paragraph{Precise statement.} On a finite lattice with compact gauge group, Haar integration by parts yields exact relations among Wilson-loop expectations. In a regulated continuum formalism, the planar Makeenko--Migdal relation has the form $\partial_\mu^x\delta W(C)/\delta\sigma_{\mu\nu}(x)=\lambda\oint_Cdy_\nu\,\delta^{(d)}(x-y)W(C_{xy})W(C_{yx})$, with lambda fixed by the action and generator normalization. The product is the planar replacement of a two-loop expectation; coincident-point and self-intersection terms require a regulator.

\paragraph{Editorial reconstruction.} This is not a rigorous construction of four-dimensional continuum Yang--Mills measure. An explicit choice, for which lambda is the displayed coefficient, is Hermitian generators with $\operatorname{Tr}T^aT^b=\delta^{ab}$ and formal action $N(4\lambda)^{-1}\int\operatorname{Tr}F_{\mu\nu}F_{\mu\nu}$, with repeated antisymmetric indices summed.

\paragraph{Literature and fate.} Makeenko--Migdal 1979 equations12,14,16--19 and the opened 2000 exposition give the earlier loop equations and their regularization/factorization caveats. The seminar illustrates an existing programme. \citep{makeenkomigdal1979exact,makeenko2000large}.

\paragraph{Present status.} Status as of 12 November 2026: the reconstructed result is established; the cited verification has the limits stated below.

\paragraph{Assessment.} Established mathematics; useful as a precise critical annotation.

\paragraph{Audit.} Literal correctness: Correct after the stated hypotheses and normalization are supplied.. Publication match: exact. Verification: readtheresult. Development: annotated. Confidence: high.

\noindent References: \citep{makeenkomigdal1979exact,makeenko2000large}.\par

\subsection{Record 616: Principal chiral fields and determinants}

\label{talk:616}

\noindent\textbf{Date:} 1983-11-28. \textbf{Speaker:} P. B. Wiegmann. \textbf{Pages:} 1147--1148.\par

Exact-solution announcement, fermionic determinants and multivalued chiral actions.

\subsubsection{1983-616.1 --- Principal chiral field and exact-solution announcement}

\label{item:1983-616.1}

\noindent\textbf{Kind:} result. \textbf{Fate:} E. \textbf{Pages:} 1147--1148.\par

\paragraph{Notebook content.} The principal chiral model is presented as exactly solvable.

\paragraph{Precise statement.} For a smooth field $g:\mathbb R^2\to SU(N)$, the Euclidean classical action is $S[g]=(2\lambda^2)^{-1}\int-\operatorname{Tr}[(g^{-1}\partial_\mu g)^2]d^2x$, with $\lambda^2>0$ and a specified trace normalization. The quantum exact-solution announcement concerns the associated renormalized two-dimensional field theory; it is stronger than merely writing its classical equation of motion.

\paragraph{Editorial reconstruction.} For compact unitary fields, $g^{-1}\partial g$ is anti-Hermitian, so the minus trace supplies the positive metric. The source does not give a regulator or a mathematical construction of the continuum measure.

\paragraph{Literature and fate.} Polyakov--Wiegmann, Theory of nonabelian Goldstone bosons, was published10November 1983, before this28November talk, and its abstract explicitly treatsO(4). The generalSU(N) paper appears28June 1984. Thus the checkedO(4) announcement is E; a generalSU(N) interpretation needs the later paper and is only a partial match here. \citep{polyakovwiegmann1983goldstone,wiegmann1984sun}.

\paragraph{Present status.} Status as of 12 November 2026: the reconstructed result is established; the cited verification has the limits stated below.

\paragraph{Assessment.} Established mathematics; useful as a precise critical annotation.

\paragraph{Audit.} Literal correctness: Correct as a qualified field-theoretic announcement; broader group claims require the later reference.. Publication match: partial. Verification: publishermetadatachecked. Development: annotated. Confidence: high.

\noindent References: \citep{polyakovwiegmann1983goldstone,wiegmann1984sun}.\par

\subsubsection{1983-616.2 --- Higher-dimensional analogy and terminology}

\label{item:1983-616.2}

\noindent\textbf{Kind:} remark. \textbf{Fate:} F. \textbf{Pages:} 1147--1148.\par

\paragraph{Notebook content.} The chiral model is compared with higher-dimensional theories and described through integrability terminology.

\paragraph{Precise statement.} The source does not specify a higher-dimensional Hamiltonian, observables, boundary conditions or a notion of complete integrability. Consequently no general theorem about four-dimensional exact solvability can be extracted.

\paragraph{Editorial reconstruction.} Terms such as exact solution, integrable and equivalent are retained as programme labels unless accompanied by a definite spectral or correlation-function statement.

\paragraph{Literature and fate.} The opened chiral-model papers are two-dimensional and do not by their titles or abstracts verify an unspecified higher-dimensional extension. This remains F. \citep{polyakovwiegmann1983goldstone,wiegmann1984sun}.

\paragraph{Present status.} Status as of 12 November 2026: see the dated fate assessment; no stronger literature-completeness claim is made.

\paragraph{Assessment.} No precise research question is assigned.

\paragraph{Audit.} Literal correctness: Underspecified programme.. Publication match: thematic. Verification: publishermetadatachecked. Development: annotated. Confidence: medium.

\noindent References: \citep{polyakovwiegmann1983goldstone,wiegmann1984sun}.\par

\subsubsection{1983-616.3 --- Chiral potentials and logarithmic fermion determinant}

\label{item:1983-616.3}

\noindent\textbf{Kind:} construction. \textbf{Fate:} A. \textbf{Pages:} 1147--1148.\par

\paragraph{Notebook content.} Chiral gauge potentials are introduced and an additive functional is called a determinant.

\paragraph{Precise statement.} Given a regulated invertible Dirac operator $D(A)$, the fermionic effective action is $W[A]=-\log\det D(A)$, up to the normalization $\log\det D(0)$ and local counterterms. At a finite regulator, $\delta\log\det D=\operatorname{Tr}(D^{-1}\delta D)$. Locally one may parametrize the two chiral connection components by group-valued fields, but their gauge covariance and anomaly determine the resulting chiral action.

\paragraph{Editorial reconstruction.} An additive local action corresponds to the logarithm of a determinant, not the determinant itself. The trace normalization, regularization, determinant phase and treatment of zero modes are omitted in the notes, so the numerical coefficient of the displayed action cannot be certified.

\paragraph{Literature and fate.} The 1983 and 1984 Polyakov--Wiegmann papers are a close historical development; the 1984 multivalued-action abstract confirms the relation to chiral fields and current algebras. Their full coefficient-level derivations were not read, hence A only for the identified programme, not an exact normalization match. \citep{polyakovwiegmann1983goldstone,polyakovwiegmann1984multivalued}.

\paragraph{Present status.} Status as of 12 November 2026: see the dated fate assessment; no stronger literature-completeness claim is made.

\paragraph{Assessment.} Established mathematics; useful as a precise critical annotation.

\paragraph{Audit.} Literal correctness: Determinant must be distinguished from its logarithm; regularization data are missing.. Publication match: partial. Verification: publishermetadatachecked. Development: annotated. Confidence: high.

\noindent References: \citep{polyakovwiegmann1983goldstone,polyakovwiegmann1984multivalued}.\par

\subsubsection{1983-616.4 --- Wess--Zumino extension and quantized ambiguity}

\label{item:1983-616.4}

\noindent\textbf{Kind:} construction. \textbf{Fate:} A. \textbf{Pages:} 1147--1148.\par

\paragraph{Notebook content.} A three-dimensional extension supplies a term whose action can be multivalued.

\paragraph{Precise statement.} For $N\ge2$, choose a closed bi-invariant3-form $\omega_3$ on $SU(N)$ normalized to have periods $2\pi\mathbb Z$. If a smooth $g:\Sigma\to SU(N)$ on a closed oriented surface extends to $\widetilde g:B\to SU(N)$ with $\partial B=\Sigma$, define $\Gamma[g]=\int_B\widetilde g^*\omega_3$ modulo $2\pi$. Two extensions differ by an integral period, so $e^{ik\Gamma[g]}$ is extension independent for integer k. The Euclidean action contains $ik\Gamma$ alongside a positive kinetic term.

\paragraph{Editorial reconstruction.} This normalization isolates the invariant mathematical content. A digit such as8pi in the notebook is not emended in isolation: changing trace and3-form conventions changes the displayed coefficient. The extension existence and period normalization must be stated.

\paragraph{Literature and fate.} Polyakov--Wiegmann 1984 is the verified near-term publication on multivalued actions. Metadata and abstract were read; the standard period argument supplies the reconstruction, without claiming the notebook coefficient was verified in the paper. \citep{polyakovwiegmann1984multivalued}.

\paragraph{Present status.} Status as of 12 November 2026: see the dated fate assessment; no stronger literature-completeness claim is made.

\paragraph{Assessment.} Established mathematics; useful as a precise critical annotation.

\paragraph{Audit.} Literal correctness: Valid after normalization and integer-level conditions are supplied.. Publication match: partial. Verification: publishermetadatachecked. Development: annotated. Confidence: high.

\noindent References: \citep{polyakovwiegmann1984multivalued}.\par

\subsection{Record 617: Fermionization and the Bethe ansatz}

\label{talk:617}

\noindent\textbf{Date:} 1983-12-05. \textbf{Speaker:} P. B. Wiegmann. \textbf{Pages:} 1149--1153.\par

Current interactions, auxiliary fields, determinant identities and the spinful contact gas.

\subsubsection{1983-617.1 --- Agenda: integrability, field theory and solvable spectra}

\label{item:1983-617.1}

\noindent\textbf{Kind:} remark. \textbf{Fate:} F. \textbf{Pages:} 1149--1153.\par

\paragraph{Notebook content.} The continuation announces several themes before the detailed constructions.

\paragraph{Precise statement.} The opening agenda contains no additional specified operator, boundary-value problem or theorem beyond the subsequent items. Its terms organize the lecture and should not be split into invented open problems.

\paragraph{Editorial reconstruction.} The agenda is kept as a distinct source remark while mathematical assertions are assigned to the later entries that actually define them.

\paragraph{Literature and fate.} No exact publication identification is justified by an agenda alone; class F. \citep{wiegmann1984sun}.

\paragraph{Present status.} Status as of 12 November 2026: see the dated fate assessment; no stronger literature-completeness claim is made.

\paragraph{Assessment.} Context only.

\paragraph{Audit.} Literal correctness: An agenda rather than a theorem.. Publication match: thematic. Verification: publishermetadatachecked. Development: annotated. Confidence: medium.

\noindent References: \citep{wiegmann1984sun}.\par

\subsubsection{1983-617.2 --- Quantum integrability of the principal chiral model}

\label{item:1983-617.2}

\noindent\textbf{Kind:} result. \textbf{Fate:} A. \textbf{Pages:} 1149--1153.\par

\paragraph{Notebook content.} The principal chiral model is described as exactly soluble through a fermionic and Bethe-ansatz construction.

\paragraph{Precise statement.} The model has field $g(x)\in SU(N)$, global $SU(N)_L\times SU(N)_R$ symmetry $g\mapsto LgR^{-1}$ and the positive Euclidean action fixed in616.1. The reported quantum solution concerns its two-dimensional spectrum and scattering structure, not merely the classical flat-current identities.

\paragraph{Editorial reconstruction.} Classical conservation laws, a quantum factorized S-matrix and a mathematically constructed quantum field theory are distinct claims. The notes do not supply a regulator or enough formulas to identify every quantum statement coefficient by coefficient.

\paragraph{Literature and fate.} Wiegmann 1984, published28June, explicitly presents a Bethe-ansatz solution for the generalSU(N) principal chiral model. Its title/abstract and institutional bibliographic entry were checked, giving A for the announced programme. The already published10November 1983 O(4) case is separately recorded at616.1. \citep{wiegmann1984sun,polyakovwiegmann1983goldstone}.

\paragraph{Present status.} Status as of 12 November 2026: see the dated fate assessment; no stronger literature-completeness claim is made.

\paragraph{Assessment.} Established mathematics; useful as a precise critical annotation.

\paragraph{Audit.} Literal correctness: Correct after the stated hypotheses and normalization are supplied.. Publication match: partial. Verification: publishermetadatachecked. Development: annotated. Confidence: high.

\noindent References: \citep{wiegmann1984sun,polyakovwiegmann1983goldstone}.\par

\subsubsection{1983-617.3 --- Colour currents and the multi-flavour fermion model}

\label{item:1983-617.3}

\noindent\textbf{Kind:} definition. \textbf{Fate:} E. \textbf{Pages:} 1149--1153.\par

\paragraph{Notebook content.} Fermions carry a colour index and a flavour index, and their currents enter a four-fermion interaction.

\paragraph{Precise statement.} Let $\psi_{a\alpha}$ be two-dimensional Dirac Grassmann fields, with colour $\alpha=1,\ldots,N$ and flavour $a=1,\ldots,N_f$. For specified Lie-algebra generators $T^A$, define $j_\mu^A=\sum_a\bar\psi_a\gamma_\mu T^A\psi_a$. A current-current model has a regulated action consisting of $\sum_a\bar\psi_a\gamma_\mu\partial_\mu\psi_a$ plus a chosen multiple of $\sum_{\mu,A}j_\mu^Aj_\mu^A$. The sign and normalization of that multiple determine the auxiliary-field contour.

\paragraph{Editorial reconstruction.} The flavour count, colour rank and spinor dimension are separate parameters. Merely defining this action does not prove equivalence to the principal chiral field.

\paragraph{Literature and fate.} The chiral-model papers verify a connection to fermionic/Thirring models at the abstract level. The source definition is a standard field-theoretic construction; exact equality of observables requires the next missing data. \citep{polyakovwiegmann1983goldstone,polyakovwiegmann1984multivalued}.

\paragraph{Present status.} Status as of 12 November 2026: the reconstructed result is established; the cited verification has the limits stated below.

\paragraph{Assessment.} Established mathematics; useful as a precise critical annotation.

\paragraph{Audit.} Literal correctness: Correct after the stated hypotheses and normalization are supplied.. Publication match: partial. Verification: publishermetadatachecked. Development: annotated. Confidence: high.

\noindent References: \citep{polyakovwiegmann1983goldstone,polyakovwiegmann1984multivalued}.\par

\subsubsection{1983-617.4 --- The infinite-flavour equivalence claim}

\label{item:1983-617.4}

\noindent\textbf{Kind:} claim. \textbf{Fate:} F. \textbf{Pages:} 1149--1153.\par

\paragraph{Notebook content.} An equivalence with the chiral model is asserted in an infinite-flavour limit.

\paragraph{Precise statement.} A precise equivalence would require a family of regulated fermion measures indexed by $N_f$, a specified scaling of the current coupling and cutoff, a map from fermionic observables to chiral observables, and a topology of convergence of their correlation functions. None of these choices is fully fixed by the notes.

\paragraph{Editorial reconstruction.} The assertion cannot be upgraded to a theorem that every current-current model converges to the principal chiral model. Formal saddle-point relations may motivate it, but they do not specify a unique limit.

\paragraph{Literature and fate.} The 1983 and 1984 papers establish closely related field-theoretic programmes, but their opened abstracts do not verify this exact infinite-flavour statement. Therefore its literal form is F; no claim of a failed or still-open theorem is made. \citep{polyakovwiegmann1983goldstone,wiegmann1984sun}.

\paragraph{Present status.} Status as of 12 November 2026: see the dated fate assessment; no stronger literature-completeness claim is made.

\paragraph{Assessment.} Recover a full primary derivation before posing a modern mathematical problem.

\paragraph{Audit.} Literal correctness: Missing regulator, coupling scaling, observable map and convergence notion.. Publication match: thematic. Verification: publishermetadatachecked. Development: annotated. Confidence: medium.

\noindent References: \citep{polyakovwiegmann1983goldstone,wiegmann1984sun}.\par

\subsubsection{1983-617.5 --- Correct Hubbard--Stratonovich sign}

\label{item:1983-617.5}

\noindent\textbf{Kind:} identity. \textbf{Fate:} E. \textbf{Pages:} 1149--1153.\par

\paragraph{Notebook content.} An auxiliary Gaussian field is introduced with a sign that would make the real Gaussian divergent.

\paragraph{Precise statement.} For $\lambda>0$ and $j\in\mathbb R^m$, $e^{\lambda|j|^2}=(4\pi\lambda)^{-m/2}\int_{\mathbb R^m}\exp[-|A|^2/(4\lambda)+A\cdot j]dA$. The same coefficientwise identity holds for commuting even nilpotent Grassmann bilinears. The version with $+|A|^2/(4\lambda)$ on a real contour diverges; a different sign requires an explicitly rotated contour or imaginary linear coupling.

\paragraph{Editorial reconstruction.} The finite-dimensional Gaussian identity fixes the sign independently of formal path-integral notation. Its functional use requires a regulator and the chosen contour.

\paragraph{Literature and fate.} This is a pre-existing Gaussian transformation. The correction is a direct complete-square calculation and is not evidence for a new publication match. \citep{polyakovwiegmann1983goldstone}.

\paragraph{Present status.} Status as of 12 November 2026: the reconstructed result is established; the cited verification has the limits stated below.

\paragraph{Assessment.} Established mathematics; useful as a precise critical annotation.

\paragraph{Audit.} Literal correctness: Divergent on the stated real contour with the written positive quadratic sign.. Publication match: partial. Verification: publishermetadatachecked. Development: annotated. Confidence: high.

\noindent References: \citep{polyakovwiegmann1983goldstone}.\par

\subsubsection{1983-617.6 --- Flavour multiplicity and the effective determinant action}

\label{item:1983-617.6}

\noindent\textbf{Kind:} identity. \textbf{Fate:} E. \textbf{Pages:} 1149--1153.\par

\paragraph{Notebook content.} Integrating out identical fermion flavours produces a determinant term in the auxiliary-field action.

\paragraph{Precise statement.} At a finite regulator, independent flavours give $\int\prod_a d\bar\psi_a d\psi_a\,e^{-\sum_a\bar\psi_aD(A)\psi_a}=\det D(A)^{N_f}$. Hence an auxiliary weight $e^{-S_{\rm aux}[A]}$ yields $S_{\rm eff}[A]=S_{\rm aux}[A]-N_f\log\det D(A)$, with a chosen logarithm and determinant normalization. The logarithm is a global functional, not a pointwise scalar to be integrated again without a defined density.

\paragraph{Editorial reconstruction.} The sign follows from the convention $e^{-S}$. A local expression for the determinant, if available, requires its own anomaly and counterterm calculation.

\paragraph{Literature and fate.} The finite-regulator identity is elementary Grassmann Gaussian integration. The relevant chiral-model papers are partial historical matches; no unregulated determinant theorem is inferred from their abstracts. \citep{polyakovwiegmann1983goldstone,polyakovwiegmann1984multivalued}.

\paragraph{Present status.} Status as of 12 November 2026: the reconstructed result is established; the cited verification has the limits stated below.

\paragraph{Assessment.} Established mathematics; useful as a precise critical annotation.

\paragraph{Audit.} Literal correctness: The effective-action sign and the determinant/log distinction require correction.. Publication match: partial. Verification: publishermetadatachecked. Development: annotated. Confidence: high.

\noindent References: \citep{polyakovwiegmann1983goldstone,polyakovwiegmann1984multivalued}.\par

\subsubsection{1983-617.7 --- Flat connections and the proposed localization}

\label{item:1983-617.7}

\noindent\textbf{Kind:} claim. \textbf{Fate:} F. \textbf{Pages:} 1149--1153.\par

\paragraph{Notebook content.} The auxiliary integral is said to reduce to pure-gauge fields, and a curvature formula is written.

\paragraph{Precise statement.} For a smooth matrix connection, $F_{\mu\nu}=\partial_\mu A_\nu-\partial_\nu A_\mu+[A_\mu,A_\nu]$. Every $A=g^{-1}dg$ is flat. Conversely a flat connection is locally pure gauge and is globally pure gauge on a trivial bundle over a simply connected domain; on general domains its holonomy need not be trivial. Localization of a functional integral onto this locus is a separate analytic assertion.

\paragraph{Editorial reconstruction.} The two derivative terms in the source curvature are reversed relative to the chosen $A=g^{-1}dg$ convention. Correcting that identity does not prove the claimed localization or remove nontrivial global holonomy.

\paragraph{Literature and fate.} The geometric statements are established and directly checked by Maurer--Cartan. The specific localization limit is not determined or verified by the opened chiral-model abstracts, so the combined source claim is F. \citep{polyakovwiegmann1983goldstone,polyakovwiegmann1984multivalued}.

\paragraph{Present status.} Status as of 12 November 2026: see the dated fate assessment; no stronger literature-completeness claim is made.

\paragraph{Assessment.} Critical correction only until a regulated statement is identified.

\paragraph{Audit.} Literal correctness: Curvature convention is inconsistent; localization and global pure-gauge conclusions lack hypotheses.. Publication match: thematic. Verification: publishermetadatachecked. Development: annotated. Confidence: medium.

\noindent References: \citep{polyakovwiegmann1983goldstone,polyakovwiegmann1984multivalued}.\par

\subsubsection{1983-617.8 --- Gauge covariance of the Dirac operator}

\label{item:1983-617.8}

\noindent\textbf{Kind:} identity. \textbf{Fate:} E. \textbf{Pages:} 1149--1153.\par

\paragraph{Notebook content.} Gauge transformations are used in manipulating the auxiliary determinant.

\paragraph{Precise statement.} For $D_A=\gamma_\mu(\partial_\mu+A_\mu)$ and a smooth invertible group-valued g, $A_\mu^g=g^{-1}A_\mu g+g^{-1}\partial_\mu g$ gives $D_{A^g}=g^{-1}D_Ag$. At a finite regulator preserving this conjugation, determinants agree. A chiral change of fermionic variables is different and can carry an anomalous Jacobian.

\paragraph{Editorial reconstruction.} Naive determinant multiplicativity in an infinite-dimensional regularization is not justified by the finite matrix identity. The nontrivial Wess--Zumino term comes from the chiral/anomaly calculation, not from ordinary vector-gauge conjugation alone.

\paragraph{Literature and fate.} The operator identity is direct and pre-existing. The opened Polyakov--Wiegmann metadata gives the relevant application context without verifying every regulator convention. \citep{polyakovwiegmann1984multivalued}.

\paragraph{Present status.} Status as of 12 November 2026: the reconstructed result is established; the cited verification has the limits stated below.

\paragraph{Assessment.} Established mathematics; useful as a precise critical annotation.

\paragraph{Audit.} Literal correctness: Correct after the stated hypotheses and normalization are supplied.. Publication match: partial. Verification: publishermetadatachecked. Development: annotated. Confidence: high.

\noindent References: \citep{polyakovwiegmann1984multivalued}.\par

\subsubsection{1983-617.9 --- The Abelian determinant and Schwinger normalization}

\label{item:1983-617.9}

\noindent\textbf{Kind:} formula. \textbf{Fate:} E. \textbf{Pages:} 1149--1153.\par

\paragraph{Notebook content.} The Abelian case is attributed to Schwinger and written as a quadratic functional of a scalar potential.

\paragraph{Precise statement.} For one massless Dirac fermion in two Euclidean dimensions with unit charge, use $D_A=\gamma_\mu(\partial_\mu+iA_\mu)$, write $A=d\chi+*d\Phi$, and assume sufficient decay, zero flux, no zero-mode obstruction and a gauge-invariant relative-determinant prescription. In the standard convention, $-\log[\det D_A/\det D_0]=(2\pi)^{-1}\int|\nabla\Phi|^2d^2x$. Charge e replaces the coefficient by $e^2/(2\pi)$.

\paragraph{Editorial reconstruction.} The factors i, the charge, the sign of the logarithm and the scalar decomposition fix the coefficient. This is an explicit conventional reconstruction; the notes do not provide enough normalization data for a literal coefficient match.

\paragraph{Literature and fate.} The verified APS record is Schwinger 1962, Gauge Invariance and Mass.II, received2July and published1December 1962. The notebook 1961 date and its exact coefficient were not verified; full coefficient derivation was not read. \citep{schwinger1962gauge}.

\paragraph{Present status.} Status as of 12 November 2026: the reconstructed result is established; the cited verification has the limits stated below.

\paragraph{Assessment.} Established mathematics; useful as a precise critical annotation.

\paragraph{Audit.} Literal correctness: Historical year and determinant conventions need qualification.. Publication match: partial. Verification: publishermetadatachecked. Development: annotated. Confidence: high.

\noindent References: \citep{schwinger1962gauge}.\par

\subsubsection{1983-617.10 --- The spinful contact-gas spectral problem}

\label{item:1983-617.10}

\noindent\textbf{Kind:} question. \textbf{Fate:} E. \textbf{Pages:} 1149--1153.\par

\paragraph{Notebook content.} The problem of finding eigenfunctions of a one-dimensional many-particle Hamiltonian is posed.

\paragraph{Precise statement.} For fixed N, circle length $\ell>0$ and real c, consider $H_N=-\sum_{j=1}^N\partial_{x_j}^2+2c\sum_{i<j}\delta(x_i-x_j)$ on wavefunctions in $L^2((\mathbb R/\ell\mathbb Z)^N;(\mathbb C^2)^{\otimes N})$ antisymmetric under simultaneous coordinate-spin exchange. Define the operator by its closed semibounded contact-interaction quadratic form, equivalently continuity and the contact derivative jump. Seek its energies and wavefunctions in fixed particle-number and spin sectors.

\paragraph{Editorial reconstruction.} The kinetic energy has a minus sign. Repulsive c>0 and attractive c<0 require different root configurations; on the line rather than a circle the spectral statement includes continuous-spectrum generalized eigenfunctions.

\paragraph{Literature and fate.} Yang 1967 gives the repulsive spinful Bethe construction before this lecture. Equations1--12 and20--27 were read. His footnote1 does not prove full completeness in every sector, so E applies to the established solution method and explicit eigenpair construction, not an unqualified modern completeness theorem. \citep{yang1967manybody,yang1968smatrix}.

\paragraph{Present status.} Status as of 12 November 2026: the reconstructed result is established; the cited verification has the limits stated below.

\paragraph{Assessment.} Established integrable model; clarify domain and completeness rather than manufacture a new open problem.

\paragraph{Audit.} Literal correctness: The written kinetic sign is wrong for the intended positive kinetic energy; boundary conditions must be fixed.. Publication match: exact. Verification: readtheresult. Development: annotated. Confidence: high.

\noindent Canonical problem: GP-1983-33.\par

\noindent References: \citep{yang1967manybody,yang1968smatrix}.\par

\subsubsection{1983-617.11 --- CAR, number and spin operators, and Fock realization}

\label{item:1983-617.11}

\noindent\textbf{Kind:} construction. \textbf{Fate:} E. \textbf{Pages:} 1149--1153.\par

\paragraph{Notebook content.} Creation and annihilation fields realize the spinful many-body space.

\paragraph{Precise statement.} The fermionic fields obey $\{\psi_\sigma(x),\psi_\tau^*(y)\}=\delta_{\sigma\tau}\delta(x-y)$ and all same-type anticommutators vanish. Set $\widehat N=\sum_\sigma\int\psi_\sigma^*\psi_\sigma$ and $S^z=\frac12(N_\uparrow-N_\downarrow)$. For antisymmetric $\Psi$, the map $\Psi\mapsto(N!)^{-1/2}\sum_{\boldsymbol\sigma}\int\Psi(\boldsymbol x,\boldsymbol\sigma)\psi_{\sigma_1}^*(x_1)\cdots\psi_{\sigma_N}^*(x_N)|0\rangle\,d\boldsymbol x$ is the fixed-N Fock realization.

\paragraph{Editorial reconstruction.} The spin projection is not generally the total-spin quantum number. The CAR enforce simultaneous coordinate-spin antisymmetry, not antisymmetry of coordinates alone at fixed ordered spin labels.

\paragraph{Literature and fate.} Standard fermionic second quantization predates the seminar; Yang 1967 supplies the model context. The normalization follows by applying the CAR. \citep{yang1967manybody}.

\paragraph{Present status.} Status as of 12 November 2026: the reconstructed result is established; the cited verification has the limits stated below.

\paragraph{Assessment.} Established mathematics; useful as a precise critical annotation.

\paragraph{Audit.} Literal correctness: Correct after the stated hypotheses and normalization are supplied.. Publication match: partial. Verification: readtheresult. Development: annotated. Confidence: high.

\noindent References: \citep{yang1967manybody}.\par

\subsubsection{1983-617.12 --- Pair potentials and normal ordering}

\label{item:1983-617.12}

\noindent\textbf{Kind:} construction. \textbf{Fate:} E. \textbf{Pages:} 1149--1153.\par

\paragraph{Notebook content.} The pair interaction is rewritten as a quartic field expression.

\paragraph{Precise statement.} For an even real pair potential V for which the operator is defined, $\sum_{i<j}V(x_i-x_j)$ corresponds to $\frac12\sum_{\sigma,\tau}\int dx\,dy\,V(x-y)\psi_\sigma^*(x)\psi_\tau^*(y)\psi_\tau(y)\psi_\sigma(x)$. Thus $V=2c\delta$ gives $c\sum_{\sigma,\tau}\int\psi_\sigma^*\psi_\tau^*\psi_\tau\psi_\sigma$. The kinetic field term is $\sum_\sigma\int\partial_x\psi_\sigma^*\partial_x\psi_\sigma$.

\paragraph{Editorial reconstruction.} The one-half prevents double counting. Coincident density products must be normal ordered or defined through the contact quadratic form; raw density-squared includes an undefined self-contraction.

\paragraph{Literature and fate.} The formula is the direct second-quantization identity for the model in Yang 1967. It is not an additional solution theorem. \citep{yang1967manybody}.

\paragraph{Present status.} Status as of 12 November 2026: the reconstructed result is established; the cited verification has the limits stated below.

\paragraph{Assessment.} Established mathematics; useful as a precise critical annotation.

\paragraph{Audit.} Literal correctness: Correct after the stated hypotheses and normalization are supplied.. Publication match: partial. Verification: readtheresult. Development: annotated. Confidence: high.

\noindent References: \citep{yang1967manybody}.\par

\subsubsection{1983-617.13 --- Coordinate Bethe waves and nondiffractive scattering}

\label{item:1983-617.13}

\noindent\textbf{Kind:} method. \textbf{Fate:} E. \textbf{Pages:} 1149--1153.\par

\paragraph{Notebook content.} In each ordering sector the wavefunction is expanded in plane waves, with coefficients exchanged at contacts.

\paragraph{Precise statement.} For the equal-mass, equal-coupling contact model, in a chamber $x_{Q1}<\cdots<x_{QN}$ write $\Psi_Q=\sum_{P\in S_N}A_Q(P)\exp(i\sum_jk_{Pj}x_{Qj})$, with spin-valued amplitudes. The contact conditions relate amplitudes across adjacent chambers by two-body matrices. Consistency of different three-particle reorderings requires the Yang--Baxter relation; it is not a consequence of energy conservation alone.

\paragraph{Editorial reconstruction.} The finite sum assumes nondiffractive scattering, so the outgoing momentum multiset is the incoming one. Arbitrary one-dimensional pair potentials do not satisfy this assumption.

\paragraph{Literature and fate.} Yang 1967 equations1--12 and Yang 1968 are pre-existing exact constructions. Earlier 1978 factorization entries are related recurrences, while the contact-gas spectral question uses GP-1983-33. \citep{yang1967manybody,yang1968smatrix}.

\paragraph{Present status.} Status as of 12 November 2026: the reconstructed result is established; the cited verification has the limits stated below.

\paragraph{Assessment.} Established mathematics; useful as a precise critical annotation.

\paragraph{Audit.} Literal correctness: Correct after the stated hypotheses and normalization are supplied.. Publication match: exact. Verification: readtheresult. Development: annotated. Confidence: high.

\noindent References: \citep{yang1967manybody,yang1968smatrix}.\par

\subsubsection{1983-617.14 --- Relativistic regularization and the O(N) extension remark}

\label{item:1983-617.14}

\noindent\textbf{Kind:} remark. \textbf{Fate:} F. \textbf{Pages:} 1149--1153.\par

\paragraph{Notebook content.} A relativistic regularization and an O(N)-type extension are mentioned at the end of the discussion.

\paragraph{Precise statement.} The notes do not specify a lattice action, cutoff relation, dispersion scaling, renormalized coupling or the exact O(N) model. No unique continuum-limit theorem or exact spectral identity follows from these phrases.

\paragraph{Editorial reconstruction.} One must distinguish an integrable regulator, its scaling limit and a model sharing a symmetry group. Agreement of broad subject or symmetry is insufficient to identify a publication.

\paragraph{Literature and fate.} The opened chiral-model papers are related but do not verify the unspecified construction as written. Class F, with no invented general regularization problem. \citep{wiegmann1984sun,polyakovwiegmann1983goldstone}.

\paragraph{Present status.} Status as of 12 November 2026: see the dated fate assessment; no stronger literature-completeness claim is made.

\paragraph{Assessment.} Identify the regulator from a fuller source before developing.

\paragraph{Audit.} Literal correctness: Underspecified.. Publication match: thematic. Verification: publishermetadatachecked. Development: annotated. Confidence: low.

\noindent References: \citep{wiegmann1984sun,polyakovwiegmann1983goldstone}.\par

\subsection{Record 618: Compactification; continued integrable models}

\label{talk:618}

\noindent\textbf{Date:} 1983-12-12. \textbf{Speaker:} V. Ginzburg; P. B. Wiegmann. \textbf{Pages:} 1154--1163.\par

V. Ginzburg opens on P001154; the explicit P. Wiegmann continuation on that page owns the remainder. The catalogue expansion V. L. is unconfirmed.

\subsubsection{1983-618.1 --- Circle compactification and metric blocks}

\label{item:1983-618.1}

\noindent\textbf{Kind:} construction. \textbf{Fate:} E. \textbf{Pages:} 1154--1163.\par

\paragraph{Notebook content.} A five-dimensional metric with a compact fifth coordinate is written in blocks, with the off-diagonal block interpreted as an electromagnetic potential.

\paragraph{Precise statement.} Let $\varphi\in\mathbb R/2\pi\mathbb Z$, let $g_{\mu\nu}$ be a four-dimensional metric, let $r>0$ and let $A=A_\mu dx^\mu$. For fields independent of $\varphi$, the Kaluza--Klein ansatz is $ds_5^2=g_{\mu\nu}dx^\mu dx^\nu+r^2(d\varphi+\kappa A_\mu dx^\mu)^2$. Hence $G_{\mu\nu}=g_{\mu\nu}+r^2\kappa^2 A_\mu A_\nu$, $G_{\mu5}=r^2\kappa A_\mu$ and $G_{55}=r^2$. The coordinate change $\varphi\mapsto\varphi-\kappa\lambda(x)$ induces $A\mapsto A+d\lambda$.

\paragraph{Editorial reconstruction.} The displayed block decomposition requires the quadratic $A_\mu A_\nu$ term; omitting it confuses the four-dimensional base metric with the upper block of the five-dimensional metric. A variable radius is an additional scalar field. The sign of the compact contribution depends on the adopted Lorentz-signature convention.

\paragraph{Literature and fate.} This is an established circle-reduction construction. The opened Witten (1981) publisher record is thematically relevant to extra dimensions, but the notebook does not identify a particular Witten theorem or paper and no exact publication match is asserted. \citep{witten1981kaluza}.

\paragraph{Present status.} Status as of 12 November 2026: the reconstructed result is established; the cited verification has the limits stated below.

\paragraph{Assessment.} Established mathematics; useful as a precise critical annotation.

\paragraph{Audit.} Literal correctness: Correct after the stated hypotheses and normalization are supplied.. Publication match: thematic. Verification: publishermetadata. Development: annotated. Confidence: high.

\noindent References: \citep{witten1981kaluza}.\par

\subsubsection{1983-618.2 --- Fourier modes give squared masses}

\label{item:1983-618.2}

\noindent\textbf{Kind:} result. \textbf{Fate:} E. \textbf{Pages:} 1154--1163.\par

\paragraph{Notebook content.} The compact direction gives Fourier modes, electric charges and a mass formula.

\paragraph{Precise statement.} For a scalar field on a product with constant radius $r$, write $\Phi(x,\varphi)=\sum_{n\in\mathbb Z}\phi_n(x)e^{in\varphi}$. If the five-dimensional mass is $m_5$, then $m_n^2=m_5^2+n^2/r^2$; for $m_5=0$ the nonnegative mass is $|n|/r$. With the metric convention of the preceding item the covariant derivative is $D_\mu\phi_n=(\partial_\mu-in\kappa A_\mu)\phi_n$.

\paragraph{Editorial reconstruction.} The angle has period $2\pi$; in the length coordinate $y=r\varphi$ of period $2\pi r$, the Fourier factor is $e^{iny/r}$. The quantity $n^2/r^2$ is a contribution to squared mass, not mass. Charge normalization also depends on the normalization of the four-dimensional gauge kinetic term.

\paragraph{Literature and fate.} The mass tower follows directly by applying the compact Laplacian to the Fourier modes. This calculation is established; the Witten metadata establishes background relevance only. \citep{witten1981kaluza}.

\paragraph{Present status.} Status as of 12 November 2026: the reconstructed result is established; the cited verification has the limits stated below.

\paragraph{Assessment.} Established mathematics; useful as a precise critical annotation.

\paragraph{Audit.} Literal correctness: The transcription conflates mass with squared mass and angle with length unless the coordinate convention is repaired.. Publication match: thematic. Verification: publishermetadata. Development: annotated. Confidence: high.

\noindent References: \citep{witten1981kaluza}.\par

\subsubsection{1983-618.3 --- Numerical mass and charge estimates}

\label{item:1983-618.3}

\noindent\textbf{Kind:} remark. \textbf{Fate:} F. \textbf{Pages:} 1154--1163.\par

\paragraph{Notebook content.} Numerical scales are associated with the compactification radius and electric charge.

\paragraph{Precise statement.} Under the preceding metric ansatz, the mathematically determined relations are $m_n^2-m_5^2=n^2/r^2$ and a charge proportional to $n\kappa$. Numerical masses or a physical relation between $r$, Newton constant and the electric charge require specified units and the normalization of the gravitational and gauge actions.

\paragraph{Editorial reconstruction.} The notebook fragment does not fix all these conventions. Its numerical estimate therefore cannot be reconstructed as a unique physical prediction; no missing number is supplied by guesswork.

\paragraph{Literature and fate.} No exact numerical statement has been identified in an opened primary source. The correct disposition is F for the incomplete estimate, alongside the established structural formula in item 2. \citep{witten1981kaluza}.

\paragraph{Present status.} Status as of 12 November 2026: see the dated fate assessment; no stronger literature-completeness claim is made.

\paragraph{Assessment.} No independent research problem is extracted from an unspecified numerical normalization.

\paragraph{Audit.} Literal correctness: Insufficient normalization data to assess the numerical claim.. Publication match: thematic. Verification: publishermetadata. Development: annotated. Confidence: low.

\noindent References: \citep{witten1981kaluza}.\par

\subsubsection{1983-618.4 --- The positive-energy theorem heading}

\label{item:1983-618.4}

\noindent\textbf{Kind:} remark. \textbf{Fate:} F. \textbf{Pages:} 1154--1163.\par

\paragraph{Notebook content.} Positive energy is mentioned in the introductory discussion associated with Witten.

\paragraph{Precise statement.} A possible intended statement is nonnegativity of ADM energy for complete asymptotically flat initial data satisfying the dominant energy condition, with the dimension, falloff, regularity, spin and boundary hypotheses needed for the chosen theorem. The notebook does not provide enough detail to select a precise formulation.

\paragraph{Editorial reconstruction.} The positivity statement must not be read as positivity for every metric in the preceding five-dimensional ansatz. A correct theorem requires a definite asymptotic problem and its full hypotheses.

\paragraph{Literature and fate.} Witten (1981) provides a primary positive-energy paper, but only its publisher metadata was checked here. This is a thematic identification, not a verified exact theorem match; the notebook heading remains F. \citep{witten1981positive}.

\paragraph{Present status.} Status as of 12 November 2026: see the dated fate assessment; no stronger literature-completeness claim is made.

\paragraph{Assessment.} Useful bibliographic orientation; no well-posed new problem follows from the heading alone.

\paragraph{Audit.} Literal correctness: A heading without enough hypotheses to constitute a theorem.. Publication match: thematic. Verification: publishermetadata. Development: annotated. Confidence: low.

\noindent References: \citep{witten1981positive}.\par

\subsubsection{1983-618.5 --- The continued contact gas: attraction and repulsion}

\label{item:1983-618.5}

\noindent\textbf{Kind:} definition. \textbf{Fate:} E. \textbf{Pages:} 1154--1163.\par

\paragraph{Notebook content.} Wiegmann resumes the one-dimensional many-particle problem and distinguishes the signs of the contact interaction.

\paragraph{Precise statement.} For fixed $N$ and circle length $\ell>0$, let $H_N=-\sum_{j=1}^N\partial_{x_j}^2+2c\sum_{i<j}\delta(x_i-x_j)$ act on square-integrable spin-$1/2$ wavefunctions antisymmetric under simultaneous exchange of a coordinate and its spin. The contact quadratic form defines a self-adjoint semibounded operator for every fixed real $c$. The cases $c>0$ and $c<0$ are respectively repulsive and attractive.

\paragraph{Editorial reconstruction.} Equal masses and the common coupling are part of this integrable model. Attractive bound clusters lead to complex rapidities, so a list restricted to real repulsive rapidities cannot describe the entire attractive spectrum.

\paragraph{Literature and fate.} Yang (1967), equations 1--12 and 20--27, supplies the precise repulsive spinful setting; Yang (1968) discusses factorized scattering. Both predate the seminar. The present annotation keeps the attractive extension separate from what was checked in the original repulsive construction. \citep{yang1967manybody,yang1968smatrix}.

\paragraph{Present status.} Status as of 12 November 2026: the reconstructed result is established; the cited verification has the limits stated below.

\paragraph{Assessment.} Established mathematics; useful as a precise critical annotation.

\paragraph{Audit.} Literal correctness: Correct after the stated hypotheses and normalization are supplied.. Publication match: exact. Verification: readtheresult. Development: annotated. Confidence: high.

\noindent References: \citep{yang1967manybody,yang1968smatrix}.\par

\subsubsection{1983-618.6 --- Classical NLS and the Bethe chronology}

\label{item:1983-618.6}

\noindent\textbf{Kind:} remark. \textbf{Fate:} E. \textbf{Pages:} 1154--1163.\par

\paragraph{Notebook content.} The notes compare a nonlinear-wave solution with a result obtained approximately forty years earlier by Bethe.

\paragraph{Precise statement.} Bethe (1931) treated a quantum spin chain. Zakharov and Shabat treated the classical focusing nonlinear Schrödinger equation, in a convention such as $iu_t+u_{xx}+\kappa|u|^2u=0$, $\kappa>0$, by the inverse scattering method; the Russian paper appeared in 1971 and its English translation in 1972. Yang (1967) concerns the continuum spinful contact gas, and the 1978 quantum inverse-scattering work is a further distinct development.

\paragraph{Editorial reconstruction.} The Russian phrase \ru{на 40 лет раньше Бете} means that Bethe did it about forty years earlier, not that the later result preceded Bethe by forty years. The analogy concerns methods and integrability, not identity of the spin-chain, continuum-gas and classical PDE problems.

\paragraph{Literature and fate.} The Bethe publisher page and the original Zakharov--Shabat translation were opened. The latter records submission on 22 December 1970, Russian publication in July 1971 and English publication in January 1972; its introductory setup and equations 3--20 were read. The Bethe page has inconsistent month metadata, so only the year 1931 is asserted. The chronology and translation correction are secure. \citep{bethe1931metalle,zakharovshabat1972nls,yang1967manybody,sklyaninfaddeev1978quantum}.

\paragraph{Present status.} Status as of 12 November 2026: the reconstructed result is established; the cited verification has the limits stated below.

\paragraph{Assessment.} Established mathematics; useful as a precise critical annotation.

\paragraph{Audit.} Literal correctness: The English direction of the forty-year comparison is reversed; the underlying models must be distinguished.. Publication match: exact. Verification: readtheresult. Development: annotated. Confidence: high.

\noindent References: \citep{bethe1931metalle,zakharovshabat1972nls,yang1967manybody,sklyaninfaddeev1978quantum}.\par

\subsubsection{1983-618.7 --- Particle number, spin projection and total spin}

\label{item:1983-618.7}

\noindent\textbf{Kind:} definition. \textbf{Fate:} E. \textbf{Pages:} 1154--1163.\par

\paragraph{Notebook content.} The particle number and a number M of particles of one spin orientation determine the displayed spin.

\paragraph{Precise statement.} If $M$ is the number of down spins among $N$ spin-$1/2$ particles, then $S_z=(N-2M)/2$. The total-spin quantum number satisfies $S\geq |S_z|$ and is not determined by $M$ alone. In a highest-weight state one has $S=S_z=N/2-M$, with $0\leq M\leq\lfloor N/2\rfloor$.

\paragraph{Editorial reconstruction.} The formula fixes the convention: M counts down spins. Calling it the number of up spins changes the sign of the projection. Identifying the projection with total spin additionally imposes a highest-weight condition.

\paragraph{Literature and fate.} These are standard SU(2) representation facts and the convention used in the nested Bethe construction. Yang (1967) and the opened algebraic Bethe-ansatz treatment support the model-specific application. \citep{yang1967manybody,faddeev1996bethe}.

\paragraph{Present status.} Status as of 12 November 2026: the reconstructed result is established; the cited verification has the limits stated below.

\paragraph{Assessment.} Established mathematics; useful as a precise critical annotation.

\paragraph{Audit.} Literal correctness: The spin-count convention and the distinction between total spin and spin projection require correction.. Publication match: exact. Verification: readtheresult. Development: annotated. Confidence: high.

\noindent References: \citep{yang1967manybody,faddeev1996bethe}.\par

\subsubsection{1983-618.8 --- Spin Young diagrams and conjugate coordinate symmetry}

\label{item:1983-618.8}

\noindent\textbf{Kind:} construction. \textbf{Fate:} E. \textbf{Pages:} 1154--1163.\par

\paragraph{Notebook content.} Young diagrams are used to describe the symmetry of spin and coordinate factors.

\paragraph{Precise statement.} In $(\mathbb C^2)^{\otimes N}$, total spin $S=N/2-r$ corresponds to the symmetric-group shape $(N-r,r)$, where $0\leq r\leq\lfloor N/2\rfloor$. To obtain a fermionic state under simultaneous coordinate-spin exchange, the coordinate symmetry is the conjugate shape $(2^r,1^{N-2r})$. A Young symmetrizer symmetrizes along rows and antisymmetrizes along columns.

\paragraph{Editorial reconstruction.} The integer $r$ equals the down-spin count only for the highest-weight vector of that irreducible spin multiplet. The lower row is not itself the antisymmetric direction; antisymmetrization is along columns. Tensor multiplicities and a basis choice are separate from specifying the shape.

\paragraph{Literature and fate.} This is the established Schur--Weyl description underlying the spin-sector decomposition. The exact diagram conventions are reconstructed directly; the hook-formula reference is supporting symmetric-group background rather than a claimed exact source for the entire physical assertion. \citep{yang1967manybody,framerobinsonthrall1954hook}.

\paragraph{Present status.} Status as of 12 November 2026: the reconstructed result is established; the cited verification has the limits stated below.

\paragraph{Assessment.} Established mathematics; useful as a precise critical annotation.

\paragraph{Audit.} Literal correctness: Correct after replacing row antisymmetry by column antisymmetry and restricting the identification r=M to highest weight.. Publication match: partial. Verification: readtheresult. Development: annotated. Confidence: high.

\noindent References: \citep{yang1967manybody,framerobinsonthrall1954hook}.\par

\subsubsection{1983-618.9 --- Antisymmetry in the Fock-to-coordinate map}

\label{item:1983-618.9}

\noindent\textbf{Kind:} construction. \textbf{Fate:} E. \textbf{Pages:} 1154--1163.\par

\paragraph{Notebook content.} The field-operator state is identified with a many-particle coordinate wavefunction.

\paragraph{Precise statement.} For fermionic creation fields, the coefficient wavefunction contributing to a fixed-N Fock vector is antisymmetric under simultaneous exchange of $(x_i,\sigma_i)$ and $(x_j,\sigma_j)$. With the standard $1/\sqrt{N!}$ convention, the antisymmetric coordinate wavefunction and its Fock vector have the same norm. An arbitrary coefficient array contributes only its antisymmetrization.

\paragraph{Editorial reconstruction.} A fixed spin component need not be totally antisymmetric in coordinates by itself; its coordinate permutation symmetry is coupled to the spin symmetry described in item 8. This is the continuation of the construction in record 617.

\paragraph{Literature and fate.} The statement follows from the canonical anticommutation relations. Yang (1967) supplies the appropriate spinful wavefunction context; no new result or independent open problem is implied. \citep{yang1967manybody}.

\paragraph{Present status.} Status as of 12 November 2026: the reconstructed result is established; the cited verification has the limits stated below.

\paragraph{Assessment.} Established mathematics; useful as a precise critical annotation.

\paragraph{Audit.} Literal correctness: Correct after the stated hypotheses and normalization are supplied.. Publication match: partial. Verification: readtheresult. Development: annotated. Confidence: high.

\noindent References: \citep{yang1967manybody}.\par

\subsubsection{1983-618.10 --- Contact matching and Bethe wave amplitudes}

\label{item:1983-618.10}

\noindent\textbf{Kind:} method. \textbf{Fate:} E. \textbf{Pages:} 1154--1163.\par

\paragraph{Notebook content.} The contact potential relates coefficients of plane waves in adjacent coordinate-ordering regions.

\paragraph{Precise statement.} For $H_N=-\sum_j\partial_{x_j}^2+2c\sum_{i<j}\delta(x_i-x_j)$, the wavefunction is continuous across $x_i=x_j$ and satisfies $[(\partial_{x_j}-\partial_{x_i})\Psi]_{(x_j-x_i)=0^-}^{0^+}=2c\,\Psi|_{x_i=x_j}$. In each collision-free chamber, Bethe wavefunctions are sums of plane waves whose coefficients are linked by these pairwise contact conditions. Consistency at triple coincidences is encoded by the Yang--Baxter relation.

\paragraph{Editorial reconstruction.} With $r=x_j-x_i$ and $R=(x_i+x_j)/2$, the relative kinetic term is $-2\partial_r^2$. Integrating $-2\partial_r^2+2c\delta(r)$ gives $[\partial_r\Psi]_{0^-}^{0^+}=c\Psi(0)$, hence the displayed derivative jump. Changing the Hamiltonian coupling convention changes every subsequent scattering and Bethe denominator.

\paragraph{Literature and fate.} Yang (1967) explicitly derives the repulsive contact-gas construction before the seminar. The recurrence belongs to GP-1983-33. The checked construction produces eigenpairs subject to admissibility; it does not by itself prove every possible completeness statement. \citep{yang1967manybody}.

\paragraph{Present status.} Status as of 12 November 2026: the reconstructed result is established; the cited verification has the limits stated below.

\paragraph{Assessment.} Established method; the useful work is correction of normalization and separation of construction from completeness.

\paragraph{Audit.} Literal correctness: Correct after the stated hypotheses and normalization are supplied.. Publication match: exact. Verification: readtheresult. Development: annotated. Confidence: high.

\noindent Canonical problem: GP-1983-33.\par

\noindent References: \citep{yang1967manybody}.\par

\subsubsection{1983-618.11 --- Conserved power sums and absence of diffraction}

\label{item:1983-618.11}

\noindent\textbf{Kind:} method. \textbf{Fate:} E. \textbf{Pages:} 1154--1163.\par

\paragraph{Notebook content.} Many conservation laws are invoked to explain elastic, nondiffractive scattering.

\paragraph{Precise statement.} Assume scattering preserves the particle number N and additive asymptotic charges $q_r=\sum_{j=1}^N k_j^r$ for $1\leq r\leq N$. Newton identities then determine the elementary symmetric functions of the momenta, and hence the polynomial $\prod_{j=1}^N(z-k_j)$. The outgoing momentum multiset therefore equals the incoming multiset.

\paragraph{Editorial reconstruction.} This is a conditional argument: the independent charges, their operator domains and their asymptotic additivity must first be established for the model. Conservation of only total momentum and energy does not determine the momentum multiset when $N\geq3$.

\paragraph{Literature and fate.} The algebraic implication is elementary and exact. Yang (1968) and the quantum inverse-scattering framework supply relevant integrable-model context, but a generic assertion that one-dimensional motion alone supplies all these charges would be false. \citep{yang1968smatrix,sklyaninfaddeev1978quantum}.

\paragraph{Present status.} Status as of 12 November 2026: the reconstructed result is established; the cited verification has the limits stated below.

\paragraph{Assessment.} Established mathematics; useful as a precise critical annotation.

\paragraph{Audit.} Literal correctness: Correct after the stated hypotheses and normalization are supplied.. Publication match: partial. Verification: readtheresult. Development: annotated. Confidence: high.

\noindent References: \citep{yang1968smatrix,sklyaninfaddeev1978quantum}.\par

\subsubsection{1983-618.12 --- Factorized many-body scattering}

\label{item:1983-618.12}

\noindent\textbf{Kind:} construction. \textbf{Fate:} E. \textbf{Pages:} 1154--1163.\par

\paragraph{Notebook content.} The many-body scattering operator is assembled from pairwise exchanges.

\paragraph{Precise statement.} For nondiffractive scattering with fixed rapidities and an internal space V per particle, assign an adjacent two-body scattering operator to each exchange. Products associated with reduced exchange words agree if the two-body operators obey the Yang--Baxter equation and operators on disjoint tensor factors commute. Inverse exchanges cancel when $R_{12}(u)R_{21}(-u)=I$.

\paragraph{Editorial reconstruction.} An overall scalar phase and the distinction between $R$ and $\check R=PR$ must be fixed consistently. Factorization is an integrability property of the specified interaction; it is not a consequence of spatial dimension alone.

\paragraph{Literature and fate.} Yang (1967, 1968) gives the original contact-gas consistency and S-matrix formulations. The product construction and its consistency are established before the lecture. \citep{yang1967manybody,yang1968smatrix}.

\paragraph{Present status.} Status as of 12 November 2026: the reconstructed result is established; the cited verification has the limits stated below.

\paragraph{Assessment.} Established mathematics; useful as a precise critical annotation.

\paragraph{Audit.} Literal correctness: Correct after the stated hypotheses and normalization are supplied.. Publication match: exact. Verification: readtheresult. Development: annotated. Confidence: high.

\noindent References: \citep{yang1967manybody,yang1968smatrix}.\par

\subsubsection{1983-618.13 --- The Yang--Baxter consistency equation}

\label{item:1983-618.13}

\noindent\textbf{Kind:} result. \textbf{Fate:} E. \textbf{Pages:} 1154--1163.\par

\paragraph{Notebook content.} Two reorderings of three-particle scattering are set equal.

\paragraph{Precise statement.} For a meromorphic family $R(u)\in\operatorname{End}(V\otimes V)$, the difference-form Yang--Baxter equation is $R_{12}(u-v)R_{13}(u-w)R_{23}(v-w)=R_{23}(v-w)R_{13}(u-w)R_{12}(u-v)$ in $\operatorname{End}(V^{\otimes3})$. It is an identity wherever the factors are defined and hence an identity of meromorphic functions.

\paragraph{Editorial reconstruction.} The subscripts identify the tensor factors acted upon. Passing to $\check R=PR$ produces a braid-form equation with correspondingly changed positions and spectral arguments; one cannot replace every R by PR while leaving the displayed equation otherwise untouched.

\paragraph{Literature and fate.} Yang (1967) supplies the consistency relation; the opened Faddeev account gives the algebraic framework. The equation is established, with the convention made explicit here. \citep{yang1967manybody,sklyaninfaddeev1978quantum,faddeev1996bethe}.

\paragraph{Present status.} Status as of 12 November 2026: the reconstructed result is established; the cited verification has the limits stated below.

\paragraph{Assessment.} Established mathematics; useful as a precise critical annotation.

\paragraph{Audit.} Literal correctness: Correct after the stated hypotheses and normalization are supplied.. Publication match: exact. Verification: readtheresult. Development: annotated. Confidence: high.

\noindent References: \citep{yang1967manybody,sklyaninfaddeev1978quantum,faddeev1996bethe}.\par

\subsubsection{1983-618.14 --- Difference parameters and the limits of a classification claim}

\label{item:1983-618.14}

\noindent\textbf{Kind:} remark. \textbf{Fate:} F. \textbf{Pages:} 1154--1163.\par

\paragraph{Notebook content.} Difference dependence, the eight-vertex form and a classification of solutions are discussed.

\paragraph{Precise statement.} A familiar two-state ansatz is $R(u)=\begin{pmatrix}a(u)&0&0&d(u)\\0&b(u)&c(u)&0\\0&c(u)&b(u)&0\\d(u)&0&0&a(u)\end{pmatrix}$. A classification theorem must specify the symmetry ansatz, regularity and nondegeneracy conditions, equivalences, and whether the spectral dependence is of difference form.

\paragraph{Editorial reconstruction.} The existence of rational, trigonometric or elliptic families does not establish exhaustiveness for arbitrary $4$ by $4$ solutions. The notebook fragment omits the restrictions required to identify a specific classification theorem.

\paragraph{Literature and fate.} Belavin (1980) was verified at publisher-metadata level. It supplies relevant primary context, but no opened exact theorem was matched to an unrestricted classification claim. The fragment is therefore F, rather than a completed classification or an invented open problem. \citep{belavin1980discrete}.

\paragraph{Present status.} Status as of 12 November 2026: see the dated fate assessment; no stronger literature-completeness claim is made.

\paragraph{Assessment.} A precise classification question requires additional hypotheses absent from the notes.

\paragraph{Audit.} Literal correctness: The classification scope is insufficiently specified.. Publication match: thematic. Verification: publishermetadata. Development: annotated. Confidence: medium.

\noindent References: \citep{belavin1980discrete}.\par

\subsubsection{1983-618.15 --- The classical r-matrix limit}

\label{item:1983-618.15}

\noindent\textbf{Kind:} method. \textbf{Fate:} E. \textbf{Pages:} 1154--1163.\par

\paragraph{Notebook content.} A first-order expansion of a quantum R-matrix leads to a classical equation.

\paragraph{Precise statement.} If $R(u;\hbar)=I+\hbar r(u)+O(\hbar^2)$ satisfies the quantum Yang--Baxter equation, its coefficient of $\hbar^2$ gives $[r_{12}(u-v),r_{13}(u-w)]+[r_{12}(u-v),r_{23}(v-w)]+[r_{13}(u-w),r_{23}(v-w)]=0$. For the rational family a representative is $r(u)=P/u$; adding a scalar multiple of the identity does not change the commutators.

\paragraph{Editorial reconstruction.} This expansion is a necessary classical-limit condition. It does not prove that every classical solution quantizes, nor does it furnish a classification without the hypotheses of an appropriate theorem.

\paragraph{Literature and fate.} The implication is verified by direct expansion. Belavin--Drinfeld (1982) is a precise primary reference for the classical subject, but its full classification proof was not read in this session; the broader publication match remains partial. \citep{belavindrinfeld1982classical,belavin1980discrete}.

\paragraph{Present status.} Status as of 12 November 2026: the reconstructed result is established; the cited verification has the limits stated below.

\paragraph{Assessment.} Established mathematics; useful as a precise critical annotation.

\paragraph{Audit.} Literal correctness: Correct after the stated hypotheses and normalization are supplied.. Publication match: partial. Verification: publishermetadata. Development: annotated. Confidence: high.

\noindent References: \citep{belavindrinfeld1982classical,belavin1980discrete}.\par

\subsubsection{1983-618.16 --- The tetrahedron equation and nontrivial solutions}

\label{item:1983-618.16}

\noindent\textbf{Kind:} question. \textbf{Fate:} E. \textbf{Pages:} 1154--1163.\par

\paragraph{Notebook content.} The three-dimensional analogue of the Yang--Baxter equation is displayed, and the existence of nontrivial solutions is questioned.

\paragraph{Precise statement.} A constant tetrahedron equation for $R\in\operatorname{End}(V^{\otimes3})$ is $R_{123}R_{145}R_{246}R_{356}=R_{356}R_{246}R_{145}R_{123}$ on $V^{\otimes6}$; spectral-parameter versions require a specified compatible parameter assignment. In the broad sense of a nontrivial solution, existence was already established before the 12 December 1983 lecture.

\paragraph{Editorial reconstruction.} Zamolodchikov proposed nontrivial solutions in 1981, and Baxter proved their validity in a paper published in June 1983. A narrower request involving a particular finite dimension, positivity, parameter dependence or classical limit cannot be supplied from the fragment. Later q-oscillator constructions use appropriate algebraic Fock spaces and restrictions on q; they are not interchangeable with every possible finite-dimensional requirement.

\paragraph{Literature and fate.} Baxter (1983) was received on 28 October 1982 and published in June 1983. Its opened publisher abstract explicitly states the proof of Zamolodchikov solutions, so the broad existence question is E. The full Baxter proof was not read. Bazhanov--Sergeev (2005 preprint; 2006 article), section 4, equations 27--33, was read for a later explicit construction. This corrects the earlier provisional attribution of the fate solely to that later construction. \citep{zamolodchikov1981tetrahedron,baxter1983tetrahedron,bazhanovsergeev2006tetrahedron}.

\paragraph{Present status.} Status as of 12 November 2026: broad nontrivial existence was proved before the 1983 lecture (E); any additional unstated restriction remains unassessable (F).

\paragraph{Assessment.} Historical correction is substantive; no new broad existence problem remains.

\paragraph{Audit.} Literal correctness: A broad assertion of nonexistence or continuing openness is historically incorrect by the lecture date; an unstated narrower question is indeterminate.. Publication match: partial. Verification: publishermetadata. Development: annotated. Confidence: high.

\noindent Canonical problem: GP-1983-34.\par

\noindent References: \citep{zamolodchikov1981tetrahedron,baxter1983tetrahedron,bazhanovsergeev2006tetrahedron}.\par

\subsubsection{1983-618.17 --- The rational SU(2)-invariant R-matrix}

\label{item:1983-618.17}

\noindent\textbf{Kind:} result. \textbf{Fate:} E. \textbf{Pages:} 1154--1163.\par

\paragraph{Notebook content.} The two-body matrix for the spin-$1/2$ delta gas is given in terms of the permutation operator.

\paragraph{Precise statement.} Let $P(v\otimes w)=w\otimes v$ and let c be real. The normalized rational family $R_c(u)=(uI+icP)/(u+ic)$ has triplet eigenvalue 1 and singlet eigenvalue $(u-ic)/(u+ic)$. It obeys the Yang--Baxter equation, $R_c(u)R_c(-u)=I$, and is unitary for real u wherever defined; apparent singularities are understood meromorphically.

\paragraph{Editorial reconstruction.} Any separate fermionic exchange sign is excluded from this R convention. The projectors $\Pi_\pm=(I\pm P)/2$ make the two eigenvalues explicit. Overall scalar normalizations do not change the Yang--Baxter equation but matter when matching a physical phase convention.

\paragraph{Literature and fate.} This is the established rational solution in the Yang/Faddeev framework. An independent exact polynomial calculation in the group algebra of $S_3$ checked the numerator Yang--Baxter identity and returned zero for every coefficient. The check confirms the reconstruction, not a new result. \citep{yang1967manybody,faddeev1996bethe}.

\paragraph{Present status.} Status as of 12 November 2026: the reconstructed result is established; the cited verification has the limits stated below.

\paragraph{Assessment.} Established mathematics; useful as a precise critical annotation.

\paragraph{Audit.} Literal correctness: Correct after the stated hypotheses and normalization are supplied.. Publication match: exact. Verification: readtheresult. Development: annotated. Confidence: high.

\noindent References: \citep{yang1967manybody,faddeev1996bethe}.\par

\subsubsection{1983-618.18 --- The contact-matrix claim and an elementary counterexample}

\label{item:1983-618.18}

\noindent\textbf{Kind:} question. \textbf{Fate:} F. \textbf{Pages:} 1154--1163.\par

\paragraph{Notebook content.} The notes suggest that choosing a contact matrix different from I and P destroys the required factorization.

\paragraph{Precise statement.} For $V=aI+bP$ with real a,b, consider the displayed scattering convention $S_V(u)=(uI+iVP)(uI+iV)^{-1}$. On $\Pi_+$ it equals 1, while on $\Pi_-$ it equals $(u-i(a-b))/(u+i(a-b))$. Thus $S_V(u)=R_{a-b}(u)$. In particular $V=2I+P$, which is neither I nor P, gives the standard rational solution with coupling 1.

\paragraph{Editorial reconstruction.} The calculation disproves the literal matrix-level exclusion. It does not exhibit a new physical interaction class: in the fermionic contact problem only the effective singlet coupling is relevant here. A classification modulo physically equivalent contact matrices would require an equivalence relation that the notebook does not specify.

\paragraph{Literature and fate.} The bounded correction follows immediately by evaluating the two projectors, and the rational Yang--Baxter identity was checked independently. The original unrestricted claim is false, while a possible narrower classification request remains F. No novelty is claimed for the counterexample or its effective-coupling interpretation. \citep{yang1967manybody}.

\paragraph{Present status.} Status as of 12 November 2026: see the dated fate assessment; no stronger literature-completeness claim is made.

\paragraph{Assessment.} Useful elementary correction, with no claim of a new integrable model or open classification result.

\paragraph{Audit.} Literal correctness: False as a literal exclusion of every V other than I and P.. Publication match: partial. Verification: readtheresult. Development: annotated. Confidence: high.

\noindent References: \citep{yang1967manybody}.\par

\subsubsection{1983-618.19 --- What finding eigenfunctions means outside the integrable case}

\label{item:1983-618.19}

\noindent\textbf{Kind:} question. \textbf{Fate:} F. \textbf{Pages:} 1154--1163.\par

\paragraph{Notebook content.} The difficulty of finding eigenfunctions for more general interactions is discussed.

\paragraph{Precise statement.} For a real bounded pair potential on a finite periodic configuration space, the Schrödinger operator is self-adjoint with compact resolvent and therefore has a complete orthonormal eigenbasis. If particle exchange symmetry is imposed, the potential must preserve that symmetry and the same conclusion holds in the corresponding invariant subspace. The absence of a Bethe factorization is a failure of that explicit method, not nonexistence of eigenfunctions.

\paragraph{Editorial reconstruction.} To turn the request into a mathematical open problem one must specify the operator, boundary conditions and the permitted meaning of explicit or closed-form solution. The notebook supplies none of the necessary algorithmic or representational restrictions.

\paragraph{Literature and fate.} Yang (1968) provides relevant integrable-scattering context. No exact formal solvability problem can be matched to this general remark, so its research-question status is F. \citep{yang1968smatrix}.

\paragraph{Present status.} Status as of 12 November 2026: see the dated fate assessment; no stronger literature-completeness claim is made.

\paragraph{Assessment.} No independent well-posed research problem is extracted without a specified solution class.

\paragraph{Audit.} Literal correctness: Correct as a statement of difficulty, but not as nonexistence of an eigenbasis on a compact domain.. Publication match: thematic. Verification: publishermetadata. Development: annotated. Confidence: medium.

\noindent References: \citep{yang1968smatrix}.\par

\subsubsection{1983-618.20 --- Periodic transport and the length of the circle}

\label{item:1983-618.20}

\noindent\textbf{Kind:} construction. \textbf{Fate:} E. \textbf{Pages:} 1154--1163.\par

\paragraph{Notebook content.} Transporting a particle once around the periodic interval yields a product of two-body matrices.

\paragraph{Precise statement.} For a circle of circumference $\ell=2L$, fix rapidities $k_1,\ldots,k_N$ and define $Z_j=R_{j,j-1}(k_j-k_{j-1})\cdots R_{j1}(k_j-k_1)R_{jN}(k_j-k_N)\cdots R_{j,j+1}(k_j-k_{j+1})$. In the stated orientation the periodic condition is $e^{ik_j\ell}Z_jv=v$. Every other particle occurs exactly once and the index j is omitted.

\paragraph{Editorial reconstruction.} If the circle is represented by $[-L,L]$, its length is $2L$, so the phase is $e^{2ik_jL}$. Reversing the transport orientation inverts the matrix product and the free propagation phase together; mixing these conventions changes the Bethe equations.

\paragraph{Literature and fate.} The construction is part of the pre-existing Yang solution. The critical correction is the circumference and the cyclic product convention, not a new quantization condition. \citep{yang1967manybody}.

\paragraph{Present status.} Status as of 12 November 2026: the reconstructed result is established; the cited verification has the limits stated below.

\paragraph{Assessment.} Established mathematics; useful as a precise critical annotation.

\paragraph{Audit.} Literal correctness: The written phase must use the full circumference, and the transport product must omit its own particle index.. Publication match: exact. Verification: readtheresult. Development: annotated. Confidence: high.

\noindent References: \citep{yang1967manybody}.\par

\subsubsection{1983-618.21 --- Commutativity of the spin transport operators}

\label{item:1983-618.21}

\noindent\textbf{Kind:} result. \textbf{Fate:} E. \textbf{Pages:} 1154--1163.\par

\paragraph{Notebook content.} The spin transport operators are asserted to commute.

\paragraph{Precise statement.} For the rational R-matrix, or a family satisfying the same Yang--Baxter, regularity $R(0)=P$ and inverse-exchange conditions, the periodic operators $Z_j$ of item 20 commute for fixed rapidities. With monodromy $T_0(u)=R_{0N}(u-k_N)\cdots R_{01}(u-k_1)$ and auxiliary trace $t(u)=\operatorname{tr}_0T_0(u)$, one has $Z_j=t(k_j)$ in the specified product convention and hence $[Z_i,Z_j]=0$.

\paragraph{Editorial reconstruction.} The operators act on the same spin space $(\mathbb C^2)^{\otimes N}$. This assertion concerns the transfer matrices and periodic transport, not commutativity of every operator-valued matrix entry of T.

\paragraph{Literature and fate.} The established Yang/Faddeev construction verifies this statement. The regularity identity identifies each special-value transfer matrix with the cyclic transport product. \citep{yang1967manybody,faddeev1996bethe}.

\paragraph{Present status.} Status as of 12 November 2026: the reconstructed result is established; the cited verification has the limits stated below.

\paragraph{Assessment.} Established mathematics; useful as a precise critical annotation.

\paragraph{Audit.} Literal correctness: Correct after the stated hypotheses and normalization are supplied.. Publication match: exact. Verification: readtheresult. Development: annotated. Confidence: high.

\noindent References: \citep{yang1967manybody,faddeev1996bethe}.\par

\subsubsection{1983-618.22 --- Pauli Lax operator and monodromy}

\label{item:1983-618.22}

\noindent\textbf{Kind:} construction. \textbf{Fate:} E. \textbf{Pages:} 1154--1163.\par

\paragraph{Notebook content.} A local two-by-two auxiliary matrix is multiplied over the particles.

\paragraph{Precise statement.} On auxiliary spin 0 and quantum spin j, $P_{0j}=\tfrac12(I+\sum_{a=1}^3\sigma_0^a\sigma_j^a)$. Therefore $L_j(u)=R_{0j}(u-k_j)=\bigl[(u-k_j+ic/2)I+(ic/2)\sum_a\sigma_0^a\sigma_j^a\bigr]/(u-k_j+ic)$. Define $T_0(u)=L_N(u)\cdots L_1(u)=\begin{pmatrix}A(u)&B(u)\\ C(u)&D(u)\end{pmatrix}_0$.

\paragraph{Editorial reconstruction.} A,B,C,D are operators on the full N-spin Hilbert space, while the displayed two-by-two indices belong to a separate auxiliary spin. The monodromy order fixes the corresponding transport and exchange conventions.

\paragraph{Literature and fate.} This is the standard quantum inverse-scattering construction. The 1978 primary record was checked at metadata level; Faddeev sections 3--4 and the relevant displayed equations were read for the explicit algebra. \citep{sklyaninfaddeev1978quantum,faddeev1996bethe}.

\paragraph{Present status.} Status as of 12 November 2026: the reconstructed result is established; the cited verification has the limits stated below.

\paragraph{Assessment.} Established mathematics; useful as a precise critical annotation.

\paragraph{Audit.} Literal correctness: Correct after the stated hypotheses and normalization are supplied.. Publication match: exact. Verification: readtheresult. Development: annotated. Confidence: high.

\noindent References: \citep{sklyaninfaddeev1978quantum,faddeev1996bethe}.\par

\subsubsection{1983-618.23 --- RTT relation and the auxiliary trace}

\label{item:1983-618.23}

\noindent\textbf{Kind:} result. \textbf{Fate:} E. \textbf{Pages:} 1154--1163.\par

\paragraph{Notebook content.} The local consistency relation is promoted to a relation for the monodromy matrix.

\paragraph{Precise statement.} On two auxiliary spaces 0 and $0^\prime$, the monodromy satisfies $R_{00^\prime}(u-v)T_0(u)T_{0^\prime}(v)=T_{0^\prime}(v)T_0(u)R_{00^\prime}(u-v)$. Taking the auxiliary traces gives $[t(u),t(v)]=0$ for $t(u)=A(u)+D(u)$. One first argues where R is invertible and then extends the equality meromorphically.

\paragraph{Editorial reconstruction.} The local relation is iterated through independent quantum sites. The cyclicity used is in the finite auxiliary tensor spaces; the operator entries on the common quantum space are not individually commutative.

\paragraph{Literature and fate.} The RTT and transfer-matrix statements are established. The opened Faddeev treatment provides the explicit derivation, while the original 1978 publication was checked at metadata level only. \citep{sklyaninfaddeev1978quantum,faddeev1996bethe}.

\paragraph{Present status.} Status as of 12 November 2026: the reconstructed result is established; the cited verification has the limits stated below.

\paragraph{Assessment.} Established mathematics; useful as a precise critical annotation.

\paragraph{Audit.} Literal correctness: Correct for the auxiliary trace; interpreting it as commutativity of all entries would be false.. Publication match: exact. Verification: readtheresult. Development: annotated. Confidence: high.

\noindent References: \citep{sklyaninfaddeev1978quantum,faddeev1996bethe}.\par

\subsubsection{1983-618.24 --- Exchange relations and the BB correction}

\label{item:1983-618.24}

\noindent\textbf{Kind:} identity. \textbf{Fate:} E. \textbf{Pages:} 1154--1163.\par

\paragraph{Notebook content.} Commutation relations among A, B and D are used to construct eigenvectors.

\paragraph{Precise statement.} In the convention $R(u)=(uI+icP)/(u+ic)$, set $w=u-v$. The relations are $[B(u),B(v)]=0$, $A(u)B(v)=\frac{w-ic}{w}B(v)A(u)+\frac{ic}{w}B(u)A(v)$, and $D(u)B(v)=\frac{w+ic}{w}B(v)D(u)-\frac{ic}{w}B(u)D(v)$, initially for generic u,v and thereafter as meromorphic identities.

\paragraph{Editorial reconstruction.} The commuting pair is B with B. A printed relation asserting that B and D commute would remove the unwanted terms needed for the Bethe equations and is incorrect. Spectral and sign conventions must be changed together when comparing sources.

\paragraph{Literature and fate.} Faddeev equations 66--104 were read for the algebraic Bethe construction. The corrected relations are established; they also determine the normalization used in items 25--27. \citep{faddeev1996bethe}.

\paragraph{Present status.} Status as of 12 November 2026: the reconstructed result is established; the cited verification has the limits stated below.

\paragraph{Assessment.} Established mathematics; useful as a precise critical annotation.

\paragraph{Audit.} Literal correctness: Replace the asserted BD commutator by the BB commutator; the remaining coefficients require a consistent R convention.. Publication match: exact. Verification: readtheresult. Development: annotated. Confidence: high.

\noindent References: \citep{faddeev1996bethe}.\par

\subsubsection{1983-618.25 --- Reference vector and its transfer eigenvalue}

\label{item:1983-618.25}

\noindent\textbf{Kind:} construction. \textbf{Fate:} E. \textbf{Pages:} 1154--1163.\par

\paragraph{Notebook content.} The all-up spin vector is used as the reference state.

\paragraph{Precise statement.} For $\Omega=|\uparrow\rangle^{\otimes N}$ and the normalized Lax operators above, $C(u)\Omega=0$, $A(u)\Omega=\Omega$, and $D(u)\Omega=d(u)\Omega$, where $d(u)=\prod_{j=1}^N(u-k_j)/(u-k_j+ic)$. Consequently $t(u)\Omega=[1+d(u)]\Omega$.

\paragraph{Editorial reconstruction.} The transfer matrix is the sum A+D. The constant contribution 1 from A cannot be omitted; it changes the eigenvalue already for the reference state.

\paragraph{Literature and fate.} This is the standard pseudovacuum computation in the opened algebraic Bethe-ansatz source. It provides a direct local check on the corrected eigenvalue. \citep{faddeev1996bethe}.

\paragraph{Present status.} Status as of 12 November 2026: the reconstructed result is established; the cited verification has the limits stated below.

\paragraph{Assessment.} Established mathematics; useful as a precise critical annotation.

\paragraph{Audit.} Literal correctness: The reference-state transfer eigenvalue requires the missing additive 1.. Publication match: exact. Verification: readtheresult. Development: annotated. Confidence: high.

\noindent References: \citep{faddeev1996bethe}.\par

\subsubsection{1983-618.26 --- Bethe vectors and cancellation of unwanted terms}

\label{item:1983-618.26}

\noindent\textbf{Kind:} method. \textbf{Fate:} E. \textbf{Pages:} 1154--1163.\par

\paragraph{Notebook content.} Products of B operators acting on the reference vector yield the candidate eigenvectors.

\paragraph{Precise statement.} Let $v=B(\lambda_1)\cdots B(\lambda_M)\Omega$ be nonzero and let the roots be generic enough that the displayed denominators are nonzero. The unwanted terms cancel if $d(\lambda_a)^{-1}=\prod_{b\ne a}\frac{\lambda_a-\lambda_b+ic}{\lambda_a-\lambda_b-ic}$ for $1\leq a\leq M$. Then $t(u)v=\Lambda(u)v$, with $\Lambda(u)=\prod_{a=1}^M\frac{u-\lambda_a-ic}{u-\lambda_a}+d(u)\prod_{a=1}^M\frac{u-\lambda_a+ic}{u-\lambda_a}$.

\paragraph{Editorial reconstruction.} There are M B factors, one for each lowered spin, not N factors independent of the spin sector. Coincident or singular roots, vanishing Bethe vectors and spin descendants require a separate limiting or representation-theoretic treatment; they are not covered by the generic cancellation calculation.

\paragraph{Literature and fate.} The construction is established and was checked against Faddeev equations 66--104. It gives eigenvectors subject to nonvanishing and admissibility, not an automatic proof that all formal roots yield distinct eigenstates. \citep{faddeev1996bethe}.

\paragraph{Present status.} Status as of 12 November 2026: the reconstructed result is established; the cited verification has the limits stated below.

\paragraph{Assessment.} Established mathematics; useful as a precise critical annotation.

\paragraph{Audit.} Literal correctness: Use M factors and impose the regularity and nonzero-vector hypotheses omitted by the shorthand.. Publication match: exact. Verification: readtheresult. Development: annotated. Confidence: high.

\noindent References: \citep{faddeev1996bethe}.\par

\subsubsection{1983-618.27 --- Nested charge and spin Bethe equations}

\label{item:1983-618.27}

\noindent\textbf{Kind:} result. \textbf{Fate:} E. \textbf{Pages:} 1154--1163.\par

\paragraph{Notebook content.} Periodic transport and the auxiliary eigenproblem produce coupled equations for the particle and spin rapidities.

\paragraph{Precise statement.} For the repulsive Hamiltonian with interaction $2c\delta$, $c>0$, circle length $\ell$, and regular admissible roots, put $\mu_a=\lambda_a+ic/2$. The coupled equations are $e^{ik_j\ell}=\prod_{a=1}^M\frac{k_j-\mu_a+ic/2}{k_j-\mu_a-ic/2}$ and $\prod_{j=1}^N\frac{\mu_a-k_j+ic/2}{\mu_a-k_j-ic/2}=\prod_{b\ne a}\frac{\mu_a-\mu_b+ic}{\mu_a-\mu_b-ic}$. The corresponding nonzero contact wavefunction has energy $E=\sum_jk_j^2$.

\paragraph{Editorial reconstruction.} The spin product omits b=a. Including that factor contributes -1 and requires an additional compensating minus sign. A notation with charge shifts ic and spin shifts 2ic is consistent only after rescaling the coupling everywhere. The periodic phase uses $\ell=2L$ for an interval $[-L,L]$.

\paragraph{Literature and fate.} The original Yang paper, equations 20--27, and the later algebraic account verify this nested construction. The normalization and admissibility conditions make the checked statement precise; a full continuum-completeness assertion is excluded from this theorem. \citep{yang1967manybody,faddeev1996bethe}.

\paragraph{Present status.} Status as of 12 November 2026: the reconstructed result is established; the cited verification has the limits stated below.

\paragraph{Assessment.} Established mathematics; useful as a precise critical annotation.

\paragraph{Audit.} Literal correctness: Correct after fixing the interaction normalization, self-root sign and full circumference.. Publication match: exact. Verification: readtheresult. Development: annotated. Confidence: high.

\noindent References: \citep{yang1967manybody,faddeev1996bethe}.\par

\subsubsection{1983-618.28 --- Why the number 2 to the N is not continuum completeness}

\label{item:1983-618.28}

\noindent\textbf{Kind:} claim. \textbf{Fate:} F. \textbf{Pages:} 1154--1163.\par

\paragraph{Notebook content.} The number of solutions is compared with $2^N$.

\paragraph{Precise statement.} At fixed particle rapidities the spin space has dimension $2^N$, and its fixed-M weight space has dimension $\binom NM$. The complete continuum Hilbert space on a circle is infinite dimensional. Already for N=1, the states $e^{2\pi imx/\ell}|\uparrow\rangle$ and $e^{2\pi imx/\ell}|\downarrow\rangle$, $m\in\mathbb Z$, give infinitely many momentum-spin eigenstates.

\paragraph{Editorial reconstruction.} Counting generic highest-weight Bethe vectors does not count all spin descendants or singular solutions. A spin-space count and completeness of the full coordinate-spin spectral problem are distinct assertions, with different quantifiers.

\paragraph{Literature and fate.} The literal identification of $2^N$ with the full continuum spectrum is false. Yang (1967), footnote 1 on page 1314, expresses a belief concerning completeness rather than furnishing a blanket proof. The unqualified notebook completeness claim is therefore F; the established finite-dimensional spin count is retained explicitly. \citep{yang1967manybody,faddeev1996bethe}.

\paragraph{Present status.} Status as of 12 November 2026: see the dated fate assessment; no stronger literature-completeness claim is made.

\paragraph{Assessment.} A precise completeness theorem needs its model, sector, admissibility and singular-root conventions; no broad unresolved problem is invented.

\paragraph{Audit.} Literal correctness: False if 2\textasciicircum{}N is asserted to count the full continuum eigenbasis; incomplete if it refers only to regular highest-weight Bethe solutions.. Publication match: partial. Verification: readtheresult. Development: annotated. Confidence: high.

\noindent References: \citep{yang1967manybody,faddeev1996bethe}.\par

\subsubsection{1983-618.29 --- Thermodynamic root-density equations}

\label{item:1983-618.29}

\noindent\textbf{Kind:} method. \textbf{Fate:} E. \textbf{Pages:} 1154--1163.\par

\paragraph{Notebook content.} Logarithms of the coupled equations are replaced by integral equations for densities.

\paragraph{Precise statement.} Let $N,\ell\to\infty$ with $N/\ell\to n>0$ and $M/\ell\to m\in[0,n]$. In a regime described by smooth charge and real spin-root densities $\rho,\sigma$ and hole densities $\rho^h,\sigma^h$, differentiation of the counting equations gives $\rho+\rho^h=\frac1{2\pi}+a_1*\sigma$ and $\sigma+\sigma^h=a_1*\rho-a_2*\sigma$, where $a_p(x)=\frac1{2\pi}\frac{pc}{(pc/2)^2+x^2}$. The normalizations are $\int\rho=N/\ell$, $\int\sigma=M/\ell$ and $E/\ell=\int k^2\rho(k)\,dk$. In the usual highest-weight parametrization the real-root spin sector has $0\leq m\leq n/2$; other magnetizations are related by spin symmetry.

\paragraph{Editorial reconstruction.} The macroscopic limit concerns N and the circumference; M need not tend to infinity when the limiting down-spin density is zero. The two displayed equations describe the stipulated real-root or ground-state regime. Finite-temperature sectors with spin strings require additional string densities and equations, so this pair is not a complete finite-temperature thermodynamic Bethe ansatz.

\paragraph{Literature and fate.} The original Yang paper supplies the thermodynamic integral-equation method in its intended regime. The reconstruction fixes the density normalization and explicitly separates it from an unrestricted thermal completeness claim. \citep{yang1967manybody}.

\paragraph{Present status.} Status as of 12 November 2026: the reconstructed result is established; the cited verification has the limits stated below.

\paragraph{Assessment.} Established mathematics; useful as a precise critical annotation.

\paragraph{Audit.} Literal correctness: Correct within the stated real-root regime and scaling; incomplete as a formula for every finite-temperature spin sector.. Publication match: exact. Verification: readtheresult. Development: annotated. Confidence: high.

\noindent References: \citep{yang1967manybody}.\par

\subsection{Record 619: YEAR 1984}

\label{talk:619}

\noindent\textbf{Date:} 1984. \textbf{Speaker:} . \textbf{Pages:} .\par

Catalogue year heading, with no separate mathematical occurrence in the supplied source.

\subsection{Record 620: Quadratic exchange algebras and finite-dimensional elliptic representations}

\label{talk:620}

\noindent\textbf{Date:} 1984-02-20. \textbf{Speaker:} I. V. Cherednik. \textbf{Pages:} 1164--1166.\par

The report itself begins on P001164 and ends on P001166. P001162–P001163 belong to the preceding Wiegmann material, and P001167 begins Zelevinsky’s report. The central new result has a close match in Cherednik’s two-page paper received 23 February 1984, three days after the seminar; the elliptic \(R\)-matrix and rank-two Sklyanin algebra have earlier sources.

\subsubsection{1983-620.1 --- PBW quantization and the classical Poisson limit}

\label{item:1983-620.1}

\noindent\textbf{Kind:} construction. \textbf{Fate:} E. \textbf{Pages:} 1164.\par

\paragraph{Notebook content.} The notes begin with the passage from a Lie algebra and its enveloping algebra to a commutative classical algebra.

\paragraph{Precise statement.} For a complex Lie algebra \(\mathfrak g\), the PBW filtration on \(U(\mathfrak g)\) has \(\operatorname{gr}U(\mathfrak g)\cong S(\mathfrak g)\). The commutator induces the Poisson bracket on \(S(\mathfrak g)\) extending \(\{x,y\}=[x,y]_{\mathfrak g}\). Equivalently the algebra with relations \(xy-yx=\hbar[x,y]\) is a flat \(\mathbb C[\hbar]\)-deformation with fibres \(S(\mathfrak g)\) at \(0\) and \(U(\mathfrak g)\) at \(1\).

\paragraph{Literature and fate.} This is PBW background. Frenkel (2007), Appendix A.2, Theorem A.2.2, explicitly states the ordered basis and associated-graded identification. The deformation formulation is its direct consequence.

\paragraph{Present status.} Established under the stated hypotheses as of 12 November 2026.

\paragraph{Assessment.} Separates the recorded algebraic construction from stronger claims not implied by its formulas.

\paragraph{Audit.} Literal correctness: correct. Publication match: exact. Verification: read the result. Development: annotated. Confidence: high.

\noindent References: \citep{frenkel2007loop}.\par

\subsubsection{1983-620.2 --- Deforming enveloping algebras beyond the Lie case}

\label{item:1983-620.2}

\noindent\textbf{Kind:} programme. \textbf{Fate:} F. \textbf{Pages:} 1164.\par

\paragraph{Notebook content.} The speaker motivates quadratic algebras as further deformations and asks for their representations.

\paragraph{Precise statement.} A general deformation programme requires a specified initial algebra, parameter space, flatness condition, and representation class. The notebook does not formulate one universal classification problem at this point; the concrete exchange and elliptic constructions are handled separately.

\paragraph{Literature and fate.} The subsequent papers provide concrete examples, but the introductory programme has no exact single theorem or open question.

\paragraph{Present status.} No uniquely formulated general deformation problem is identified.

\paragraph{Assessment.} Separates the recorded algebraic construction from stronger claims not implied by its formulas.

\paragraph{Audit.} Literal correctness: not determinable. Publication match: thematic only. Verification: unresolved lead. Development: annotated. Confidence: medium.

\noindent References: \citep{sklyanin1982algebraic,cherednik1985finite}.\par

\subsubsection{1983-620.3 --- The RTT exchange algebra}

\label{item:1983-620.3}

\noindent\textbf{Kind:} definition. \textbf{Fate:} E. \textbf{Pages:} 1164--1165.\par

\paragraph{Notebook content.} A matrix of noncommuting entries is constrained by an \(R\)-matrix exchange relation.

\paragraph{Precise statement.} Let \(V\) be finite dimensional and \(R(u,v)\in\operatorname{End}(V\otimes V)\) be a specified meromorphic matrix. The RTT algebra is generated by the chosen coefficients of \(L(u)\), modulo \(R_{12}(u,v)L_1(u)L_2(v)=L_2(v)L_1(u)R_{12}(u,v)\), where \(L_1=L\otimes1\) and \(L_2=1\otimes L\). The allowed coefficient space and the values at poles are part of the definition.

\paragraph{Literature and fate.} The exchange relation is the algebraic framework in Sklyanin (1982), \S2, and Cherednik (1985), \S1. It defines an associative quotient without, by itself, asserting a PBW basis or flatness.

\paragraph{Present status.} Established under the stated hypotheses as of 12 November 2026.

\paragraph{Assessment.} Separates the recorded algebraic construction from stronger claims not implied by its formulas.

\paragraph{Audit.} Literal correctness: correct. Publication match: exact. Verification: read the result. Development: annotated. Confidence: high.

\noindent References: \citep{sklyanin1982algebraic,cherednik1985finite}.\par

\subsubsection{1983-620.4 --- Yang–Baxter consistency does not supply every PBW assertion}

\label{item:1983-620.4}

\noindent\textbf{Kind:} claim. \textbf{Fate:} E. \textbf{Pages:} 1165.\par

\paragraph{Notebook content.} The Yang–Baxter equation is presented as the consistency condition behind an ordered-monomial statement.

\paragraph{Precise statement.} The equation \(R_{12}(u,v)R_{13}(u,w)R_{23}(v,w)=R_{23}(v,w)R_{13}(u,w)R_{12}(u,v)\) equates the two three-factor reordering routes. A basis theorem still requires specified ordering, parameters, and a confluence or flatness argument; the equation alone does not establish arbitrary monomial independence. For instance \(R=1\) gives commuting RTT entries, so different orderings of the same generators represent the same element.

\paragraph{Editorial reconstruction.} The precise intended basis statement is not recoverable without its missing ordering and parameter restrictions.

\paragraph{Literature and fate.} Sklyanin’s exchange-algebra construction explains the consistency mechanism. The final example is our elementary check against an unrestricted reading of the notebook’s basis wording, not a counterexample to a correctly formulated PBW theorem for a particular algebra.

\paragraph{Present status.} Established under the stated hypotheses as of 12 November 2026.

\paragraph{Assessment.} Separates the recorded algebraic construction from stronger claims not implied by its formulas.

\paragraph{Audit.} Literal correctness: false as written. Publication match: partial. Verification: read the result. Development: annotated. Confidence: high.

\noindent References: \citep{sklyanin1982algebraic}.\par

\subsubsection{1983-620.5 --- Unitarity is an additional normalization condition}

\label{item:1983-620.5}

\noindent\textbf{Kind:} claim. \textbf{Fate:} E. \textbf{Pages:} 1165.\par

\paragraph{Notebook content.} An arrow in the notes suggests that Yang–Baxter implies \(R_{12}(u)R_{21}(-u)=1\).

\paragraph{Precise statement.} Yang–Baxter does not imply unitarity. For \(D=\operatorname{diag}(1,2)\), the constant matrix \(R=D\otimes D\) solves Yang–Baxter because the three leg operators commute, whereas \(R_{12}R_{21}=D^2\otimes D^2=\operatorname{diag}(1,4,4,16)\ne1\). Unitarity, or scalar unitarity followed by an admissible normalization, must be imposed or proved for the specific family.

\paragraph{Literature and fate.} This is a direct countercheck. Cherednik (1985), \S1, Proposition 1, proves inversion properties for his particular elliptic matrices; those properties do not follow from the Yang–Baxter equation alone.

\paragraph{Present status.} Established under the stated hypotheses as of 12 November 2026.

\paragraph{Assessment.} Separates the recorded algebraic construction from stronger claims not implied by its formulas.

\paragraph{Audit.} Literal correctness: false as written. Publication match: partial. Verification: read the result. Development: annotated. Confidence: high.

\noindent References: \citep{cherednik1985finite}.\par

\subsubsection{1983-620.6 --- Elliptic functions with torsion multipliers}

\label{item:1983-620.6}

\noindent\textbf{Kind:} construction. \textbf{Fate:} E. \textbf{Pages:} 1165.\par

\paragraph{Notebook content.} Functions \(w_\alpha\) with prescribed elliptic multipliers and a normalization are introduced.

\paragraph{Precise statement.} Let \(E=\mathbb C/(\mathbb Z+\tau\mathbb Z)\), \(\operatorname{Im}\tau>0\), \(n\ge2\), and \(\omega_\alpha=(\alpha_1+\alpha_2\tau)/n\), \(\alpha\in(\mathbb Z/n)^2\). For generic \(\eta\), define \(w_\alpha(u)=e^{2\pi i\alpha_2u/n}\theta(u+\eta+\omega_\alpha)\theta(\eta)/(\theta(\eta+\omega_\alpha)\theta(u+\eta))\), with \(\theta\) the odd theta function. These are meromorphic sections of flat torsion line bundles, normalized by \(w_\alpha(0)=1\); \(w_0=1\). Denominators must be nonzero as parameters, and the allowed simple pole is part of the construction.

\paragraph{Editorial reconstruction.} The explicit theta formula is a modern normalization of the recorded multiplier characterization, not a claim that the notebook used the same parameter scale.

\paragraph{Literature and fate.} Belavin (1980) gives the elliptic torsion construction; Cherednik (1985), \S1, uses its pole and multiplier version. The parameter is rescaled in different sources, so the pole may be written \(-\eta\) or \(-\eta/n\).

\paragraph{Present status.} Established under the stated hypotheses as of 12 November 2026.

\paragraph{Assessment.} Separates the recorded algebraic construction from stronger claims not implied by its formulas.

\paragraph{Audit.} Literal correctness: correct with missing hypotheses. Publication match: exact. Verification: read the result. Development: annotated. Confidence: high.

\noindent References: \citep{belavin1980discrete,cherednik1985finite}.\par

\subsubsection{1983-620.7 --- Clock and shift matrices and the permutation operator}

\label{item:1983-620.7}

\noindent\textbf{Kind:} construction. \textbf{Fate:} E. \textbf{Pages:} 1165.\par

\paragraph{Notebook content.} Finite Heisenberg matrices are used as a basis of the matrix algebra.

\paragraph{Precise statement.} On the basis \(e_j\), \(j\in\mathbb Z/n\), let \(Qe_j=\zeta^je_j\), \(\Lambda e_j=e_{j+1}\), \(\zeta=e^{2\pi i/n}\), and \(I_\alpha=Q^{\alpha_1}\Lambda^{\alpha_2}\). They form a basis of \(\operatorname{End}(\mathbb C^n)\), orthogonal for the trace pairing with the inverse basis, and \(P=\frac1n\sum_\alpha I_\alpha\otimes I_\alpha^{-1}\), where \(P(x\otimes y)=y\otimes x\).

\paragraph{Literature and fate.} This is the finite Heisenberg ingredient in Belavin (1980) and Cherednik (1985), \S1. The permutation identity follows by applying the sum to \(e_j\otimes e_k\).

\paragraph{Present status.} Established under the stated hypotheses as of 12 November 2026.

\paragraph{Assessment.} Separates the recorded algebraic construction from stronger claims not implied by its formulas.

\paragraph{Audit.} Literal correctness: correct. Publication match: exact. Verification: read the result. Development: annotated. Confidence: high.

\noindent References: \citep{belavin1980discrete,cherednik1985finite}.\par

\subsubsection{1983-620.8 --- The Belavin elliptic R-matrix}

\label{item:1983-620.8}

\noindent\textbf{Kind:} construction. \textbf{Fate:} E. \textbf{Pages:} 1165--1166.\par

\paragraph{Notebook content.} The elliptic coefficients and matrix basis are combined into an \(R\)-matrix.

\paragraph{Precise statement.} With compatible choices of the torsion multipliers and Heisenberg basis above, \(R(u)=\frac1n\sum_\alpha w_\alpha(u)I_\alpha\otimes I_\alpha^{-1}\) is the normalized Belavin elliptic solution of the difference-form Yang–Baxter equation, with \(R(0)=P\). Scalar rescalings preserve Yang–Baxter, but alter the scalar in any inversion relation.

\paragraph{Editorial reconstruction.} The matrix-basis and multiplier indices must be chosen compatibly; the formula fixes the regularity normalization, not an unproved unitarity normalization.

\paragraph{Literature and fate.} Belavin (1980), Discrete groups and integrability of quantum systems, is the earlier primary source. Cherednik’s 1985 \S1 recalls the same elliptic family before constructing its representations.

\paragraph{Present status.} Established under the stated hypotheses as of 12 November 2026.

\paragraph{Assessment.} Separates the recorded algebraic construction from stronger claims not implied by its formulas.

\paragraph{Audit.} Literal correctness: correct with missing hypotheses. Publication match: exact. Verification: read the result. Development: annotated. Confidence: high.

\noindent References: \citep{belavin1980discrete,cherednik1985finite}.\par

\subsubsection{1983-620.9 --- The rational permutation solution}

\label{item:1983-620.9}

\noindent\textbf{Kind:} example. \textbf{Fate:} E. \textbf{Pages:} 1166.\par

\paragraph{Notebook content.} A rational \(R\)-matrix built from the identity and permutation is exhibited.

\paragraph{Precise statement.} For a complex parameter \(\eta\), \(R(u)=(uI+\eta P)/(u+\eta)\), away from its pole, satisfies the difference-form Yang–Baxter equation. It also satisfies \(R(u)R_{21}(-u)=I\) wherever both sides are defined, by \(P^2=I\); this inversion identity is a calculation for this family.

\paragraph{Literature and fate.} This is the rational specialization of the standard exchange-algebra framework. The Yang–Baxter identity reduces to the permutation relations, and inversion follows by multiplying the two numerators.

\paragraph{Present status.} Established under the stated hypotheses as of 12 November 2026.

\paragraph{Assessment.} Separates the recorded algebraic construction from stronger claims not implied by its formulas.

\paragraph{Audit.} Literal correctness: correct with missing hypotheses. Publication match: partial. Verification: read the result. Development: annotated. Confidence: high.

\noindent References: \citep{sklyanin1982algebraic}.\par

\subsubsection{1983-620.10 --- The fundamental evaluation representation}

\label{item:1983-620.10}

\noindent\textbf{Kind:} construction. \textbf{Fate:} E. \textbf{Pages:} 1166.\par

\paragraph{Notebook content.} The \(R\)-matrix itself is used as an operator-valued \(L\)-matrix.

\paragraph{Precise statement.} If \(R\) solves Yang–Baxter, then \(L_0(u)=R_{0a}(u-a)\), acting on an auxiliary factor \(0\) and a fixed quantum factor \(a\), satisfies the RTT relation with the same \(R\). This supplies a representation of the corresponding RTT quotient whenever the chosen \(L\)-coefficient ansatz is respected.

\paragraph{Literature and fate.} The assertion is the Yang–Baxter equation with one leg designated as the representation space. Cherednik (1985), \S1, uses this construction as the starting point for fusion.

\paragraph{Present status.} Established under the stated hypotheses as of 12 November 2026.

\paragraph{Assessment.} Separates the recorded algebraic construction from stronger claims not implied by its formulas.

\paragraph{Audit.} Literal correctness: correct. Publication match: exact. Verification: read the result. Development: annotated. Confidence: high.

\noindent References: \citep{cherednik1985finite}.\par

\subsubsection{1983-620.11 --- Finite-generator elliptic exchange algebras}

\label{item:1983-620.11}

\noindent\textbf{Kind:} definition. \textbf{Fate:} A. \textbf{Pages:} 1166.\par

\paragraph{Notebook content.} The spectral dependence is separated from finitely many noncommuting generators.

\paragraph{Precise statement.} Choose the finite-dimensional elliptic coefficient space used by the \(R\)-matrix and set \(L(u)=\sum_\alpha S_\alpha w_\alpha(u)I_\alpha\), with the appropriate common scalar normalization. Expanding RTT in this coefficient basis defines homogeneous quadratic relations among the \(n^2\) generators \(S_\alpha\). A PBW or Hilbert-series assertion is an additional theorem and is not included in this definition.

\paragraph{Editorial reconstruction.} The ansatz is presented as a quotient definition; no generic flatness claim is added.

\paragraph{Literature and fate.} Sklyanin’s rank-two construction appears in 1982; Cherednik (1985), \S1 and subsequent fusion statements, provides the general elliptic setting used in the report. The correspondence is close, but the notebook’s compressed coefficient conventions require reconstruction.

\paragraph{Present status.} Established under the stated hypotheses as of 12 November 2026.

\paragraph{Assessment.} Separates the recorded algebraic construction from stronger claims not implied by its formulas.

\paragraph{Audit.} Literal correctness: correct with missing hypotheses. Publication match: partial. Verification: read the result. Development: annotated. Confidence: high.

\noindent References: \citep{sklyanin1982algebraic,cherednik1985finite}.\par

\subsubsection{1983-620.12 --- Rank-two Sklyanin relations and finite theta modules}

\label{item:1983-620.12}

\noindent\textbf{Kind:} example. \textbf{Fate:} E. \textbf{Pages:} 1166.\par

\paragraph{Notebook content.} For \(n=2\), four generators and cyclic quadratic relations are written, followed by finite-dimensional representations.

\paragraph{Precise statement.} In a standard normalization, for cyclic \((\alpha,\beta,\gamma)\) of \((1,2,3)\), \([S_\alpha,S_\beta]=i(S_0S_\gamma+S_\gamma S_0)\) and \([S_0,S_\alpha]=iJ_{\beta\gamma}(S_\beta S_\gamma+S_\gamma S_\beta)\). The parameters obey \(J_{12}+J_{23}+J_{31}+J_{12}J_{23}J_{31}=0\), with degenerate parameter values treated separately. The elliptic difference-operator construction acts on finite spaces of even theta functions of order \(4\ell\), of dimension \(2\ell+1\), for \(2\ell\in\mathbb Z_{\ge0}\) under the usual genericity conditions.

\paragraph{Editorial reconstruction.} The (i)-factors and cyclic parameter labels depend on generator normalization; they are fixed here rather than silently copied across conventions.

\paragraph{Literature and fate.} Sklyanin (1982), \S2, equations (27) and (32), gives the algebra; Sklyanin (1983), \S1, Theorem 4 and the paragraph following it, Russian pp. 37–38, gives invariant theta spaces and specializes to order \(4\ell\), dimension \(2\ell+1\). Both publications predate February 1984.

\paragraph{Present status.} Established under the stated hypotheses as of 12 November 2026.

\paragraph{Assessment.} Separates the recorded algebraic construction from stronger claims not implied by its formulas.

\paragraph{Audit.} Literal correctness: correct with missing hypotheses. Publication match: exact. Verification: read the result. Development: annotated. Confidence: medium.

\noindent References: \citep{sklyanin1982algebraic,sklyanin1983representations}.\par

\subsubsection{1983-620.13 --- Fusion and finite-dimensional elliptic representations}

\label{item:1983-620.13}

\noindent\textbf{Kind:} theorem. \textbf{Fate:} A. \textbf{Pages:} 1166.\par

\paragraph{Notebook content.} The report constructs finite-dimensional representations beyond the defining module.

\paragraph{Precise statement.} For the generic Belavin elliptic \(R\)-matrix, fusion of shifted tensor products yields well-defined fused operators on the symmetric and exterior powers of the defining \(n\)-dimensional space. After the scalar normalization specified in the construction, their coefficients give finite-dimensional representations of the elliptic exchange algebra, with dimensions \(\binom{n+q-1}{q}\) on \(S^q\mathbb C^n\) and \(\binom nq\) on \(\Lambda^q\mathbb C^n\), \(0\le q\le n\) in the exterior case.

\paragraph{Editorial reconstruction.} Generic parameters exclude the poles and degeneracies in the fusion normalization. The statement concerns existence and dimension, not an unrecorded classification of every irreducible representation.

\paragraph{Literature and fate.} Cherednik (1985), Theorems 2–3 and Corollary 4, is the exact near-contemporary publication; received 23 February 1984, published in Functional Analysis and Its Applications 19(1), 77–79. This is distinct from the earlier rank-two theta modules.

\paragraph{Present status.} Established under the stated hypotheses as of 12 November 2026.

\paragraph{Assessment.} Separates the recorded algebraic construction from stronger claims not implied by its formulas.

\paragraph{Audit.} Literal correctness: correct with missing hypotheses. Publication match: exact. Verification: read the result. Development: annotated. Confidence: high.

\noindent References: \citep{cherednik1985finite}.\par

\subsection{Record 621: Models of classical-group representations and hidden symmetries}

\label{talk:621}

\noindent\textbf{Date:} 1984-02-27. \textbf{Speaker:} A. V. Zelevinsky, joint work with I. M. Gelfand. \textbf{Pages:} 1167--1170 upper.\par

The actual report begins on P001167 and ends in the upper part of P001170; Nikiforov begins below it on that page. Its model tables, basic Grassmannians, and hidden actions match Gelfand–Zelevinsky’s 1984 paper, received 15 March 1984. The paper explicitly separates proved even constructions from a supergroup formulation whose proof is deferred. Gelfand’s brief research suggestions are recorded separately without invented conjectures.

\subsubsection{1983-621.1 --- Multiplicity-one models and the sphere example}

\label{item:1983-621.1}

\noindent\textbf{Kind:} definition. \textbf{Fate:} E. \textbf{Pages:} 1167.\par

\paragraph{Notebook content.} The representation of rotations on functions on a sphere motivates the definition of a model.

\paragraph{Precise statement.} For a complex reductive algebraic group \(G\), a model is the algebraic direct sum \(\bigoplus_{\lambda\in\operatorname{Irr}_{\mathrm{alg}}(G)}V_\lambda\), with each finite-dimensional irreducible once. For compact \(SO(3)\), \(L^2(S^2)\) is the Hilbert direct sum of spherical-harmonic spaces of degrees \(\ell\ge0\), each irreducible of dimension \(2\ell+1\). Its finite \(SO(3)\)-type subspace is the corresponding algebraic model.

\paragraph{Literature and fate.} Gelfand–Zelevinsky (1984), \S0, Russian p. 14, gives this definition and example. Hilbert completion and algebraic direct sum are different realizations of the same finite-type spectrum.

\paragraph{Present status.} Established with the stated algebraic and parity conventions as of 12 November 2026.

\paragraph{Assessment.} Identifies the exact model and the scope of the hidden Lie-algebra action.

\paragraph{Audit.} Literal correctness: correct. Publication match: exact. Verification: read the result. Development: annotated. Confidence: high.

\noindent References: \citep{gelfandzelevinsky1984models}.\par

\subsubsection{1983-621.2 --- The hidden so8 action on the sl3 model}

\label{item:1983-621.2}

\noindent\textbf{Kind:} example. \textbf{Fate:} E. \textbf{Pages:} 1167.\par

\paragraph{Notebook content.} A model for \(SU(3)\) or \(\mathfrak{sl}_3\) carries a larger orthogonal symmetry.

\paragraph{Precise statement.} The direct sum of all finite-dimensional irreducible \(\mathfrak{sl}_3(\mathbb C)\)-modules, each once, has a realization whose \(\mathfrak{sl}_3\)-action extends to \(\mathfrak{so}_8(\mathbb C)\). This is an action of the larger Lie algebra; it need not integrate to an algebraic action of \(SO(8)\) on the same function space.

\paragraph{Editorial reconstruction.} The group-level and larger-algebra irreducibility qualifications are not silently inferred from the phrase hidden symmetry.

\paragraph{Literature and fate.} Gelfand–Zelevinsky (1984), opening construction 2 on p. 14, attributes the example to Biedenharn and Flath. The exact extension is identified; an additional assertion of irreducibility for the larger algebra was not verified from the checked passage.

\paragraph{Present status.} Established with the stated algebraic and parity conventions as of 12 November 2026.

\paragraph{Assessment.} Identifies the exact model and the scope of the hidden Lie-algebra action.

\paragraph{Audit.} Literal correctness: correct with missing hypotheses. Publication match: partial. Verification: read the result. Development: annotated. Confidence: high.

\noindent References: \citep{gelfandzelevinsky1984models}.\par

\subsubsection{1983-621.3 --- Complexified Minkowski space as a Grassmannian chart}

\label{item:1983-621.3}

\noindent\textbf{Kind:} construction. \textbf{Fate:} E. \textbf{Pages:} 1167.\par

\paragraph{Notebook content.} The twistor compactification of Minkowski space is cited as a second geometric motivation.

\paragraph{Precise statement.} The open chart in \(\operatorname{Gr}(2,\mathbb C^4)\) consisting of planes transverse to a fixed complementary two-plane is the affine space \(\operatorname{Hom}(\mathbb C^2,\mathbb C^2)\cong\mathbb C^4\), via graphs. The determinant quadratic form on \(2\times2\) matrices gives the complexified Minkowski conformal structure. Real Minkowski space requires an additional real form.

\paragraph{Editorial reconstruction.} The distinction between a complex fourfold and a chosen real Minkowski slice is made explicit.

\paragraph{Literature and fate.} Gelfand–Zelevinsky (1984), opening construction 3, p. 14, identifies the Grassmannian compactification. The graph-chart statement is the precise algebraic construction.

\paragraph{Present status.} Established with the stated algebraic and parity conventions as of 12 November 2026.

\paragraph{Assessment.} Identifies the exact model and the scope of the hidden Lie-algebra action.

\paragraph{Audit.} Literal correctness: correct with missing hypotheses. Publication match: partial. Verification: read the result. Development: annotated. Confidence: high.

\noindent References: \citep{gelfandzelevinsky1984models}.\par

\subsubsection{1983-621.4 --- The earlier exterior-algebra induction models}

\label{item:1983-621.4}

\noindent\textbf{Kind:} theorem. \textbf{Fate:} E. \textbf{Pages:} 1167--1168.\par

\paragraph{Notebook content.} An induction table uses exterior algebras of defining representations.

\paragraph{Precise statement.} Write \(\Lambda(r)=\Lambda^\bullet\mathbb C^r\). For the pairs \((G,M,\tau)\) equal to \((GL(r),O(r),\Lambda(r))\), \((O(2r+1),O(r+1)\times O(r),\Lambda(r+1)\otimes1)\), \((O(2r),O(r)\times O(r),\Lambda(r)\otimes1)\), and \((Sp(2r),GL(r),\Lambda(r))\), algebraic induction \(\operatorname{Ind}_M^G\tau\) is a model of \(G\). Use the displayed full orthogonal groups; restricting a disconnected-group model to its Lie algebra can change multiplicities.

\paragraph{Literature and fate.} Gelfand–Zelevinsky (1984), \S0, Theorem 1, p. 15, explicitly presents this as the earlier model construction from its reference 3. It is class E, not a new theorem of the 1984 talk.

\paragraph{Present status.} Established with the stated algebraic and parity conventions as of 12 November 2026.

\paragraph{Assessment.} Identifies the exact model and the scope of the hidden Lie-algebra action.

\paragraph{Audit.} Literal correctness: correct with missing hypotheses. Publication match: exact. Verification: read the result. Development: annotated. Confidence: high.

\noindent References: \citep{gelfandzelevinsky1984models}.\par

\subsubsection{1983-621.5 --- The symmetric-algebra induction models}

\label{item:1983-621.5}

\noindent\textbf{Kind:} theorem. \textbf{Fate:} A. \textbf{Pages:} 1168.\par

\paragraph{Notebook content.} A dual-looking table replaces exterior algebras and orthogonal factors by symmetric algebras and symplectic factors.

\paragraph{Precise statement.} Write \(S(r)=S^\bullet\mathbb C^r\). The triples \((GL(2r),Sp(2r),S(2r))\), \((O(2r),GL(r),S(r))\), \((Sp(4r),Sp(2r)\times Sp(2r),S(2r)\otimes1)\), and \((Sp(4r+2),Sp(2r+2)\times Sp(2r),S(2r+2)\otimes1)\) give algebraic induced models \(\operatorname{Ind}_M^G\tau\). The inducing representations are infinite-dimensional, interpreted as their algebraic direct sums.

\paragraph{Literature and fate.} Gelfand–Zelevinsky (1984), Theorem 2, p. 15, with proof in \S1, is an exact match to the new table. The omitted odd \(GL\) and odd \(O\) cases are supplied by the function-model version rather than by inventing missing table entries.

\paragraph{Present status.} Established with the stated algebraic and parity conventions as of 12 November 2026.

\paragraph{Assessment.} Identifies the exact model and the scope of the hidden Lie-algebra action.

\paragraph{Audit.} Literal correctness: correct with missing hypotheses. Publication match: exact. Verification: read the result. Development: annotated. Confidence: high.

\noindent References: \citep{gelfandzelevinsky1984models}.\par

\subsubsection{1983-621.6 --- The exterior–symmetric parallelism}

\label{item:1983-621.6}

\noindent\textbf{Kind:} remark. \textbf{Fate:} A. \textbf{Pages:} 1168.\par

\paragraph{Notebook content.} The two model tables are compared by a formal replacement of classical groups and functors.

\paragraph{Precise statement.} At the level of these model tables, the replacements \(GL(r)\mapsto GL(2r)\), \(O(r)\mapsto Sp(2r)\), \(Sp(2r)\mapsto O(4r)\), and \(\Lambda(r)\mapsto S(2r)\) carry the exterior table to the symmetric one. This is a correspondence among the stated constructions, not a general equivalence of their representation categories.

\paragraph{Literature and fate.} Gelfand–Zelevinsky (1984), paragraph following Theorem 2, p. 15, explicitly says that the geometric meaning of this doubling procedure is unclear there. Later notions of super duality cannot be substituted without an exact result match.

\paragraph{Present status.} Established with the stated algebraic and parity conventions as of 12 November 2026.

\paragraph{Assessment.} Identifies the exact model and the scope of the hidden Lie-algebra action.

\paragraph{Audit.} Literal correctness: correct. Publication match: exact. Verification: read the result. Development: annotated. Confidence: high.

\noindent References: \citep{gelfandzelevinsky1984models}.\par

\subsubsection{1983-621.7 --- Gelfand’s polyhedral-model suggestion}

\label{item:1983-621.7}

\noindent\textbf{Kind:} programme. \textbf{Fate:} F. \textbf{Pages:} 1168.\par

\paragraph{Notebook content.} Gelfand suggests a connection with polyhedra.

\paragraph{Precise statement.} No polyhedron, lattice, group action, or proposed correspondence is specified. A precise theorem would require these data; the word “polyhedra” alone does not select a representation-theoretic conjecture.

\paragraph{Literature and fate.} Related earlier Grassmannian and geometric-model records were checked, including 1980-508, but no exact earlier problem occurrence or exact publication match was identified.

\paragraph{Present status.} Insufficiently specified for a mathematical status assignment.

\paragraph{Assessment.} Identifies the exact model and the scope of the hidden Lie-algebra action.

\paragraph{Audit.} Literal correctness: not determinable. Publication match: none found. Verification: unresolved lead. Development: annotated. Confidence: medium.

\subsubsection{1983-621.8 --- Basic Grassmannians for classical groups}

\label{item:1983-621.8}

\noindent\textbf{Kind:} definition. \textbf{Fate:} E. \textbf{Pages:} 1168--1169.\par

\paragraph{Notebook content.} A family of ordinary, isotropic, and Lagrangian Grassmannians is introduced.

\paragraph{Precise statement.} The basic spaces are \(\operatorname{Gr}(r,\mathbb C^{2r})\) for \(GL(2r)\), \(\operatorname{Gr}(r+1,\mathbb C^{2r+1})\) for \(GL(2r+1)\), maximal isotropic \(r\)-planes for a nondegenerate symmetric form in dimensions \(2r\) and \(2r+1\), and Lagrangian \(r\)-planes in a symplectic \(2r\)-space. The even orthogonal space has two connected components; a connected \(SO(2r)\) acts transitively on each component. For a product group, take the product of the basic spaces.

\paragraph{Literature and fate.} Gelfand–Zelevinsky (1984), \S0, pp. 16–17, defines exactly this list. Component choices cannot be dropped from the orthogonal construction.

\paragraph{Present status.} Established with the stated algebraic and parity conventions as of 12 November 2026.

\paragraph{Assessment.} Identifies the exact model and the scope of the hidden Lie-algebra action.

\paragraph{Audit.} Literal correctness: correct with missing hypotheses. Publication match: exact. Verification: read the result. Development: annotated. Confidence: high.

\noindent References: \citep{gelfandzelevinsky1984models}.\par

\subsubsection{1983-621.9 --- The odd–even orthogonal Grassmannian isomorphism}

\label{item:1983-621.9}

\noindent\textbf{Kind:} exercise. \textbf{Fate:} E. \textbf{Pages:} 1169.\par

\paragraph{Notebook content.} An identification of a component of an even orthogonal Grassmannian with an odd orthogonal one is set as an exercise.

\paragraph{Precise statement.} Let \(V\) be a complex quadratic space of dimension \(2r+2\), and let \(H\subset V\) be a nondegenerate hyperplane. For either component of \(\operatorname{OG}(r+1,V)\), the map \(W\mapsto W\cap H\) is an isomorphism onto \(\operatorname{OG}(r,H)\). Each isotropic \(r\)-plane \(U\subset H\) has two maximal isotropic extensions in \(V\), one in each component; they are the two isotropic lines in the nondegenerate quadratic plane \(U^\perp/U\).

\paragraph{Literature and fate.} Gelfand–Zelevinsky (1984), p. 16 and \S1, identifies \(\Gamma(SO(2r+2))\) with \(\Gamma(O(2r+1))\). The quotient-plane argument is the standard elementary proof.

\paragraph{Present status.} Established with the stated algebraic and parity conventions as of 12 November 2026.

\paragraph{Assessment.} Identifies the exact model and the scope of the hidden Lie-algebra action.

\paragraph{Audit.} Literal correctness: correct with missing hypotheses. Publication match: exact. Verification: read the result. Development: annotated. Confidence: high.

\noindent Canonical problem: GP-1983-60.\par

\noindent References: \citep{gelfandzelevinsky1984models}.\par

\subsubsection{1983-621.10 --- Hidden Lie-algebra actions from an open orbit}

\label{item:1983-621.10}

\noindent\textbf{Kind:} method. \textbf{Fate:} A. \textbf{Pages:} 1169.\par

\paragraph{Notebook content.} A model is realized on an open orbit in a homogeneous space for a larger group.

\paragraph{Precise statement.} Let \(Y=L/P\), let \(\mathcal E\) be an \(L\)-equivariant algebraic vector bundle, and let \(\Omega\subset Y\) be an open \(G\)-orbit for \(G\subset L\). Global infinitesimal \(L\)-actions restrict to differential operators on \(\Gamma(\Omega,\mathcal E)\), extending the \(G\)-action even if \(\Omega\) is not stable under the whole group \(L\). If the stabilizer is \(M\) and the fibre representation is \(\tau\), this space is \(\operatorname{Ind}_M^G\tau\).

\paragraph{Literature and fate.} Gelfand–Zelevinsky (1984), \S0, p. 16, defines these restriction models and explains the relation to induction. Theorem 3 gives the concrete exterior-bundle examples.

\paragraph{Present status.} Established with the stated algebraic and parity conventions as of 12 November 2026.

\paragraph{Assessment.} Identifies the exact model and the scope of the hidden Lie-algebra action.

\paragraph{Audit.} Literal correctness: correct with missing hypotheses. Publication match: exact. Verification: read the result. Development: annotated. Confidence: high.

\noindent References: \citep{gelfandzelevinsky1984models}.\par

\subsubsection{1983-621.11 --- The mentioned even differential}

\label{item:1983-621.11}

\noindent\textbf{Kind:} remark. \textbf{Fate:} F. \textbf{Pages:} 1169.\par

\paragraph{Notebook content.} An even differential is mentioned in the supergeometric discussion.

\paragraph{Precise statement.} The surviving remark does not determine the operator, its domain, its square, or a cohomology group. It therefore does not specify a differential complex whose properties can be checked.

\paragraph{Literature and fate.} The matched model paper verifies the surrounding bundle and supergroup constructions but does not identify an exact theorem from this isolated formulation.

\paragraph{Present status.} Unidentifiable as a separate precise differential construction.

\paragraph{Assessment.} Identifies the exact model and the scope of the hidden Lie-algebra action.

\paragraph{Audit.} Literal correctness: not determinable. Publication match: thematic only. Verification: unresolved lead. Development: annotated. Confidence: medium.

\noindent References: \citep{gelfandzelevinsky1984models}.\par

\subsubsection{1983-621.12 --- Exterior-bundle functions and supergroup models}

\label{item:1983-621.12}

\noindent\textbf{Kind:} construction. \textbf{Fate:} A. \textbf{Pages:} 1169--1170.\par

\paragraph{Notebook content.} Sections of an exterior algebra bundle are interpreted as functions on a superspace, with larger supergroup symmetries.

\paragraph{Precise statement.} For a vector bundle \(E\to Y\), the split superspace with odd bundle \(E\) has function sheaf \(\Lambda^\bullet E^*\). Thus a model realized by sections of \(\Lambda^\bullet U\) is represented by the odd bundle \(U^*\). The corresponding overgroups in the recorded classical cases are \(OSp(2|2r)\) over \(GL(r)\), \(GL(1|r)\) over \(O(r)\), and \(OSp(2|2r)\times Sp(2r)\) over \(Sp(2r)\), with the basic super-Grassmannians and their open \(G\)-orbits.

\paragraph{Editorial reconstruction.} The dual on the odd bundle fixes the function-sheaf convention. Verification here means the exact published statement was read, not that a deferred proof was inspected.

\paragraph{Literature and fate.} Gelfand–Zelevinsky (1984), Theorem 3′, pp. 17–18, states this supergroup reformulation and its hidden-superalgebra corollary. The paper explicitly defers the proof of Theorem 3′ to another work; reading this announcement is not reading that deferred proof.

\paragraph{Present status.} Established with the stated algebraic and parity conventions as of 12 November 2026.

\paragraph{Assessment.} Identifies the exact model and the scope of the hidden Lie-algebra action.

\paragraph{Audit.} Literal correctness: correct with missing hypotheses. Publication match: exact. Verification: read the result. Development: annotated. Confidence: high.

\noindent References: \citep{gelfandzelevinsky1984models}.\par

\subsubsection{1983-621.13 --- The phrase almost irreducible}

\label{item:1983-621.13}

\noindent\textbf{Kind:} claim. \textbf{Fate:} F. \textbf{Pages:} 1170.\par

\paragraph{Notebook content.} The larger representation is described as almost irreducible.

\paragraph{Precise statement.} The record gives no definition of “almost irreducible”: it could refer to finite length, a finite-dimensional exception, a quotient, or a parameter-generic condition. None of those meanings is selected by the phrase alone.

\paragraph{Literature and fate.} The checked Gelfand–Zelevinsky paper constructs the hidden actions and states that the resulting larger representations will be studied separately. It does not justify replacing the notebook phrase by a particular irreducibility theorem.

\paragraph{Present status.} The literal phrase has no uniquely determined modern assertion.

\paragraph{Assessment.} Identifies the exact model and the scope of the hidden Lie-algebra action.

\paragraph{Audit.} Literal correctness: not determinable. Publication match: thematic only. Verification: unresolved lead. Development: annotated. Confidence: medium.

\noindent References: \citep{gelfandzelevinsky1984models}.\par

\subsubsection{1983-621.14 --- Function models and the intentional half-integer subgroup notation}

\label{item:1983-621.14}

\noindent\textbf{Kind:} theorem. \textbf{Fate:} A. \textbf{Pages:} 1170.\par

\paragraph{Notebook content.} Models on functions are listed using \(GL(r-\tfrac12)\) and odd-indexed symplectic groups.

\paragraph{Precise statement.} Define \(GL(r-\tfrac12)\subset GL(r)\) and \(Sp(2r-1)\subset Sp(2r)\) as stabilizers of a nonzero linear functional in the defining module. Then the function-model stabilizers are \(M_0=Sp(n-1)\) for \(G=GL(n)\), \(M_0=GL((n-1)/2)\) for \(G=O(n)\), and \(M_0=Sp(n-1)\times Sp(n)\) for \(G=Sp(2n)\), using those stabilizer conventions. In each case \(\mathcal O(G/M_0)=\operatorname{Ind}_{M_0}^G1\) is a model.

\paragraph{Literature and fate.} Gelfand–Zelevinsky (1984), Theorem 2′, pp. 15–16, defines exactly this notation; it is not an OCR error or a new reductive group of fractional rank. Theorem 4 realizes the function models using overgroups \(SO(2n+2)\), orthogonal products, and \(GL(2n+1)\), respectively.

\paragraph{Present status.} Established with the stated algebraic and parity conventions as of 12 November 2026.

\paragraph{Assessment.} Identifies the exact model and the scope of the hidden Lie-algebra action.

\paragraph{Audit.} Literal correctness: correct. Publication match: exact. Verification: read the result. Development: annotated. Confidence: high.

\noindent References: \citep{gelfandzelevinsky1984models}.\par

\subsubsection{1983-621.15 --- Induction in stages and the distinction from the super construction}

\label{item:1983-621.15}

\noindent\textbf{Kind:} method. \textbf{Fate:} A. \textbf{Pages:} 1170.\par

\paragraph{Notebook content.} The report passes from an induced representation with a fibre to a model on functions.

\paragraph{Precise statement.} If \(M_0\subset M\subset G\) and \(\tau\cong\operatorname{Ind}_{M_0}^{M}1\) in the algebraic induction setting, then \(\operatorname{Ind}_{M}^{G}\tau\cong\operatorname{Ind}_{M_0}^{G}1\). This accounts for the function realization of the symmetric-algebra models. It is separate from interpreting a finite exterior-algebra fibre as the odd function algebra of a superspace.

\paragraph{Literature and fate.} Gelfand–Zelevinsky (1984), paragraph before Theorem 2′, and the discussion following Theorems 3–4, explicitly distinguish the two constructions.

\paragraph{Present status.} Established with the stated algebraic and parity conventions as of 12 November 2026.

\paragraph{Assessment.} Identifies the exact model and the scope of the hidden Lie-algebra action.

\paragraph{Audit.} Literal correctness: correct. Publication match: exact. Verification: read the result. Development: annotated. Confidence: high.

\noindent References: \citep{gelfandzelevinsky1984models}.\par

\subsubsection{1983-621.16 --- Further hidden symmetries}

\label{item:1983-621.16}

\noindent\textbf{Kind:} question. \textbf{Fate:} F. \textbf{Pages:} 1170.\par

\paragraph{Notebook content.} The discussion asks what other symmetries may occur.

\paragraph{Precise statement.} A classification question would need to specify the model, the class of allowable differential operators or overgroups, and an equivalence relation. The notebook does not supply these restrictions for the general question.

\paragraph{Literature and fate.} The paper constructs the displayed families, but its examples do not answer every possible enlargement. No exact later result matching an unrestricted “other symmetries” problem was verified.

\paragraph{Present status.} Insufficiently specified to label open or solved as a single problem.

\paragraph{Assessment.} Identifies the exact model and the scope of the hidden Lie-algebra action.

\paragraph{Audit.} Literal correctness: not determinable. Publication match: thematic only. Verification: unresolved lead. Development: annotated. Confidence: medium.

\noindent References: \citep{gelfandzelevinsky1984models}.\par

\subsubsection{1983-621.17 --- Gelfand’s volume suggestion}

\label{item:1983-621.17}

\noindent\textbf{Kind:} programme. \textbf{Fate:} F. \textbf{Pages:} 1168.\par

\paragraph{Notebook content.} A connection with volume is suggested alongside the model constructions.

\paragraph{Precise statement.} No region, metric, measure, or formula relating volume to the representation model is specified.

\paragraph{Literature and fate.} Earlier geometric-volume discussions, including 1980-508, are thematic precedents only. No exact problem occurrence or result match was established.

\paragraph{Present status.} The suggestion does not determine a numerical or geometric proposition.

\paragraph{Assessment.} Identifies the exact model and the scope of the hidden Lie-algebra action.

\paragraph{Audit.} Literal correctness: not determinable. Publication match: none found. Verification: unresolved lead. Development: annotated. Confidence: medium.

\subsubsection{1983-621.18 --- Gelfand’s dilogarithm suggestion}

\label{item:1983-621.18}

\noindent\textbf{Kind:} programme. \textbf{Fate:} F. \textbf{Pages:} 1168.\par

\paragraph{Notebook content.} Dilogarithms are mentioned as a possible connection.

\paragraph{Precise statement.} No dilogarithm identity, arguments, branch convention, or representation-theoretic construction is recorded.

\paragraph{Literature and fate.} The earlier index contains related Grassmannian and dilogarithm discussions, but these do not identify the intended assertion. No new problem ID is assigned.

\paragraph{Present status.} The short suggestion remains mathematically underspecified.

\paragraph{Assessment.} Identifies the exact model and the scope of the hidden Lie-algebra action.

\paragraph{Audit.} Literal correctness: not determinable. Publication match: thematic only. Verification: unresolved lead. Development: annotated. Confidence: medium.

\subsubsection{1983-621.19 --- Gelfand’s localization suggestion}

\label{item:1983-621.19}

\noindent\textbf{Kind:} programme. \textbf{Fate:} F. \textbf{Pages:} 1168.\par

\paragraph{Notebook content.} Localization is mentioned as a further direction.

\paragraph{Precise statement.} A localization statement needs a category, a space or algebra to localize, a parameter, and a comparison functor. The record does not identify these data.

\paragraph{Literature and fate.} The geometric model paper supplies a setting for infinitesimal actions, but that construction is not automatically a localization equivalence. No exact independent localization theorem was matched to this remark.

\paragraph{Present status.} No uniquely stated localization problem is recoverable.

\paragraph{Assessment.} Identifies the exact model and the scope of the hidden Lie-algebra action.

\paragraph{Audit.} Literal correctness: not determinable. Publication match: thematic only. Verification: unresolved lead. Development: annotated. Confidence: medium.

\noindent References: \citep{gelfandzelevinsky1984models}.\par

\subsection{Record 622: Hypergeometric equations, discrete orthogonal polynomials, and angular momentum}

\label{talk:622}

\noindent\textbf{Date:} 1984-02-27. \textbf{Speaker:} A. F. Nikiforov. \textbf{Pages:} 1170 lower--1171.\par

Nikiforov’s heading begins in the lower part of P001170 and the report ends on P001171. The catalogue’s “F. Nikiforov” is identified as A. F. Nikiforov. The material matches Nikiforov–Suslov–Uvarov’s 1985 paper and 1991 English monograph; their cited preprints from 1982 and 1983 show that the main discrete-polynomial correspondences already predated this seminar.

\subsubsection{1983-622.1 --- Polynomial eigenfunctions of hypergeometric type}

\label{item:1983-622.1}

\noindent\textbf{Kind:} result. \textbf{Fate:} E. \textbf{Pages:} 1170.\par

\paragraph{Notebook content.} The equation with quadratic and linear coefficients is used to characterize polynomial solutions.

\paragraph{Precise statement.} Let \(\deg\sigma\le2\), \(\deg\tau\le1\), and \(L=\sigma(x)d^2/dx^2+\tau(x)d/dx\). It preserves each polynomial degree filtration. A degree-\(n\) eigenpolynomial satisfying \(Ly_n+\lambda_ny_n=0\) must have \(\lambda_n=-n\tau'-n(n-1)\sigma''/2\). If this eigenvalue differs from those in smaller degrees, there is a unique monic degree-\(n\) eigenpolynomial. Orthogonality additionally requires a positive Pearson weight and vanishing boundary terms; polynomial solvability alone does not imply it.

\paragraph{Editorial reconstruction.} The eigenvalue is obtained from the coefficient of the highest monomial, and uniqueness follows by solving the lower triangular coefficient system. The read preface supplies the framework, not this full hypothesis-complete statement.

\paragraph{Literature and fate.} Nikiforov–Suslov–Uvarov (1991), preface p. V, equation (0.1), gives the hypergeometric differential equation. The displayed eigenvalue and monic uniqueness criterion are our elementary reconstruction from its triangular action on monomials; no Chapter 1 theorem or proof was read.

\paragraph{Present status.} Established under the stated regularity and nondegeneracy hypotheses as of 12 November 2026.

\paragraph{Assessment.} Fixes the operator normalization and the hypotheses needed for polynomial or orthogonal solutions.

\paragraph{Audit.} Literal correctness: correct with missing hypotheses. Publication match: partial. Verification: indirect. Development: annotated. Confidence: high.

\noindent References: \citep{nikiforovsuslovuvarov1991classical}.\par

\subsubsection{1983-622.2 --- Pearson weights and Rodrigues’ formula}

\label{item:1983-622.2}

\noindent\textbf{Kind:} construction. \textbf{Fate:} E. \textbf{Pages:} 1170.\par

\paragraph{Notebook content.} A weight and a Rodrigues-type formula are written.

\paragraph{Precise statement.} If \((\sigma\rho)'=\tau\rho\), then, when the expression is nonzero of degree \(n\), \(y_n(x)=B_n\rho(x)^{-1}\frac{d^n}{dx^n}(\sigma(x)^n\rho(x))\) solves the corresponding polynomial equation. To obtain an orthogonal family on an interval, require \(\rho>0\), the relevant finite moments, and boundary terms that vanish under the integrations by parts.

\paragraph{Editorial reconstruction.} This is the standard Pearson–Rodrigues construction, stated with explicit nonvanishing and integration-by-parts conditions. The cited read passage identifies the construction but does not print the full formula or proof.

\paragraph{Literature and fate.} Nikiforov–Suslov–Uvarov (1991), preface p. V, in the paragraph following equation (0.1), announces the Rodrigues construction and subsequent orthogonality argument. The explicit formula and its endpoint hypotheses were reconstructed from the Pearson equation; their Chapter 1 presentation was not read.

\paragraph{Present status.} Established under the stated regularity and nondegeneracy hypotheses as of 12 November 2026.

\paragraph{Assessment.} Fixes the operator normalization and the hypotheses needed for polynomial or orthogonal solutions.

\paragraph{Audit.} Literal correctness: correct with missing hypotheses. Publication match: partial. Verification: indirect. Development: annotated. Confidence: high.

\noindent References: \citep{nikiforovsuslovuvarov1991classical}.\par

\subsubsection{1983-622.3 --- The noninteger Rodrigues contour expression}

\label{item:1983-622.3}

\noindent\textbf{Kind:} method. \textbf{Fate:} F. \textbf{Pages:} 1170.\par

\paragraph{Notebook content.} The order of differentiation is allowed to become nonintegral and a contour-integral expression is suggested.

\paragraph{Precise statement.} A contour realization for nonintegral order needs a branch of the powers, a contour or cycle, endpoint conditions, and convergence assumptions. These data are not fixed in the notebook, so there is no unique formula to which a universal solution statement can be attached.

\paragraph{Literature and fate.} Earlier index entries 1968-160.16 and 1981-556.10 discuss related fractional Rodrigues or Euler-integral ideas. They are thematic precedents, not an exact repeated problem. No specific cycle formula was verified here.

\paragraph{Present status.} The nonintegral construction is underspecified as recorded.

\paragraph{Assessment.} Fixes the operator normalization and the hypotheses needed for polynomial or orthogonal solutions.

\paragraph{Audit.} Literal correctness: not determinable. Publication match: thematic only. Verification: unresolved lead. Development: annotated. Confidence: medium.

\noindent References: \citep{nikiforovsuslovuvarov1991classical}.\par

\subsubsection{1983-622.4 --- Stability under differentiation}

\label{item:1983-622.4}

\noindent\textbf{Kind:} result. \textbf{Fate:} E. \textbf{Pages:} 1170.\par

\paragraph{Notebook content.} Differentiating a polynomial solution is said to preserve the same kind of equation.

\paragraph{Precise statement.} If \(y\) solves \(\sigma y''+\tau y'+\lambda y=0\), with \(\deg\sigma\le2\), \(\deg\tau\le1\), then \(v=y^{(m)}\) satisfies \(\sigma v''+(\tau+m\sigma')v'+(\lambda+m\tau'+m(m-1)\sigma''/2)v=0\). The transformed Pearson weight is \(\rho_m=\sigma^m\rho\) where the original Pearson equation holds. For a polynomial of degree \(n\), take \(0\le m\le n\).

\paragraph{Editorial reconstruction.} Differentiating the equation m times leaves only the first two derivatives of the quadratic coefficient and the first derivative of the linear coefficient. Differentiating sigma times the transformed weight verifies its Pearson equation.

\paragraph{Literature and fate.} Nikiforov–Suslov–Uvarov (1991), preface p. V, states that derivatives remain of hypergeometric type. The displayed shifted coefficients follow from our direct Leibniz-rule calculation; no Chapter 1 proof or displayed parameter formula was read.

\paragraph{Present status.} Established under the stated regularity and nondegeneracy hypotheses as of 12 November 2026.

\paragraph{Assessment.} Fixes the operator normalization and the hypotheses needed for polynomial or orthogonal solutions.

\paragraph{Audit.} Literal correctness: correct. Publication match: partial. Verification: indirect. Development: annotated. Confidence: high.

\noindent References: \citep{nikiforovsuslovuvarov1991classical}.\par

\subsubsection{1983-622.5 --- Several-variable hypergeometric analogues}

\label{item:1983-622.5}

\noindent\textbf{Kind:} programme. \textbf{Fate:} F. \textbf{Pages:} 1170.\par

\paragraph{Notebook content.} The discussion asks about analogues in several variables.

\paragraph{Precise statement.} No system of partial differential or difference equations, parameter configuration, integration cycle, or desired classification is specified. Classical several-variable functions and later generalized hypergeometric systems are different possible targets.

\paragraph{Literature and fate.} The earlier index entry 1981-556.11 contains a similarly broad programme, but no exact problem identity can be certified. Later hypergeometric theories are not treated as an answer to this unformulated question.

\paragraph{Present status.} Insufficiently specified to assign one solved/open status.

\paragraph{Assessment.} Fixes the operator normalization and the hypotheses needed for polynomial or orthogonal solutions.

\paragraph{Audit.} Literal correctness: not determinable. Publication match: thematic only. Verification: unresolved lead. Development: annotated. Confidence: medium.

\subsubsection{1983-622.6 --- The discrete hypergeometric operator and step size}

\label{item:1983-622.6}

\noindent\textbf{Kind:} construction. \textbf{Fate:} E. \textbf{Pages:} 1171.\par

\paragraph{Notebook content.} The differential equation is replaced by a second-order difference equation.

\paragraph{Precise statement.} For step \(h\ne0\), set \(\Delta_hf=(f(x+h)-f(x))/h\) and \(\nabla_hf=(f(x)-f(x-h))/h\). The operator \(\sigma\Delta_h\nabla_h+\tau\Delta_h\), with degrees at most two and one, preserves polynomial degree. If it is written in the symmetric form \(\widetilde\sigma\Delta_h\nabla_h+\frac12\widetilde\tau(\Delta_h+\nabla_h)\), then \(\sigma=\widetilde\sigma-\frac h2\widetilde\tau\) and \(\tau=\widetilde\tau\). The factor \(h\) disappears only for unit step.

\paragraph{Editorial reconstruction.} The read preface directly verifies the symmetric equation. The explicit normalized-step convention and conversion are an elementary annotation, not a claim to have read Chapter 2.

\paragraph{Literature and fate.} Nikiforov–Suslov–Uvarov (1991), preface p. V, equation (0.2), displays the symmetric step-\(h\) difference equation and was read. The forward-difference conversion is our algebraic calculation from \(\Delta_h-\nabla_h=h\Delta_h\nabla_h\). No Chapter 2 theorem or proof was read.

\paragraph{Present status.} Established under the stated regularity and nondegeneracy hypotheses as of 12 November 2026.

\paragraph{Assessment.} Fixes the operator normalization and the hypotheses needed for polynomial or orthogonal solutions.

\paragraph{Audit.} Literal correctness: correct with missing hypotheses. Publication match: partial. Verification: read the result. Development: annotated. Confidence: high.

\noindent References: \citep{nikiforovsuslovuvarov1991classical}.\par

\subsubsection{1983-622.7 --- Clebsch–Gordan coefficients and Hahn polynomials}

\label{item:1983-622.7}

\noindent\textbf{Kind:} theorem. \textbf{Fate:} E. \textbf{Pages:} 1171.\par

\paragraph{Notebook content.} The report identifies angular-momentum coupling coefficients with discrete orthogonal polynomials.

\paragraph{Precise statement.} Let \(j_1,j_2\) be nonnegative half-integers and fix an admissible \(m\ge|j_1-j_2|\). Set \(N=j_1+j_2-m+1\), \(x=j_1-m_1\), \(\alpha=m+j_1-j_2\), \(\beta=m-j_1+j_2\), and \(n=j-m\). Then \(x,n\in\{0,\ldots,N-1\}\). For the Hahn polynomial \(y_n\) orthogonal with weight \(\rho(x)=\Gamma(\beta+x+1)\Gamma(\alpha+N-x)/(\Gamma(x+1)\Gamma(N-x))\), normalized by \(d_n^2=\sum_x\rho(x)y_n(x)^2\), one has \(\left|\langle j_1m_1,j_2(m-m_1)\mid jm\rangle\right|^2=\rho(x)y_n(x)^2/d_n^2\). Squaring removes the phase convention.

\paragraph{Literature and fate.} Nikiforov–Suslov–Uvarov (1985), \S2.2, equation (40), Russian pp. 31–32, is the exact match. Its reference 27 identifies the corresponding 1982 preprint, so the result is class E despite the 1985 journal publication.

\paragraph{Present status.} Established under the stated regularity and nondegeneracy hypotheses as of 12 November 2026.

\paragraph{Assessment.} Fixes the operator normalization and the hypotheses needed for polynomial or orthogonal solutions.

\paragraph{Audit.} Literal correctness: correct with missing hypotheses. Publication match: exact. Verification: read the result. Development: annotated. Confidence: high.

\noindent References: \citep{nikiforovsuslovuvarov1985rotation,nikiforovsuslovuvarov1991classical}.\par

\subsubsection{1983-622.8 --- The Hahn weight and its missing variable}

\label{item:1983-622.8}

\noindent\textbf{Kind:} formula. \textbf{Fate:} E. \textbf{Pages:} 1171.\par

\paragraph{Notebook content.} The displayed gamma-function weight has a term independent of the lattice variable where the symmetric endpoint factor should occur.

\paragraph{Precise statement.} For \(N\in\mathbb Z_{>0}\), \(\alpha,\beta>-1\), and \(x=0,\ldots,N-1\), the positive Hahn weight is proportional to \(\rho(x)=\Gamma(\beta+x+1)\Gamma(\alpha+N-x)/(\Gamma(x+1)\Gamma(N-x))\). Replacing the second numerator by \(\Gamma(\alpha+N-1)\) changes the weight and does not give this Hahn family.

\paragraph{Literature and fate.} Nikiforov–Suslov–Uvarov (1985), table 1, Russian p. 29, gives the exact endpoint factor \(\alpha+N-x\). The P001171 scan was inspected and the suspect term was checked against it.

\paragraph{Present status.} Established under the stated regularity and nondegeneracy hypotheses as of 12 November 2026.

\paragraph{Assessment.} Fixes the operator normalization and the hypotheses needed for polynomial or orthogonal solutions.

\paragraph{Audit.} Literal correctness: false as written. Publication match: exact. Verification: read the result. Development: annotated. Confidence: high.

\noindent References: \citep{nikiforovsuslovuvarov1985rotation}.\par

\subsubsection{1983-622.9 --- Quadratic lattices and Racah recoupling coefficients}

\label{item:1983-622.9}

\noindent\textbf{Kind:} construction. \textbf{Fate:} E. \textbf{Pages:} 1171.\par

\paragraph{Notebook content.} The discussion passes to nonuniform lattices and the coefficients for recoupling three angular momenta.

\paragraph{Precise statement.} For admissible half-integers \(j_1,j_2,j_3,J\), the change between the two coupling bases has entries \(U_{j_{12},j_{23}}=(-1)^{j_1+j_2+j_3+J}\sqrt{(2j_{12}+1)(2j_{23}+1)}\left\{\begin{smallmatrix}j_1&j_2&j_{12}\\j_3&J&j_{23}\end{smallmatrix}\right\}\). This is an orthogonal matrix after standard phase choices. As the intermediate label varies over its finite triangle-inequality interval, its rows are normalized Racah polynomials on a quadratic lattice \(x(x+\gamma+\delta+1)\), with parameters determined by the four fixed angular momenta.

\paragraph{Editorial reconstruction.} The recoupling matrix is stated with all admissibility data; no single parameter table is imposed without choosing an ordering of the four angular momenta.

\paragraph{Literature and fate.} Nikiforov–Suslov–Uvarov (1985), \S2.3, equation (44), Russian p. 33, gives the exact Racah correspondence. References 12 and 14 identify the 1982 Racah preprint and the 1983 nonuniform-lattice preprint. The 1991 monograph develops the latter in Chapter 3 and recoupling in \S5.3.

\paragraph{Present status.} Established under the stated regularity and nondegeneracy hypotheses as of 12 November 2026.

\paragraph{Assessment.} Fixes the operator normalization and the hypotheses needed for polynomial or orthogonal solutions.

\paragraph{Audit.} Literal correctness: correct with missing hypotheses. Publication match: exact. Verification: read the result. Development: annotated. Confidence: high.

\noindent References: \citep{nikiforovsuslovuvarov1985rotation,nikiforovsuslovuvarov1991classical}.\par

\subsection{Record 623: The half-line Toda chain by inverse spectral reconstruction}

\label{talk:623}

\noindent\textbf{Date:} 1984-03-05. \textbf{Speaker:} Yu. M. Berezanskii. \textbf{Pages:} 1172--1173 upper.\par

The report begins on P001172 and ends in the upper part of P001173; Fradkin begins below it. Its bounded half-line inverse-spectral theorem closely matches Berezanskii’s 1985 note, received 9 April 1984. That note says the approach had already been presented in Kyiv in October 1983, so the seminar is not asserted to be the first public announcement.

\subsubsection{1983-623.1 --- The normalized half-line Toda system}

\label{item:1983-623.1}

\noindent\textbf{Kind:} definition. \textbf{Fate:} E. \textbf{Pages:} 1172.\par

\paragraph{Notebook content.} The report writes a Toda-type differential system for Jacobi coefficients.

\paragraph{Precise statement.} For \(n\ge0\), impose \(a_n(t)>0\), \(b_n(t)\in\mathbb R\), \(a_{-1}(t)=0\), and \(\dot a_n=\frac12a_n(b_{n+1}-b_n)\), \(\dot b_n=a_n^2-a_{n-1}^2\). This normalization corresponds to multiplying the spectral measure by \(e^{t\lambda}\). The convention without the factor \(1/2\) in the first equation simultaneously doubles the second equation and uses \(e^{2t\lambda}\).

\paragraph{Literature and fate.} Berezanskii (1985), \S2, equation (4), with its time factor \(f=1\), gives the precise system. The incomplete or inconsistent notebook formula must be corrected as a pair of equations, not by changing only one coefficient.

\paragraph{Present status.} Established under the stated regularity and nondegeneracy hypotheses as of 12 November 2026.

\paragraph{Assessment.} Fixes the operator normalization and the hypotheses needed for polynomial or orthogonal solutions.

\paragraph{Audit.} Literal correctness: false as written. Publication match: exact. Verification: read the result. Development: annotated. Confidence: high.

\noindent References: \citep{berezanskii1985integration}.\par

\subsubsection{1983-623.2 --- Bounded Jacobi matrices and their spectral measure}

\label{item:1983-623.2}

\noindent\textbf{Kind:} construction. \textbf{Fate:} E. \textbf{Pages:} 1172.\par

\paragraph{Notebook content.} The coefficients are assembled into a tridiagonal operator and a spectral measure.

\paragraph{Precise statement.} For bounded real sequences \(b_n\) and bounded positive \(a_n\), the Jacobi matrix \(Je_n=a_{n-1}e_{n-1}+b_ne_n+a_ne_{n+1}\) is a bounded self-adjoint operator on \(\ell^2(\mathbb Z_{\ge0})\). The vector \(e_0\) is cyclic, and its spectral measure is a compactly supported probability measure with infinite support. Conversely such a measure, with positive leading-coefficient orthonormal polynomials, determines the Jacobi coefficients.

\paragraph{Literature and fate.} Berezanskii (1985), \S1, equations (1)–(2), gives the inverse spectral construction; boundedness makes self-adjointness and moment determinacy automatic here.

\paragraph{Present status.} Established under the stated regularity and nondegeneracy hypotheses as of 12 November 2026.

\paragraph{Assessment.} Fixes the operator normalization and the hypotheses needed for polynomial or orthogonal solutions.

\paragraph{Audit.} Literal correctness: correct with missing hypotheses. Publication match: exact. Verification: read the result. Development: annotated. Confidence: high.

\noindent References: \citep{berezanskii1985integration}.\par

\subsubsection{1983-623.3 --- Exponential spectral evolution and the bounded solution}

\label{item:1983-623.3}

\noindent\textbf{Kind:} theorem. \textbf{Fate:} A. \textbf{Pages:} 1172--1173.\par

\paragraph{Notebook content.} Evolution is reduced to multiplying the measure by an exponential and then reconstructing the coefficients.

\paragraph{Precise statement.} For bounded initial Jacobi coefficients with \(a_n(0)>0\), let \(\rho_0\) be the spectral probability measure. Put \(d\rho_t(\lambda)=e^{t\lambda}d\rho_0(\lambda)/\int e^{t\mu}d\rho_0(\mu)\). Orthogonalizing \(1,\lambda,\lambda^2,\ldots\) for \(\rho_t\) and reconstructing \(a_n(t),b_n(t)\) gives the unique \(C^1\) solution of the normalized half-line Toda system in the class bounded in \(n\), uniformly on compact time intervals. Compact support makes the construction valid for every real \(t\).

\paragraph{Editorial reconstruction.} The classification records the near-contemporary publication by the speaker, while explicitly declining a first-announcement claim.

\paragraph{Literature and fate.} Berezanskii (1985), \S3, equation (7), and \S5, Theorem, are an exact match. Received 9 April 1984, published in Doklady 281(1), 16–19 (1985); the introduction records earlier Kyiv circulation in October 1983.

\paragraph{Present status.} Established under the stated regularity and nondegeneracy hypotheses as of 12 November 2026.

\paragraph{Assessment.} Fixes the operator normalization and the hypotheses needed for polynomial or orthogonal solutions.

\paragraph{Audit.} Literal correctness: correct with missing hypotheses. Publication match: exact. Verification: read the result. Development: annotated. Confidence: high.

\noindent References: \citep{berezanskii1985integration}.\par

\subsubsection{1983-623.4 --- Recovering the coefficients from orthogonal polynomials}

\label{item:1983-623.4}

\noindent\textbf{Kind:} method. \textbf{Fate:} E. \textbf{Pages:} 1173.\par

\paragraph{Notebook content.} Integral formulas recover the two Jacobi coefficient sequences.

\paragraph{Precise statement.} With \(P_n(\lambda;t)\) orthonormal for \(\rho_t\), \(P_0=1\), and positive leading coefficients, \(a_n(t)=\int\lambda P_n(\lambda;t)P_{n+1}(\lambda;t)\,d\rho_t(\lambda)\) and \(b_n(t)=\int\lambda P_n(\lambda;t)^2\,d\rho_t(\lambda)\). If \(D_n=\det(s_{j+k})_{j,k=0}^n\), \(D_{-1}=1\), then \(a_n=\sqrt{D_{n-1}D_{n+1}}/D_n\), provided all these Hankel determinants are positive.

\paragraph{Literature and fate.} Berezanskii (1985), \S1, equation (2), and \S4, gives these reconstruction formulas and the determinant method used to verify the Toda equations.

\paragraph{Present status.} Established under the stated regularity and nondegeneracy hypotheses as of 12 November 2026.

\paragraph{Assessment.} Fixes the operator normalization and the hypotheses needed for polynomial or orthogonal solutions.

\paragraph{Audit.} Literal correctness: correct. Publication match: exact. Verification: read the result. Development: annotated. Confidence: high.

\noindent References: \citep{berezanskii1985integration}.\par

\subsubsection{1983-623.5 --- Hankel positivity and infinite support}

\label{item:1983-623.5}

\noindent\textbf{Kind:} criterion. \textbf{Fate:} E. \textbf{Pages:} 1172.\par

\paragraph{Notebook content.} Moments and positivity of their Hankel matrices are discussed.

\paragraph{Precise statement.} For a positive Borel measure on \(\mathbb R\) with all moments \(s_k\), every Hankel matrix \((s_{j+k})_{j,k=0}^n\) is positive semidefinite because its quadratic form is \(\int|\sum_{j=0}^nc_j\lambda^j|^2d\rho\). It is positive definite for every \(n\) exactly when the support is infinite. A normalized strictly positive moment sequence admits the Jacobi reconstruction, but uniqueness of a representing measure is a further determinacy issue.

\paragraph{Literature and fate.} The moment and determinant setup is used explicitly in Berezanskii (1985), \S1. The strict-versus-semidefinite distinction is the elementary polynomial-zero argument; finite support leads to a terminating recurrence.

\paragraph{Present status.} Established under the stated regularity and nondegeneracy hypotheses as of 12 November 2026.

\paragraph{Assessment.} Fixes the operator normalization and the hypotheses needed for polynomial or orthogonal solutions.

\paragraph{Audit.} Literal correctness: correct with missing hypotheses. Publication match: partial. Verification: read the result. Development: annotated. Confidence: high.

\noindent References: \citep{berezanskii1985integration}.\par

\subsubsection{1983-623.6 --- Determinacy, self-adjointness, and growth}

\label{item:1983-623.6}

\noindent\textbf{Kind:} criterion. \textbf{Fate:} E. \textbf{Pages:} 1172--1173.\par

\paragraph{Notebook content.} The text links uniqueness of the measure to growth of the Jacobi coefficients.

\paragraph{Precise statement.} For a Jacobi matrix with \(a_n>0\), \(b_n\in\mathbb R\), the Hamburger moment problem is determinate precisely when the minimal Jacobi operator on finitely supported sequences is essentially self-adjoint. A standard sufficient condition is \(\sum_{n\ge0}a_n^{-1}=\infty\). Bounded \(a_n\) satisfies it. This is a sufficient slow-growth condition, not an equivalence with an arbitrary informal growth bound.

\paragraph{Editorial reconstruction.} The phrase about growth is interpreted as sufficient control. A reversed “not too slowly” reading would contradict the inspected handwritten text and the operator-theoretic condition.

\paragraph{Literature and fate.} Berezanskii (1985), \S1, explicitly assumes self-adjointness and describes the coefficients as not growing too rapidly. The full general moment-theory proof and the sharp necessity of any particular growth test were not read in that short note; this assessment is therefore indirect beyond its stated hypothesis.

\paragraph{Present status.} Established under the stated regularity and nondegeneracy hypotheses as of 12 November 2026.

\paragraph{Assessment.} Fixes the operator normalization and the hypotheses needed for polynomial or orthogonal solutions.

\paragraph{Audit.} Literal correctness: correct with missing hypotheses. Publication match: partial. Verification: indirect. Development: annotated. Confidence: high.

\noindent References: \citep{berezanskii1985integration}.\par

\subsubsection{1983-623.7 --- Extending the inverse-spectral evolution to unbounded coefficients}

\label{item:1983-623.7}

\noindent\textbf{Kind:} programme. \textbf{Fate:} F. \textbf{Pages:} 1173.\par

\paragraph{Notebook content.} The speaker indicates that the boundedness restriction can be weakened.

\paragraph{Precise statement.} For an unbounded Jacobi operator, self-adjointness, existence of exponential moments \(\int e^{t\lambda}d\rho_0\), persistence of the operator domain, and the uniqueness class must be specified. The record does not give a maximal admissible class or a theorem valid for every unbounded initial sequence.

\paragraph{Literature and fate.} Berezanskii (1985), the paragraph after the \S5 theorem, explicitly says the assertion extends with changes to unbounded self-adjoint operators, but does not supply those full hypotheses there. This verifies the direction of the remark without identifying a uniquely stated open problem.

\paragraph{Present status.} The unbounded extension is announced, but the maximal scope cannot be determined from the recorded formulation.

\paragraph{Assessment.} Fixes the operator normalization and the hypotheses needed for polynomial or orthogonal solutions.

\paragraph{Audit.} Literal correctness: not determinable. Publication match: partial. Verification: read the result. Development: annotated. Confidence: medium.

\noindent References: \citep{berezanskii1985integration}.\par

\subsection{Record 624: Gauge unification and supergravity}

\label{talk:624}

\noindent\textbf{Date:} 1984-03-05. \textbf{Speaker:} E. S. Fradkin. \textbf{Pages:} 1173 lower--1174.\par

The actual report begins in the lower part of P001173, after Berezanskii, and continues through P001174. Catalogue overlap does not assign the Toda discussion to Fradkin. The report combines standard gauge constructions, a tentative \(E_8\) unification programme, bibliographic pointers, and Gelfand’s question about an infinite-dimensional geometric space.

\subsubsection{1983-624.1 --- Unification by a larger gauge group}

\label{item:1983-624.1}

\noindent\textbf{Kind:} programme. \textbf{Fate:} F. \textbf{Pages:} 1173.\par

\paragraph{Notebook content.} The report discusses combining strong, electromagnetic, and weak interactions and then including gravitation.

\paragraph{Precise statement.} The identifiable mathematical construction is an inclusion of gauge Lie algebras into a larger algebra, together with a choice of matter representation and symmetry-breaking vacuum. The notes do not specify a complete action, matter content, vacuum, or quantum theory sufficient to define a unique unification problem.

\paragraph{Literature and fate.} This is a programme, not a determinate conjecture. The \(SU(5)\), supergravity, and \(E_8\) examples are assessed separately below.

\paragraph{Present status.} No unique mathematical problem can be extracted from this programme as recorded; no claim of unresolved physics is inferred.

\paragraph{Assessment.} Useful for distinguishing the recorded mathematical assertion from broader physical expectations.

\paragraph{Audit.} Literal correctness: not determinable. Publication match: thematic only. Verification: unresolved lead. Development: annotated. Confidence: medium.

\subsubsection{1983-624.2 --- Electroweak gauge symmetry and the massless photon}

\label{item:1983-624.2}

\noindent\textbf{Kind:} construction. \textbf{Fate:} E. \textbf{Pages:} 1173.\par

\paragraph{Notebook content.} The notes associate \(SU(2)\times U(1)\) with electroweak unification.

\paragraph{Precise statement.} In the four-dimensional classical electroweak model with gauge algebra \(\mathfrak{su}(2)\oplus\mathfrak u(1)\), positive couplings \(g,g'\), and a scalar doublet of hypercharge \(1/2\), a vacuum \(\langle\phi\rangle=(0,v/\sqrt2)\), \(v>0\), has unbroken generator \(Q=T_3+Y\). The quadratic gauge-boson masses are \(m_W=gv/2\), \(m_Z=v\sqrt{g^2+g'^2}/2\), and \(m_\gamma=0\).

\paragraph{Editorial reconstruction.} The vacuum, normalization, and mass formulas make the recorded group identification precise; the notebook does not spell out all these conventions.

\paragraph{Literature and fate.} Weinberg’s 1967 paper gives the gauge and scalar construction, equations (4), (7), and (10)–(13). Its last page raises the renormalizability question; that question is not proved by merely writing the classical model.

\paragraph{Present status.} Established background as of 12 November 2026; the statement below is limited to the explicitly specified mathematical model.

\paragraph{Assessment.} Useful for distinguishing the recorded mathematical assertion from broader physical expectations.

\paragraph{Audit.} Literal correctness: correct with missing hypotheses. Publication match: exact. Verification: read the result. Development: annotated. Confidence: high.

\noindent References: \citep{weinberg1967leptons}.\par

\subsubsection{1983-624.3 --- The global form of the standard subgroup of SU(5)}

\label{item:1983-624.3}

\noindent\textbf{Kind:} construction. \textbf{Fate:} E. \textbf{Pages:} 1173.\par

\paragraph{Notebook content.} The report passes from electroweak symmetry to \(SU(5)\).

\paragraph{Precise statement.} The block-diagonal subgroup of \(SU(5)\) is \(S(U(3)\times U(2))\), whose Lie algebra is \(\mathfrak{su}(3)\oplus\mathfrak{su}(2)\oplus\mathfrak u(1)\). Globally it is \((SU(3)\times SU(2)\times U(1))/\mathbb Z_6\), rather than a direct product. The adjoint representation of \(SU(5)\) has dimension \(5^2-1=24\).

\paragraph{Editorial reconstruction.} The finite central quotient is supplied to distinguish the global group from its Lie algebra; the notes use the customary local group notation.

\paragraph{Literature and fate.} Georgi and Glashow’s 1974 article, pp. 438–441, is the identifiable original \(SU(5)\) proposal. The publisher abstract and bibliographic record were opened; the full article was not read. The subgroup and dimension statements are direct matrix calculations.

\paragraph{Present status.} Established background as of 12 November 2026; the statement below is limited to the explicitly specified mathematical model.

\paragraph{Assessment.} Useful for distinguishing the recorded mathematical assertion from broader physical expectations.

\paragraph{Audit.} Literal correctness: correct with missing hypotheses. Publication match: exact. Verification: publisher metadata checked. Development: annotated. Confidence: high.

\noindent References: \citep{georgiglashow1974unity}.\par

\subsubsection{1983-624.4 --- Supersymmetry as an odd extension of the Poincare algebra}

\label{item:1983-624.4}

\noindent\textbf{Kind:} definition. \textbf{Fate:} E. \textbf{Pages:} 1173--1174.\par

\paragraph{Notebook content.} Fradkin describes odd generators as a way to combine bosonic and fermionic symmetries.

\paragraph{Precise statement.} For four-dimensional \(N=1\) super-Poincare symmetry without central charges, the odd Weyl-spinor generators satisfy \(\{Q_\alpha,\overline Q_{\dot\beta}\}=2\sigma^\mu_{\alpha\dot\beta}P_\mu\), \(\{Q_\alpha,Q_\beta\}=0\), and \([P_\mu,Q_\alpha]=0\), with the usual Lorentz covariance. The bracket is a Lie superbracket; odd-odd brackets are symmetric.

\paragraph{Editorial reconstruction.} The displayed (N=1), zero-central-charge version is a concrete standard example, not a reconstruction of an unspecified extended supersymmetry algebra.

\paragraph{Literature and fate.} Leites (1980), Chapter I, \S1.9, supplies the precise superalgebra convention and its historical introduction discusses physical supersymmetry. The notebook’s review citation is most plausibly Fayet and Ferrara (1977), not “Ferrari and Fayet, 1978”; that review’s full text was not obtained.

\paragraph{Present status.} Established background as of 12 November 2026; the statement below is limited to the explicitly specified mathematical model.

\paragraph{Assessment.} Useful for distinguishing the recorded mathematical assertion from broader physical expectations.

\paragraph{Audit.} Literal correctness: correct with missing hypotheses. Publication match: partial. Verification: read the result. Development: annotated. Confidence: high.

\noindent References: \citep{leites1980introduction,fayetferrara1977supersymmetry}.\par

\subsubsection{1983-624.5 --- Gauging the superalgebra and the role of the vierbein}

\label{item:1983-624.5}

\noindent\textbf{Kind:} construction. \textbf{Fate:} E. \textbf{Pages:} 1174.\par

\paragraph{Notebook content.} A connection valued in a Poincare-type superalgebra is written, with gravitational and fermionic components.

\paragraph{Precise statement.} A superalgebra-valued one-form can be written \(\mathcal A=e^aP_a+\frac12\omega^{ab}M_{ab}+\overline\psi Q+\cdots\), with curvature \(\mathcal F=d\mathcal A+\mathcal A\wedge\mathcal A\). To interpret \(e^a\) as a vierbein, it must be an invertible coframe. This geometric interpretation and any torsion constraint are additional assumptions; an arbitrary gauge connection alone is not a gravitational metric.

\paragraph{Editorial reconstruction.} Invertibility and the distinction between a gauge translation and a spacetime diffeomorphism are added explicitly; no full supergravity action is reconstructed.

\paragraph{Literature and fate.} This is standard supergravity geometry. Van Nieuwenhuizen (1981), introductory sections on gauge theory, the Palatini formulation, and torsion, is the identifiable review; its full text was not obtained, so the bibliographic match remains indirect.

\paragraph{Present status.} Established background as of 12 November 2026; the statement below is limited to the explicitly specified mathematical model.

\paragraph{Assessment.} Useful for distinguishing the recorded mathematical assertion from broader physical expectations.

\paragraph{Audit.} Literal correctness: correct with missing hypotheses. Publication match: partial. Verification: indirect. Development: annotated. Confidence: high.

\noindent References: \citep{vannieuwenhuizen1981supergravity}.\par

\subsubsection{1983-624.6 --- The minimal massless supergravity multiplet}

\label{item:1983-624.6}

\noindent\textbf{Kind:} result. \textbf{Fate:} E. \textbf{Pages:} 1174.\par

\paragraph{Notebook content.} The report groups gravitational and spinor fields into a supermultiplet.

\paragraph{Precise statement.} In four-dimensional Minkowski space, the massless \(N=1\) supergravity multiplet contains helicities \(+2,-2,+3/2,-3/2\): two physical bosonic and two physical fermionic states. These are on-shell helicity counts, not the numbers of components of unconstrained tensor fields.

\paragraph{Literature and fate.} This particle-content statement predates the report and belongs to the standard supergravity construction summarized by Van Nieuwenhuizen (1981), introductory discussion of the multiplet. The review was identified bibliographically, not read in full.

\paragraph{Present status.} Established background as of 12 November 2026; the statement below is limited to the explicitly specified mathematical model.

\paragraph{Assessment.} Useful for distinguishing the recorded mathematical assertion from broader physical expectations.

\paragraph{Audit.} Literal correctness: correct with missing hypotheses. Publication match: partial. Verification: indirect. Development: annotated. Confidence: high.

\noindent References: \citep{vannieuwenhuizen1981supergravity}.\par

\subsubsection{1983-624.7 --- The 248-dimensional representation of E8}

\label{item:1983-624.7}

\noindent\textbf{Kind:} remark. \textbf{Fate:} E. \textbf{Pages:} 1174.\par

\paragraph{Notebook content.} The \(E_8\) representation of dimension \(248\) is called fundamental in the notes.

\paragraph{Precise statement.} For the complex simple Lie algebra \(\mathfrak e_8\), the adjoint representation has dimension \(248\) and is its smallest nontrivial irreducible representation. Thus “fundamental” here can refer to its fundamental highest weight, but the representation should be identified explicitly as adjoint. An embedding of a unification subgroup does not by itself supply a consistent particle model.

\paragraph{Editorial reconstruction.} The representation-theoretic identification is precise; the notebook’s physical model remains underspecified.

\paragraph{Literature and fate.} Stepanyantz (2024), introduction immediately after equation (8), explicitly identifies the smallest representation as the adjoint \(248\); reference 109 identifies the Konshtein–Fradkin \(E_8\) announcement of 1980. The original announcement was not read.

\paragraph{Present status.} Established background as of 12 November 2026; the statement below is limited to the explicitly specified mathematical model.

\paragraph{Assessment.} Useful for distinguishing the recorded mathematical assertion from broader physical expectations.

\paragraph{Audit.} Literal correctness: correct with missing hypotheses. Publication match: partial. Verification: read the result. Development: annotated. Confidence: high.

\noindent References: \citep{stepanyantz2024flipped}.\par

\subsubsection{1983-624.8 --- Finiteness and the cited instanton calculation}

\label{item:1983-624.8}

\noindent\textbf{Kind:} claim. \textbf{Fate:} F. \textbf{Pages:} 1174.\par

\paragraph{Notebook content.} The report links extended supersymmetry, unification, and finiteness, with a Fradkin–Tseytlin citation.

\paragraph{Precise statement.} A finiteness assertion requires a specified action, dimension, gauge group, matter content, regularization, loop order, and treatment of anomalies. These data are not supplied for the proposed \(E_8\) theory. Consequently the notes do not determine an all-orders finiteness proposition.

\paragraph{Literature and fate.} The closest bibliographic match is Fradkin and Tseytlin, Physics Letters B \textbf{134} (1984), 301–306, on instanton zero modes and beta-functions in gauged supergravity. Its publisher-indexed abstract concerns a one-loop calculation, not an arbitrary all-orders \(E_8\) statement. The original full text could not be opened.

\paragraph{Present status.} The notebook assertion lacks enough hypotheses to assign a single finiteness theorem or open problem.

\paragraph{Assessment.} Useful for distinguishing the recorded mathematical assertion from broader physical expectations.

\paragraph{Audit.} Literal correctness: not determinable. Publication match: partial. Verification: indirect. Development: annotated. Confidence: medium.

\noindent References: \citep{fradkintseytlin1984instanton}.\par

\subsubsection{1983-624.9 --- Gelfand’s infinite-dimensional geometric space}

\label{item:1983-624.9}

\noindent\textbf{Kind:} question. \textbf{Fate:} F. \textbf{Pages:} 1174.\par

\paragraph{Notebook content.} Gelfand asks about an infinite-dimensional space behind the physical construction.

\paragraph{Precise statement.} The recorded question does not specify the space, its category, its geometric structure, or a desired theorem. Spaces of fields, connections, and gauge equivalence classes are different possible interpretations; none is selected by the text.

\paragraph{Literature and fate.} No exact problem match was found in searches combining Gelfand, Fradkin, supergravity, and infinite-dimensional geometry. The thematic relationship to moduli of gauge fields does not identify the intended problem.

\paragraph{Present status.} Unidentifiable as a precise problem from the surviving formulation.

\paragraph{Assessment.} Useful for distinguishing the recorded mathematical assertion from broader physical expectations.

\paragraph{Audit.} Literal correctness: not determinable. Publication match: none found. Verification: unresolved lead. Development: annotated. Confidence: medium.

\subsection{Record 625: Gauge mass generation, extended supergravity, and open gauge algebras}

\label{talk:625}

\noindent\textbf{Date:} 1984-03-12. \textbf{Speaker:} E. S. Fradkin. \textbf{Pages:} 1175--1176.\par

The dated report occupies P001175–P001176. It revisits gauge curvature, renormalization, spontaneous symmetry breaking, extended supermultiplets, conformal gravity, functional integration, and quantization of gauge algebras that close modulo equations of motion. These are separate assertions with different degrees of precision.

\subsubsection{1983-625.1 --- Gauge connection, curvature, and transformation conventions}

\label{item:1983-625.1}

\noindent\textbf{Kind:} definition. \textbf{Fate:} E. \textbf{Pages:} 1175.\par

\paragraph{Notebook content.} The report writes a matrix gauge potential, a curvature, and an infinitesimal gauge transformation.

\paragraph{Precise statement.} For \(D_\mu=\partial_\mu+A_\mu\), define \(F_{\mu\nu}=[D_\mu,D_\nu]=\partial_\mu A_\nu-\partial_\nu A_\mu+[A_\mu,A_\nu]\). Under \(D_\mu\mapsto g^{-1}D_\mu g\), \(A_\mu\mapsto g^{-1}A_\mu g+g^{-1}\partial_\mu g\), and \(F_{\mu\nu}\mapsto g^{-1}F_{\mu\nu}g\). Infinitesimally, \(g=1+\epsilon\) gives \(\delta A_\mu=\partial_\mu\epsilon+[A_\mu,\epsilon]\).

\paragraph{Editorial reconstruction.} The modern statement fixes a single internally consistent convention; an overall sign convention is legitimate only when used throughout.

\paragraph{Literature and fate.} These identities follow by expanding the commutator. They agree with the gauge-theoretic setup in ’t Hooft (1971), \S2. The derivative order in the recorded curvature is inconsistent with the other displayed convention unless the whole curvature convention is changed.

\paragraph{Present status.} Established under the explicitly stated assumptions as of 12 November 2026.

\paragraph{Assessment.} Clarifies the hypotheses and the distinction between a classical construction and a quantum assertion.

\paragraph{Audit.} Literal correctness: false as written. Publication match: exact. Verification: read the result. Development: annotated. Confidence: high.

\noindent References: \citep{thooft1971massive}.\par

\subsubsection{1983-625.2 --- Explicit vector masses and renormalizability}

\label{item:1983-625.2}

\noindent\textbf{Kind:} claim. \textbf{Fate:} E. \textbf{Pages:} 1175.\par

\paragraph{Notebook content.} A mass term is described as destroying gauge invariance and causing a renormalization problem.

\paragraph{Precise statement.} A term \(\frac12m^2 A_\mu^aA^{a\mu}\) is not invariant under the full local connection transformation for \(m\ne0\). In an interacting nonabelian vector theory, simply adding this Proca term does not give the standard renormalizable Higgs construction. The assertion needs this interacting nonabelian setting: a free massive Abelian vector field does not furnish the claimed obstruction.

\paragraph{Literature and fate.} ’t Hooft (1971), introduction and \S\S3–5, explains how spontaneous breaking preserves the identities needed for renormalization and unitarity. The notebook’s unqualified formulation is therefore too broad.

\paragraph{Present status.} Established under the explicitly stated assumptions as of 12 November 2026.

\paragraph{Assessment.} Clarifies the hypotheses and the distinction between a classical construction and a quantum assertion.

\paragraph{Audit.} Literal correctness: correct with missing hypotheses. Publication match: exact. Verification: read the result. Development: annotated. Confidence: high.

\noindent References: \citep{thooft1971massive}.\par

\subsubsection{1983-625.3 --- What perturbative renormalizability asserts}

\label{item:1983-625.3}

\noindent\textbf{Kind:} definition. \textbf{Fate:} E. \textbf{Pages:} 1175.\par

\paragraph{Notebook content.} The discussion asks what renormalizability means and links it to gauge invariance.

\paragraph{Precise statement.} For a specified perturbative local quantum field theory, renormalizability means that ultraviolet divergences can be absorbed into a finite set of local counterterms already allowed by the theory’s field content and symmetries, together with field and parameter redefinitions, to every perturbative order under consideration. Gauge invariance by itself is insufficient without suitable dimension, interactions, and absence of obstructing anomalies.

\paragraph{Editorial reconstruction.} This is a precise formulation of the concept discussed, not an assertion that the notes prove an all-orders theorem for every later model.

\paragraph{Literature and fate.} ’t Hooft (1971), \S4 and the introductory qualifications concerning Ward identities, supplies the relevant earlier proof in the spontaneously broken gauge setting. A general statement about every gauge-invariant action is not implied.

\paragraph{Present status.} Established under the explicitly stated assumptions as of 12 November 2026.

\paragraph{Assessment.} Clarifies the hypotheses and the distinction between a classical construction and a quantum assertion.

\paragraph{Audit.} Literal correctness: correct with missing hypotheses. Publication match: partial. Verification: read the result. Development: annotated. Confidence: high.

\noindent References: \citep{thooft1971massive}.\par

\subsubsection{1983-625.4 --- Vacuum displacement and the scalar quadratic term}

\label{item:1983-625.4}

\noindent\textbf{Kind:} construction. \textbf{Fate:} E. \textbf{Pages:} 1175.\par

\paragraph{Notebook content.} The scalar field is displaced by a constant, and a condensate and an excitation are discussed.

\paragraph{Precise statement.} Write \(\phi=v+\eta\), with constant vacuum \(v\) satisfying \(V'(v)=0\), and fluctuation \(\eta\). For a canonically normalized real scalar and \(V(\phi)=a\phi^2+b\phi^4\), \(b>0\), \(a<0\), a nonzero minimum has \(v^2=-a/(2b)\), and \(V(v+\eta)=V(v)+\frac12(-4a)\eta^2+O(\eta^3)\). The squared scalar mass is \(V''(v)=-4a\), not the coefficient of \(\eta^2\) without the factor \(1/2\).

\paragraph{Editorial reconstruction.} The polynomial potential here is an explicitly normalized representative used to explain the mass coefficient; the symmetry-breaking gauge model uses a charged multiplet.

\paragraph{Literature and fate.} Higgs (1964), equations (1)–(4), provides the earlier classical mechanism in a charged scalar model. The notebook’s labels for the constant and the variable are interchanged. The elementary Taylor expansion fixes the normalization issue separately.

\paragraph{Present status.} Established under the explicitly stated assumptions as of 12 November 2026.

\paragraph{Assessment.} Clarifies the hypotheses and the distinction between a classical construction and a quantum assertion.

\paragraph{Audit.} Literal correctness: false as written. Publication match: exact. Verification: read the result. Development: annotated. Confidence: high.

\noindent References: \citep{higgs1964broken}.\par

\subsubsection{1983-625.5 --- The gauge-boson mass matrix and unbroken symmetries}

\label{item:1983-625.5}

\noindent\textbf{Kind:} result. \textbf{Fate:} E. \textbf{Pages:} 1175.\par

\paragraph{Notebook content.} The report says that the displaced scalar generates masses.

\paragraph{Precise statement.} Let a compact gauge group act unitarily on a finite-dimensional scalar space, with anti-Hermitian generators \(T_a\), couplings included in \(T_a\), and vacuum \(v\). The quadratic form from \(\lVert D_\mu\phi\rVert^2\) is, up to the common kinetic normalization, \(M_{ab}=\operatorname{Re}\langle T_av,T_bv\rangle\). It is positive semidefinite; its kernel consists of Lie-algebra elements \(X\) satisfying \(Xv=0\). Thus exactly the unbroken infinitesimal generators remain massless at this classical level.

\paragraph{Literature and fate.} Higgs (1964), p. 509, explicitly distinguishes massive directions from surviving massless gauge fields; Weinberg (1967) exhibits the photon. The formula above is the direct quadratic expansion. It does not imply that every matter particle acquires a mass.

\paragraph{Present status.} Established under the explicitly stated assumptions as of 12 November 2026.

\paragraph{Assessment.} Clarifies the hypotheses and the distinction between a classical construction and a quantum assertion.

\paragraph{Audit.} Literal correctness: correct with missing hypotheses. Publication match: partial. Verification: read the result. Development: annotated. Confidence: high.

\noindent References: \citep{higgs1964broken,weinberg1967leptons}.\par

\subsubsection{1983-625.6 --- The proposed unification energy scale}

\label{item:1983-625.6}

\noindent\textbf{Kind:} remark. \textbf{Fate:} F. \textbf{Pages:} 1175.\par

\paragraph{Notebook content.} A scale of order \(10^{14}\) is indicated in the physical discussion.

\paragraph{Precise statement.} A numerical unification scale is the output of a model-dependent running-coupling calculation, including a renormalization scheme, particle spectrum, thresholds, and matching conditions. None of these are supplied sufficiently to identify a unique numerical prediction.

\paragraph{Literature and fate.} The \(SU(5)\) literature supplies thematic context, but no exact calculation matching this isolated number was verified. It should remain a historical estimate, not be promoted to a mathematical theorem or a current prediction.

\paragraph{Present status.} The recorded number cannot be assigned a uniquely checkable modern status.

\paragraph{Assessment.} Clarifies the hypotheses and the distinction between a classical construction and a quantum assertion.

\paragraph{Audit.} Literal correctness: not determinable. Publication match: thematic only. Verification: unresolved lead. Development: annotated. Confidence: medium.

\noindent References: \citep{georgiglashow1974unity}.\par

\subsubsection{1983-625.7 --- The four-dimensional N=8 massless multiplet}

\label{item:1983-625.7}

\noindent\textbf{Kind:} result. \textbf{Fate:} E. \textbf{Pages:} 1176.\par

\paragraph{Notebook content.} The table lists multiplicities \(1,8,28,56,70\) at decreasing spins.

\paragraph{Precise statement.} For a massless four-dimensional \(N=8\) supergravity multiplet with highest helicity \(2\), the helicity multiplicities are the binomial coefficients \(\binom8k\), \(0\le k\le8\), at helicity \(2-k/2\). Equivalently there are one graviton, eight gravitini, twenty-eight vectors, fifty-six spin-\(1/2\) fields, and seventy real scalars, with \(128\) bosonic and \(128\) fermionic physical states. Under the assumptions of a finite massless multiplet, CPT completion, and no spin above \(2\), one has \(N\le8\).

\paragraph{Literature and fate.} The original \(N=8\) Lagrangian is Cremmer and Julia (1978), pp. 48–51. Its bibliographic record was located but its full text was not opened. The multiplicities are checked directly from eight fermionic raising operators and CPT symmetry.

\paragraph{Present status.} Established under the explicitly stated assumptions as of 12 November 2026.

\paragraph{Assessment.} Clarifies the hypotheses and the distinction between a classical construction and a quantum assertion.

\paragraph{Audit.} Literal correctness: correct with missing hypotheses. Publication match: exact. Verification: indirect. Development: annotated. Confidence: high.

\noindent References: \citep{cremmerjulia1978eight}.\par

\subsubsection{1983-625.8 --- The eleven-dimensional supergravity field content}

\label{item:1983-625.8}

\noindent\textbf{Kind:} result. \textbf{Fate:} E. \textbf{Pages:} 1176.\par

\paragraph{Notebook content.} Eleven-dimensional supergravity and the absence of certain anomalies are mentioned.

\paragraph{Precise statement.} The eleven-dimensional massless supergravity multiplet has a metric, a three-form potential, and a Majorana gravitino. The little group is \(SO(9)\): the bosonic counts are \(11(11-3)/2=44\) and \(\binom93=84\), and the gamma-traceless vector-spinor has \(9\cdot16-16=128\) states. Hence the physical degrees of freedom are \(128+128\).

\paragraph{Editorial reconstruction.} The field-count assertion is precise. The anomaly remark is retained as an unresolved qualification rather than interpreted as a theorem about all quantum anomalies.

\paragraph{Literature and fate.} Cremmer, Julia, and Scherk (1978), pp. 409–412, is the original construction, identified bibliographically. The notebook’s anomaly phrase does not specify an anomaly polynomial, a quantum definition, or boundary conditions, and does not imply perturbative finiteness.

\paragraph{Present status.} Established under the explicitly stated assumptions as of 12 November 2026.

\paragraph{Assessment.} Clarifies the hypotheses and the distinction between a classical construction and a quantum assertion.

\paragraph{Audit.} Literal correctness: correct with missing hypotheses. Publication match: partial. Verification: indirect. Development: annotated. Confidence: high.

\noindent References: \citep{cremmerjuliascherk1978eleven}.\par

\subsubsection{1983-625.9 --- Conformal curvature-squared gravity}

\label{item:1983-625.9}

\noindent\textbf{Kind:} construction. \textbf{Fate:} E. \textbf{Pages:} 1176.\par

\paragraph{Notebook content.} A four-index curvature expression and a conformal quadratic action are written.

\paragraph{Precise statement.} For a smooth four-dimensional pseudo-Riemannian metric, the Weyl tensor is the trace-free part of the Riemann tensor. Its squared norm obeys \(C_{\mu\nu\rho\sigma}C^{\mu\nu\rho\sigma}=R_{\mu\nu\rho\sigma}R^{\mu\nu\rho\sigma}-2R_{\mu\nu}R^{\mu\nu}+\frac13R^2\). Modulo the Euler density \(E_4=R_{\mu\nu\rho\sigma}^2-4R_{\mu\nu}^2+R^2\), the conformal action is \(2\int(R_{\mu\nu}R^{\mu\nu}-R^2/3)\sqrt{|g|}\,d^4x\).

\paragraph{Editorial reconstruction.} The statement gives the invariant identity independently of the notebook’s compressed index expression; no unrecorded supersymmetric completion is asserted.

\paragraph{Literature and fate.} The displayed identity is the standard four-dimensional curvature decomposition. Van Nieuwenhuizen’s review discusses conformal supergravity, but that review was not read in full. The recorded malformed Kronecker-symbol expression must not be treated as a valid definition of the Weyl tensor, and the purely bosonic quadratic action is not by itself a supergravity action.

\paragraph{Present status.} Established under the explicitly stated assumptions as of 12 November 2026.

\paragraph{Assessment.} Clarifies the hypotheses and the distinction between a classical construction and a quantum assertion.

\paragraph{Audit.} Literal correctness: false as written. Publication match: partial. Verification: indirect. Development: annotated. Confidence: high.

\noindent References: \citep{vannieuwenhuizen1981supergravity}.\par

\subsubsection{1983-625.10 --- Formal functional integration over fields}

\label{item:1983-625.10}

\noindent\textbf{Kind:} method. \textbf{Fate:} F. \textbf{Pages:} 1176.\par

\paragraph{Notebook content.} The discussion introduces a functional integral for the gravitational or gauge theory.

\paragraph{Precise statement.} An expression \(\int \exp(iS[\phi])\,D\phi\) is formal until the field space, measure or perturbative prescription, gauge quotient, boundary conditions, and regularization are specified. The notes supply no such construction sufficient to define a unique mathematical integral.

\paragraph{Literature and fate.} The earlier index contains related formal gravitational-integral discussion at 1981-546.5, but this does not identify an identical precise problem. No claim of a rigorous construction or a new open problem follows from the notation.

\paragraph{Present status.} The formal expression remains mathematically underspecified in this record.

\paragraph{Assessment.} Clarifies the hypotheses and the distinction between a classical construction and a quantum assertion.

\paragraph{Audit.} Literal correctness: not determinable. Publication match: thematic only. Verification: unresolved lead. Development: annotated. Confidence: medium.

\subsubsection{1983-625.11 --- Quantizing gauge algebras that close modulo field equations}

\label{item:1983-625.11}

\noindent\textbf{Kind:} method. \textbf{Fate:} E. \textbf{Pages:} 1176.\par

\paragraph{Notebook content.} The speaker mentions gauge transformations whose commutators close only after using equations of motion, and a Vilkovisky-related method.

\paragraph{Precise statement.} For fields \(\phi^i\) and gauge generators \(R^i_\alpha\), an open algebra has \([R_\alpha,R_\beta]^i=R^i_\gamma f^\gamma_{\alpha\beta}+M^{ij}_{\alpha\beta}\,\delta S/\delta\phi^j\). The Batalin–Vilkovisky method introduces ghosts and antifields and seeks an extended action satisfying the classical master equation \((S_{\mathrm{BV}},S_{\mathrm{BV}})=0\), with suitable properness and reducibility data. Open closure occurs in supergravity; ordinary spacetime diffeomorphisms already close off shell and should not be used as the universal example.

\paragraph{Editorial reconstruction.} The master-equation formulation identifies the method; no unrecorded quantization theorem for the unspecified model is attributed to the seminar.

\paragraph{Literature and fate.} Batalin and Vilkovisky (1983), abstract and pp. 2567–2582, explicitly address dependent generators and open algebras; the publisher abstract was opened. The earlier Fradkin–Vilkovisky (1975), pp. 224–226, constraint-quantization paper is another relevant bibliographic pointer, but its full text was not read.

\paragraph{Present status.} Established under the explicitly stated assumptions as of 12 November 2026.

\paragraph{Assessment.} Clarifies the hypotheses and the distinction between a classical construction and a quantum assertion.

\paragraph{Audit.} Literal correctness: correct with missing hypotheses. Publication match: partial. Verification: publisher metadata checked. Development: annotated. Confidence: high.

\noindent References: \citep{batalinvilkovisky1983dependent,fradkinvilkovisky1975constraints}.\par

\subsection{Record 626: Lie superalgebras and exceptional families}

\label{talk:626}

\noindent\textbf{Date:} 1984-03-19. \textbf{Speaker:} D. A. Leites. \textbf{Pages:} 1177--1178.\par

The heading on P001177 begins Leites’s report; it continues through P001178, before Penkov’s heading on P001179. It introduces superspaces, signs, matrix families, invariant forms, traces, and two exceptional constructions. A short question about Jordan form remains underspecified. The remaining matrix formulas were read in the supplied transcription; the outage prevented a final visual recheck of P001178.

\subsubsection{1983-626.1 --- Superspaces, homogeneous maps, and block matrices}

\label{item:1983-626.1}

\noindent\textbf{Kind:} definition. \textbf{Fate:} E. \textbf{Pages:} 1177.\par

\paragraph{Notebook content.} The report starts with a vector space decomposed into two parity components and the corresponding parity of linear maps.

\paragraph{Precise statement.} Over a field of characteristic zero, a superspace is \(V=V_{\overline0}\oplus V_{\overline1}\). A homogeneous map \(T\) has parity \(|T|\in\mathbb Z/2\) when \(T(V_i)\subset V_{i+|T|}\). For \(V=\mathbb C^{m|n}\), even endomorphisms are block diagonal and odd endomorphisms block off-diagonal; all endomorphisms are sums of these components.

\paragraph{Literature and fate.} Leites (1980), Chapter I, \S\S1.1–1.3, gives these definitions. They are background rather than a new classification result.

\paragraph{Present status.} Classical superalgebraic background, already known before the 1984 report and still valid with the stated conventions on 12 November 2026.

\paragraph{Assessment.} Makes parity, characteristic, and matrix conventions explicit.

\paragraph{Audit.} Literal correctness: correct with missing hypotheses. Publication match: exact. Verification: read the result. Development: annotated. Confidence: high.

\noindent References: \citep{leites1980introduction}.\par

\subsubsection{1983-626.2 --- Koszul signs and the Lie superalgebra identities}

\label{item:1983-626.2}

\noindent\textbf{Kind:} definition. \textbf{Fate:} E. \textbf{Pages:} 1177.\par

\paragraph{Notebook content.} The signs in skew symmetry and the Jacobi identity are discussed, with a reference to Quillen’s sign rule.

\paragraph{Precise statement.} For homogeneous \(x,y,z\), a Lie superbracket has parity \(|[x,y]|=|x|+|y|\), satisfies \([x,y]=-(-1)^{|x||y|}[y,x]\), and obeys \((-1)^{|x||z|}[x,[y,z]]+(-1)^{|y||x|}[y,[z,x]]+(-1)^{|z||y|}[z,[x,y]]=0\). Interchanging homogeneous tensor factors contributes \((-1)^{|x||y|}\).

\paragraph{Literature and fate.} Leites (1980), Chapter I, \S1.9.1, states the identities, and the historical introduction discusses the sign convention. Earlier occurrences include 1976-404.1.

\paragraph{Present status.} Classical superalgebraic background, already known before the 1984 report and still valid with the stated conventions on 12 November 2026.

\paragraph{Assessment.} Makes parity, characteristic, and matrix conventions explicit.

\paragraph{Audit.} Literal correctness: correct with missing hypotheses. Publication match: exact. Verification: read the result. Development: annotated. Confidence: high.

\noindent References: \citep{leites1980introduction}.\par

\subsubsection{1983-626.3 --- The associative supercommutator}

\label{item:1983-626.3}

\noindent\textbf{Kind:} construction. \textbf{Fate:} E. \textbf{Pages:} 1177.\par

\paragraph{Notebook content.} An associative superalgebra is turned into a Lie superalgebra by a modified commutator.

\paragraph{Precise statement.} If \(A=A_{\overline0}\oplus A_{\overline1}\) is an associative algebra whose multiplication preserves parity, then \([x,y]=xy-(-1)^{|x||y|}yx\) for homogeneous elements extends bilinearly to a Lie superbracket. For \(A=\operatorname{End}(V)\) this gives \(\mathfrak{gl}(V)\). Associativity proves the super-Jacobi identity by cancellation of the six ordered triple products.

\paragraph{Literature and fate.} This construction is explicit in Leites (1980), Chapter I, \S1.9.1, and is already implicit in earlier superalgebra records.

\paragraph{Present status.} Classical superalgebraic background, already known before the 1984 report and still valid with the stated conventions on 12 November 2026.

\paragraph{Assessment.} Makes parity, characteristic, and matrix conventions explicit.

\paragraph{Audit.} Literal correctness: correct with missing hypotheses. Publication match: exact. Verification: read the result. Development: annotated. Confidence: high.

\noindent References: \citep{leites1980introduction}.\par

\subsubsection{1983-626.4 --- Supertrace, special linear superalgebra, and Euler characteristic}

\label{item:1983-626.4}

\noindent\textbf{Kind:} definition. \textbf{Fate:} E. \textbf{Pages:} 1177.\par

\paragraph{Notebook content.} The difference of the two block traces is introduced and compared with an Euler characteristic.

\paragraph{Precise statement.} For \(X=\begin{pmatrix}A&B\\C&D\end{pmatrix}\), define \(\operatorname{str}X=\operatorname{tr}A-\operatorname{tr}D\). It vanishes on supercommutators; hence \(\mathfrak{sl}(m|n)=\ker\operatorname{str}\) is a Lie subsuperalgebra. For a finite complex with even chain endomorphism \(T\), the alternating trace on chain spaces equals the alternating trace induced on cohomology; for \(T=1\) this is the Euler characteristic.

\paragraph{Editorial reconstruction.} Finiteness and the chain-map hypothesis are added to make the Euler-characteristic analogy an actual trace identity.

\paragraph{Literature and fate.} Leites (1980), Chapter I, \S1.6, supplies supertrace; the cohomological equality follows by splitting cycles and boundaries in a finite complex. The simplicity of \(\mathfrak{sl}(m|n)\) is not automatic when \(m=n\), where scalar matrices lie in its center.

\paragraph{Present status.} Classical superalgebraic background, already known before the 1984 report and still valid with the stated conventions on 12 November 2026.

\paragraph{Assessment.} Makes parity, characteristic, and matrix conventions explicit.

\paragraph{Audit.} Literal correctness: correct with missing hypotheses. Publication match: exact. Verification: read the result. Development: annotated. Confidence: high.

\noindent References: \citep{leites1980introduction}.\par

\subsubsection{1983-626.5 --- Classical and vector-field families}

\label{item:1983-626.5}

\noindent\textbf{Kind:} remark. \textbf{Fate:} E. \textbf{Pages:} 1177.\par

\paragraph{Notebook content.} The list includes linear, orthogonal, symplectic, and vector-field families denoted \(W,S,H,K\).

\paragraph{Precise statement.} The list names several different constructions. For example \(W(n)=\operatorname{Der}\Lambda(\xi_1,\ldots,\xi_n)\) is finite dimensional when all \(n\) coordinates are odd; allowing even polynomial or formal coordinates produces infinite-dimensional vector-field superalgebras. Divergence, symplectic, or contact preservation yields the corresponding special, Hamiltonian, or contact families, with simplicity requiring the appropriate derived algebra and rank hypotheses.

\paragraph{Editorial reconstruction.} The example (W(n)) fixes one category; no uniform unqualified simplicity claim is attached to the entire list.

\paragraph{Literature and fate.} Kac (1977), A sketch of Lie superalgebra theory, pp. 31–64, classifies simple finite-dimensional examples; its publisher abstract was opened, but the full classification proof was not read. The notebook list is not an exhaustive theorem across all finite, polynomial, and formal categories.

\paragraph{Present status.} Classical superalgebraic background, already known before the 1984 report and still valid with the stated conventions on 12 November 2026.

\paragraph{Assessment.} Makes parity, characteristic, and matrix conventions explicit.

\paragraph{Audit.} Literal correctness: correct with missing hypotheses. Publication match: partial. Verification: publisher metadata checked. Development: annotated. Confidence: high.

\noindent References: \citep{kac1977sketch}.\par

\subsubsection{1983-626.6 --- Orthosymplectic algebras from an even form}

\label{item:1983-626.6}

\noindent\textbf{Kind:} construction. \textbf{Fate:} E. \textbf{Pages:} 1177--1178.\par

\paragraph{Notebook content.} Preservation of a nondegenerate graded form is used to define an orthosymplectic family.

\paragraph{Precise statement.} Let \(V=\mathbb C^{m|2n}\) carry a nondegenerate even supersymmetric bilinear form \(B\): it is symmetric on \(V_{\overline0}\), alternating on \(V_{\overline1}\), and vanishes on mixed parity. The homogeneous endomorphisms satisfying \(B(Xv,w)+(-1)^{|X||v|}B(v,Xw)=0\) form \(\mathfrak{osp}(m|2n)\), whose even part is \(\mathfrak{so}(m)\oplus\mathfrak{sp}(2n)\).

\paragraph{Literature and fate.} Leites (1980), Chapter I, \S1.5, gives the graded bilinear-form conventions. The intrinsic identity avoids ambiguity among supertranspose conventions in block notation.

\paragraph{Present status.} Classical superalgebraic background, already known before the 1984 report and still valid with the stated conventions on 12 November 2026.

\paragraph{Assessment.} Makes parity, characteristic, and matrix conventions explicit.

\paragraph{Audit.} Literal correctness: correct with missing hypotheses. Publication match: exact. Verification: read the result. Development: annotated. Confidence: high.

\noindent References: \citep{leites1980introduction}.\par

\subsubsection{1983-626.7 --- The periplectic algebra}

\label{item:1983-626.7}

\noindent\textbf{Kind:} construction. \textbf{Fate:} E. \textbf{Pages:} 1178.\par

\paragraph{Notebook content.} The family denoted \(\Pi(n)\) is displayed in blocks with one symmetric and one skew-symmetric off-diagonal block.

\paragraph{Precise statement.} With one standard convention, the periplectic algebra preserving a nondegenerate odd symmetric form on \(\mathbb C^{n|n}\) is \(\mathfrak{pe}(n)=\{\begin{pmatrix}A&B\\C&-A^{t}\end{pmatrix}:B^{t}=B,\ C^{t}=-C\}\). Its special subalgebra imposes \(\operatorname{tr}A=0\); this special algebra is simple for \(n\ge3\). A parity reversal of the form interchanges the symmetric and alternating block conventions.

\paragraph{Editorial reconstruction.} The notation \(\Pi(n)\) is normalized to \(\mathfrak{pe}(n)\); the special-algebra and rank qualifications are supplied explicitly.

\paragraph{Literature and fate.} This is one of the classical matrix families in Kac’s 1977 classification. The classification citation was verified at publisher-metadata level; the displayed block condition itself is checked by preserving the odd form.

\paragraph{Present status.} Classical superalgebraic background, already known before the 1984 report and still valid with the stated conventions on 12 November 2026.

\paragraph{Assessment.} Makes parity, characteristic, and matrix conventions explicit.

\paragraph{Audit.} Literal correctness: correct with missing hypotheses. Publication match: partial. Verification: indirect. Development: annotated. Confidence: medium.

\noindent References: \citep{kac1977sketch,leites1980introduction}.\par

\subsubsection{1983-626.8 --- The queer algebra as a supercentralizer}

\label{item:1983-626.8}

\noindent\textbf{Kind:} construction. \textbf{Fate:} E. \textbf{Pages:} 1178.\par

\paragraph{Notebook content.} An odd complex structure \(J\) and matrices with equal diagonal and off-diagonal blocks are written.

\paragraph{Precise statement.} Let \(J=\begin{pmatrix}0&I\\-I&0\end{pmatrix}\) on \(\mathbb C^{n|n}\). Its supercentralizer consists of \(\mathfrak q(n)=\{\begin{pmatrix}A&B\\B&A\end{pmatrix}\}\): an even matrix commutes with \(J\), while an odd one anticommutes with it. The bracket is the supercommutator, not an ordinary commutator on arbitrary inhomogeneous matrices.

\paragraph{Literature and fate.} The queer family is classical background included in Kac (1977); Leites’s general matrix and parity conventions explain the centralizer condition. The block calculation is immediate from the displayed \(J\).

\paragraph{Present status.} Classical superalgebraic background, already known before the 1984 report and still valid with the stated conventions on 12 November 2026.

\paragraph{Assessment.} Makes parity, characteristic, and matrix conventions explicit.

\paragraph{Audit.} Literal correctness: correct with missing hypotheses. Publication match: partial. Verification: indirect. Development: annotated. Confidence: high.

\noindent References: \citep{kac1977sketch,leites1980introduction}.\par

\subsubsection{1983-626.9 --- Queer trace and the central quotient}

\label{item:1983-626.9}

\noindent\textbf{Kind:} definition. \textbf{Fate:} E. \textbf{Pages:} 1178.\par

\paragraph{Notebook content.} A trace of the off-diagonal block is defined, with a purported equality to a supertrace involving \(J\).

\paragraph{Precise statement.} For \(X=\begin{pmatrix}A&B\\B&A\end{pmatrix}\), the odd linear functional \(\operatorname{qtr}X=\operatorname{tr}B\) vanishes on supercommutators. Thus \(\mathfrak{sq}(n)=\ker\operatorname{qtr}\) is a subsuperalgebra. For \(n\ge3\), \(\mathfrak{psq}(n)=\mathfrak{sq}(n)/\mathbb CI\) is simple. With the preceding \(J\), however, \(\operatorname{str}(XJ)=-2\operatorname{tr}B\), so \(\operatorname{qtr}X=-\frac12\operatorname{str}(XJ)\).

\paragraph{Literature and fate.} The family and its special central quotient are standard classification data. The normalization correction follows by multiplying the two displayed block matrices and does not depend on later literature. A final scan check of the handwritten normalization was prevented by the outage, so the correction is against the supplied transcription.

\paragraph{Present status.} Classical superalgebraic background, already known before the 1984 report and still valid with the stated conventions on 12 November 2026.

\paragraph{Assessment.} Makes parity, characteristic, and matrix conventions explicit.

\paragraph{Audit.} Literal correctness: false as written. Publication match: partial. Verification: indirect. Development: annotated. Confidence: medium.

\noindent References: \citep{kac1977sketch}.\par

\subsubsection{1983-626.10 --- Gelfand’s question about Jordan form}

\label{item:1983-626.10}

\noindent\textbf{Kind:} question. \textbf{Fate:} F. \textbf{Pages:} 1178.\par

\paragraph{Notebook content.} Gelfand asks whether there is a Jordan-form theorem in this setting.

\paragraph{Precise statement.} A classification problem would need to specify whether the operator is even or odd, whether its entries lie in \(\mathbb C\) or a supercommutative base algebra, and which parity-preserving or supergroup conjugations are allowed. These alternatives produce different problems; none is fixed in the surviving question.

\paragraph{Literature and fate.} Searches for Leites, Gelfand, Jordan form, and supermatrices did not identify an exact published answer to this particular formulation. Work on Jordan superalgebras is a different subject and was not accepted as a match.

\paragraph{Present status.} The question remains insufficiently specified for an exact problem/status assignment.

\paragraph{Assessment.} Makes parity, characteristic, and matrix conventions explicit.

\paragraph{Audit.} Literal correctness: not determinable. Publication match: none found. Verification: unresolved lead. Development: annotated. Confidence: medium.

\subsubsection{1983-626.11 --- The exceptional algebra F(4)}

\label{item:1983-626.11}

\noindent\textbf{Kind:} construction. \textbf{Fate:} E. \textbf{Pages:} 1178.\par

\paragraph{Notebook content.} An exceptional algebra is described using an orthogonal algebra, an \(\mathfrak{sl}_2\) factor, and a spin representation.

\paragraph{Precise statement.} The complex exceptional simple Lie superalgebra \(F(4)\) has even part \(\mathfrak{so}_7\oplus\mathfrak{sl}_2\) and odd part \(\operatorname{Spin}_7\otimes\mathbb C^2\), where \(\operatorname{Spin}_7\) denotes the eight-dimensional spin module. Its superdimension is \(24|16\). The equivariant odd-odd bracket is normalized so that the super-Jacobi identity holds.

\paragraph{Literature and fate.} This is the exceptional family in Kac’s 1977 classification; the old notation in the notes is identified by its even part and odd module, not by typography alone. The classification article was checked through its publisher metadata, not read in full.

\paragraph{Present status.} Classical superalgebraic background, already known before the 1984 report and still valid with the stated conventions on 12 November 2026.

\paragraph{Assessment.} Makes parity, characteristic, and matrix conventions explicit.

\paragraph{Audit.} Literal correctness: correct with missing hypotheses. Publication match: partial. Verification: indirect. Development: annotated. Confidence: medium.

\noindent References: \citep{kac1977sketch}.\par

\subsubsection{1983-626.12 --- The one-parameter exceptional family D(2,1;alpha)}

\label{item:1983-626.12}

\noindent\textbf{Kind:} construction. \textbf{Fate:} E. \textbf{Pages:} 1178.\par

\paragraph{Notebook content.} Three two-dimensional symplectic spaces and coefficients whose sum is zero are used to construct an exceptional family.

\paragraph{Precise statement.} Let \((V_i,\psi_i)\), \(i=1,2,3\), be two-dimensional complex symplectic spaces. Put \(\mathfrak g_{\overline0}=\bigoplus_i\mathfrak{sp}(V_i)\), \(\mathfrak g_{\overline1}=V_1\otimes V_2\otimes V_3\). Define \(P_i(x,y)z=\psi_i(y,z)x+\psi_i(x,z)y\). The odd bracket is \(\sum_i\sigma_i\psi_j(x_j,y_j)\psi_k(x_k,y_k)P_i(x_i,y_i)\), where \(\{i,j,k\}=\{1,2,3\}\). It satisfies Jacobi exactly when \(\sigma_1+\sigma_2+\sigma_3=0\); it is simple when every \(\sigma_i\ne0\). Common rescaling and permutation of the triple yield the isomorphisms; a normalized ratio gives \(D(2,1;\alpha)\), \(\alpha\notin\{0,-1\}\).

\paragraph{Editorial reconstruction.} The bracket is written once for each of the three factors, avoiding an unnormalized sum over all six permutations.

\paragraph{Literature and fate.} Poletaeva (2010), \S2, equations (1)–(5), explicitly records this construction, Jacobi condition, simplicity criterion, and parameter equivalence, attributing the equivalent presentation to Scheunert. This verifies the exact mathematical content; it is not evidence that the family first appeared in 2010. Kac’s classification predates the seminar.

\paragraph{Present status.} Classical superalgebraic background, already known before the 1984 report and still valid with the stated conventions on 12 November 2026.

\paragraph{Assessment.} Makes parity, characteristic, and matrix conventions explicit.

\paragraph{Audit.} Literal correctness: correct with missing hypotheses. Publication match: exact. Verification: read the result. Development: annotated. Confidence: high.

\noindent References: \citep{poletaeva2010matrix,kac1977sketch}.\par

\subsection{Record 627: Typical Borel–Weil–Bott theory for special linear supergroups}

\label{talk:627}

\noindent\textbf{Date:} 1984-03-19. \textbf{Speaker:} I. B. Penkov. \textbf{Pages:} 1179.\par

Penkov’s heading is on P001179, after Leites. The report concerns \(SL(n|m)\), parity-dependent full flags, line bundles, and the typical-weight analogue of Borel–Weil–Bott. The exact result is documented in Penkov’s later full treatment; the short 1984 announcement is distinguished from the 1988/1990 exposition.

\subsubsection{1983-627.1 --- Flag stabilizers and parity sequences}

\label{item:1983-627.1}

\noindent\textbf{Kind:} definition. \textbf{Fate:} E. \textbf{Pages:} 1179.\par

\paragraph{Notebook content.} For \(SL(n|m)\), the report describes different Borels by strings of \(n\) plus signs and \(m\) minus signs.

\paragraph{Precise statement.} Let \(G=SL(n|m)\) over \(\mathbb C\), \(n,m>0\), \(n\ne m\). A complete flag in the free superspace has successive quotients of rank \(1|0\) or \(0|1\). Its parity sequence \(J\) has \(n\) even and \(m\) odd entries. The stabilizer \(B_J\) is a triangular Borel subsupergroup; up to conjugacy by the even group, these choices are classified by their parity sequences. This definition must not be replaced without qualification by “maximal solvable subsuperalgebra.”

\paragraph{Literature and fate.} Penkov (1988), \S1.2, develops the Borel and flag conventions. Manin and Voronov (1988/1990), introduction and \S1, explicitly warn that the ordinary maximal-solvable characterization does not transfer unchanged to supergroups.

\paragraph{Present status.} Established with the stated hypotheses as of 12 November 2026.

\paragraph{Assessment.} Separates the precise typical-weight theorem from statements that fail or require modification in supergeometry.

\paragraph{Audit.} Literal correctness: correct with missing hypotheses. Publication match: exact. Verification: read the result. Development: annotated. Confidence: high.

\noindent References: \citep{penkov1988borel,maninvoronov1988supercell}.\par

\subsubsection{1983-627.2 --- Dimension and reduced space of a complete superflag}

\label{item:1983-627.2}

\noindent\textbf{Kind:} construction. \textbf{Fate:} E. \textbf{Pages:} 1179.\par

\paragraph{Notebook content.} The flag space \(G/B_J\) is used as the geometric carrier of the representation.

\paragraph{Precise statement.} The reduced variety of the complete flag superspace is \(\operatorname{Fl}(\mathbb C^n)\times\operatorname{Fl}(\mathbb C^m)\). Its superdimension is \(\bigl(n(n-1)/2+m(m-1)/2\,\big|\,nm\bigr)\), independently of the ordering \(J\), although the full superspace and its sheaf of functions depend on \(J\).

\paragraph{Literature and fate.} The flag construction is given in Manin and Voronov, \S1; the dimension follows by counting positive even and odd roots. Equality of reduced spaces does not imply isomorphism of the superspaces.

\paragraph{Present status.} Established with the stated hypotheses as of 12 November 2026.

\paragraph{Assessment.} Separates the precise typical-weight theorem from statements that fail or require modification in supergeometry.

\paragraph{Audit.} Literal correctness: correct with missing hypotheses. Publication match: partial. Verification: read the result. Development: annotated. Confidence: high.

\noindent References: \citep{maninvoronov1988supercell}.\par

\subsubsection{1983-627.3 --- Characters and equivariant line bundles}

\label{item:1983-627.3}

\noindent\textbf{Kind:} construction. \textbf{Fate:} E. \textbf{Pages:} 1179.\par

\paragraph{Notebook content.} Line bundles are identified with characters of a Cartan or Borel subgroup.

\paragraph{Precise statement.} For a fixed \(B_J\) and its diagonal torus \(T\), even \(G\)-equivariant line bundles on \(G/B_J\) are classified by algebraic characters of \(B_J\), equivalently characters of \(T\), through the associated-bundle construction. One must specify equivariance and parity. This statement alone does not identify the entire ordinary Picard group of an arbitrary supermanifold.

\paragraph{Editorial reconstruction.} Only the justified equivariant Picard assertion is used; no unverified claim about all non-equivariant super line bundles is added.

\paragraph{Literature and fate.} Penkov (1988), \S1, constructs the homogeneous sheaves from inducing representations. The precise character statement is standard homogeneous-bundle theory; the notebook’s bare Picard notation requires the equivariance qualification.

\paragraph{Present status.} Established with the stated hypotheses as of 12 November 2026.

\paragraph{Assessment.} Separates the precise typical-weight theorem from statements that fail or require modification in supergeometry.

\paragraph{Audit.} Literal correctness: correct with missing hypotheses. Publication match: partial. Verification: read the result. Development: annotated. Confidence: high.

\noindent References: \citep{penkov1988borel}.\par

\subsubsection{1983-627.4 --- The typical Borel–Weil–Bott theorem}

\label{item:1983-627.4}

\noindent\textbf{Kind:} theorem. \textbf{Fate:} A. \textbf{Pages:} 1179.\par

\paragraph{Notebook content.} The displayed theorem excludes odd-root hyperplanes and places cohomology in one Weyl-group degree.

\paragraph{Precise statement.} Let \(B\) be a triangular Borel of \(SL(n|m)\), \(n\ne m\), and \(\mathcal L_\lambda=G\times_B\mathbb C_{-\lambda}\) an even homogeneous line bundle for an integral character. Put \(\rho_B=\frac12\sum_{\alpha\in\Delta^+_{\overline0}}\alpha-\frac12\sum_{\beta\in\Delta^+_{\overline1}}\beta\) and \(\chi=\lambda+\rho_B\). Assume \((\chi,\beta)\ne0\) for every odd root \(\beta\). If \(\chi\) is even-singular, all cohomology vanishes. If \(\chi\) is even-regular, let \(w\in S_n\times S_m\) send it into the dominant even chamber. Exactly degree \(\ell(w)\) can be nonzero, and that cohomology is irreducible. For the distinguished Borel, its contragredient has highest weight \(w\chi-\rho_B\); for other parity sequences one transports the highest weight and parity by odd reflections.

\paragraph{Editorial reconstruction.} The bundle sign and contragredient convention are stated explicitly. Atypical weights are outside the theorem; no vanishing statement is inferred for them. The simplified highest-weight formula is restricted to the distinguished Borel.

\paragraph{Literature and fate.} Penkov (1988), \S2.1, Theorem 1, Russian p. 93, states and proves the typical theorem. The early Penkov–Skornyakov announcement is C. R. Acad. Sci. Paris 299 (1984), 1005–1008; that original note was identified but not read. The 1988 Russian treatment appeared in English in 1990.

\paragraph{Present status.} Established with the stated hypotheses as of 12 November 2026.

\paragraph{Assessment.} Separates the precise typical-weight theorem from statements that fail or require modification in supergeometry.

\paragraph{Audit.} Literal correctness: correct with missing hypotheses. Publication match: exact. Verification: read the result. Development: annotated. Confidence: high.

\noindent References: \citep{penkov1988borel,penkovskornyakov1984borel}.\par

\subsubsection{1983-627.5 --- Changing a Borel by an odd reflection}

\label{item:1983-627.5}

\noindent\textbf{Kind:} method. \textbf{Fate:} A. \textbf{Pages:} 1179.\par

\paragraph{Notebook content.} The weight notation is shifted by a parity-sequence-dependent \(\rho_J\), and the different flags are compared.

\paragraph{Precise statement.} If an isotropic odd simple root \(\beta\) is reflected, then the positive-root system changes by \(\beta\mapsto-\beta\) and \(\rho_{B'}=\rho_B+\beta\). For a typical irreducible module, the highest weight changes from \(\lambda\) to \(\lambda-\beta\), and a highest vector changes parity. Thus \(\lambda+\rho_B\) is preserved; unqualified equality of graded representations across all flags omits a parity choice.

\paragraph{Literature and fate.} Penkov (1988), \S1.2 and the odd-reflection steps in the proof of \S2.1, gives the mechanism underlying the flag-dependent notation. This is a method in the proved typical theory, not a claim that all superflags are isomorphic.

\paragraph{Present status.} Established with the stated hypotheses as of 12 November 2026.

\paragraph{Assessment.} Separates the precise typical-weight theorem from statements that fail or require modification in supergeometry.

\paragraph{Audit.} Literal correctness: correct with missing hypotheses. Publication match: partial. Verification: read the result. Development: annotated. Confidence: high.

\noindent References: \citep{penkov1988borel}.\par

\subsection{Record 628: Superflags, homogeneous line bundles, and affine flag constructions}

\label{talk:628}

\noindent\textbf{Date:} 1984-03-26. \textbf{Speaker:} I. B. Penkov. \textbf{Pages:} 1180--1181.\par

The continuation occupies P001180–P001181. It repeats the typical theorem and character construction, mentions a differential-operator version, and illustrates flags and additional volume data. Here “affine flags” means flags with determinant data or the homogeneous space \(G/N\), not the affine flag variety of a loop group.

\subsubsection{1983-628.1 --- The repeated typical cohomology statement}

\label{item:1983-628.1}

\noindent\textbf{Kind:} theorem. \textbf{Fate:} A. \textbf{Pages:} 1180.\par

\paragraph{Notebook content.} The line-bundle theorem from 19 March is restated with the shifted weight \(\chi-\rho_J\).

\paragraph{Precise statement.} With the hypotheses and conventions of item 1983-627.4, the cohomology of the corresponding homogeneous line bundle is zero for even-singular \(\chi\), and otherwise is concentrated in degree \(\ell(w)\), with an irreducible typical module. Character integrality, the odd-root nonvanishing condition, and the even Weyl-group regularity are separate requirements.

\paragraph{Editorial reconstruction.} The earlier item supplies the precise convention; the repeated notebook formula is not used to erase parity or contragredient issues.

\paragraph{Literature and fate.} This is a second occurrence of the same typical theorem, matched to Penkov (1988), \S2.1, Theorem 1. It is not a new open question and receives no new problem identifier.

\paragraph{Present status.} Established with the stated hypotheses as of 12 November 2026.

\paragraph{Assessment.} Separates the precise typical-weight theorem from statements that fail or require modification in supergeometry.

\paragraph{Audit.} Literal correctness: correct with missing hypotheses. Publication match: exact. Verification: read the result. Development: annotated. Confidence: high.

\noindent References: \citep{penkov1988borel}.\par

\subsubsection{1983-628.2 --- The announced differential-operator version}

\label{item:1983-628.2}

\noindent\textbf{Kind:} method. \textbf{Fate:} A. \textbf{Pages:} 1180.\par

\paragraph{Notebook content.} The speaker says there is also a version involving \(D\)-modules.

\paragraph{Precise statement.} The identifiable direction is to relate representations to sheaves of differential operators on superflag varieties. A precise equivalence would require a twisting parameter, a specified module category, and regularity or typicality assumptions. The record does not state these data or an equivalence theorem.

\paragraph{Editorial reconstruction.} The defined direction is matched bibliographically; the notebook does not state enough hypotheses to reproduce the full later theorem.

\paragraph{Literature and fate.} Penkov (1988), bibliography item 22, explicitly identifies Penkov–Skornyakov, Cohomologie des D-modules tordus typiques sur les supervarietes de drapeaux, C. R. Acad. Sci. Paris 299 (1984), 1005–1008. Its title matches this announcement; the original note was not obtained. Bibliography item 21 separately identifies Penkov’s 1987 localization note. Thus the direction has a concrete published continuation, while its exact category-equivalence hypotheses remain unverified here.

\paragraph{Present status.} An exact contemporary paper on typical twisted D-module cohomology is identified, but its full theorem was not read; no unverified localization equivalence is asserted.

\paragraph{Assessment.} Separates the precise typical-weight theorem from statements that fail or require modification in supergeometry.

\paragraph{Audit.} Literal correctness: not determinable. Publication match: partial. Verification: indirect. Development: annotated. Confidence: medium.

\noindent References: \citep{penkov1988borel,penkovskornyakov1984borel}.\par

\subsubsection{1983-628.3 --- Flags with volume data and the space G/N}

\label{item:1983-628.3}

\noindent\textbf{Kind:} construction. \textbf{Fate:} E. \textbf{Pages:} 1180--1181.\par

\paragraph{Notebook content.} Flags are supplemented with forms of volume, and an affine flag space is described.

\paragraph{Precise statement.} For an ordinary connected reductive group \(G\), a Borel \(B=TN\), and unipotent radical \(N\), the quotient \(G/N\to G/B\) is a principal \(T\)-bundle; for \(SL_r\), it can be interpreted as complete flags with compatible nonzero determinant data on their subspaces. Its analogous super construction must use the appropriate Berezinian lines and parity conventions. This homogeneous space is distinct from an affine flag variety attached to a loop group.

\paragraph{Editorial reconstruction.} The ordinary quotient is stated completely. The notebook’s super volume terminology is normalized to Berezinian data without claiming an unrecorded global theorem.

\paragraph{Literature and fate.} The quotient construction is standard and supplies the precise ordinary case. Penkov’s and Manin–Voronov’s flag formalism provides the supergeometric setting; an exact unrestricted identification using ordinary volume forms in every odd direction is not asserted.

\paragraph{Present status.} Established with the stated hypotheses as of 12 November 2026.

\paragraph{Assessment.} Separates the precise typical-weight theorem from statements that fail or require modification in supergeometry.

\paragraph{Audit.} Literal correctness: correct with missing hypotheses. Publication match: partial. Verification: read the result. Development: annotated. Confidence: high.

\noindent References: \citep{penkov1988borel,maninvoronov1988supercell}.\par

\subsubsection{1983-628.4 --- Monotone ranks in a flag and the underlying ordinary flags}

\label{item:1983-628.4}

\noindent\textbf{Kind:} example. \textbf{Fate:} E. \textbf{Pages:} 1181.\par

\paragraph{Notebook content.} An example chain contains a step whose even rank decreases, and the even and odd flag components are illustrated.

\paragraph{Precise statement.} In a flag of locally free direct summands, both ranks are nondecreasing: \(V^{a|b}\subset V^{c|d}\) requires \(a\le c\) and \(b\le d\). Thus a chain \(T^{1|0}\subset T^{2|0}\subset T^{1|1}\) cannot be the intended flag; \(T^{2|1}\) is the natural next rank if the last step adds one odd dimension. The reduced full flag remains the product of the even and odd ordinary full flags.

\paragraph{Editorial reconstruction.} The correction (2|1) is the minimal rank-compatible reconstruction; the original image should decide whether the error is scribal or transcriptional.

\paragraph{Literature and fate.} This is a rank check, not a new theorem. The example is interpreted using the general flag definition in Manin and Voronov, \S1. The proposed corrected superscript is conditional on the supplied transcription because the final scan recheck was blocked.

\paragraph{Present status.} Established with the stated hypotheses as of 12 November 2026.

\paragraph{Assessment.} Separates the precise typical-weight theorem from statements that fail or require modification in supergeometry.

\paragraph{Audit.} Literal correctness: false as written. Publication match: exact. Verification: read the result. Development: annotated. Confidence: medium.

\noindent References: \citep{maninvoronov1988supercell}.\par

\subsection{Record 629: Verma modules and singular vectors for affine sl2}

\label{talk:629}

\noindent\textbf{Date:} 1984-04-02. \textbf{Speaker:} D. B. Fuchs. \textbf{Pages:} 1182--1183 upper.\par

Fuchs’s dated report begins on P001182 and ends in the upper part of P001183, before Beilinson’s heading. It defines the derived affine algebra, Verma modules and reducibility, asks about Casimirs, and introduces the formal-power method later published by Malikov, Feigin, and Fuchs. The notes mix a lowest-weight definition with a highest-weight reducibility formula; both conventions are separated below.

\subsubsection{1983-629.1 --- The six-generator derived affine algebra}

\label{item:1983-629.1}

\noindent\textbf{Kind:} definition. \textbf{Fate:} E. \textbf{Pages:} 1182.\par

\paragraph{Notebook content.} Six Chevalley generators and the affine \(2\times2\) Cartan matrix are written.

\paragraph{Precise statement.} Let \(A=(a_{ij})=\begin{pmatrix}2&-2\\-2&2\end{pmatrix}\). Over \(\mathbb C\), impose \([h_i,h_j]=0\), \([h_i,e_j]=a_{ij}e_j\), \([h_i,f_j]=-a_{ij}f_j\), \([e_i,f_j]=\delta_{ij}h_i\), and \((\operatorname{ad}e_i)^3e_j=(\operatorname{ad}f_i)^3f_j=0\) for \(i\ne j\). This gives \(\widehat{\mathfrak{sl}}_2'=(\mathfrak{sl}_2\otimes\mathbb C[t,t^{-1}])\oplus\mathbb Cc\), with \(c=h_1+h_2\). The loop-degree derivation is an additional generator, absent from this presentation.

\paragraph{Literature and fate.} Malikov, Feigin, and Fuchs (1986), \S1 and opening of \S3, use this exact presentation. The distinction between the derived algebra and its extension by a derivation is necessary.

\paragraph{Present status.} Established with the stated algebraic conventions as of 12 November 2026.

\paragraph{Assessment.} Resolves the highest-versus-lowest-weight convention and distinguishes finite algebraic expressions from formal completions.

\paragraph{Audit.} Literal correctness: correct with missing hypotheses. Publication match: exact. Verification: read the result. Development: annotated. Confidence: high.

\noindent References: \citep{malikovfeiginfuchs1986singular}.\par

\subsubsection{1983-629.2 --- Highest- and lowest-weight Verma modules}

\label{item:1983-629.2}

\noindent\textbf{Kind:} definition. \textbf{Fate:} E. \textbf{Pages:} 1182.\par

\paragraph{Notebook content.} The generating vector in the notebook is killed by the negative generators.

\paragraph{Precise statement.} For a highest-weight convention, \(M(\mu)=U(\widehat{\mathfrak{sl}}_2')\otimes_{U(\mathfrak h+\mathfrak n_+)}\mathbb C_\mu\), where \(e_i v_\mu=0\) and \(h_i v_\mu=\mu_i v_\mu\). As a vector space it is \(U(\mathfrak n_-)v_\mu\). The notebook’s convention \(f_i v=0\), \(h_i v=\lambda_i v\) is lowest weight; the Chevalley involution \(e_i\leftrightarrow f_i,\ h_i\mapsto-h_i\) converts it to the highest weight \(\mu=-\lambda\).

\paragraph{Editorial reconstruction.} The two definitions are both legitimate; only mixing their parameter signs is erroneous.

\paragraph{Literature and fate.} Malikov–Feigin–Fuchs (1986), \S1, defines the highest-weight convention used in its formulas. The conversion is algebraic and must precede comparison with the notebook.

\paragraph{Present status.} Established with the stated algebraic conventions as of 12 November 2026.

\paragraph{Assessment.} Resolves the highest-versus-lowest-weight convention and distinguishes finite algebraic expressions from formal completions.

\paragraph{Audit.} Literal correctness: correct. Publication match: exact. Verification: read the result. Development: annotated. Confidence: high.

\noindent References: \citep{malikovfeiginfuchs1986singular}.\par

\subsubsection{1983-629.3 --- Reducibility and nontrivial singular vectors}

\label{item:1983-629.3}

\noindent\textbf{Kind:} criterion. \textbf{Fate:} E. \textbf{Pages:} 1182.\par

\paragraph{Notebook content.} Reducibility is related to a vector annihilated by the same nilpotent subalgebra as the vacuum.

\paragraph{Precise statement.} A root-graded Verma module is reducible if and only if it contains a nonzero singular vector not proportional to its generating vector. In highest-weight notation this means \(e_1w=e_2w=0\), \(w\notin\mathbb Cv_\mu\); in the lowest-weight convention the \(f_i\) replace the \(e_i\). The nontriviality condition is essential because the vacuum is always singular.

\paragraph{Literature and fate.} Malikov–Feigin–Fuchs (1986), \S1, immediately before Theorem 1.1, states this criterion. The root grading, or an extension by the degree derivation, keeps the positive-height argument distinct from the degenerate affine Cartan weight alone.

\paragraph{Present status.} Established with the stated algebraic conventions as of 12 November 2026.

\paragraph{Assessment.} Resolves the highest-versus-lowest-weight convention and distinguishes finite algebraic expressions from formal completions.

\paragraph{Audit.} Literal correctness: correct with missing hypotheses. Publication match: exact. Verification: read the result. Development: annotated. Confidence: high.

\noindent References: \citep{malikovfeiginfuchs1986singular}.\par

\subsubsection{1983-629.4 --- The affine sl2 reducibility hyperplanes}

\label{item:1983-629.4}

\noindent\textbf{Kind:} theorem. \textbf{Fate:} E. \textbf{Pages:} 1182.\par

\paragraph{Notebook content.} The notebook gives two families of affine-linear equations and a critical-level equation, but permits \(n=0\).

\paragraph{Precise statement.} In the highest-weight convention, \(M(\mu_1,\mu_2)\) is reducible precisely if, for some integer \(k\ge0\), either \((k+1)(\mu_1+1)+k(\mu_2+1)\in\mathbb Z_{>0}\), or \(k(\mu_1+1)+(k+1)(\mu_2+1)\in\mathbb Z_{>0}\), or \(\mu_1+\mu_2+2=0\). In the notebook’s lowest-weight convention substitute \(\mu_i=-\lambda_i\); the critical equation becomes \(\lambda_1+\lambda_2=2\). Zero is excluded from the real-root integer condition.

\paragraph{Editorial reconstruction.} The notebook’s signs and its \(n\ge0\) condition cannot both be retained with its lowest-weight definition.

\paragraph{Literature and fate.} Malikov–Feigin–Fuchs (1986), \S3, Proposition 3.1, Russian p. 29, states the exact specialization of the earlier Kac–Kazhdan criterion. Its indexing uses \(k>0\); replacing \(k\) by \(k+1\) gives the displayed form. The paper’s \S1, Theorem 1.1, identifies the earlier result.

\paragraph{Present status.} Established with the stated algebraic conventions as of 12 November 2026.

\paragraph{Assessment.} Resolves the highest-versus-lowest-weight convention and distinguishes finite algebraic expressions from formal completions.

\paragraph{Audit.} Literal correctness: false as written. Publication match: exact. Verification: read the result. Development: annotated. Confidence: high.

\noindent References: \citep{malikovfeiginfuchs1986singular}.\par

\subsubsection{1983-629.5 --- Casimir polynomials and the completion issue}

\label{item:1983-629.5}

\noindent\textbf{Kind:} question. \textbf{Fate:} F. \textbf{Pages:} 1182.\par

\paragraph{Notebook content.} Gelfand asks about Casimir expressions in the generators.

\paragraph{Precise statement.} On the literal finite-polynomial reading, the question is the center of \(U(\widehat{\mathfrak{sl}}_2')\), and the answer is \(\mathbb C[c]\). If infinite normally ordered expressions or their action at fixed level are intended, a topology and level must instead be specified. These are different algebraic questions.

\paragraph{Editorial reconstruction.} The PBW argument is our annotation, not an attribution of a new theorem or of the later completed-center theory to the seminar. No new open-problem ID is assigned to the ambiguous remainder.

\paragraph{Literature and fate.} The finite-polynomial answer has an elementary PBW proof: the highest symbol of a central element is a polynomial in finitely many loop modes. Commuting with a mode of sufficiently large degree gives independent new mode variables; their coefficients force every derivative with respect to the original \(\mathfrak{sl}_2\) modes to vanish, since \(\mathfrak{sl}_2\) has zero center. The symbol depends only on \(c\), and induction on PBW degree gives \(\mathbb C[c]\). Frenkel (2007), \S2.1.2 and \S3.5, concerns the different completed question; \S4.3, Proposition 4.3.9, gives the scalar center away from critical level.

\paragraph{Present status.} The finite-polynomial interpretation is answered by the elementary argument; the precise intended completed or representation-level Casimir question is not determined by the short record.

\paragraph{Assessment.} Resolves the highest-versus-lowest-weight convention and distinguishes finite algebraic expressions from formal completions.

\paragraph{Audit.} Literal correctness: correct with missing hypotheses. Publication match: partial. Verification: read the result. Development: annotated. Confidence: medium.

\noindent References: \citep{frenkel2007loop}.\par

\subsubsection{1983-629.6 --- The Malikov–Feigin–Fuchs singular-vector product}

\label{item:1983-629.6}

\noindent\textbf{Kind:} construction. \textbf{Fate:} A. \textbf{Pages:} 1183.\par

\paragraph{Notebook content.} Products with nonintegral exponents are proposed as a way to write singular vectors.

\paragraph{Precise statement.} For integers \(n,k>0\) and \(t\in\mathbb C\), let \(F(n,k,t)\) be the alternating product of \(2k-1\) factors beginning and ending with \(f_1\), with exponents \(n+(k-1)t,n+(k-2)t,\ldots,n-(k-1)t\). For \(\mu_1=n-1-(k-1)t\), \(\mu_2=-n-1+kt\), this formal product has a well-defined value in \(U(\mathfrak n_-)\), and \(F(n,k,t)v_\mu\) is a singular vector of degree \(n(k\alpha_1+(k-1)\alpha_2)\). Interchanging the two simple roots gives the other family.

\paragraph{Editorial reconstruction.} The precise product length, parameter restrictions, and highest-weight convention are supplied from the matched theorem.

\paragraph{Literature and fate.} Malikov, Feigin, and Fuchs (1986), \S3, Theorem 3.2, equations (6)–(7), is an exact match; the paper was received 26 December 1984. It proves that the formal expression becomes an ordinary enveloping-algebra element rather than treating arbitrary complex powers as analytic operators.

\paragraph{Present status.} Established with the stated algebraic conventions as of 12 November 2026.

\paragraph{Assessment.} Resolves the highest-versus-lowest-weight convention and distinguishes finite algebraic expressions from formal completions.

\paragraph{Audit.} Literal correctness: correct with missing hypotheses. Publication match: exact. Verification: read the result. Development: annotated. Confidence: high.

\noindent References: \citep{malikovfeiginfuchs1986singular}.\par

\subsubsection{1983-629.7 --- Generalized binomial commutation}

\label{item:1983-629.7}

\noindent\textbf{Kind:} method. \textbf{Fate:} A. \textbf{Pages:} 1183.\par

\paragraph{Notebook content.} A generalized commutator formula is used to manipulate nonintegral powers.

\paragraph{Precise statement.} If the iterated right commutators \(D_y^j(x)=[\cdots[[x,y],y],\ldots,y]\) eventually vanish, the formal rule \([x,y^z]=\sum_{j\ge1}\binom zj y^{z-j}D_y^j(x)\) is finite. For integer \(z\ge0\) it follows by induction; polynomial continuation of the coefficients defines the formal-power calculation in a suitable localization or ordered-symbol algebra. It does not define a complex power of every operator in every representation.

\paragraph{Literature and fate.} Malikov–Feigin–Fuchs (1986), \S2, Propositions 2.1–2.2, and the universal formula (8) in the proof of Theorem 3.2 provide the exact computational method.

\paragraph{Present status.} Established with the stated algebraic conventions as of 12 November 2026.

\paragraph{Assessment.} Resolves the highest-versus-lowest-weight convention and distinguishes finite algebraic expressions from formal completions.

\paragraph{Audit.} Literal correctness: correct with missing hypotheses. Publication match: exact. Verification: read the result. Development: annotated. Confidence: high.

\noindent References: \citep{malikovfeiginfuchs1986singular}.\par

\subsection{Record 630: Mixed Hodge structures and absolute Hodge cohomology}

\label{talk:630}

\noindent\textbf{Date:} 1984-04-02. \textbf{Speaker:} A. A. Beilinson. \textbf{Pages:} 1183 lower--1185.\par

Beilinson’s heading is in the lower part of P001183, after Fuchs, and the report continues through P001185. It introduces pure and mixed real Hodge structures, canonical complex splittings, an operator comparing conjugate splittings, Tate objects, and Ext groups, ending with a derived-cohomology programme. The category and coefficient ring must be kept explicit; the checked Ext formula below is for real mixed Hodge structures.

\subsubsection{1983-630.1 --- Pure real Hodge structures}

\label{item:1983-630.1}

\noindent\textbf{Kind:} definition. \textbf{Fate:} E. \textbf{Pages:} 1183.\par

\paragraph{Notebook content.} The report defines a decomposition indexed by two integers and a conjugation condition.

\paragraph{Precise statement.} A pure real Hodge structure of weight \(n\) is a finite-dimensional real vector space \(H_{\mathbb R}\) with \(H_{\mathbb C}=\bigoplus_{p+q=n}H^{p,q}\) and \(\overline{H^{p,q}}=H^{q,p}\). Equivalently, a finite decreasing filtration \(F\) satisfies \(H_{\mathbb C}=F^pH_{\mathbb C}\oplus\overline{F^{n-p+1}H_{\mathbb C}}\) for every \(p\).

\paragraph{Literature and fate.} Deligne (1971), \S2.1, supplies the classical definition. Bunke, Nikolaus, and Tamme (2018), Definition 2.1, states the real version with these conventions.

\paragraph{Present status.} Established with real coefficients and the stated filtration conventions as of 12 November 2026.

\paragraph{Assessment.} Makes coefficient, weight, and derived-category conventions explicit without turning an exercise into an open problem.

\paragraph{Audit.} Literal correctness: correct. Publication match: exact. Verification: read the result. Development: annotated. Confidence: high.

\noindent References: \citep{deligne1971hodge,bunkenikolaustamme2018regulator}.\par

\subsubsection{1983-630.2 --- Mixed real Hodge structures}

\label{item:1983-630.2}

\noindent\textbf{Kind:} definition. \textbf{Fate:} E. \textbf{Pages:} 1183--1184.\par

\paragraph{Notebook content.} The direct-sum picture is modified by a weight filtration and a conjugation condition modulo lower terms.

\paragraph{Precise statement.} A mixed real Hodge structure is a finite-dimensional \(H_{\mathbb R}\), a finite increasing real filtration \(W\), and a finite decreasing filtration \(F\) on \(H_{\mathbb C}\), such that the induced \(F\) on \(\operatorname{gr}^W_nH_{\mathbb C}\) defines a pure real Hodge structure of weight \(n\) for every \(n\). Exhaustiveness and separatedness are included in “finite filtration.”

\paragraph{Literature and fate.} Deligne (1971), \S2.3, defines mixed Hodge structures. Bunke–Nikolaus–Tamme (2018), Definition 2.2, gives the same real-coefficient convention.

\paragraph{Present status.} Established with real coefficients and the stated filtration conventions as of 12 November 2026.

\paragraph{Assessment.} Makes coefficient, weight, and derived-category conventions explicit without turning an exercise into an open problem.

\paragraph{Audit.} Literal correctness: correct. Publication match: exact. Verification: read the result. Development: annotated. Confidence: high.

\noindent References: \citep{deligne1971hodge,bunkenikolaustamme2018regulator}.\par

\subsubsection{1983-630.3 --- Deligne’s canonical complex bigrading}

\label{item:1983-630.3}

\noindent\textbf{Kind:} construction. \textbf{Fate:} E. \textbf{Pages:} 1184.\par

\paragraph{Notebook content.} A bigrading is used to represent the mixed structure and to compare it with its conjugate.

\paragraph{Precise statement.} There is a unique functorial decomposition \(H_{\mathbb C}=\bigoplus I^{p,q}\) with \(F^p=\bigoplus_{r\ge p}I^{r,s}\), \(W_n=\bigoplus_{r+s\le n}I^{r,s}\), and \(\overline{I^{p,q}}\equiv I^{q,p}\pmod{\bigoplus_{r<q,\ s<p}I^{r,s}}\). This splits \(W\) over \(\mathbb C\) compatibly with \(F\). The summands need not be conjugate exactly, so this decomposition does not identify the original mixed Hodge structure with a direct sum of pure real structures.

\paragraph{Editorial reconstruction.} The strict inequalities in the conjugation remainder specify the otherwise compressed phrase “lower terms.”

\paragraph{Literature and fate.} Deligne (1971), \S1.2, Lemma 1.2.11, constructs the two functorial complex bigradings used in the mixed-Hodge setting of \S2.3. The conjugation refinement is the standard Deligne splitting convention; it must not be confused with a real splitting preserving the original Hodge filtration.

\paragraph{Present status.} Established with real coefficients and the stated filtration conventions as of 12 November 2026.

\paragraph{Assessment.} Makes coefficient, weight, and derived-category conventions explicit without turning an exercise into an open problem.

\paragraph{Audit.} Literal correctness: correct with missing hypotheses. Publication match: exact. Verification: read the result. Development: annotated. Confidence: high.

\noindent References: \citep{deligne1971hodge}.\par

\subsubsection{1983-630.4 --- Real splitting and exactness in the category}

\label{item:1983-630.4}

\noindent\textbf{Kind:} result. \textbf{Fate:} E. \textbf{Pages:} 1184.\par

\paragraph{Notebook content.} Split structures are contrasted with general mixed structures, and categorical properties are discussed.

\paragraph{Precise statement.} A mixed real Hodge structure is split over \(\mathbb R\) precisely when its canonical bigrading satisfies \(\overline{I^{p,q}}=I^{q,p}\); then it is a direct sum of pure real Hodge structures. Morphisms preserving \(W\) and \(F\) are strict for both filtrations, and kernels and cokernels with induced filtrations make the category \(\mathrm{MHS}_{\mathbb R}\) abelian. A functorial real splitting obtained by changing \(F\) must not be confused with an isomorphism preserving the original \(F\).

\paragraph{Literature and fate.} Deligne (1971), \S2.3, Corollary 2.3.5, gives strictness and the abelian-category statement; the splitting condition follows from the bigrading.

\paragraph{Present status.} Established with real coefficients and the stated filtration conventions as of 12 November 2026.

\paragraph{Assessment.} Makes coefficient, weight, and derived-category conventions explicit without turning an exercise into an open problem.

\paragraph{Audit.} Literal correctness: correct with missing hypotheses. Publication match: exact. Verification: read the result. Development: annotated. Confidence: high.

\noindent References: \citep{deligne1971hodge}.\par

\subsubsection{1983-630.5 --- Mixed Hodge structures on algebraic cohomology}

\label{item:1983-630.5}

\noindent\textbf{Kind:} theorem. \textbf{Fate:} E. \textbf{Pages:} 1184.\par

\paragraph{Notebook content.} Cohomology of algebraic varieties is given as the main source of mixed Hodge structures.

\paragraph{Precise statement.} For a separated complex algebraic variety of finite type, each singular cohomology group \(H^j(X,\mathbb R)\) has a functorial mixed real Hodge structure. If \(X\) is smooth and projective, this structure is pure of weight \(j\). Singularities and nonproperness generally require mixed weights.

\paragraph{Literature and fate.} Deligne (1971), Theorem 3.2.5(iii) and Definition 3.2.12, treats the smooth nonproper case; Deligne (1974), Proposition 8.2.2, proves the general separated finite-type construction and functoriality. These results predate the seminar.

\paragraph{Present status.} Established with real coefficients and the stated filtration conventions as of 12 November 2026.

\paragraph{Assessment.} Makes coefficient, weight, and derived-category conventions explicit without turning an exercise into an open problem.

\paragraph{Audit.} Literal correctness: correct with missing hypotheses. Publication match: exact. Verification: read the result. Development: annotated. Confidence: high.

\noindent References: \citep{deligne1971hodge,deligne1974hodge}.\par

\subsubsection{1983-630.6 --- The operator comparing conjugate splittings}

\label{item:1983-630.6}

\noindent\textbf{Kind:} construction. \textbf{Fate:} E. \textbf{Pages:} 1184.\par

\paragraph{Notebook content.} An operator called arithmetic monodromy compares a splitting with its conjugate and is said to become trivial in the split case.

\paragraph{Precise statement.} Put \(J^{p,q}=\overline{I^{q,p}}\). Both \(I\) and \(J\) split the same weight-graded Hodge spaces. Define \(T\) to carry \(I^{p,q}\) to \(J^{p,q}\) by the identity on \(\operatorname{gr}^W_{p+q}\). Then \(T\) is unipotent, \((T-1)I^{p,q}\subset\bigoplus_{r<p,\ s<q}I^{r,s}\), and \(\overline T=T^{-1}\). In particular \((T-1)W_n\subset W_{n-2}\), and \(T=1\) exactly for a structure split over \(\mathbb R\).

\paragraph{Editorial reconstruction.} The operator is defined explicitly from the two splittings. This supplies a precise interpretation of the notebook terminology without asserting that Deligne uses that name.

\paragraph{Literature and fate.} These assertions follow directly from the canonical bigrading and its conjugation condition in Deligne’s splitting formalism. The term “arithmetic monodromy” in the notebook is not identified with topological monodromy of a family or with the nilpotent operator of a degeneration.

\paragraph{Present status.} Established with real coefficients and the stated filtration conventions as of 12 November 2026.

\paragraph{Assessment.} Makes coefficient, weight, and derived-category conventions explicit without turning an exercise into an open problem.

\paragraph{Audit.} Literal correctness: correct with missing hypotheses. Publication match: partial. Verification: read the result. Development: annotated. Confidence: medium.

\noindent References: \citep{deligne1971hodge}.\par

\subsubsection{1983-630.7 --- Rational and integral coefficient structures}

\label{item:1983-630.7}

\noindent\textbf{Kind:} definition. \textbf{Fate:} E. \textbf{Pages:} 1184--1185.\par

\paragraph{Notebook content.} The coefficient ring is allowed to be \(\mathbb Z\), \(\mathbb Q\), or \(\mathbb R\).

\paragraph{Precise statement.} For rational coefficients, replace \(H_{\mathbb R}\) by a finite-dimensional rational vector space with rational \(W\), and impose the Hodge condition after tensoring with \(\mathbb R\) and \(\mathbb C\). An integral structure additionally specifies a finitely generated lattice; if only free lattices are allowed, the resulting category is not abelian under arbitrary cokernels. Statements about integral Ext therefore require an explicit convention about torsion and the exact category.

\paragraph{Editorial reconstruction.} The torsion qualification is added because the notebook varies the coefficient ring without separately defining each homological category.

\paragraph{Literature and fate.} Deligne’s coefficient conventions explain the rational and integral forms. The real Ext computation later in this report is not silently asserted unchanged in every integral category.

\paragraph{Present status.} Established with real coefficients and the stated filtration conventions as of 12 November 2026.

\paragraph{Assessment.} Makes coefficient, weight, and derived-category conventions explicit without turning an exercise into an open problem.

\paragraph{Audit.} Literal correctness: correct with missing hypotheses. Publication match: partial. Verification: read the result. Development: annotated. Confidence: high.

\noindent References: \citep{deligne1971hodge}.\par

\subsubsection{1983-630.8 --- Tate objects and the sign of the Hodge type}

\label{item:1983-630.8}

\noindent\textbf{Kind:} definition. \textbf{Fate:} E. \textbf{Pages:} 1185.\par

\paragraph{Notebook content.} The rank-one object \(A(i)=(2\pi i)^iA\) is written with a Hodge bidegree.

\paragraph{Precise statement.} With the standard cohomological Tate convention, \(A(i)\) has Hodge type \((-i,-i)\) and weight \(-2i\), and tensoring by it changes bidegrees by \((p,q)\mapsto(p-i,q-i)\). A convention using type \((i,i)\) is the opposite twist \(A(-i)\); the period factor alone does not determine a sign convention unless the filtration is stated.

\paragraph{Literature and fate.} Bunke–Nikolaus–Tamme (2018), Example 2.3, explicitly fixes the standard real Tate object and weight. The notebook’s opposite-looking bidegree is recorded as a convention ambiguity rather than automatically attributed to OCR.

\paragraph{Present status.} Established with real coefficients and the stated filtration conventions as of 12 November 2026.

\paragraph{Assessment.} Makes coefficient, weight, and derived-category conventions explicit without turning an exercise into an open problem.

\paragraph{Audit.} Literal correctness: not determinable. Publication match: partial. Verification: read the result. Development: annotated. Confidence: medium.

\noindent References: \citep{bunkenikolaustamme2018regulator}.\par

\subsubsection{1983-630.9 --- Absolute Hodge cohomology of a mixed structure}

\label{item:1983-630.9}

\noindent\textbf{Kind:} definition. \textbf{Fate:} A. \textbf{Pages:} 1185.\par

\paragraph{Notebook content.} Cohomology is defined using derived Hom from the trivial Hodge structure.

\paragraph{Precise statement.} For \(M\in\mathrm{MHS}_{\mathbb R}\), set \(H^i_{\mathrm H}(M)=\operatorname{Ext}^i_{\mathrm{MHS}_{\mathbb R}}(\mathbb R(0),M)\). The zeroth group is \(W_0M_{\mathbb R}\cap F^0M_{\mathbb C}\), viewed inside \(M_{\mathbb C}\). For a bounded complex \(K\) of mixed Hodge structures, use \(\operatorname{Hom}_{D^b(\mathrm{MHS}_{\mathbb R})}(\mathbb R(0),K[i])\).

\paragraph{Literature and fate.} Beilinson’s absolute-Hodge construction was published in 1986. Bunke–Nikolaus–Tamme (2018), introduction and \S2, explicitly identifies this derived formulation and cites Beilinson’s Theorem 3.4. The original 1986 chapter was verified bibliographically but not read.

\paragraph{Present status.} Established with real coefficients and the stated filtration conventions as of 12 November 2026.

\paragraph{Assessment.} Makes coefficient, weight, and derived-category conventions explicit without turning an exercise into an open problem.

\paragraph{Audit.} Literal correctness: correct with missing hypotheses. Publication match: exact. Verification: indirect. Development: annotated. Confidence: high.

\noindent References: \citep{beilinson1986absolute,bunkenikolaustamme2018regulator}.\par

\subsubsection{1983-630.10 --- The first Ext quotient and vanishing in higher degrees}

\label{item:1983-630.10}

\noindent\textbf{Kind:} exercise. \textbf{Fate:} A. \textbf{Pages:} 1185.\par

\paragraph{Notebook content.} The quotient formula for \(H^1\) is set as an exercise, with vanishing above degree one.

\paragraph{Precise statement.} For every finite-dimensional real mixed Hodge structure \(M\), \(\operatorname{Ext}^1(\mathbb R(0),M)\cong W_0M_{\mathbb C}/(W_0M_{\mathbb R}+(F^0M_{\mathbb C}\cap W_0M_{\mathbb C}))\), as a real vector space, and \(\operatorname{Ext}^i(\mathbb R(0),M)=0\) for \(i>1\). These groups are computed by the two-term complex \(W_0M_{\mathbb R}\oplus(F^0M_{\mathbb C}\cap W_0M_{\mathbb C})\to W_0M_{\mathbb C}\), \((a,b)\mapsto a-b\), in degrees \(0,1\).

\paragraph{Editorial reconstruction.} The quotient is real-linear, not a quotient by complex vector subspaces. The exact checked assertion uses real coefficients.

\paragraph{Literature and fate.} Bunke–Nikolaus–Tamme (2018), Construction 2.8, equation (8), and Proposition 2.10, give the mapping-complex model; specialize its first argument to \(\mathbb R(0)\) and the second to a degree-zero structure. This directly verifies both the quotient and the higher vanishing. The original Beilinson chapter’s corresponding proof was not read.

\paragraph{Present status.} Established with real coefficients and the stated filtration conventions as of 12 November 2026.

\paragraph{Assessment.} Makes coefficient, weight, and derived-category conventions explicit without turning an exercise into an open problem.

\paragraph{Audit.} Literal correctness: correct with missing hypotheses. Publication match: exact. Verification: read the result. Development: annotated. Confidence: high.

\noindent Canonical problem: GP-1983-61.\par

\noindent References: \citep{beilinson1986absolute,bunkenikolaustamme2018regulator}.\par

\subsubsection{1983-630.11 --- A derived Hodge object for an algebraic variety}

\label{item:1983-630.11}

\noindent\textbf{Kind:} programme. \textbf{Fate:} A. \textbf{Pages:} 1185.\par

\paragraph{Notebook content.} The report aims to attach complexes with Hodge coefficients to varieties, in analogy with other cohomological formalisms.

\paragraph{Precise statement.} For a smooth complex algebraic variety \(X\), the absolute-Hodge formalism supplies an object \(R\Gamma_{\mathrm H}(X)\) in the derived category of mixed real Hodge structures whose cohomology is the ordinary cohomology of \(X\) with its mixed Hodge structure. One defines \(H^n_{\mathrm{abs}}(X,\mathbb R(p))=\operatorname{Hom}_{D^b(\mathrm{MHS}_{\mathbb R})}(\mathbb R(0),R\Gamma_{\mathrm H}(X)(p)[n])\). A chosen chain-level model is additional data.

\paragraph{Editorial reconstruction.} The smooth-variety case is stated precisely; the notebook’s wider analogy is not expanded into unrecorded comparison theorems.

\paragraph{Literature and fate.} Beilinson (1986), Théorem 3.4, is identified by Bunke–Nikolaus–Tamme (2018), Proposition 2.7, as the derived equivalence between complexes of mixed Hodge structures and mixed Hodge complexes. Their \S3 constructs a concrete absolute-Hodge model for smooth varieties. The original chapter was not read, so its exact result is verified indirectly through this explicit primary reuse.

\paragraph{Present status.} Established with real coefficients and the stated filtration conventions as of 12 November 2026.

\paragraph{Assessment.} Makes coefficient, weight, and derived-category conventions explicit without turning an exercise into an open problem.

\paragraph{Audit.} Literal correctness: correct with missing hypotheses. Publication match: partial. Verification: indirect. Development: annotated. Confidence: high.

\noindent References: \citep{beilinson1986absolute,bunkenikolaustamme2018regulator}.\par

\subsection{Record 631: Critical Ising fields and renormalization; a perverse-sheaf aside}

\label{talk:631}

\noindent\textbf{Date:} 1984-04-16. \textbf{Speaker:} Ya. G. Sinai. \textbf{Pages:} 1186--1187.\par

Separates the finite-volume model, phase regimes, critical block limits, and universality. A final MacPherson allusion supplies no reconstructible theorem.

\subsubsection{1983-631.1 --- Finite-volume Ising Gibbs measure}

\label{item:1983-631.1}

\noindent\textbf{Kind:} definition. \textbf{Fate:} E. \textbf{Pages:} 1186--1187.\par

\paragraph{Notebook content.} The notes introduce spins \(\sigma_x=\pm1\) and a Gibbs weight determined by pair interactions and inverse temperature.

\paragraph{Precise statement.} On finite \(\Lambda\), with specified boundary condition, \(P_\beta(\sigma)=Z_{\Lambda,\beta}^{-1}\exp(\beta\sum_{\{x,y\}}J_{xy}\sigma_x\sigma_y)\); a field term may be added but is set to zero in the subsequent critical discussion.

\paragraph{Literature and fate.} Classical model definition, long before the talk. Infinite-volume statements require an indicated Gibbs state.

\paragraph{Present status.} As of 12 November 2026: Classical model definition, long before the talk. Infinite-volume statements require an indicated Gibbs state.

\paragraph{Audit.} Literal correctness: correct with missing hypotheses. Publication match: exact. Verification: read the result. Development: none. Confidence: high.

\noindent References: \citep{onsager1944crystal}.\par

\subsubsection{1983-631.2 --- Subcritical exponential decay}

\label{item:1983-631.2}

\noindent\textbf{Kind:} claim. \textbf{Fate:} B. \textbf{Pages:} 1186--1187.\par

\paragraph{Notebook content.} At high temperature the spin correlation is said to decay exponentially.

\paragraph{Precise statement.} For nonnegative translation-invariant finite-range couplings on \(\mathbb Z^d\), the zero-field two-point function satisfies \(\langle\sigma_0\sigma_x\rangle_\beta\le e^{-c(\beta)|x|}\) for \(\beta<\beta_c\), with a harmless metric-dependent constant convention.

\paragraph{Literature and fate.} The general-dimensional sharpness theorem was proved after the talk (Aizenman--Barsky--Fernandez, 1987, as identified by Duminil-Copin--Tassion). The latter paper's Ising Theorem 2.1(3), not its percolation Theorem 1.1, is the verified modern statement. Planar exact-model antecedents do not establish the general version.

\paragraph{Present status.} As of 12 November 2026: The general-dimensional sharpness theorem was proved after the talk (Aizenman--Barsky--Fernandez, 1987, as identified by Duminil-Copin--Tassion). The latter paper's Ising Theorem 2.1(3), not its percolation Theorem 1.1, is the verified modern statement. Planar exact-model antecedents do not establish the general version.

\paragraph{Audit.} Literal correctness: correct with missing hypotheses. Publication match: exact. Verification: read the result. Development: none. Confidence: high.

\noindent References: \citep{duminilcopintassion2016sharpness}.\par

\subsubsection{1983-631.3 --- Ordered phase and spontaneous magnetization}

\label{item:1983-631.3}

\noindent\textbf{Kind:} claim. \textbf{Fate:} E. \textbf{Pages:} 1186--1187.\par

\paragraph{Notebook content.} The low-temperature correlation has a nonzero long-distance limit; the critical inequality is compressed.

\paragraph{Precise statement.} For the square-lattice ferromagnetic nearest-neighbour model in the ordered phase \(\beta>\beta_c\), the long-distance two-point limit is \(m(\beta)^2>0\). At the critical point the spontaneous magnetization vanishes.

\paragraph{Literature and fate.} Yang's 1952 magnetization formula is earlier than the seminar. The ordered-phase statement requires the strict inequality and a specified infinite-volume state.

\paragraph{Present status.} As of 12 November 2026: Yang's 1952 magnetization formula is earlier than the seminar. The ordered-phase statement requires the strict inequality and a specified infinite-volume state.

\paragraph{Audit.} Literal correctness: correct with missing hypotheses. Publication match: exact. Verification: publisher metadata checked. Development: none. Confidence: high.

\noindent References: \citep{yang1952magnetization}.\par

\subsubsection{1983-631.4 --- Critical exponent and singular free energy}

\label{item:1983-631.4}

\noindent\textbf{Kind:} claim. \textbf{Fate:} E. \textbf{Pages:} 1186--1187.\par

\paragraph{Notebook content.} A power law with exponent \(1/4\) and the distinction from the ordered/disordered regimes are stated.

\paragraph{Precise statement.} At criticality in the planar nearest-neighbour model the two-point correlation decays as an amplitude times \(|x-y|^{-1/4}\); the exact free energy has its classical critical singularity.

\paragraph{Literature and fate.} The exact square-lattice results predate the talk (Onsager; Wu--McCoy--Tracy--Barouch). This does not justify the same exponent for an arbitrary interaction.

\paragraph{Present status.} As of 12 November 2026: The exact square-lattice results predate the talk (Onsager; Wu--McCoy--Tracy--Barouch). This does not justify the same exponent for an arbitrary interaction.

\paragraph{Audit.} Literal correctness: correct with missing hypotheses. Publication match: exact. Verification: indirect. Development: none. Confidence: high.

\noindent References: \citep{onsager1944crystal,wumccoytracybarouch1976ising}.\par

\subsubsection{1983-631.5 --- Joint block-magnetization limit}

\label{item:1983-631.5}

\noindent\textbf{Kind:} conjecture. \textbf{Fate:} B. \textbf{Pages:} 1186--1187.\par

\paragraph{Notebook content.} The seminar asks for simultaneous limits of rescaled sums over several regions.

\paragraph{Precise statement.} For bounded regular blocks \(B_j\) in the critical planar nearest-neighbour model, seek a joint limit of \(L^{-15/8}\sum_{x/L\in B_j}\sigma_x\), with consistent lattice amplitude and boundary conditions.

\paragraph{Literature and fate.} Camia--Garban--Newman constructed the magnetization field and compatible block limits: 2012 preprint, 2015 journal, Theorem 1.3, Corollary 1.6 and Appendix A. Generalized-field convergence and block observables are distinguished.

\paragraph{Present status.} As of 12 November 2026: Camia--Garban--Newman constructed the magnetization field and compatible block limits: 2012 preprint, 2015 journal, Theorem 1.3, Corollary 1.6 and Appendix A. Generalized-field convergence and block observables are distinguished.

\paragraph{Audit.} Literal correctness: correct with missing hypotheses. Publication match: exact. Verification: read the result. Development: none. Confidence: high.

\noindent Canonical problem: GP-1983-90.\par

\noindent References: \citep{camia2015magnetization}.\par

\subsubsection{1983-631.6 --- Universality of the critical limiting field}

\label{item:1983-631.6}

\noindent\textbf{Kind:} question. \textbf{Fate:} C. \textbf{Pages:} 1186--1187.\par

\paragraph{Notebook content.} The joint limit is expected not to depend on the initial model; its class and normalization are unstated.

\paragraph{Precise statement.} For zero-field, translation-invariant, square-symmetric, finite-range ferromagnetic pair interactions on \(\mathbb Z^2\) with connected interaction graph, establish the critical bulk magnetization-field limit and its identification with the nearest-neighbour field after a model-dependent multiplicative normalization. The notebook does not specify this class; it is an explicit editorial reconstruction.

\paragraph{Editorial reconstruction.} The interaction class and amplitude convention are editorial choices made explicit in the development, not wording supplied by the notebook.

\paragraph{Literature and fate.} Partially developed (C). Nearest-neighbour field convergence is established. Verified 2025 results treat boundary correlations under weak perturbations. A September 2026 weak-perturbation bulk preprint was located but its primary text could not be opened; it is an unverified lead, not evidence for a new theorem. No verified result found in this run covering the stated arbitrary-strength family. This chapter proves only conditional covariance scaling from uniform two-point asymptotics.

\paragraph{Present status.} As of 12 November 2026: Partially developed (C). Nearest-neighbour field convergence is established. Verified 2025 results treat boundary correlations under weak perturbations. A September 2026 weak-perturbation bulk preprint was located but its primary text could not be opened; it is an unverified lead, not evidence for a new theorem. No verified result found in this run covering the stated arbitrary-strength family. This chapter proves only conditional covariance scaling from uniform two-point asymptotics.

\paragraph{Assessment.} Worth a bounded follow-up: verify exact perturbative bulk hypotheses and tightness/uniform-integrability estimates; a covariance computation alone cannot settle arbitrary-strength universality.

\paragraph{Audit.} Literal correctness: correct with missing hypotheses. Publication match: partial. Verification: read the result. Development: developed in this chapter. Confidence: high.

\noindent Canonical problem: GP-1983-91.\par

\noindent References: \citep{camia2015magnetization,cavagiulianigreenblatt2025boundary,openai2026isingunverified}.\par

\subsubsection{1983-631.7 --- Renormalization as a pushforward transformation}

\label{item:1983-631.7}

\noindent\textbf{Kind:} method. \textbf{Fate:} E. \textbf{Pages:} 1186--1187.\par

\paragraph{Notebook content.} Block variables and repeated changes of scale are used to compare models.

\paragraph{Precise statement.} A block map sends a microscopic measure to its pushforward distribution on coarse variables; iteration and rescaling suggest fixed points and stable manifolds. Its image need not remain in a finite-dimensional family of nearest-neighbour Gibbs measures.

\paragraph{Literature and fate.} The method and epsilon-expansion programme were already established by Wilson and collaborators. The 1974 review is a bibliographic match, but direct full-text access failed in this run.

\paragraph{Present status.} As of 12 November 2026: The method and epsilon-expansion programme were already established by Wilson and collaborators. The 1974 review is a bibliographic match, but direct full-text access failed in this run.

\paragraph{Audit.} Literal correctness: correct with missing hypotheses. Publication match: partial. Verification: indirect. Development: none. Confidence: high.

\noindent References: \citep{wilsonkogut1974renormalization}.\par

\subsubsection{1983-631.8 --- Four-dimensional critical Gaussianity}

\label{item:1983-631.8}

\noindent\textbf{Kind:} conjecture. \textbf{Fate:} B. \textbf{Pages:} 1186--1187.\par

\paragraph{Notebook content.} Dimension four is singled out at the boundary between nonclassical and mean-field behaviour.

\paragraph{Precise statement.} For the specified critical four-dimensional Ising and scalar \(\varphi^4\) models, nontrivial scaling limits are Gaussian after the appropriate normalization; logarithmic effects must not be erased.

\paragraph{Literature and fate.} Aizenman--Duminil-Copin established marginal triviality (2019 preprint; 2021 publication). This is later development, not a theorem available in April 1984.

\paragraph{Present status.} As of 12 November 2026: Aizenman--Duminil-Copin established marginal triviality (2019 preprint; 2021 publication). This is later development, not a theorem available in April 1984.

\paragraph{Audit.} Literal correctness: correct with missing hypotheses. Publication match: partial. Verification: publisher metadata checked. Development: none. Confidence: high.

\noindent References: \citep{aizenmanduminilcopin2021marginal}.\par

\subsubsection{1983-631.9 --- The epsilon-expansion strategy}

\label{item:1983-631.9}

\noindent\textbf{Kind:} method. \textbf{Fate:} E. \textbf{Pages:} 1186--1187.\par

\paragraph{Notebook content.} The discussion invokes a change of dimension near four.

\paragraph{Precise statement.} The expansion about \(d=4\) organizes perturbative critical exponents in \(\epsilon=4-d\); it is not by itself a rigorous convergence theorem for fixed \(\epsilon=1\).

\paragraph{Literature and fate.} Wilson--Fisher (1972) and Wilson--Kogut (1974) precede the talk. No additional precise research question is present.

\paragraph{Present status.} As of 12 November 2026: Wilson--Fisher (1972) and Wilson--Kogut (1974) precede the talk. No additional precise research question is present.

\paragraph{Audit.} Literal correctness: correct with missing hypotheses. Publication match: partial. Verification: publisher metadata checked. Development: none. Confidence: high.

\noindent References: \citep{wilsonfisher1972critical,wilsonkogut1974renormalization}.\par

\subsubsection{1983-631.10 --- MacPherson and perverse sheaves: isolated reference}

\label{item:1983-631.10}

\noindent\textbf{Kind:} reference. \textbf{Fate:} F. \textbf{Pages:} 1186--1187.\par

\paragraph{Notebook content.} A brief aside refers to MacPherson and perverse sheaves.

\paragraph{Precise statement.} The notes do not identify a category, stratification, functor, theorem, or problem to be reconstructed.

\paragraph{Literature and fate.} No unique claim can be identified from the surviving words. No generic perverse-sheaf theorem is substituted.

\paragraph{Present status.} As of 12 November 2026: No unique claim can be identified from the surviving words. No generic perverse-sheaf theorem is substituted.

\paragraph{Audit.} Literal correctness: correct with missing hypotheses. Publication match: none found. Verification: unresolved lead. Development: none. Confidence: low.

\subsection{Record 633: Centres, fixed points, and property (FA)}

\label{talk:633}

\noindent\textbf{Date:} 1984-04-23. \textbf{Speaker:} Jean-Pierre Serre. \textbf{Pages:} 1187.\par

The Serre material occurs above the Margulis note on P001187, despite catalogue numbering.

\subsubsection{1983-633.1 --- Unique centre of a bounded Hilbert set}

\label{item:1983-633.1}

\noindent\textbf{Kind:} theorem. \textbf{Fate:} E. \textbf{Pages:} 1187.\par

\paragraph{Notebook content.} The notes use the centre of a bounded set in a Hilbert space.

\paragraph{Precise statement.} Every nonempty bounded subset \(A\) of a real or complex Hilbert space is contained in a unique closed ball of minimum radius; no compactness of \(A\) is needed.

\paragraph{Literature and fate.} An older centre lemma. Bekka--de la Harpe--Valette, Lemma 2.2.7, gives a precise proof using weak compactness and the parallelogram identity.

\paragraph{Present status.} As of 12 November 2026: An older centre lemma. Bekka--de la Harpe--Valette, Lemma 2.2.7, gives a precise proof using weak compactness and the parallelogram identity.

\paragraph{Audit.} Literal correctness: correct with missing hypotheses. Publication match: exact. Verification: read the result. Development: none. Confidence: high.

\noindent References: \citep{bekkadelaharpevalette2008kazhdan}.\par

\subsubsection{1983-633.2 --- Bounded affine orbit implies a fixed point}

\label{item:1983-633.2}

\noindent\textbf{Kind:} theorem. \textbf{Fate:} E. \textbf{Pages:} 1187.\par

\paragraph{Notebook content.} The centre is applied to an invariant bounded orbit.

\paragraph{Precise statement.} An action by affine isometries of a Hilbert space has a fixed point if it has one bounded orbit. The centre of that orbit is fixed.

\paragraph{Literature and fate.} Classical before the talk; the equivalences are stated in Bekka--de la Harpe--Valette, Proposition 2.2.9.

\paragraph{Present status.} As of 12 November 2026: Classical before the talk; the equivalences are stated in Bekka--de la Harpe--Valette, Proposition 2.2.9.

\paragraph{Audit.} Literal correctness: correct with missing hypotheses. Publication match: exact. Verification: read the result. Development: none. Confidence: high.

\noindent References: \citep{bekkadelaharpevalette2008kazhdan}.\par

\subsubsection{1983-633.3 --- Property (FA)}

\label{item:1983-633.3}

\noindent\textbf{Kind:} definition. \textbf{Fate:} E. \textbf{Pages:} 1187.\par

\paragraph{Notebook content.} A group has a fixed point for every action on a tree.

\paragraph{Precise statement.} Property (FA) requires that every simplicial action without edge inversions on a tree fixes a vertex. Without this convention an edge flip gives a misleading counterexample.

\paragraph{Literature and fate.} Serre's 1977 tree theory is the prior source. Real-tree fixed-point properties and Hilbert-space fixed-point properties are different notions.

\paragraph{Present status.} As of 12 November 2026: Serre's 1977 tree theory is the prior source. Real-tree fixed-point properties and Hilbert-space fixed-point properties are different notions.

\paragraph{Audit.} Literal correctness: correct with missing hypotheses. Publication match: exact. Verification: read the result. Development: none. Confidence: high.

\noindent References: \citep{serre1977arbres,bekkadelaharpevalette2008kazhdan}.\par

\subsubsection{1983-633.4 --- Finite groups have property (FA)}

\label{item:1983-633.4}

\noindent\textbf{Kind:} example. \textbf{Fate:} E. \textbf{Pages:} 1187.\par

\paragraph{Notebook content.} Finite groups are examples.

\paragraph{Precise statement.} A finite group acting without inversions on a tree has a bounded orbit and fixes a vertex; with inversions it still fixes the midpoint of a preserved edge geometrically.

\paragraph{Literature and fate.} Known before the talk; follows from the centre of the finite invariant subtree.

\paragraph{Present status.} As of 12 November 2026: Known before the talk; follows from the centre of the finite invariant subtree.

\paragraph{Audit.} Literal correctness: correct with missing hypotheses. Publication match: exact. Verification: read the result. Development: none. Confidence: high.

\noindent References: \citep{serre1977arbres,bekkadelaharpevalette2008kazhdan}.\par

\subsubsection{1983-633.5 --- The group SL(3,Z) has property (FA)}

\label{item:1983-633.5}

\noindent\textbf{Kind:} example. \textbf{Fate:} E. \textbf{Pages:} 1187.\par

\paragraph{Notebook content.} \(\mathrm{SL}(3,\mathbb Z)\) is named as an infinite example.

\paragraph{Precise statement.} \(\mathrm{SL}(3,\mathbb Z)\), viewed as a discrete group, has property (FA); its higher-rank fixed-point rigidity is not shared by every lattice in every rank-one group.

\paragraph{Literature and fate.} This predates the seminar and is covered by the established higher-rank/property-(T) framework and the tree fixed-point implication.

\paragraph{Present status.} As of 12 November 2026: This predates the seminar and is covered by the established higher-rank/property-(T) framework and the tree fixed-point implication.

\paragraph{Audit.} Literal correctness: correct with missing hypotheses. Publication match: exact. Verification: indirect. Development: none. Confidence: high.

\noindent References: \citep{serre1977arbres,bekkadelaharpevalette2008kazhdan}.\par

\subsection{Record 632: Higher-rank lattices and actions on trees}

\label{talk:632}

\noindent\textbf{Date:} 1984-04-23. \textbf{Speaker:} G. A. Margulis. \textbf{Pages:} 1187.\par

A short generalization follows Serre's examples in the physical notebook order.

\subsubsection{1983-632.1 --- Property (FA) for irreducible higher-real-rank lattices}

\label{item:1983-632.1}

\noindent\textbf{Kind:} theorem. \textbf{Fate:} E. \textbf{Pages:} 1187.\par

\paragraph{Notebook content.} The note extends the preceding tree fixed-point statement to irreducible lattices.

\paragraph{Precise statement.} An irreducible lattice in a connected semisimple real Lie group with finite centre, no compact factors, and total real rank at least two has property (FA). Irreducibility is essential when rank-one factors are present.

\paragraph{Literature and fate.} Margulis's 1981 paper is the pre-seminar attribution. Cornulier--Kar's introduction and reference 15 verify the precise historical identification; the original Margulis proof was not retrieved.

\paragraph{Present status.} As of 12 November 2026: Margulis's 1981 paper is the pre-seminar attribution. Cornulier--Kar's introduction and reference 15 verify the precise historical identification; the original Margulis proof was not retrieved.

\paragraph{Audit.} Literal correctness: correct with missing hypotheses. Publication match: exact. Verification: indirect. Development: none. Confidence: high.

\noindent References: \citep{margulis1981amalgams,cornulierkar2011fa}.\par

\subsection{Record 634: Academic-year heading 1984--85}

\label{talk:634}

\noindent\textbf{Date:} 1984-85. \textbf{Speaker:} —. \textbf{Pages:} 1188.\par

Documentary heading only; it is not a mathematical talk.

\subsection{Record 635: Lipschitz structures, signature operators, and slice areas; Grothendieck programme}

\label{talk:635}

\noindent\textbf{Date:} 1984-09-03. \textbf{Speaker:} I. M. Gelfand. \textbf{Pages:} 1188--1190.\par

Separates the earlier Lipschitz index theorem from later local characteristic formulas, checks two operator displays, and gives a precise form of the slice-area question.

\subsubsection{1984-635.1 --- Bi-Lipschitz equivalence and the cube--sphere example}

\label{item:1984-635.1}

\noindent\textbf{Kind:} definition. \textbf{Fate:} E. \textbf{Pages:} 1188--1190.\par

\paragraph{Notebook content.} The discussion uses Lipschitz equivalence and compares piecewise-smooth and smooth shapes.

\paragraph{Precise statement.} A bi-Lipschitz map \(f\) satisfies \(c\,d(x,y)\le d(fx,fy)\le C\,d(x,y)\) for constants \(0<c\le C<\infty\). Compact smooth manifolds have mutually bi-Lipschitz smooth Riemannian metrics; this is a compactness statement and does not hold for arbitrary complete metrics on a noncompact manifold.

\paragraph{Literature and fate.} Standard geometry before the talk. The rough boundary of a cube does not prevent a bi-Lipschitz equivalence of its boundary with a sphere.

\paragraph{Present status.} As of 12 November 2026: Standard geometry before the talk. The rough boundary of a cube does not prevent a bi-Lipschitz equivalence of its boundary with a sphere.

\paragraph{Audit.} Literal correctness: correct with missing hypotheses. Publication match: partial. Verification: read the result. Development: none. Confidence: high.

\noindent References: \citep{sullivan1979hyperbolic}.\par

\subsubsection{1984-635.2 --- Existence and uniqueness of Lipschitz structures away from dimension four}

\label{item:1984-635.2}

\noindent\textbf{Kind:} theorem. \textbf{Fate:} E. \textbf{Pages:} 1188--1190.\par

\paragraph{Notebook content.} Sullivan's Lipschitz-structure theorem is reported in an abbreviated form.

\paragraph{Precise statement.} A topological \(n\)-manifold, \(n\ne4\), admits a locally Euclidean Lipschitz structure; any two such structures are related by a suitable homeomorphism arbitrarily close to the identity. Uniqueness is not literal equality of atlases.

\paragraph{Literature and fate.} Sullivan's January 1978 preprint, Corollary 3, and 1979 publication predate the 1984 report. The omitted dimensional exception is substantive.

\paragraph{Present status.} As of 12 November 2026: Sullivan's January 1978 preprint, Corollary 3, and 1979 publication predate the 1984 report. The omitted dimensional exception is substantive.

\paragraph{Audit.} Literal correctness: correct with missing hypotheses. Publication match: exact. Verification: read the result. Development: annotated. Confidence: high.

\noindent References: \citep{sullivan1979hyperbolic}.\par

\subsubsection{1984-635.3 --- Characteristic information from measurable metric data}

\label{item:1984-635.3}

\noindent\textbf{Kind:} question. \textbf{Fate:} B. \textbf{Pages:} 1188--1190.\par

\paragraph{Notebook content.} The question asks for characteristic-number formulas without differentiating a measurable metric.

\paragraph{Precise statement.} Construct characteristic-number information on an oriented compact Lipschitz or quasiconformal manifold from sufficiently bounded measurable metric/conformal data, without assuming a classical curvature tensor. Distinguish global index formulas from local characteristic-cycle formulas.

\paragraph{Literature and fate.} Teleman's signature-index theorem was published in 1983 before the talk. The later local formula programme is realized in Connes--Sullivan--Teleman, Theorems 1.1 and 3.2 (1993 preprint, 1994 publication). Global index and local cycle versions are different; no classical curvature tensor is claimed for arbitrary measurable data.

\paragraph{Present status.} As of 12 November 2026: Teleman's signature-index theorem was published in 1983 before the talk. The later local formula programme is realized in Connes--Sullivan--Teleman, Theorems 1.1 and 3.2 (1993 preprint, 1994 publication). Global index and local cycle versions are different; no classical curvature tensor is claimed for arbitrary measurable data.

\paragraph{Audit.} Literal correctness: correct with missing hypotheses. Publication match: partial. Verification: read the result. Development: none. Confidence: high.

\noindent Canonical problem: GP-1984-01.\par

\noindent References: \citep{teleman1983signature,sullivanteleman1983novikov,connessullivanteleman1994quasiconformal}.\par

\subsubsection{1984-635.4 --- Analytic proof of rational Pontryagin invariance}

\label{item:1984-635.4}

\noindent\textbf{Kind:} theorem. \textbf{Fate:} E. \textbf{Pages:} 1188--1190.\par

\paragraph{Notebook content.} The Lipschitz operator construction is connected to Novikov's theorem.

\paragraph{Precise statement.} Rational Pontryagin classes of smooth manifolds are invariants under homeomorphism. The Lipschitz signature-index machinery gives an analytic route to the theorem after the necessary localization and index arguments.

\paragraph{Literature and fate.} Sullivan--Teleman's three-page proof was received in 1982 and published in December 1983, before this report. It is an alternative proof of an older theorem, not a new 1984 invariance claim.

\paragraph{Present status.} As of 12 November 2026: Sullivan--Teleman's three-page proof was received in 1982 and published in December 1983, before this report. It is an alternative proof of an older theorem, not a new 1984 invariance claim.

\paragraph{Audit.} Literal correctness: correct with missing hypotheses. Publication match: exact. Verification: read the result. Development: none. Confidence: high.

\noindent References: \citep{sullivanteleman1983novikov}.\par

\subsubsection{1984-635.5 --- The displayed two-by-two Dirac-type square}

\label{item:1984-635.5}

\noindent\textbf{Kind:} formula. \textbf{Fate:} E. \textbf{Pages:} 1188--1190.\par

\paragraph{Notebook content.} A \(2\times2\) first-order matrix is asserted to square to a scalar operator, but its two diagonal signs agree.

\paragraph{Precise statement.} For commuting symbols \(a,b\), \(M=\begin{pmatrix}a&ib\\-ib&a\end{pmatrix}\) has \(M^2=\begin{pmatrix}a^2+b^2&2iab\\-2iab&a^2+b^2\end{pmatrix}\). Replacing the lower-right \(a\) by \(-a\) gives \(M^2=(a^2+b^2)I\).

\paragraph{Literature and fate.} The notebook display is false as written. Direct multiplication identifies a minimal sign repair; the chapter supplies the explicit check rather than silently altering the diplomatic text.

\paragraph{Present status.} As of 12 November 2026: The notebook display is false as written. Direct multiplication identifies a minimal sign repair; the chapter supplies the explicit check rather than silently altering the diplomatic text.

\paragraph{Audit.} Literal correctness: false as written. Publication match: thematic only. Verification: read the result. Development: annotated. Confidence: high.

\noindent References: \citep{teleman1983signature}.\par

\subsubsection{1984-635.6 --- A grad--div--curl square root of the Laplacian}

\label{item:1984-635.6}

\noindent\textbf{Kind:} formula. \textbf{Fate:} E. \textbf{Pages:} 1188--1190.\par

\paragraph{Notebook content.} The vector-calculus display mixes dimensions and component counts.

\paragraph{Precise statement.} On a scalar plus a three-vector over \(\mathbb R^3\), define \(Q(f,v)=(-\operatorname{div}v,\nabla f+\operatorname{curl}v)\). Then \(Q^2=-\Delta I_4\). There are four fibre components and three base coordinates.

\paragraph{Literature and fate.} The repaired identity follows from the standard vector-calculus identities. The four-component symbol is checked independently in the program.

\paragraph{Present status.} As of 12 November 2026: The repaired identity follows from the standard vector-calculus identities. The four-component symbol is checked independently in the program.

\paragraph{Audit.} Literal correctness: false as written. Publication match: thematic only. Verification: read the result. Development: annotated. Confidence: high.

\noindent References: \citep{teleman1983signature}.\par

\subsubsection{1984-635.7 --- Hodge--de Rham operator and signature index}

\label{item:1984-635.7}

\noindent\textbf{Kind:} construction. \textbf{Fate:} E. \textbf{Pages:} 1188--1190.\par

\paragraph{Notebook content.} The notes replace a coordinate display by an operator on forms.

\paragraph{Precise statement.} For a compact oriented smooth Riemannian manifold, \(D=d+d^*\) is elliptic and \(D^2=\Delta\); with the appropriate chirality in dimension divisible by four, its signature component has index equal to the signature. Lipschitz versions require domains and closed-operator estimates.

\paragraph{Literature and fate.} The smooth construction predates the talk. Teleman's 1983 theorem supplies the Lipschitz analytic framework; a formal coordinate square alone is not its proof.

\paragraph{Present status.} As of 12 November 2026: The smooth construction predates the talk. Teleman's 1983 theorem supplies the Lipschitz analytic framework; a formal coordinate square alone is not its proof.

\paragraph{Audit.} Literal correctness: correct with missing hypotheses. Publication match: exact. Verification: read the result. Development: none. Confidence: high.

\noindent References: \citep{teleman1983signature}.\par

\subsubsection{1984-635.8 --- Directional average of slice-derivative energy}

\label{item:1984-635.8}

\noindent\textbf{Kind:} formula. \textbf{Fate:} E. \textbf{Pages:} 1188--1190.\par

\paragraph{Notebook content.} The slice-area integral is compared to \(\pi V\), without fully explicit averaging conventions.

\paragraph{Precise statement.} For measurable \(E\subset\mathbb R^3\) of finite volume, \(S_\theta(t)=R1_E(\theta,t)\) satisfies \((4\pi)^{-1}\int_{S^2}\int_{\mathbb R}|S_\theta'(t)|^2dt\,d\theta=2\pi|E|\), with weak derivatives and normalized spherical average.

\paragraph{Literature and fate.} Classical Radon--Plancherel theory fixes the constant. The chapter gives a complete proof and an independent ball calculation; (pi V) is false with this normalized convention.

\paragraph{Present status.} As of 12 November 2026: Classical Radon--Plancherel theory fixes the constant. The chapter gives a complete proof and an independent ball calculation; (pi V) is false with this normalized convention.

\paragraph{Audit.} Literal correctness: false as written. Publication match: exact. Verification: read the result. Development: annotated. Confidence: high.

\noindent Canonical problem: GP-1984-02.\par

\noindent References: \citep{gelfandgraevvilenkin1966integral}.\par

\subsubsection{1984-635.9 --- Differentiability of slices for a measurable finite-volume body}

\label{item:1984-635.9}

\noindent\textbf{Kind:} question. \textbf{Fate:} E. \textbf{Pages:} 1188--1190.\par

\paragraph{Notebook content.} The notes ask about differentiability for general measurable sets.

\paragraph{Precise statement.} If \(E\) is measurable and \(|E|<\infty\), then \(S_\theta\in H^1(\mathbb R)\) for almost every direction; the representative is locally absolutely continuous and its weak derivative agrees with its classical derivative almost everywhere. No finite-perimeter assumption is needed.

\paragraph{Literature and fate.} This follows from the finite-measure Radon--Plancherel argument proved in the chapter. The theorem does not give all directions or every raw plane section; null-set changes can alter exceptional sections.

\paragraph{Present status.} As of 12 November 2026: This follows from the finite-measure Radon--Plancherel argument proved in the chapter. The theorem does not give all directions or every raw plane section; null-set changes can alter exceptional sections.

\paragraph{Audit.} Literal correctness: correct with missing hypotheses. Publication match: exact. Verification: read the result. Development: annotated. Confidence: high.

\noindent Canonical problem: GP-1984-02.\par

\noindent References: \citep{gelfandgraevvilenkin1966integral}.\par

\subsubsection{1984-635.10 --- Grothendieck's programme: headings and manuscript extent}

\label{item:1984-635.10}

\noindent\textbf{Kind:} reference. \textbf{Fate:} F. \textbf{Pages:} 1188--1190.\par

\paragraph{Notebook content.} A contents-style description invokes Grothendieck, moduli, anabelian ideas, dessins and a long manuscript.

\paragraph{Precise statement.} The headings can be matched to the 1984 Esquisse d'un programme, but no unique individual theorem or testable problem is stated here. French champs refers to stacks in this context, not scalar fields. The roughly 1300-page manuscript allusion must not be identified with the short Esquisse.

\paragraph{Literature and fate.} The original French Esquisse was opened and compared. The programme is identifiable; the mathematical assertion is not. More specific anabelian or dessin conjectures are not inserted into the notebook.

\paragraph{Present status.} As of 12 November 2026: The original French Esquisse was opened and compared. The programme is identifiable; the mathematical assertion is not. More specific anabelian or dessin conjectures are not inserted into the notebook.

\paragraph{Audit.} Literal correctness: not determinable. Publication match: partial. Verification: read the result. Development: none. Confidence: medium.

\noindent References: \citep{grothendieck1984esquisse}.\par

\subsection{Record 636: Catalan numbers, special linear systems, and finite SU(2) groups}

\label{talk:636}

\noindent\textbf{Date:} 1984-09-03. \textbf{Speaker:} V. I. Arnold. \textbf{Pages:} 1190--1191.\par

The enumerative count requires a general curve; the finite-group discussion requires binary lifts and distinguishes the quotient surface, Milnor fibre, and Cartan form. An unattributed groupoid note follows.

\subsubsection{1984-636.1 --- Catalan number}

\label{item:1984-636.1}

\noindent\textbf{Kind:} definition. \textbf{Fate:} E. \textbf{Pages:} 1190--1191.\par

\paragraph{Notebook content.} The integer \((2k)!/(k!(k+1)!)\) occurs as the enumerative answer.

\paragraph{Precise statement.} For \(k\ge0\), \(C_k=\frac1{k+1}\binom{2k}{k}\) is the \(k\)-th Catalan number.

\paragraph{Literature and fate.} Classical before the talk; its occurrence does not by itself identify a new enumerative theorem.

\paragraph{Present status.} As of 12 November 2026: Classical before the talk; its occurrence does not by itself identify a new enumerative theorem.

\paragraph{Audit.} Literal correctness: correct with missing hypotheses. Publication match: exact. Verification: read the result. Development: none. Confidence: high.

\noindent References: \citep{griffithsharris1980special}.\par

\subsubsection{1984-636.2 --- Pencils of minimal generic degree on a genus-two-k curve}

\label{item:1984-636.2}

\noindent\textbf{Kind:} theorem. \textbf{Fate:} E. \textbf{Pages:} 1190--1191.\par

\paragraph{Notebook content.} A curve of genus \(2k\) is said to have \(C_k\) functions of degree \(k+1\).

\paragraph{Precise statement.} A general smooth complex projective curve of genus \(2k\) has \(C_k\) distinct \(g^1_{k+1}\)'s, counted as linear series; their maps to \(\mathbb P^1\) are taken modulo target automorphisms. The count is not valid for every curve or for unquotiented functions.

\paragraph{Literature and fate.} Griffiths--Harris (1980), Main Theorem II(b), specializes to this Brill--Noether number-zero count. Generality and the quotient by target automorphisms are indispensable.

\paragraph{Present status.} As of 12 November 2026: Griffiths--Harris (1980), Main Theorem II(b), specializes to this Brill--Noether number-zero count. Generality and the quotient by target automorphisms are indispensable.

\paragraph{Audit.} Literal correctness: correct with missing hypotheses. Publication match: exact. Verification: read the result. Development: annotated. Confidence: high.

\noindent References: \citep{griffithsharris1980special}.\par

\subsubsection{1984-636.3 --- Branch-parameter dimension count}

\label{item:1984-636.3}

\noindent\textbf{Kind:} argument. \textbf{Fate:} E. \textbf{Pages:} 1190--1191.\par

\paragraph{Notebook content.} A comparison with \(3g-3\) supports the proposed degree and finiteness.

\paragraph{Precise statement.} For a simply branched degree-\(d\) cover of \(\mathbb P^1\) by a genus-\(g\) curve, Riemann--Hurwitz gives \(2g+2d-2\) branch points. Quotienting their positions by \(\mathrm{PGL}_2\) gives dimension \(2g+2d-5\); for \(g=2k,d=k+1\), this equals \(3g-3\).

\paragraph{Literature and fate.} This is a necessary dimension comparison, not a proof of existence, dominance, transversality, or the Catalan degree.

\paragraph{Present status.} As of 12 November 2026: This is a necessary dimension comparison, not a proof of existence, dominance, transversality, or the Catalan degree.

\paragraph{Audit.} Literal correctness: correct with missing hypotheses. Publication match: partial. Verification: read the result. Development: none. Confidence: high.

\noindent References: \citep{griffithsharris1980special}.\par

\subsubsection{1984-636.4 --- Finite subgroups of SU(2) and binary polyhedral groups}

\label{item:1984-636.4}

\noindent\textbf{Kind:} classification. \textbf{Fate:} E. \textbf{Pages:} 1190--1191.\par

\paragraph{Notebook content.} Finite rotation groups and finite subgroups of \(\mathrm{SU}(2)\) are discussed together.

\paragraph{Precise statement.} Under \(\mathrm{SU}(2)\to\mathrm{SO}(3)\), the noncyclic finite subgroups are binary dihedral and binary tetrahedral, octahedral and icosahedral groups. The binary icosahedral group has order120 and is not the rotation group of order60.

\paragraph{Literature and fate.} The classification is old. The source's unqualified identification of the two group lists requires the binary lift correction.

\paragraph{Present status.} As of 12 November 2026: The classification is old. The source's unqualified identification of the two group lists requires the binary lift correction.

\paragraph{Audit.} Literal correctness: false as written. Publication match: exact. Verification: read the result. Development: annotated. Confidence: high.

\noindent References: \citep{mckay1980graphs}.\par

\subsubsection{1984-636.5 --- McKay tensor graph and the affine Cartan matrix}

\label{item:1984-636.5}

\noindent\textbf{Kind:} theorem. \textbf{Fate:} E. \textbf{Pages:} 1190--1191.\par

\paragraph{Notebook content.} Tensoring irreducible representations gives a graph and a Cartan matrix.

\paragraph{Precise statement.} For a finite \(\Gamma\subset\mathrm{SU}(2)\), fix its defining two-dimensional representation \(V\). With \(V\otimes V_i=\sum_jm_{ij}V_j\), the matrix \(2I-(m_{ij})\) is the corresponding affine ADE Cartan matrix, with standard low-rank conventions.

\paragraph{Literature and fate.} McKay's 1980 article, especially Propositions 3--6, predates the report. The fixed tensor factor cannot be replaced by arbitrary irreducible representations.

\paragraph{Present status.} As of 12 November 2026: McKay's 1980 article, especially Propositions 3--6, predates the report. The fixed tensor factor cannot be replaced by arbitrary irreducible representations.

\paragraph{Audit.} Literal correctness: correct with missing hypotheses. Publication match: exact. Verification: read the result. Development: none. Confidence: high.

\noindent References: \citep{mckay1980graphs}.\par

\subsubsection{1984-636.6 --- Monster and modular-function allusion}

\label{item:1984-636.6}

\noindent\textbf{Kind:} reference. \textbf{Fate:} F. \textbf{Pages:} 1190--1191.\par

\paragraph{Notebook content.} A numerical relationship between sporadic groups and modular quantities is mentioned.

\paragraph{Precise statement.} No particular graded module, coefficient identity, Hauptmodul assertion, or conjecture is stated unambiguously.

\paragraph{Literature and fate.} Borcherds's 1992 theorem is relevant later context, but only a thematic match. The notebook is not assigned the full moonshine theorem by inference.

\paragraph{Present status.} As of 12 November 2026: Borcherds's 1992 theorem is relevant later context, but only a thematic match. The notebook is not assigned the full moonshine theorem by inference.

\paragraph{Audit.} Literal correctness: correct with missing hypotheses. Publication match: thematic only. Verification: read the result. Development: none. Confidence: low.

\noindent References: \citep{borcherds1992moonshine}.\par

\subsubsection{1984-636.7 --- Invariant ring and quotient surface}

\label{item:1984-636.7}

\noindent\textbf{Kind:} theorem. \textbf{Fate:} E. \textbf{Pages:} 1190--1191.\par

\paragraph{Notebook content.} Three generators and a relation are associated to the finite subgroup.

\paragraph{Precise statement.} For a finite \(\Gamma\subset\mathrm{SL}_2(\mathbb C)\), the quotient \(\mathbb C^2/\Gamma\) is a two-dimensional surface; in the ADE description its invariant ring is generated by three quantities with one relation, embedding the surface as a hypersurface in \(\mathbb C^3\).

\paragraph{Literature and fate.} The ambient affine three-space does not make the quotient three-dimensional. The classical invariant-surface statement is matched by McKay's article.

\paragraph{Present status.} As of 12 November 2026: The ambient affine three-space does not make the quotient three-dimensional. The classical invariant-surface statement is matched by McKay's article.

\paragraph{Audit.} Literal correctness: false as written. Publication match: exact. Verification: read the result. Development: annotated. Confidence: high.

\noindent References: \citep{mckay1980graphs}.\par

\subsubsection{1984-636.8 --- Milnor lattice and the ADE Cartan form}

\label{item:1984-636.8}

\noindent\textbf{Kind:} theorem. \textbf{Fate:} E. \textbf{Pages:} 1190--1191.\par

\paragraph{Notebook content.} Vanishing cycles are compared with a Coxeter matrix.

\paragraph{Precise statement.} For the simple surface singularity in its local Milnor fibre, the vanishing-cycle intersection lattice is the negative ADE Cartan lattice, up to the orientation/sign convention. The Coxeter transformation is a different automorphism.

\paragraph{Literature and fate.} Arnold's simple-singularity classification and McKay's singularity discussion predate the talk. The exact local Milnor-fibre hypothesis and the Cartan/Coxeter distinction must be restored.

\paragraph{Present status.} As of 12 November 2026: Arnold's simple-singularity classification and McKay's singularity discussion predate the talk. The exact local Milnor-fibre hypothesis and the Cartan/Coxeter distinction must be restored.

\paragraph{Audit.} Literal correctness: false as written. Publication match: exact. Verification: indirect. Development: annotated. Confidence: high.

\noindent References: \citep{arnold1972normal,mckay1980graphs}.\par

\subsubsection{1984-636.9 --- Fundamental groupoid on a selected set of objects}

\label{item:1984-636.9}

\noindent\textbf{Kind:} definition. \textbf{Fate:} E. \textbf{Pages:} 1190--1191.\par

\paragraph{Notebook content.} An unattributed note describes paths associated with a set \(A\).

\paragraph{Precise statement.} \(\Pi_1(X,A)\) has object set \(A\subset X\) and homotopy classes relative endpoints of paths whose two endpoints both lie in \(A\); composition is endpoint-matching concatenation.

\paragraph{Literature and fate.} Brown's 1967 groupoid form of van Kampen is earlier context. The lower note has no independent speaker heading; it remains under this record only as a source-boundary convention.

\paragraph{Present status.} As of 12 November 2026: Brown's 1967 groupoid form of van Kampen is earlier context. The lower note has no independent speaker heading; it remains under this record only as a source-boundary convention.

\paragraph{Audit.} Literal correctness: correct with missing hypotheses. Publication match: exact. Verification: read the result. Development: annotated. Confidence: high.

\noindent References: \citep{brown1967groupoids}.\par

\subsection{Record 637: Algebraic complete integrability, Laurent series, and Poisson--Lie geometry}

\label{talk:637}

\noindent\textbf{Date:} 1984-09-10. \textbf{Speaker:} M. A. Semenov-Tian-Shansky. \textbf{Pages:} 1191--1194.\par

The local pole-count criterion is separated from compact abelian completion; the SO(4) formula remains partly unreadable. The later portion anticipates the 1985 dressing-action paper.

\subsubsection{1984-637.1 --- Polynomial Poisson systems and the count of integrals}

\label{item:1984-637.1}

\noindent\textbf{Kind:} setup. \textbf{Fate:} E. \textbf{Pages:} 1191--1194.\par

\paragraph{Notebook content.} A polynomial Poisson tensor, Casimirs, and commuting Hamiltonians define the integrability problem.

\paragraph{Precise statement.} On an \(n\)-dimensional Poisson variety of generic rank \(2m\), take \(k=n-2m\) independent local Casimirs and \(m\) further independent pairwise commuting Hamiltonians on a regular leaf. Independence is required on a dense regular set.

\paragraph{Literature and fate.} This is the classical Poisson formulation of Liouville integrability; singular strata require separate treatment.

\paragraph{Present status.} As of 12 November 2026: This is the classical Poisson formulation of Liouville integrability; singular strata require separate treatment.

\paragraph{Audit.} Literal correctness: correct with missing hypotheses. Publication match: exact. Verification: read the result. Development: none. Confidence: high.

\noindent References: \citep{adler1982kowalewski,semenov1983classical}.\par

\subsubsection{1984-637.2 --- Liouville tori and their dimension}

\label{item:1984-637.2}

\noindent\textbf{Kind:} theorem. \textbf{Fate:} E. \textbf{Pages:} 1191--1194.\par

\paragraph{Notebook content.} The number of integrals and the dimension of invariant tori are compressed inconsistently.

\paragraph{Precise statement.} On a real regular symplectic leaf of dimension \(2m\), a connected compact common level of \(m\) independent commuting Hamiltonians is an \(m\)-torus, provided the associated vector fields are complete on that level. Including Casimirs gives \(k+m\), not \(k+n\), independent level conditions.

\paragraph{Literature and fate.} The standard Liouville theorem is earlier than the talk. The notebook's dimension count is not a separate theorem.

\paragraph{Present status.} As of 12 November 2026: The standard Liouville theorem is earlier than the talk. The notebook's dimension count is not a separate theorem.

\paragraph{Audit.} Literal correctness: false as written. Publication match: exact. Verification: indirect. Development: annotated. Confidence: high.

\noindent References: \citep{adler1982kowalewski}.\par

\subsubsection{1984-637.3 --- Algebraic complete integrability}

\label{item:1984-637.3}

\noindent\textbf{Kind:} definition. \textbf{Fate:} E. \textbf{Pages:} 1191--1194.\par

\paragraph{Notebook content.} A complex invariant variety is compactified to a torus.

\paragraph{Precise statement.} In the compact-abelian meaning used here, a generic complex invariant manifold of dimension \(m\) embeds as a Zariski-open subset of an \(m\)-dimensional abelian variety, and the commuting flows extend as translations. The phase dimension \(n\) is not the abelian dimension.

\paragraph{Literature and fate.} Adler--van Moerbeke's framework predates the report. This is stronger than merely having (m) first integrals or a Lax representation.

\paragraph{Present status.} As of 12 November 2026: Adler--van Moerbeke's framework predates the report. This is stronger than merely having (m) first integrals or a Lax representation.

\paragraph{Audit.} Literal correctness: false as written. Publication match: exact. Verification: read the result. Development: annotated. Confidence: high.

\noindent References: \citep{adler1982kowalewski}.\par

\subsubsection{1984-637.4 --- Necessary Laurent-family count}

\label{item:1984-637.4}

\noindent\textbf{Kind:} criterion. \textbf{Fate:} E. \textbf{Pages:} 1191--1194.\par

\paragraph{Notebook content.} A pole expansion with \(n-1\) parameters is proposed as an integrability test.

\paragraph{Precise statement.} If an algebraically completely integrable system has a divisor reached transversely by a flow and meromorphic coordinate functions with the prescribed local poles, fixing the pole time leaves an \((n-1)\)-parameter Laurent family in phase space. Pole time restores the final integration constant.

\paragraph{Literature and fate.} Adler--van Moerbeke (1982), pp.83--84, explains the pole-family programme and separates necessity from sufficiency. The transversality and meromorphic-extension assumptions are essential.

\paragraph{Present status.} As of 12 November 2026: Adler--van Moerbeke (1982), pp.83--84, explains the pole-family programme and separates necessity from sufficiency. The transversality and meromorphic-extension assumptions are essential.

\paragraph{Audit.} Literal correctness: correct with missing hypotheses. Publication match: exact. Verification: read the result. Development: none. Confidence: high.

\noindent Canonical problem: GP-1984-03.\par

\noindent References: \citep{adler1982kowalewski}.\par

\subsubsection{1984-637.5 --- Reconstructing a compactification from pole divisors}

\label{item:1984-637.5}

\noindent\textbf{Kind:} method. \textbf{Fate:} E. \textbf{Pages:} 1191--1194.\par

\paragraph{Notebook content.} The Laurent coefficients are used to add a boundary divisor to the invariant affine variety.

\paragraph{Precise statement.} One may compute families of Laurent solutions, identify their parameter varieties as candidate divisors, glue boundary charts, and then prove smooth projectivity and extension/linearization of the flows. These final global properties are additional tasks, not consequences of parameter counting alone.

\paragraph{Literature and fate.} The cited 1982 construction succeeds for specified integrable systems. It does not supply a universal theorem for every polynomial ODE satisfying a local pole test.

\paragraph{Present status.} As of 12 November 2026: The cited 1982 construction succeeds for specified integrable systems. It does not supply a universal theorem for every polynomial ODE satisfying a local pole test.

\paragraph{Audit.} Literal correctness: correct with missing hypotheses. Publication match: partial. Verification: read the result. Development: none. Confidence: high.

\noindent References: \citep{adler1982kowalewski,adler1987quadrics}.\par

\subsubsection{1984-637.6 --- Failure of the bare Laurent-to-abelian converse}

\label{item:1984-637.6}

\noindent\textbf{Kind:} claim. \textbf{Fate:} E. \textbf{Pages:} 1191--1194.\par

\paragraph{Notebook content.} The local count is stated as if it guaranteed algebraic complete integrability.

\paragraph{Precise statement.} The count-only implication is false: \(H=pq^2\) on \(\mathbb C^2\), with \(\{q,p\}=1\), has \(q=-1/(t-t_0),p=h(t-t_0)^2\). At fixed \(t_0\) this has \(n-1=1\) parameter, while for \(h\ne0\) its energy curve is \(\mathbb C^*\) with projective genus zero.

\paragraph{Literature and fate.} The published earlier programme does not prove the literal converse. The chapter gives a complete counterexample, a canonical coordinate change making both coordinates polar, and an independent genus check. Class E records the prior necessary-test framework; the editorial disproof is not assigned historical priority.

\paragraph{Present status.} As of 12 November 2026: The published earlier programme does not prove the literal converse. The chapter gives a complete counterexample, a canonical coordinate change making both coordinates polar, and an independent genus check. Class E records the prior necessary-test framework; the editorial disproof is not assigned historical priority.

\paragraph{Audit.} Literal correctness: false as written. Publication match: partial. Verification: read the result. Development: annotated. Confidence: high.

\noindent Canonical problem: GP-1984-03.\par

\noindent References: \citep{adler1982kowalewski,adler1987quadrics}.\par

\subsubsection{1984-637.7 --- Abelian varieties, Jacobians, and finite quotients}

\label{item:1984-637.7}

\noindent\textbf{Kind:} remark. \textbf{Fate:} F. \textbf{Pages:} 1191--1194.\par

\paragraph{Notebook content.} The compactified invariant variety is compared with a Jacobian or a quotient of one.

\paragraph{Precise statement.} The source does not specify which system, curve, polarization, isogeny, or finite group is intended. Every abelian variety is related to a Jacobian by a quotient in a broad sense, but that does not identify the concrete invariant variety here.

\paragraph{Literature and fate.} No unique theorem can be recovered from this isolated comparison. The 1987 SO(4) paper offers relevant models, not a decisive match.

\paragraph{Present status.} As of 12 November 2026: No unique theorem can be recovered from this isolated comparison. The 1987 SO(4) paper offers relevant models, not a decisive match.

\paragraph{Audit.} Literal correctness: correct with missing hypotheses. Publication match: thematic only. Verification: read the result. Development: none. Confidence: low.

\noindent References: \citep{adler1987quadrics}.\par

\subsubsection{1984-637.8 --- The Painleve-property classification question}

\label{item:1984-637.8}

\noindent\textbf{Kind:} question. \textbf{Fate:} F. \textbf{Pages:} 1191--1194.\par

\paragraph{Notebook content.} Gelfand asks about differential equations whose solutions have no problematic moving singularities.

\paragraph{Precise statement.} The Painleve property concerns the absence of movable critical singularities of solutions in a specified complex differential-equation class; fixed singularities of the equation may remain. No order, degree, equivalence group, or analytic category is fixed here.

\paragraph{Literature and fate.} The literal prohibition of all singularities needs correction, while the proposed classification remains underdetermined. Earlier indexed discussions 1977-439.10--12 supply context, not a single reusable problem with identical hypotheses.

\paragraph{Present status.} As of 12 November 2026: The literal prohibition of all singularities needs correction, while the proposed classification remains underdetermined. Earlier indexed discussions 1977-439.10--12 supply context, not a single reusable problem with identical hypotheses.

\paragraph{Audit.} Literal correctness: false as written. Publication match: none found. Verification: unresolved lead. Development: annotated. Confidence: low.

\subsubsection{1984-637.9 --- Euler flow on SO(4) and its Casimirs}

\label{item:1984-637.9}

\noindent\textbf{Kind:} example. \textbf{Fate:} E. \textbf{Pages:} 1191--1194.\par

\paragraph{Notebook content.} A quadratic rigid-body Hamiltonian is described using two three-vectors.

\paragraph{Precise statement.} Using \(\mathfrak{so}(4)\cong\mathfrak{so}(3)\oplus\mathfrak{so}(3)\), coordinates \(u,v\) have independent cross-product Lie--Poisson brackets and Casimirs \(|u|^2,|v|^2\). A positive quadratic inertia Hamiltonian gives the real geodesic-flow problem; extra integrability requires restrictions on its coefficients.

\paragraph{Literature and fate.} The polynomial Poisson setup and classical rigid-body interpretation predate the seminar. A general positive quadratic form is not thereby algebraically completely integrable.

\paragraph{Present status.} As of 12 November 2026: The polynomial Poisson setup and classical rigid-body interpretation predate the seminar. A general positive quadratic form is not thereby algebraically completely integrable.

\paragraph{Audit.} Literal correctness: correct with missing hypotheses. Publication match: partial. Verification: indirect. Development: none. Confidence: high.

\noindent References: \citep{adler1982kowalewski,adler1984geodesic}.\par

\subsubsection{1984-637.10 --- Unresolved SO(4) parameter constraint}

\label{item:1984-637.10}

\noindent\textbf{Kind:} formula. \textbf{Fate:} F. \textbf{Pages:} 1191--1194.\par

\paragraph{Notebook content.} A long relation among coefficient vectors contains an ambiguous separator or component grouping.

\paragraph{Precise statement.} The surviving display cannot be converted into a unique polynomial constraint without choosing between materially different readings. No such choice is made here.

\paragraph{Literature and fate.} The July 1984 SO(4) paper is a likely contextual lead and the 1987 paper was read, but neither establishes an exact match to this damaged display. The formula remains F and the ambiguity is retained in the corrections ledger.

\paragraph{Present status.} As of 12 November 2026: The July 1984 SO(4) paper is a likely contextual lead and the 1987 paper was read, but neither establishes an exact match to this damaged display. The formula remains F and the ambiguity is retained in the corrections ledger.

\paragraph{Audit.} Literal correctness: correct with missing hypotheses. Publication match: partial. Verification: unresolved lead. Development: none. Confidence: low.

\noindent References: \citep{adler1984geodesic,adler1987quadrics}.\par

\subsubsection{1984-637.11 --- Four hyperelliptic boundary curves}

\label{item:1984-637.11}

\noindent\textbf{Kind:} claim. \textbf{Fate:} F. \textbf{Pages:} 1191--1194.\par

\paragraph{Notebook content.} Four curves at infinity and a hyperelliptic description are mentioned.

\paragraph{Precise statement.} A precise assertion would need the family of invariant surfaces, defining parameters, genera, intersections, and the relation to the pole divisor. These are not sufficiently specified in the notes.

\paragraph{Literature and fate.} The 1987 quadrics paper is relevant but does not justify importing its entire family or genus into the 1984 fragment. No open problem is inferred from an unidentified claim.

\paragraph{Present status.} As of 12 November 2026: The 1987 quadrics paper is relevant but does not justify importing its entire family or genus into the 1984 fragment. No open problem is inferred from an unidentified claim.

\paragraph{Audit.} Literal correctness: correct with missing hypotheses. Publication match: partial. Verification: read the result. Development: none. Confidence: low.

\noindent References: \citep{adler1987quadrics}.\par

\subsubsection{1984-637.12 --- Lax representation in a loop algebra}

\label{item:1984-637.12}

\noindent\textbf{Kind:} construction. \textbf{Fate:} E. \textbf{Pages:} 1191--1194.\par

\paragraph{Notebook content.} The algebra of spectral-parameter loops is introduced with a notation resembling a rational-function field.

\paragraph{Precise statement.} The algebraic Laurent-polynomial loop algebra is \(\mathfrak g\otimes\mathbb C[\lambda,\lambda^{-1}]\), with bracket on coefficients. A Lax equation \(\dot L=[L,M]\) preserves suitable spectral invariants but does not alone prove a complete set of independent commuting integrals.

\paragraph{Literature and fate.} The loop and r-matrix framework is earlier than the talk, including Semenov-Tian-Shansky (1983). Laurent polynomials and all rational functions are different algebras.

\paragraph{Present status.} As of 12 November 2026: The loop and r-matrix framework is earlier than the talk, including Semenov-Tian-Shansky (1983). Laurent polynomials and all rational functions are different algebras.

\paragraph{Audit.} Literal correctness: correct with missing hypotheses. Publication match: exact. Verification: read the result. Development: annotated. Confidence: high.

\noindent References: \citep{semenov1983classical}.\par

\subsubsection{1984-637.13 --- Quadratic r-matrix Poisson bracket}

\label{item:1984-637.13}

\noindent\textbf{Kind:} theorem. \textbf{Fate:} E. \textbf{Pages:} 1191--1194.\par

\paragraph{Notebook content.} A second Poisson bracket quadratic in the matrix variable is used.

\paragraph{Precise statement.} Under the stated associative-algebra, invariant trace pairing, and skew r-matrix/modified Yang--Baxter hypotheses, the quadratic r-bracket is Poisson and supplies the corresponding commuting spectral Hamiltonians. Closure and the r-matrix identity must be checked.

\paragraph{Literature and fate.} Semenov-Tian-Shansky (1983), Section 5 and Theorem 13, is a verified prior publication match.

\paragraph{Present status.} As of 12 November 2026: Semenov-Tian-Shansky (1983), Section 5 and Theorem 13, is a verified prior publication match.

\paragraph{Audit.} Literal correctness: correct with missing hypotheses. Publication match: exact. Verification: read the result. Development: none. Confidence: high.

\noindent References: \citep{semenov1983classical}.\par

\subsubsection{1984-637.14 --- Why the associative construction does not directly restrict to so(n)}

\label{item:1984-637.14}

\noindent\textbf{Kind:} obstruction. \textbf{Fate:} E. \textbf{Pages:} 1191--1194.\par

\paragraph{Notebook content.} A difficulty with the second bracket on an orthogonal Lie algebra is noted.

\paragraph{Precise statement.} Skew-symmetric matrices are closed under commutators but not ordinary matrix multiplication. Therefore an associative quadratic-bracket construction cannot be transferred to \(\mathfrak{so}(n)\) merely by replacing the ambient algebra without checking closure or reduction.

\paragraph{Literature and fate.} This is an obstruction to a particular construction, not a theorem that no compatible second bracket can exist on any orthogonal system.

\paragraph{Present status.} As of 12 November 2026: This is an obstruction to a particular construction, not a theorem that no compatible second bracket can exist on any orthogonal system.

\paragraph{Audit.} Literal correctness: correct with missing hypotheses. Publication match: partial. Verification: read the result. Development: none. Confidence: high.

\noindent References: \citep{semenov1983classical}.\par

\subsubsection{1984-637.15 --- Symplectic leaves from double-coset geometry}

\label{item:1984-637.15}

\noindent\textbf{Kind:} theorem. \textbf{Fate:} A. \textbf{Pages:} 1191--1194.\par

\paragraph{Notebook content.} The new Poisson-group bracket is connected with double cosets.

\paragraph{Precise statement.} In the factorizable/split Poisson--Lie setting, the relevant intersections of double cosets, with appropriate connected-component and local-factorization conditions, describe symplectic leaves.

\paragraph{Literature and fate.} The precise match is Semenov-Tian-Shansky (1985), Theorem 7. Drinfeld (1983) supplies the double and Lie-bialgebra framework but does not contain this leaf theorem; class A reflects the speaker's later paper.

\paragraph{Present status.} As of 12 November 2026: The precise match is Semenov-Tian-Shansky (1985), Theorem 7. Drinfeld (1983) supplies the double and Lie-bialgebra framework but does not contain this leaf theorem; class A reflects the speaker's later paper.

\paragraph{Audit.} Literal correctness: correct with missing hypotheses. Publication match: partial. Verification: read the result. Development: none. Confidence: high.

\noindent References: \citep{drinfeld1983hamiltonian,semenov1985dressing}.\par

\subsubsection{1984-637.16 --- Infinitesimal defect of a Poisson group action}

\label{item:1984-637.16}

\noindent\textbf{Kind:} identity. \textbf{Fate:} A. \textbf{Pages:} 1191--1194.\par

\paragraph{Notebook content.} The action does not preserve brackets pointwise in the ordinary symplectic-action sense.

\paragraph{Precise statement.} For a Poisson--Lie action, differentiating the Poisson action-map condition gives \(X_M\{f,g\}-\{X_Mf,g\}-\{f,X_Mg\}=\langle[\alpha_f,\alpha_g]_*,X\rangle\), up to a fixed infinitesimal sign convention; \(\alpha_f\) is the differential along the orbit at the group identity.

\paragraph{Literature and fate.} Semenov-Tian-Shansky (1985), Proposition 5 and equation (22), gives the exact near-term formulation. The right-hand side need not vanish.

\paragraph{Present status.} As of 12 November 2026: Semenov-Tian-Shansky (1985), Proposition 5 and equation (22), gives the exact near-term formulation. The right-hand side need not vanish.

\paragraph{Audit.} Literal correctness: correct with missing hypotheses. Publication match: exact. Verification: read the result. Development: none. Confidence: high.

\noindent References: \citep{semenov1985dressing}.\par

\subsubsection{1984-637.17 --- Invariant reduction without an ordinary moment map}

\label{item:1984-637.17}

\noindent\textbf{Kind:} theorem. \textbf{Fate:} A. \textbf{Pages:} 1191--1194.\par

\paragraph{Notebook content.} A subgroup condition makes invariant functions closed under the Poisson bracket.

\paragraph{Precise statement.} For a connected subgroup \(H\), if \(\mathfrak h^\perp\subset\mathfrak g^*\) is a Lie subalgebra, the \(H\)-invariant functions of a Poisson--Lie action form a Poisson algebra. A manifold quotient additionally requires regularity; an ordinary \(\mathfrak g^*\)-valued moment map is not assumed.

\paragraph{Literature and fate.} Semenov-Tian-Shansky (1985), Theorem 6, proves this sufficient condition. The later record 639 repeats it with an unjustified converse, checked separately.

\paragraph{Present status.} As of 12 November 2026: Semenov-Tian-Shansky (1985), Theorem 6, proves this sufficient condition. The later record 639 repeats it with an unjustified converse, checked separately.

\paragraph{Audit.} Literal correctness: correct with missing hypotheses. Publication match: exact. Verification: read the result. Development: none. Confidence: high.

\noindent References: \citep{semenov1985dressing}.\par

\subsubsection{1984-637.18 --- Sosinsky on Grothendieck and dessins}

\label{item:1984-637.18}

\noindent\textbf{Kind:} reference. \textbf{Fate:} F. \textbf{Pages:} 1191--1194.\par

\paragraph{Notebook content.} A short heading or proposed report names Sosinsky and Grothendieck's programme.

\paragraph{Precise statement.} No explicit dessin, Galois action, correspondence, or problem is stated in the fragment.

\paragraph{Literature and fate.} The Esquisse provides context only. The heading is preserved without converting a seminar plan into a theorem.

\paragraph{Present status.} As of 12 November 2026: The Esquisse provides context only. The heading is preserved without converting a seminar plan into a theorem.

\paragraph{Audit.} Literal correctness: correct with missing hypotheses. Publication match: thematic only. Verification: read the result. Development: none. Confidence: low.

\noindent References: \citep{grothendieck1984esquisse}.\par

\subsection{Record 638: Atiyah--Singer: families of Dirac operators and anomalies}

\label{talk:638}

\noindent\textbf{Date:} 1984-10-08. \textbf{Speaker:} A. A. Beilinson. \textbf{Pages:} 1194--1195.\par

The contractible connection space models a based gauge-group classifying space; the family index is a map into K-theory, and the anomaly involves a complex chiral determinant.

\subsubsection{1984-638.1 --- Connections modulo based gauge transformations}

\label{item:1984-638.1}

\noindent\textbf{Kind:} construction. \textbf{Fate:} E. \textbf{Pages:} 1194--1195.\par

\paragraph{Notebook content.} The affine space of connections and a gauge quotient are used as a classifying space.

\paragraph{Precise statement.} After suitable Sobolev completion, the affine connection space \(\mathcal A\) is contractible and the based gauge group \(\mathcal G_0\) acts freely; under the standard principal-bundle hypotheses \(\mathcal A/\mathcal G_0\simeq B\mathcal G_0\). Unbased stabilizers would invalidate the free-action argument.

\paragraph{Literature and fate.} Atiyah--Singer's April 1984 publication predates the October report. Singer's 1985 exposition, Section 2, supplies the verified detailed construction.

\paragraph{Present status.} As of 12 November 2026: Atiyah--Singer's April 1984 publication predates the October report. Singer's 1985 exposition, Section 2, supplies the verified detailed construction.

\paragraph{Audit.} Literal correctness: correct with missing hypotheses. Publication match: exact. Verification: read the result. Development: none. Confidence: high.

\noindent References: \citep{atiyahsinger1984dirac,singer1985families}.\par

\subsubsection{1984-638.2 --- The family-index map is not a homotopy equivalence}

\label{item:1984-638.2}

\noindent\textbf{Kind:} claim. \textbf{Fate:} E. \textbf{Pages:} 1194--1195.\par

\paragraph{Notebook content.} A classifying-space arrow is written in a way that suggests equivalence.

\paragraph{Precise statement.} A family of coupled chiral Dirac operators defines an index class and a classifying map \(B\mathcal G_0\to\mathbb Z\times BU\), subject to the family and base hypotheses. The construction does not imply that this map is a homotopy equivalence.

\paragraph{Literature and fate.} The family-index construction is prior material. Singer (1985), Section 7, identifies the intended map; the stronger equivalence is unsupported and generally false.

\paragraph{Present status.} As of 12 November 2026: The family-index construction is prior material. Singer (1985), Section 7, identifies the intended map; the stronger equivalence is unsupported and generally false.

\paragraph{Audit.} Literal correctness: false as written. Publication match: exact. Verification: read the result. Development: annotated. Confidence: high.

\noindent References: \citep{atiyahsinger1984dirac,singer1985families}.\par

\subsubsection{1984-638.3 --- The chiral determinant and gauge anomaly}

\label{item:1984-638.3}

\noindent\textbf{Kind:} formula. \textbf{Fate:} E. \textbf{Pages:} 1194--1195.\par

\paragraph{Notebook content.} An anomaly is represented by a logarithm of a determinant; an adjoint product would erase relevant phase information.

\paragraph{Precise statement.} A determinant-line/trivialization construction uses complex chiral operators relative to a fixed reference. On a suitable index-zero invertible locus, the regularized determinant of a reference-composed chiral operator retains its phase; the positive operator \(D_A^*D_A\) alone does not determine the chiral anomaly.

\paragraph{Literature and fate.} Singer (1985), Sections 4 and 7, gives the verified construction corresponding to the earlier Atiyah--Singer report. Invertibility, index, regularization, and a local choice of logarithm must be stated.

\paragraph{Present status.} As of 12 November 2026: Singer (1985), Sections 4 and 7, gives the verified construction corresponding to the earlier Atiyah--Singer report. Invertibility, index, regularization, and a local choice of logarithm must be stated.

\paragraph{Audit.} Literal correctness: false as written. Publication match: partial. Verification: read the result. Development: annotated. Confidence: high.

\noindent References: \citep{atiyahsinger1984dirac,singer1985families}.\par

\subsection{Record 639: Poisson--Lie groups, doubles, dressing actions, and reduction}

\label{talk:639}

\noindent\textbf{Date:} 1984-10-29. \textbf{Speaker:} M. A. Semenov-Tian-Shansky. \textbf{Pages:} 1195--1198.\par

Definitions and older r-matrix constructions lead to the speaker's 1985 action, leaf and reduction results. Repeated labels, a missing star, and an unjustified converse are retained as notebook errors and explicitly repaired.

\subsubsection{1984-639.1 --- Poisson manifold}

\label{item:1984-639.1}

\noindent\textbf{Kind:} definition. \textbf{Fate:} E. \textbf{Pages:} 1195--1198.\par

\paragraph{Notebook content.} A bivector and the Jacobi identity define a Poisson structure.

\paragraph{Precise statement.} A smooth bivector \(\pi\) defines \(\{f,g\}=\pi(df,dg)\); the Leibniz identity is automatic and Jacobi is equivalent to \([\pi,\pi]_{\mathrm{Sch}}=0\). Nonconstant rank is allowed.

\paragraph{Literature and fate.} Standard before the talk; Drinfeld (1983), Section 1, uses this definition.

\paragraph{Present status.} As of 12 November 2026: Standard before the talk; Drinfeld (1983), Section 1, uses this definition.

\paragraph{Audit.} Literal correctness: correct with missing hypotheses. Publication match: exact. Verification: read the result. Development: none. Confidence: high.

\noindent References: \citep{drinfeld1983hamiltonian}.\par

\subsubsection{1984-639.2 --- Product of Poisson manifolds}

\label{item:1984-639.2}

\noindent\textbf{Kind:} construction. \textbf{Fate:} E. \textbf{Pages:} 1195--1198.\par

\paragraph{Notebook content.} The product bracket is written using gradients with inconsistent labels.

\paragraph{Precise statement.} On \(M_1\times M_2\), set \(\pi=\operatorname{pr}_1^*\pi_1+\operatorname{pr}_2^*\pi_2\) in the natural lifted-tensor sense: brackets within each factor are its original bracket, and functions pulled back from different factors commute.

\paragraph{Literature and fate.} The direct-product construction is classical. Both differentials in each summand must refer to that same factor; the repeated source labels are a local notation error.

\paragraph{Present status.} As of 12 November 2026: The direct-product construction is classical. Both differentials in each summand must refer to that same factor; the repeated source labels are a local notation error.

\paragraph{Audit.} Literal correctness: correct with missing hypotheses. Publication match: exact. Verification: read the result. Development: annotated. Confidence: high.

\noindent References: \citep{drinfeld1983hamiltonian}.\par

\subsubsection{1984-639.3 --- Multiplicative Poisson tensor}

\label{item:1984-639.3}

\noindent\textbf{Kind:} definition. \textbf{Fate:} E. \textbf{Pages:} 1195--1198.\par

\paragraph{Notebook content.} Group multiplication is required to be Poisson.

\paragraph{Precise statement.} A Poisson--Lie group satisfies \(\pi(gh)=(L_g)_*\pi(h)+(R_h)_*\pi(g)\). In particular \(\pi(e)=0\), and differentiation at the identity gives a Lie bialgebra.

\paragraph{Literature and fate.} Drinfeld's 1983 definition and infinitesimal correspondence predate the talk.

\paragraph{Present status.} As of 12 November 2026: Drinfeld's 1983 definition and infinitesimal correspondence predate the talk.

\paragraph{Audit.} Literal correctness: correct with missing hypotheses. Publication match: exact. Verification: read the result. Development: none. Confidence: high.

\noindent References: \citep{drinfeld1983hamiltonian}.\par

\subsubsection{1984-639.4 --- The zero Poisson structure}

\label{item:1984-639.4}

\noindent\textbf{Kind:} example. \textbf{Fate:} E. \textbf{Pages:} 1195--1198.\par

\paragraph{Notebook content.} The zero bracket is the elementary example.

\paragraph{Precise statement.} Any Lie group with \(\pi=0\) is Poisson--Lie, with zero cobracket and abelian dual Lie algebra. An action of this group is Poisson exactly when each acting transformation preserves the target Poisson tensor.

\paragraph{Literature and fate.} Classical and contained in the defining identities; no open problem is present.

\paragraph{Present status.} As of 12 November 2026: Classical and contained in the defining identities; no open problem is present.

\paragraph{Audit.} Literal correctness: correct with missing hypotheses. Publication match: exact. Verification: read the result. Development: none. Confidence: high.

\noindent References: \citep{drinfeld1983hamiltonian}.\par

\subsubsection{1984-639.5 --- A dual Lie algebra as a linear Poisson vector group}

\label{item:1984-639.5}

\noindent\textbf{Kind:} example. \textbf{Fate:} E. \textbf{Pages:} 1195--1198.\par

\paragraph{Notebook content.} The additive group of a dual vector space carries the Lie--Poisson bracket.

\paragraph{Precise statement.} For a finite-dimensional Lie algebra \(\mathfrak a\), \(\{\ell_X,\ell_Y\}(\xi)=\langle\xi,[X,Y]\rangle\) extends to a linear Poisson tensor on \(\mathfrak a^*\); linearity makes it multiplicative for addition.

\paragraph{Literature and fate.} An older construction used in Drinfeld's 1983 examples.

\paragraph{Present status.} As of 12 November 2026: An older construction used in Drinfeld's 1983 examples.

\paragraph{Audit.} Literal correctness: correct with missing hypotheses. Publication match: exact. Verification: read the result. Development: none. Confidence: high.

\noindent References: \citep{drinfeld1983hamiltonian}.\par

\subsubsection{1984-639.6 --- Manin triples and the double}

\label{item:1984-639.6}

\noindent\textbf{Kind:} definition. \textbf{Fate:} E. \textbf{Pages:} 1195--1198.\par

\paragraph{Notebook content.} Two isotropic subalgebras are paired by a nondegenerate form.

\paragraph{Precise statement.} A Manin triple is a Lie algebra \(\mathfrak d\) with invariant nondegenerate symmetric bilinear form and complementary isotropic Lie subalgebras \(\mathfrak g,\mathfrak g^*\). The form is not a positive-definite inner product; its isotropic summands are essential.

\paragraph{Literature and fate.} Drinfeld (1983) constructs this Lie-bialgebra double. Positivity would make a nonzero isotropic decomposition impossible.

\paragraph{Present status.} As of 12 November 2026: Drinfeld (1983) constructs this Lie-bialgebra double. Positivity would make a nonzero isotropic decomposition impossible.

\paragraph{Audit.} Literal correctness: correct with missing hypotheses. Publication match: exact. Verification: read the result. Development: none. Confidence: high.

\noindent References: \citep{drinfeld1983hamiltonian}.\par

\subsubsection{1984-639.7 --- Gelfand's exercise: multiplicativity on a vector group forces linearity}

\label{item:1984-639.7}

\noindent\textbf{Kind:} exercise. \textbf{Fate:} E. \textbf{Pages:} 1195--1198.\par

\paragraph{Notebook content.} The seminar asks to identify Poisson--Lie brackets on a vector group.

\paragraph{Precise statement.} For \((V,+)\), multiplicativity gives \(\pi(x+y)=\pi(x)+\pi(y)\). Smoothness implies real linearity, and Jacobi on linear functions is precisely the Jacobi identity for a Lie bracket on \(V^*\).

\paragraph{Literature and fate.} This closes the exercise directly and fits the pre-existing Drinfeld framework. No unproved nonlinear multiplicative examples are needed.

\paragraph{Present status.} As of 12 November 2026: This closes the exercise directly and fits the pre-existing Drinfeld framework. No unproved nonlinear multiplicative examples are needed.

\paragraph{Audit.} Literal correctness: correct with missing hypotheses. Publication match: exact. Verification: read the result. Development: annotated. Confidence: high.

\noindent Canonical problem: GP-1984-04.\par

\noindent References: \citep{drinfeld1983hamiltonian}.\par

\subsubsection{1984-639.8 --- The r-matrix/Sklyanin group bracket}

\label{item:1984-639.8}

\noindent\textbf{Kind:} construction. \textbf{Fate:} E. \textbf{Pages:} 1195--1198.\par

\paragraph{Notebook content.} A splitting operator gives a quadratic group bracket.

\paragraph{Precise statement.} For a suitable skew r-matrix satisfying the modified classical Yang--Baxter condition, the left-minus-right translated bivector defines a multiplicative Poisson bracket. In matrix notation the familiar bracket is of Sklyanin type; the invariant pairing and normalization must be fixed.

\paragraph{Literature and fate.} The construction is already in Drinfeld (1983) and Semenov-Tian-Shansky (1983), including the quadratic bracket theorem.

\paragraph{Present status.} As of 12 November 2026: The construction is already in Drinfeld (1983) and Semenov-Tian-Shansky (1983), including the quadratic bracket theorem.

\paragraph{Audit.} Literal correctness: correct with missing hypotheses. Publication match: exact. Verification: read the result. Development: none. Confidence: high.

\noindent References: \citep{drinfeld1983hamiltonian,semenov1983classical}.\par

\subsubsection{1984-639.9 --- Triangular r-matrix and the missing sign}

\label{item:1984-639.9}

\noindent\textbf{Kind:} example. \textbf{Fate:} E. \textbf{Pages:} 1195--1198.\par

\paragraph{Notebook content.} The triangular matrix formula has an unclear or missing minus sign.

\paragraph{Precise statement.} For the standard trace-skew triangular operator on \(\mathfrak{gl}_n\), \(R(X)=X_{>0}-X_{<0}\), with the diagonal component mapped to zero. Equivalent split conventions redistribute half-diagonal terms; the upper/lower signs cannot both be positive.

\paragraph{Literature and fate.} This is the standard prior r-matrix example. The source sign must be compared with the stated projection convention rather than guessed from a later formula.

\paragraph{Present status.} As of 12 November 2026: This is the standard prior r-matrix example. The source sign must be compared with the stated projection convention rather than guessed from a later formula.

\paragraph{Audit.} Literal correctness: false as written. Publication match: exact. Verification: read the result. Development: annotated. Confidence: high.

\noindent References: \citep{semenov1983classical}.\par

\subsubsection{1984-639.10 --- Leaf geometry from double cosets}

\label{item:1984-639.10}

\noindent\textbf{Kind:} theorem. \textbf{Fate:} A. \textbf{Pages:} 1195--1198.\par

\paragraph{Notebook content.} Symplectic leaves are identified with double-coset intersections.

\paragraph{Precise statement.} Within the factorizable Poisson--Lie setting of the cited theorem, take the prescribed double-coset intersections and their relevant connected components, on the appropriate factorization domain, to obtain symplectic leaves. An arbitrary Bruhat cell is not automatically a leaf.

\paragraph{Literature and fate.} The exact developed result is Semenov-Tian-Shansky (1985), Theorem 7, following this report. Drinfeld (1983) is background, not a source for this precise leaf assertion.

\paragraph{Present status.} As of 12 November 2026: The exact developed result is Semenov-Tian-Shansky (1985), Theorem 7, following this report. Drinfeld (1983) is background, not a source for this precise leaf assertion.

\paragraph{Audit.} Literal correctness: correct with missing hypotheses. Publication match: partial. Verification: read the result. Development: none. Confidence: high.

\noindent References: \citep{semenov1985dressing}.\par

\subsubsection{1984-639.11 --- Symplectic foliation of a Poisson manifold}

\label{item:1984-639.11}

\noindent\textbf{Kind:} theorem. \textbf{Fate:} E. \textbf{Pages:} 1195--1198.\par

\paragraph{Notebook content.} A Poisson manifold is decomposed into symplectic leaves.

\paragraph{Precise statement.} The distribution \(\pi^\sharp(T^*M)\), though its rank may vary, integrates into immersed symplectic leaves. The local structure near a point splits into a symplectic factor and a transverse Poisson structure vanishing at that point.

\paragraph{Literature and fate.} Weinstein's 1983 local-structure theorem predates the talk. Its precise statement is checked indirectly through the author's 1985 survey; the original full proof was not readable in this run.

\paragraph{Present status.} As of 12 November 2026: Weinstein's 1983 local-structure theorem predates the talk. Its precise statement is checked indirectly through the author's 1985 survey; the original full proof was not readable in this run.

\paragraph{Audit.} Literal correctness: correct with missing hypotheses. Publication match: exact. Verification: indirect. Development: none. Confidence: high.

\noindent References: \citep{weinstein1983local,weinstein1985poisson}.\par

\subsubsection{1984-639.12 --- Poisson actions and the infinitesimal defect identity}

\label{item:1984-639.12}

\noindent\textbf{Kind:} identity. \textbf{Fate:} A. \textbf{Pages:} 1195--1198.\par

\paragraph{Notebook content.} An action-map definition is differentiated to compare brackets and infinitesimal generators.

\paragraph{Precise statement.} For a Poisson--Lie action \(G\times M\to M\), the failure of an infinitesimal generator to be a Poisson derivation is the contraction of the Lie-bialgebra cobracket with the two orbit differentials. With the chapter's signs it is the dual-bracket pairing displayed in item 637.16.

\paragraph{Literature and fate.} Semenov-Tian-Shansky (1985), Proposition 5, equation (22), is the exact near-term publication match. This repeats the September occurrence.

\paragraph{Present status.} As of 12 November 2026: Semenov-Tian-Shansky (1985), Proposition 5, equation (22), is the exact near-term publication match. This repeats the September occurrence.

\paragraph{Audit.} Literal correctness: correct with missing hypotheses. Publication match: exact. Verification: read the result. Development: none. Confidence: high.

\noindent References: \citep{semenov1985dressing}.\par

\subsubsection{1984-639.13 --- Left and right cotangent moment maps}

\label{item:1984-639.13}

\noindent\textbf{Kind:} construction. \textbf{Fate:} E. \textbf{Pages:} 1195--1198.\par

\paragraph{Notebook content.} The cotangent bundle of a group maps to two copies of its dual algebra; diagram labels repeat.

\paragraph{Precise statement.} The cotangent lifts of left and right translations on \(T^*G\) have commuting momentum functions. With the raw left/right translation conventions, the two targets carry opposite Lie--Poisson signs; alternatively insert a sign in one momentum map.

\paragraph{Literature and fate.} The construction is older. Weinstein's 1985 survey, Section 5, verifies the distinction and cites the 1983 source. Repeated source labels are not copied into the reconstructed diagram.

\paragraph{Present status.} As of 12 November 2026: The construction is older. Weinstein's 1985 survey, Section 5, verifies the distinction and cites the 1983 source. Repeated source labels are not copied into the reconstructed diagram.

\paragraph{Audit.} Literal correctness: correct with missing hypotheses. Publication match: exact. Verification: indirect. Development: annotated. Confidence: high.

\noindent References: \citep{weinstein1983local,weinstein1985poisson}.\par

\subsubsection{1984-639.14 --- Full dual pairs versus commuting subalgebras}

\label{item:1984-639.14}

\noindent\textbf{Kind:} definition. \textbf{Fate:} E. \textbf{Pages:} 1195--1198.\par

\paragraph{Notebook content.} Mutually commuting pullbacks are called a dual pair.

\paragraph{Precise statement.} For Poisson submersions from a symplectic manifold, a full dual pair requires \((\ker dJ_1)^\omega=\ker dJ_2\). Merely requiring all pullbacks from the two targets to commute gives only an inclusion; global mutual-centralizer conclusions require suitable connected-fibre hypotheses.

\paragraph{Literature and fate.} This distinction belongs to the earlier dual-pair theory, checked through Weinstein's survey and the speaker's 1985 use of it. The notebook omits necessary rank and fullness qualifications.

\paragraph{Present status.} As of 12 November 2026: This distinction belongs to the earlier dual-pair theory, checked through Weinstein's survey and the speaker's 1985 use of it. The notebook omits necessary rank and fullness qualifications.

\paragraph{Audit.} Literal correctness: correct with missing hypotheses. Publication match: exact. Verification: indirect. Development: annotated. Confidence: high.

\noindent References: \citep{weinstein1983local,weinstein1985poisson,semenov1985dressing}.\par

\subsubsection{1984-639.15 --- Local symplectic-leaf correspondence in a dual pair}

\label{item:1984-639.15}

\noindent\textbf{Kind:} theorem. \textbf{Fate:} E. \textbf{Pages:} 1195--1198.\par

\paragraph{Notebook content.} A preimage under one arrow maps to a leaf under the other.

\paragraph{Precise statement.} For a full dual pair with the required submersion/rank assumptions, the image of a connected local preimage of a leaf under one map lies in and locally fills a leaf of the other. Identifying complete global leaves requires further completeness and fibre hypotheses.

\paragraph{Literature and fate.} The older local dual-pair result must not be enlarged into an unrestricted global correspondence. The verified source discussions are Weinstein (1985), Section 5, and Semenov-Tian-Shansky (1985).

\paragraph{Present status.} As of 12 November 2026: The older local dual-pair result must not be enlarged into an unrestricted global correspondence. The verified source discussions are Weinstein (1985), Section 5, and Semenov-Tian-Shansky (1985).

\paragraph{Audit.} Literal correctness: correct with missing hypotheses. Publication match: exact. Verification: indirect. Development: annotated. Confidence: high.

\noindent References: \citep{weinstein1983local,weinstein1985poisson,semenov1985dressing}.\par

\subsubsection{1984-639.16 --- The Lie-bialgebra double and its group}

\label{item:1984-639.16}

\noindent\textbf{Kind:} construction. \textbf{Fate:} E. \textbf{Pages:} 1195--1198.\par

\paragraph{Notebook content.} The vector sum of a Lie algebra and its dual is given a compatible bracket.

\paragraph{Precise statement.} The double \(\mathfrak d=\mathfrak g\oplus\mathfrak g^*\) has the canonical split pairing and the bracket determined by invariance together with the brackets on its two isotropic summands. Passing to groups requires choosing integrations and a factorization domain.

\paragraph{Literature and fate.} Drinfeld (1983), Theorem 2 and its integration discussion, predates the talk. A Lie-algebra direct sum as vector spaces need not be a direct product of Lie algebras.

\paragraph{Present status.} As of 12 November 2026: Drinfeld (1983), Theorem 2 and its integration discussion, predates the talk. A Lie-algebra direct sum as vector spaces need not be a direct product of Lie algebras.

\paragraph{Audit.} Literal correctness: correct with missing hypotheses. Publication match: exact. Verification: read the result. Development: none. Confidence: high.

\noindent References: \citep{drinfeld1983hamiltonian}.\par

\subsubsection{1984-639.17 --- The plus and minus brackets on the double}

\label{item:1984-639.17}

\noindent\textbf{Kind:} construction. \textbf{Fate:} A. \textbf{Pages:} 1195--1198.\par

\paragraph{Notebook content.} Two double brackets and four quotient arrows are written with repeated indices.

\paragraph{Precise statement.} With the canonical double r-tensor, left-minus-right translation gives the multiplicative bracket, while the left-plus-right version yields the Heisenberg-double bracket. The latter is nondegenerate near the identity and on its suitable factorization locus; global nondegeneracy is not automatic.

\paragraph{Literature and fate.} Semenov-Tian-Shansky (1985), the double and dual-pair constructions, supplies the near-term match. All four arrows must have distinct source/target and sign labels; the notebook repeats them.

\paragraph{Present status.} As of 12 November 2026: Semenov-Tian-Shansky (1985), the double and dual-pair constructions, supplies the near-term match. All four arrows must have distinct source/target and sign labels; the notebook repeats them.

\paragraph{Audit.} Literal correctness: false as written. Publication match: partial. Verification: read the result. Development: annotated. Confidence: high.

\noindent References: \citep{semenov1985dressing}.\par

\subsubsection{1984-639.18 --- The complementary quotient needs a dual-group star}

\label{item:1984-639.18}

\noindent\textbf{Kind:} formula. \textbf{Fate:} A. \textbf{Pages:} 1195--1198.\par

\paragraph{Notebook content.} A quotient diagram repeats \(G\) where the complementary subgroup is required.

\paragraph{Precise statement.} The double has complementary subgroups integrating \(\mathfrak g\) and \(\mathfrak g^*\). Quotients by \(G\) and by \(G^*\) play different roles; locally a factorization identifies them with the complementary factors, with the indicated Poisson signs.

\paragraph{Literature and fate.} The missing dual-group marker is a source-level diagram error. The chapter repairs the mathematical distinction without inferring a global group factorization.

\paragraph{Present status.} As of 12 November 2026: The missing dual-group marker is a source-level diagram error. The chapter repairs the mathematical distinction without inferring a global group factorization.

\paragraph{Audit.} Literal correctness: false as written. Publication match: partial. Verification: read the result. Development: annotated. Confidence: high.

\noindent References: \citep{semenov1985dressing}.\par

\subsubsection{1984-639.19 --- Coisotropic invariant functions: sufficiency but not a single-action converse}

\label{item:1984-639.19}

\noindent\textbf{Kind:} theorem. \textbf{Fate:} A. \textbf{Pages:} 1195--1198.\par

\paragraph{Notebook content.} The invariant-subalgebra criterion is stated with an if-and-only-if.

\paragraph{Precise statement.} If \(H\) is connected and \(\mathfrak h^\perp\) is a Lie subalgebra of \(\mathfrak g^*\), then \(C^\infty(M)^H\) is Poisson-closed for every Poisson--Lie action. Necessity fails for a fixed action: a trivial action has every function invariant, regardless of \(H\).

\paragraph{Literature and fate.} The sufficient direction is Semenov-Tian-Shansky (1985), Theorem 6. The chapter supplies a full noncoisotropic connected-subgroup counterexample using the additive group with \(\pi=z\partial_x\wedge\partial_y\), independently checked.

\paragraph{Present status.} As of 12 November 2026: The sufficient direction is Semenov-Tian-Shansky (1985), Theorem 6. The chapter supplies a full noncoisotropic connected-subgroup counterexample using the additive group with \(\pi=z\partial_x\wedge\partial_y\), independently checked.

\paragraph{Audit.} Literal correctness: false as written. Publication match: exact. Verification: read the result. Development: annotated. Confidence: high.

\noindent References: \citep{semenov1985dressing}.\par

\subsubsection{1984-639.20 --- Dressing action and the unit test}

\label{item:1984-639.20}

\noindent\textbf{Kind:} formula. \textbf{Fate:} A. \textbf{Pages:} 1195--1198.\par

\paragraph{Notebook content.} The displayed factorization formula for dressing fails a direct identity-element check under its stated reading.

\paragraph{Precise statement.} On a factorization domain in the double, write \(u g=g' u'\) with \(u,u'\in G^*\) and \(g,g'\in G\); the corresponding local dressing transformation sends \(g\) to \(g'\). Left/right conventions may differ but must send the identity acting element to the identity transformation.

\paragraph{Literature and fate.} Semenov-Tian-Shansky (1985), Section 6, Proposition 17 and Theorem 13, gives the precise dressing results. The abstract factorization statement is safe; an inverse is not inserted into an ambiguous source expression by guesswork.

\paragraph{Present status.} As of 12 November 2026: Semenov-Tian-Shansky (1985), Section 6, Proposition 17 and Theorem 13, gives the precise dressing results. The abstract factorization statement is safe; an inverse is not inserted into an ambiguous source expression by guesswork.

\paragraph{Audit.} Literal correctness: false as written. Publication match: partial. Verification: read the result. Development: annotated. Confidence: high.

\noindent References: \citep{semenov1985dressing}.\par

\subsubsection{1984-639.21 --- Quantization of Poisson--Lie/Lie-bialgebra structures}

\label{item:1984-639.21}

\noindent\textbf{Kind:} programme. \textbf{Fate:} B. \textbf{Pages:} 1195--1198.\par

\paragraph{Notebook content.} The lecture ends with a quantum or deformation direction.

\paragraph{Precise statement.} A precise later formulation asks for a formal Hopf-algebra quantization of a Lie bialgebra over a characteristic-zero field. This differs from convergence of analytic quantum-group representations or quantization on a prescribed global real group.

\paragraph{Literature and fate.} Etingof--Kazhdan proved formal quantization (1995 preprint; 1996 publication). This is a thematic later development of the notebook programme, not a theorem stated with those full hypotheses in 1984.

\paragraph{Present status.} As of 12 November 2026: Etingof--Kazhdan proved formal quantization (1995 preprint; 1996 publication). This is a thematic later development of the notebook programme, not a theorem stated with those full hypotheses in 1984.

\paragraph{Audit.} Literal correctness: correct with missing hypotheses. Publication match: thematic only. Verification: read the result. Development: none. Confidence: high.

\noindent References: \citep{etingofkazhdan1996quantization}.\par

\subsubsection{1984-639.22 --- Sklyanin's 1983 representation paper}

\label{item:1984-639.22}

\noindent\textbf{Kind:} reference. \textbf{Fate:} E. \textbf{Pages:} 1195--1198.\par

\paragraph{Notebook content.} The notes give a Functional Analysis and Its Applications 17(4),1983 pointer.

\paragraph{Precise statement.} The pointer identifies Sklyanin's representation work on the quadratic algebra; the exact mathematical results are indexed under Cherednik's related record 620 rather than duplicated from a bare citation.

\paragraph{Literature and fate.} A real earlier publication reference, not an open problem. Bibliographic identity is checked by the section responsible for record 620.

\paragraph{Present status.} As of 12 November 2026: A real earlier publication reference, not an open problem. Bibliographic identity is checked by the section responsible for record 620.

\paragraph{Audit.} Literal correctness: correct with missing hypotheses. Publication match: exact. Verification: publisher metadata checked. Development: none. Confidence: high.

\noindent References: \citep{sklyanin1983representations}.\par

\section{Canonical problem index}

\subsection{GP-1983-01 --- Characterize the face/h-vectors of simple or simplicial convex polytopes}

\label{problem:GP-1983-01}

Determine which integer h-vectors arise from simplicial convex polytopes, equivalently from their simple polar duals. The g-theorem imposes h-symmetry and the Macaulay M-sequence condition on successive differences.

\paragraph{Occurrences.} 1983-600.11.

\paragraph{Status.} E: solved before the talk by Billera--Lee sufficiency and Stanley necessity; the two contributions are complementary.

\noindent\citep{stanley1980faces,stanley1984lefschetz}.

\subsection{GP-1983-02 --- Density of finite-critical-orbit quadratic parameters: ambient set unspecified}

\label{problem:GP-1983-02}

The notebook asks whether parameters \(c\) for which the critical orbit of \(z^2+c\) is finite are dense, without an unambiguous ambient set.

\paragraph{Occurrences.} 1983-600.16.

\paragraph{Status.} F: the whole-plane interpretation is false because |c|>2 is an open escape region. The intended restricted formulation is unidentified, and no density-of-hyperbolicity conjecture is substituted.

\noindent\citep{douadyhubbard1985polynomes}.

\subsection{GP-1983-03 --- A rigorous operator-family framework for the anomaly sign}

\label{problem:GP-1983-03}

Editorial specialization of Gelfand's broad request: construct determinant/Pfaffian lines for elliptic families, gauge-compatible transport, and a spectral formula for mapping-torus holonomy.

\paragraph{Occurrences.} 1983-604.10.

\paragraph{Status.} B: later determinant-line and eta-invariant theory supplies a rigorous substantial answer to this specialization; the undefined general field-theory program is not claimed solved.

\noindent\citep{daifreed1994eta,wangwenwitten2019new}.

\subsection{GP-1983-04 --- Transition matrices and multiplication rules for standard symmetric functions}

\label{problem:GP-1983-04}

Determine the change-of-basis matrices among \(m_\lambda,e_\lambda,h_\lambda,p_\lambda,s_\lambda\), and the Schur multiplication coefficients, degree by degree.

\paragraph{Occurrences.} 1983-605.9, 1983-605.12, 1983-608.8, 1983-608.9, 1983-608.10.

\paragraph{Status.} E: classical formulas available before1983; exact matrix-counting and tableau rules verified in the cited reference.

\noindent\citep{macdonald1995symmetric}.

\subsection{GP-1983-05 --- Compute the normalized formal wave coefficients recursively}

\label{problem:GP-1983-05}

For a normalized scalar operator \(L_n\), recursively determine the coefficients of \(\Psi=e^{kx}(1+\sum_{r\ge1}c_rk^{-r})\) in \(L_n\Psi=k^n\Psi\), specifying integration constants.

\paragraph{Occurrences.} 1983-606.2.

\paragraph{Status.} E: the formal recursive construction was established before1983; the n=2 recurrence is supplied and its first coefficients checked.

\noindent\citep{cherednik1978differential}.

\subsection{GP-1983-06 --- Verify the complete-symmetric-function generating product}

\label{problem:GP-1983-06}

Show that \(\sum_{r\ge0}h_rt^r=\prod_i(1-tx_i)^{-1}\) as a formal series in the graded symmetric-function ring.

\paragraph{Occurrences.} 1983-605.4.

\paragraph{Status.} E: elementary exercise resolved by expanding each geometric-series factor and counting variable multiplicities.

\noindent\citep{macdonald1995symmetric}.

\subsection{GP-1983-07 --- Find the width of the rectangular partition poset}

\label{problem:GP-1983-07}

For diagrams in an \(m\)-by-\(n\) rectangle ordered by inclusion, determine the largest number of pairwise incomparable diagrams.

\paragraph{Occurrences.} 1983-600.4, 1983-600.5.

\paragraph{Status.} E: Stanley's pre1983 hard-Lefschetz/Sperner theorem identifies the width with the middle Gaussian-binomial coefficient and proves the strong Sperner property.

\noindent\citep{stanley1980weyl,stanley1984lefschetz}.

\subsection{GP-1973-53 --- Recursive odd-order KdV Lax operators (earlier stable problem)}

\label{problem:GP-1973-53}

For \(L=\partial^2+u\), construct the odd differential parts \((L^{(2r+1)/2})_+\) and the associated KdV Lax flows. The1983 occurrence treats a broader scalar order n and links only through this n=2 specialization.

\paragraph{Occurrences.} 1973-327.1, 1983-606.4.

\paragraph{Status.} E: established formal pseudodifferential construction. Earlier identifier reused for the exact KdV specialization; earlier cumulative occurrence history must be retained when merging.

\noindent\citep{gelfanddikii1976fractional,cherednik1978differential}.

\subsection{GP-1983-30 --- Direct Dyson-rank bijections for the congruences modulo 5 and 7}

\label{problem:GP-1983-30}

For $m=5$ with $n=5q+4$, and for $m=7$ with $n=7q+5$, seek a uniform direct operation on Young diagrams partitioning n that gives an m-cycle action and increases Dyson rank modulo m by 1. The rank is largest part minus number of parts. A construction obtained by separately enumerating each finite rank class and matching their enumeration indices is excluded from the intended combinatorial explanation.

\paragraph{Occurrences.} 1983-611.4.

\paragraph{Status.} Status as of 12 November 2026: C for the intended direct combinatorial construction, with an explicit opened open-problem source from 2019; no claim that every later publication has been exhaustively checked.

\paragraph{Scope.} The intended direct combinatorial construction, separated from abstract existence of bijections.

\paragraph{Assessment.} A meaningful open combinatorial direction. The negative-coefficient obstruction for one rank-polynomial route is already in Stanton and is not a new result of this companion.

\noindent\citep{atkinswinnertondyer1954partitions,stanton2019problems,stanton2019rankcrank}.

\subsection{GP-1983-31 --- The maximal irreducible dimension exponent and its constant}

\label{problem:GP-1983-31}

Let $d_n=\max_{\lambda\vdash n}f^\lambda$ and $a_n=-n^{-1/2}\log(d_n/\sqrt{n!})$. Determine the limiting exponent and its value. Aggarwal--Elboim Theorem 1 establishes a constant $d>0$ such that $|a_n-d|\leq C/\log n$ for sufficiently large n. An explicit evaluation of d and stronger asymptotics are separate questions not answered by that theorem.

\paragraph{Occurrences.} 1983-614.1.

\paragraph{Status.} Status as of 12 November 2026: the convergence component has a 2026 preprint theorem; the retained C assessment concerns the explicit constant and stronger asymptotic information. Preprint theorem status is distinguished from peer-reviewed publication.

\paragraph{Scope.} Separate existence of the limit, explicit evaluation of its value and sharper error estimates.

\paragraph{Assessment.} Potentially valuable for precise evaluation or sharper asymptotics. Numerical fits and conditional dimension bounds do not prove the constant and are not advertised as a solution.

\noindent\citep{vershikkerov1985dimensions,aggarwalelboim2026maximal}.

\subsection{GP-1983-32 --- The ordered K-group and eventual Schur positivity}

\label{problem:GP-1983-32}

For the Young-graph dimension group, identify the underlying ordered group with $\Lambda_{\mathbb Z}/(s_{(1)}-1)$. To test positivity, first represent a class by a homogeneous symmetric function f, homogenizing terms with powers of $s_{(1)}$ if necessary. The class is positive if and only if $s_{(1)}^k f$ is a nonnegative integral linear combination of Schur functions for some $k\geq0$.

\paragraph{Occurrences.} 1983-614.17.

\paragraph{Status.} Status as of 12 November 2026: E for the reconstructed ordered-group statement.

\paragraph{Scope.} The positive cone in the symmetric-function description of the Young-graph dimension group.

\paragraph{Assessment.} Useful precise annotation and historical cross-reference; no new open question is created.

\noindent\citep{vershikkerov1983kfunctor}.

\subsection{GP-1983-33 --- The spinful one-dimensional contact-gas spectral construction}

\label{problem:GP-1983-33}

For fixed N, circle length $\ell>0$ and real c, consider $H_N=-\sum_{j=1}^N\partial_{x_j}^2+2c\sum_{i<j}\delta(x_i-x_j)$ on wavefunctions in $L^2((\mathbb R/\ell\mathbb Z)^N;(\mathbb C^2)^{\otimes N})$ antisymmetric under simultaneous coordinate-spin exchange. Define the operator by its closed semibounded contact-interaction quadratic form, equivalently continuity and the contact derivative jump. Seek its energies and wavefunctions in fixed particle-number and spin sectors.

\paragraph{Occurrences.} 1983-617.10, 1983-618.10.

\paragraph{Status.} Status as of 12 November 2026: the reconstructed result is established; the cited verification has the limits stated below.

\paragraph{Scope.} The explicit repulsive Bethe eigenpair construction, with operator domain, periodic boundary conditions and completeness kept separate.

\paragraph{Assessment.} Established integrable model; clarify domain and completeness rather than manufacture a new open problem.

\noindent\citep{yang1967manybody,yang1968smatrix}.

\subsection{GP-1983-34 --- Nontrivial solutions of the tetrahedron equation}

\label{problem:GP-1983-34}

A constant tetrahedron equation for $R\in\operatorname{End}(V^{\otimes3})$ is $R_{123}R_{145}R_{246}R_{356}=R_{356}R_{246}R_{145}R_{123}$ on $V^{\otimes6}$; spectral-parameter versions require a specified compatible parameter assignment. In the broad sense of a nontrivial solution, existence was already established before the 12 December 1983 lecture.

\paragraph{Occurrences.} 1983-618.16.

\paragraph{Status.} Status as of 12 November 2026: broad nontrivial existence was proved before the 1983 lecture (E); any additional unstated restriction remains unassessable (F).

\paragraph{Scope.} Broad nontrivial existence; a narrower unstated requirement cannot be reconstructed.

\paragraph{Assessment.} Historical correction is substantive; no new broad existence problem remains.

\noindent\citep{zamolodchikov1981tetrahedron,baxter1983tetrahedron,bazhanovsergeev2006tetrahedron}.

\subsection{GP-1983-60 --- Odd–even orthogonal Grassmannian identification}

\label{problem:GP-1983-60}

For a nondegenerate hyperplane \(H\) in a complex quadratic space \(V\) of dimension \(2r+2\), prove that \(W\mapsto W\cap H\) identifies either component of \(\operatorname{OG}(r+1,V)\) with \(\operatorname{OG}(r,H)\).

\paragraph{Occurrences.} 1983-621.9.

\paragraph{Status.} Solved classical exercise. Each isotropic r-plane has two maximal extensions, one in each component. No exact earlier occurrence was identified in the supplied index.

\noindent\citep{gelfandzelevinsky1984models}.

\subsection{GP-1983-61 --- Computing Ext from the trivial real mixed Hodge structure}

\label{problem:GP-1983-61}

For a finite-dimensional real mixed Hodge structure \(M\), prove the quotient formula for \(\operatorname{Ext}^1(\mathbb R(0),M)\) and the vanishing \(\operatorname{Ext}^i(\mathbb R(0),M)=0\), \(i>1\), as stated in item 1983-630.10.

\paragraph{Occurrences.} 1983-630.10.

\paragraph{Status.} Solved exercise; follows from the two-term real mixed-Hodge mapping complex. No exact earlier problem occurrence was identified in the supplied cumulative index.

\noindent\citep{beilinson1986absolute,bunkenikolaustamme2018regulator}.

\subsection{GP-1983-90 --- Joint critical block-magnetization limits}

\label{problem:GP-1983-90}

For the critical planar nearest-neighbour Ising model, do the jointly rescaled block sums \(L^{-15/8}\sum_{x/L\in B_j}\sigma_x\) converge to a nontrivial limiting random field, with a consistent convention for lattice amplitude?

\paragraph{Occurrences.} 1983-631.5.

\paragraph{Status.} Closed for the nearest-neighbour planar model by Camia, Garban and Newman (2012 preprint; 2015 journal). The general universality question is separately indexed as GP-1983-91.

\noindent\citep{camia2015magnetization}.

\subsection{GP-1983-91 --- Universality of the planar critical bulk magnetization field}

\label{problem:GP-1983-91}

For zero-field, translation-invariant, square-symmetric, finite-range ferromagnetic pair interactions on \(\mathbb Z^2\) with connected interaction graph, establish the critical bulk magnetization-field limit and its identification with the nearest-neighbour field after a model-dependent multiplicative normalization. The notebook does not specify this class; it is an explicit editorial reconstruction.

\paragraph{Occurrences.} 1983-631.6.

\paragraph{Status.} Partially developed (C). Nearest-neighbour field convergence is established. Verified 2025 results treat boundary correlations under weak perturbations. A September 2026 weak-perturbation bulk preprint was located but its primary text could not be opened; it is an unverified lead, not evidence for a new theorem. No verified result found in this run covering the stated arbitrary-strength family. This chapter proves only conditional covariance scaling from uniform two-point asymptotics.

\noindent\citep{camia2015magnetization,cavagiulianigreenblatt2025boundary,openai2026isingunverified}.

\subsection{GP-1984-01 --- Characteristic numbers without derivatives of a measurable metric}

\label{problem:GP-1984-01}

Construct characteristic-number information on an oriented compact Lipschitz or quasiconformal manifold from sufficiently bounded measurable metric/conformal data, without assuming a classical curvature tensor. Distinguish global index formulas from local characteristic-cycle formulas.

\paragraph{Occurrences.} 1984-635.3.

\paragraph{Status.} Developed: Teleman (1983) and Sullivan--Teleman (1983) already supply the signature-index and rational-Pontryagin framework before the report; Connes--Sullivan--Teleman (1994, preprint 1993) provide later local quasiconformal constructions. The literal phrase 'curvature of every measurable metric' is not established.

\noindent\citep{teleman1983signature,sullivanteleman1983novikov,connessullivanteleman1994quasiconformal}.

\subsection{GP-1984-02 --- Differentiability and energy of directional slice areas}

\label{problem:GP-1984-02}

For measurable \(E\subset\mathbb R^3\) of finite volume, determine the regularity of its plane-slice area \(S_\theta(t)\) and the exact normalized spherical average of \(\int|S_\theta'(t)|^2dt\).

\paragraph{Occurrences.} 1965-86.2, 1984-635.8, 1984-635.9.

\paragraph{Status.} Closed as a classical Radon--Plancherel consequence: \(S_\theta\in H^1(\mathbb R)\) for almost every direction and the normalized spherical average equals \(2\pi|E|\). The chapter gives a full finite-measure proof and a separate ball check. The earlier occurrence has no problem ID and does not unambiguously specify the same averaging.

\noindent\citep{gelfandgraevvilenkin1966integral}.

\subsection{GP-1984-03 --- Scope of a Laurent-series criterion for algebraic complete integrability}

\label{problem:GP-1984-03}

Does a polynomial Hamiltonian system with a full \((n-1)\)-parameter Laurent family at a fixed movable pole necessarily admit compact abelian completion of generic invariant manifolds? Here \(n\) is phase-space dimension.

\paragraph{Occurrences.} 1984-637.4, 1984-637.6.

\paragraph{Status.} The count-only implication is false. The chapter's explicit polynomial example \(H=pq^2\) has the required Laurent family and, after a canonical linear change, poles in both coordinates, while generic energy curves have rational function field and projective genus zero. This does not refute ordinary Liouville integrability or a criterion with stronger global hypotheses.

\noindent\citep{adler1982kowalewski,adler1987quadrics}.

\subsection{GP-1984-04 --- Multiplicative Poisson tensors on a vector group are linear}

\label{problem:GP-1984-04}

For a finite-dimensional real vector group \((V,+)\), show that a smooth multiplicative Poisson tensor is linear in the base point and that its Jacobi identity is equivalent to a Lie bracket on \(V^*\).

\paragraph{Occurrences.} 1984-639.7.

\paragraph{Status.} Closed: multiplicativity gives \(\pi(x+y)=\pi(x)+\pi(y)\), and smooth additivity implies linearity. Evaluating Jacobi on linear coordinate functions gives the dual Lie-algebra Jacobi identity. This is the notebook's exercise, already within Drinfeld's 1983 framework.

\noindent\citep{drinfeld1983hamiltonian}.

\section{Mathematical developments}

\subsection{A bounded development of the magnetization normalization question}
\label{dev:ising}

The source of this development is item 1983-631.6 (GP-1983-91), read together with the block-limit conjecture 1983-631.5. The notebook asks for a joint limiting law independent of the starting model but gives neither a class of interactions nor a normalization convention. The strongest literal assertion, independence across arbitrary models without an amplitude adjustment, is not a well-posed universality statement. A defensible modern question fixes zero field, translation invariance, square symmetry, finite range, ferromagnetic pair interactions with a connected interaction graph, and a critical state on \(\mathbb Z^2\). It then asks for the critical magnetization field after a model-dependent multiplicative normalization.

For the nearest-neighbour square-lattice model, \citet[Theorem 1.3, Corollary 1.6, and Appendix A]{camia2015magnetization} constructs the field and block limits. The preprint circulated in May 2012; journal publication followed in 2015. \citet[Theorem 1.1]{cavagiulianigreenblatt2025boundary} treats boundary spin correlations for sufficiently weak finite-range perturbations in the half-plane. This is a different observable and a restricted perturbative regime. A September 2026 manuscript entitled \emph{Conformal universality of bulk Ising correlations under weak interactions} was returned by the current search. Its primary repository could not be opened in this run, so its proof is not verified here; it is recorded as an unresolved lead, not silently omitted or promoted to an established theorem. Its advertised weak-perturbation scope would still not cover arbitrary-strength finite-range ferromagnetic interactions. The residual question indexed here is precisely that broader bulk-field statement, not the already established nearest-neighbour limit.

The bounded contribution below is conditional and elementary. It identifies what a two-point estimate would imply, and exactly what it would fail to imply. No new universality theorem or priority claim is made.

\begin{proposition}[Covariance scaling from a uniform power-law tail]\ours
Let \(d\ge1\), \(0<\eta<d\), and let \((X_x)_{x\in\mathbb Z^d}\) be a stationary real mean-zero field with finite second moments. Write
\(\rho(h)=\mathbb E[X_0X_h]\).
Assume, for some \(A>0\) and \(C<\infty\),
\[
 \lim_{R\to\infty}\ \sup_{|h|\ge R}
 \left|\frac{\rho(h)}{A|h|^{-\eta}}-1\right|=0,
 \qquad
 |\rho(h)|\le C(1+|h|)^{-\eta}.
\]
For bounded Jordan-measurable sets \(B,C_0\subset\mathbb R^d\), put
\(S_L(B)=\sum_{x/L\in B}X_x\).
Then
\[
 L^{-2d+\eta}\operatorname{Cov}(S_L(B),S_L(C_0))
 \longrightarrow A\int_B\int_{C_0}|u-v|^{-\eta}\,du\,dv.
 \tag{D1}
\]
If \(B\) has positive volume, the only power normalization \(L^{-\alpha}S_L(B)\) whose variance has a finite strictly positive limit is
\(\alpha=d-\eta/2\).
\end{proposition}
\begin{proof}
All lattice points involved lie in a fixed box of side \(O(L)\). There is a constant \(K\), depending only on the dimension and that box, such that
\[
 \sum_{|h|\le r}(1+|h|)^{-\eta}\le K(1+r)^{d-\eta}.
\]
One proves this by counting lattice points in unit shells, or in dyadic annuli; \(d-\eta>0\) makes the resulting geometric sum finite. Thus the part of the normalized covariance sum with \(|x-y|\le\delta L\) is at most
\[
 K' L^{-d+\eta}(1+\delta L)^{d-\eta}
 \le K''\bigl(L^{-(d-\eta)}+\delta^{d-\eta}\bigr).
 \tag{D2}
\]
Absolute values were used, so no positivity assumption on each covariance is needed. The same near-diagonal estimate holds for the continuum integral by integrating \(|w|^{-\eta}\) in a ball.

For \(|x-y|>\delta L\), uniform tail asymptotics replace \(L^\eta\rho(x-y)\) by
\(A|x/L-y/L|^{-\eta}\), with an error tending to zero uniformly on this bounded region. After multiplication by \(L^{-2d}\), the sum is a Riemann sum. The integrand away from the diagonal is bounded, and its discontinuities are contained in the boundaries of \(B\), \(C_0\), and the surface \(|u-v|=\delta\), which have measure zero. The off-diagonal sum therefore converges to its integral. Take \(L\to\infty\), use (D2), and then let \(\delta\downarrow0\).

For \(B=C_0\) of positive volume, the limiting integral is finite and strictly positive. Hence
\[
 \operatorname{Var}(L^{-\alpha}S_L(B))
 \sim A I_B L^{2d-\eta-2\alpha},
 \qquad I_B=\int_B\int_B|u-v|^{-\eta}\,du\,dv>0,
\]
which proves the last assertion.
\end{proof}

For \(d=2,\eta=1/4\), the exponent is \(15/8\). The amplitude \(A\) remains in every limiting covariance; multiplying the field by \(A^{-1/2}\) removes it. This is a necessary normalization at the covariance level. The proposition does not prove tightness, non-Gaussianity, higher correlations, moment determinacy, or convergence in distribution. In particular, divergent second moments alone would not rule out a weak limit without uniform-integrability information. These are the substantive obstructions to promoting a two-point calculation to the notebook's law-level universality claim.

\begin{remark}[Reproducible numerical check]\oursComputed
The accompanying program evaluates
\[
 \rho_\eta(h)=(1+|h|^2)^{-\eta/2},
 \qquad
 L^{-4+\eta}\sum_{a,b=-(L-1)}^{L-1}
 (L-|a|)(L-|b|)\rho_\eta(a,b).
 \tag{D3}
\]
For the unit square, the continuum target is independently reduced to
\[
 8\int_0^1 (1+t^2)^{-\eta/2}
 \left(\frac1{2-\eta}-\frac{1+t}{3-\eta}
             +\frac{t}{4-\eta}\right)\,dt.
 \tag{D4}
\]
To obtain (D4), write the difference-vector integral as
\(4\int_0^1\int_0^1(1-u)(1-v)(u^2+v^2)^{-\eta/2}\,du\,dv\),
split into two triangles, and set \(v=ut\) in one triangle.
The computation uses several exponents, a direct four-index summation at small \(L\), and two quadrature resolutions. Numerical agreement is a check of normalization and implementation; convergence for all \(L\) rests on the proof.
\end{remark}

The independent review re-derived the lattice-shell estimate, the singular-diagonal bound, and the exponent. It also checked that the result is stated only about covariances and that no weak-convergence assertion follows from the variance calculation alone. A further bounded research run could verify a candidate perturbative bulk theorem against its exact hypotheses and derive uniform-integrability estimates for block observables in that family. It could not responsibly settle arbitrary-strength universality from the calculation above. The notebook does not record why the problem was dropped; no sociological explanation is inferred.

\subsection{The slice-area problem: finite measure, Sobolev regularity, and the constant}
\label{dev:slices}

Items 1984-635.8--9 refer to a question labelled unfinished in the notes. The earlier index contains the related occurrence 1965-86.2, where the directional averaging was itself underdetermined; it has no earlier problem ID. The present source makes the averaging question much clearer. This is an annotation of class E material, not a manufactured C/D research problem. The classical integral-transform source is \citet[Chapter I, Section 1.5]{gelfandgraevvilenkin1966integral}. The proof below supplies the exact convention and the finite-measure extension without claiming novelty.

\begin{theorem}[Directional slice energy for a measurable set]\ours
Let \(E\subset\mathbb R^3\) be Lebesgue measurable with \(|E|<\infty\). Define its plane-slice area almost everywhere by
\[
 S_\theta(t)=\int_{\theta^\perp}1_E(t\theta+y)\,dy,
 \qquad \theta\in S^2.
\]
For almost every \(\theta\), this function has an \(H^1(\mathbb R)\) representative and is locally absolutely continuous. With surface measure \(d\theta\) of total mass \(4\pi\),
\[
 \frac1{4\pi}\int_{S^2}\int_{\mathbb R}|S_\theta'(t)|^2\,dt\,d\theta
       =2\pi |E|.
 \tag{D5}
\]
The derivative is the weak derivative, which equals the classical derivative almost everywhere for that representative.
\end{theorem}
\begin{proof}
Use \(\widehat f(\xi)=\int e^{-ix\cdot\xi}f(x)\,dx\), so Plancherel in \(\mathbb R^n\) reads
\(\int|\widehat f|^2=(2\pi)^n\int|f|^2\).
For a Schwartz function \(f\), Fubini gives the Fourier-slice formula
\(\widehat{Rf}(\theta,\sigma)=\widehat f(\sigma\theta)\).
Consequently
\[
 \begin{split}
 \int_{S^2}\int_{\mathbb R}|\partial_tRf|^2\,dt\,d\theta
 &=\frac1{2\pi}\int_{S^2}\int_{\mathbb R}
                  \sigma^2|\widehat f(\sigma\theta)|^2\,d\sigma\,d\theta\\
 &=\frac1{\pi}\int_{\mathbb R^3}|\widehat f(\xi)|^2\,d\xi
 =8\pi^2\|f\|_2^2.
 \end{split}
 \tag{D6}
\]
The factor two in the second equality is essential: positive and negative radial parameters cover each nonzero frequency twice.

There is also an unweighted estimate. Polar coordinates give
\[
 \int_{S^2}\int_{\mathbb R}|\widehat f(\sigma\theta)|^2\,d\sigma\,d\theta
 =2\int_{\mathbb R^3}\frac{|\widehat f(\xi)|^2}{|\xi|^2}\,d\xi.
\]
On \(|\xi|\le1\), use \(|\widehat f|\le\|f\|_1\); the weight \(|\xi|^{-2}\) is locally integrable in dimension three. On \(|\xi|>1\), use Plancherel and \(|\xi|^{-2}\le1\). This bounds \(\|Rf\|_{L^2(S^2\times\mathbb R)}\) by a constant times \(\|f\|_1+\|f\|_2\).

Choose Schwartz \(f_j\to1_E\) in both \(L^1\) and \(L^2\). The last estimate and (D6) imply convergence of \(Rf_j\) and \(\partial_tRf_j\) in their respective \(L^2(S^2\times\mathbb R)\) spaces. On the other hand, for every direction for which slices are defined,
\[
 \int_{\mathbb R}|R(f_j-1_E)(\theta,t)|\,dt
       \le\|f_j-1_E\|_1 .
\]
Integration over \(S^2\) identifies the \(L^2\) limit with the slice function in \(L^1\), hence as a distribution. Integration by parts against compactly supported smooth test functions identifies the second limit as its weak \(t\)-derivative. Fubini for this Sobolev norm now yields
\(S_\theta\in H^1(\mathbb R)\) for almost every \(\theta\). Passing to the norm limit in (D6), using \(\|1_E\|_2^2=|E|\), and dividing by \(4\pi\) proves (D5).
\end{proof}

\begin{example}[Independent check on a ball]\ours
For a radius-\(R\) ball,
\(S_\theta(t)=\pi(R^2-t^2)_+\), independently of \(\theta\).
It is continuous at \(t=\pm R\), so its weak derivative has no delta contribution and equals
\(-2\pi t\,1_{\{|t|<R\}}\). Therefore
\[
 \int_{\mathbb R}|S_\theta'(t)|^2\,dt
 =4\pi^2\int_{-R}^{R}t^2\,dt
 =\frac{8\pi^2R^3}{3}
 =2\pi\,\frac{4\pi R^3}{3}.
\]
This independently rules out the notebook's \(\pi V\) under normalized directional averaging.
\end{example}

The finite-volume hypothesis cannot be omitted: \(\mathbb R^3\) has infinite plane-slice area everywhere. Even with finite volume, the theorem gives almost every direction and almost every slice parameter, not pointwise differentiability of the raw slice area for all parameters. Null-set changes of \(E\) can alter exceptional plane sections without changing its volume or the Sobolev class. The independent check kept this distinction, re-derived the Fourier normalization, and tested multiple radii directly.

\subsection{A Laurent family does not force compact abelian completion}
\label{dev:laurent}

The reconstruction tested here is the bare implication in item 1984-637.6. \citet[p.84]{adler1982kowalewski} already distinguishes a necessary pole-family test from an unproved sufficiency statement. Its framework contains more geometric information than a parameter count. We test exactly the count-only assertion with the compact-abelian definition used immediately before it in the notebook.

\begin{proposition}[A polynomial one-degree-of-freedom counterexample]\ours
A polynomial Hamiltonian system on \(\mathbb C^2\) can have a one-parameter family of Laurent solutions with a fixed movable pole, even with both chosen canonical coordinates having poles, while its generic energy curve has no smooth projective abelian completion.
\end{proposition}
\begin{proof}
For \(\{q,p\}=1\), take \(H=pq^2\). Hamilton's equations are
\[
 \dot q=q^2,\qquad \dot p=-2pq.
\]
For arbitrary \(h,t_0\in\mathbb C\),
\[
 q(t)=-\frac1{t-t_0},\qquad p(t)=h(t-t_0)^2
 \tag{D7}
\]
solves both equations and has energy \(h\). Fixing the pole at \(t_0=0\) leaves one freely variable parameter, equal to \(n-1\) for phase-space dimension \(n=2\).
If poles in every coordinate are desired, take
\[
 Q=q+p,\qquad P=q+2p .
\]
The change has determinant one and
\(\{Q,P\}=2\{q,p\}+\{p,q\}=1\).
Its inverse is \(q=2Q-P,\ p=P-Q\); the polynomial Hamiltonian becomes
\((P-Q)(2Q-P)^2\), while
\[
 Q=-t^{-1}+ht^2,\qquad P=-t^{-1}+2ht^2.
\]
Both coordinates now have a genuine pole for every \(h\).

For \(h\ne0\), the energy curve \(pq^2=h\) is isomorphic to
\(\mathbb C^*\) via \(q\mapsto(q,h/q^2)\). Its function field is
\(\mathbb C(q)\); therefore its smooth projective completion has genus zero and is \(\mathbb P^1\). An abelian variety of dimension one has genus one, so no such abelian completion exists.
\end{proof}

This is not a counterexample to ordinary Liouville integrability: one autonomous Hamiltonian is the required integral in two-dimensional phase space. The curve is isomorphic to the algebraic group \(\mathbb C^*\), but the vector field \(q^2\partial_q\) is not invariant under its multiplicative translations. Thus no linear-flow assertion on a generalized abelian group is being made. The exact failure is the passage from local Laurent families to a compact abelian invariant variety. The independent review checked the symplectic coordinate change and the function-field obstruction separately from substitution into the ODE. A plausible next investigation would impose and test explicit completeness and boundary-extension conditions; no unrestricted improved converse is claimed here.

\subsection{The quantifier missing from the coisotropic-invariants criterion}
\label{dev:coisotropic}

Item 1984-639.19 is a near-term publication match with a false literal converse. The sufficient direction is \citet[Theorem 6]{semenov1985dressing}. We reproduce its short mechanism in the notebook's notation and give an explicit counterexample to necessity for a single action.

\begin{proposition}[Sufficient condition for invariant functions]\ours
Let a connected Poisson--Lie group \(G\) act Poissonly on \(M\), and let \(H\subset G\) be a connected Lie subgroup. If its annihilator
\(\mathfrak h^\perp\subset\mathfrak g^*\) is closed under the dual Lie bracket, then \(C^\infty(M)^H\) is closed under the Poisson bracket.
\end{proposition}
\begin{proof}
For \(f,g\) invariant under \(H\), define the orbit differentials at \(x\) by
\(\alpha_f=d_e[u\mapsto f(u\cdot x)]\) and similarly for \(g\).
They annihilate \(\mathfrak h\). The differentiated Poisson-action identity gives, with a consistent infinitesimal sign convention,
\[
 X_M\{f,g\}-\{X_Mf,g\}-\{f,X_Mg\}
       =\langle[\alpha_f,\alpha_g]_*,X\rangle .
\]
For \(X\in\mathfrak h\), the last pairing is zero by the annihilator hypothesis, and the two middle terms vanish by invariance. Thus all infinitesimal \(H\)-derivatives of \(\{f,g\}\) vanish. Connectedness of \(H\) gives global \(H\)-invariance.
\end{proof}

\begin{example}[Necessity fails for a fixed action]\ours
Take the additive group \(G=\mathbb R^3\), with coordinates \(x,y,z\) and
\[
 \pi_G=z\,\partial_x\wedge\partial_y .
\]
It is Poisson: the only nonzero coordinate bracket is \(\{x,y\}=z\), and \(z\) is a Casimir, so the Jacobi identities hold. It is multiplicative because its coefficients are linear and addition has identity derivative. At the Lie-bialgebra level, \(\mathfrak g\) is abelian and \(\mathfrak g^*\) is Heisenberg:
\[
 [dx,dy]_*=dz,\qquad [dx,dz]_*=[dy,dz]_*=0.
\]
Let \(H\) be the \(z\)-axis. It is a closed connected subgroup, but
\(\mathfrak h^\perp=\operatorname{span}\{dx,dy\}\) is not a Lie subalgebra.

Now let \(G\) act trivially on any Poisson manifold \(M\), for instance the symplectic plane. The action map is projection \(G\times M\to M\), hence Poisson. All functions on \(M\) are \(H\)-invariant and form a Poisson algebra. This checks every hypothesis while disproving the necessity asserted for an arbitrary single action.
\end{example}

The example also shows why merely adding connectedness would not repair the converse. One must change the quantifier over actions or use an action rich enough to test all annihilator covectors, as in the regular homogeneous-quotient problem. The independent review checked both the multiplicativity of the linear tensor and the dual Heisenberg bracket. No new general reduction theorem is asserted.

\subsection{Two elementary symbol checks used in the correction ledger}
\label{dev:symbols}

\begin{remark}[Matrix and vector-calculus repairs]\oursComputed
For commuting symbols \(a,b\), the notebook's matrix
\[
 \begin{pmatrix}a&ib\\-ib&a\end{pmatrix}^{\!2}
 =
 \begin{pmatrix}a^2+b^2&2iab\\-2iab&a^2+b^2\end{pmatrix}.
\]
Replacing its lower-right entry by \(-a\) removes the off-diagonal terms. Separately, for a scalar \(f\) and a three-vector field \(v\), the operator
\[
 Q(f,v)=(-\operatorname{div}v,\nabla f+\operatorname{curl}v)
\]
satisfies \(Q^2=-\Delta I_4\), by
\(\operatorname{div}\operatorname{curl}=0\),
\(\operatorname{curl}\nabla=0\), and
\(\operatorname{curl}^2=\nabla\operatorname{div}-\Delta\).
These are four fibre components over \(\mathbb R^3\). The program independently multiplies the \(4\times4\) symbol at several vectors of commuting differential symbols and compares it with \(-|k|^2I_4\). These checks justify the local emendations and are not research novelty claims.
\end{remark}

